\documentclass[10pt, letterpaper]{article}

% Inhaltsverzeichnis für Pakettypen (nur für Übersicht im Header, wird nicht im Dokument angezeigt)
% 1. Seitenlayout und Ränder
% 2. Sprache und Zeichensatz
% 3. Mathematik und Theorem-Umgebungen
% 4. Eigene Makros
% 5. Diagramme und Grafiken
% 6. Tabellen und Aufzählungen
% 7. Inhaltsverzeichnis
% 8. Abschnittsüberschriften
% 9. Abstrakt-Umgebung
% 10. Todos/Notizen
% 11. Rahmen/Box-Umgebungen
% 12. Python-Integration
% 13. Literaturverwaltung
% 14. Hyperlinks
% 15. Absatzeinstellungen
% 16. Umgebungen
% 17  Graphik
% 18  Extra
% 00. Titel und Autor

\usepackage{slashed}

% --- 1. Seitenlayout und Ränder ---
\usepackage[margin=3cm]{geometry}

% --- 2. Sprache und Zeichensatz ---
\usepackage[english]{babel}
\usepackage[T1]{fontenc}
\usepackage[utf8]{inputenc}

% --- 3. Mathematik und Theorem-Umgebungen ---
\usepackage{amsmath, amssymb, amsthm}
\usepackage{mathrsfs}
\DeclareMathOperator{\WF}{WF}


% --- 4. Eigene Makros ---
\usepackage{xcolor}
\newcommand{\SKP}{\langle\cdot,\cdot\rangle}
\newcommand{\R}{\mathbb{R}}
\newcommand{\N}{\mathbb{N}}
\newcommand{\Q}{\mathbb{Q}}
\newcommand{\Z}{\mathbb{Z}}
\newcommand{\C}{\mathbb{C}}
\newcommand{\entwurf}[1]{\textcolor{red}{#1}}

% --- 5. Diagramme und Grafiken ---
\usepackage{graphicx}
\graphicspath{{../../Library/Images/}}
\usepackage[export]{adjustbox}
\usepackage{tikz}
\usetikzlibrary{decorations.pathreplacing, arrows.meta, positioning}
\usepackage{tikz-cd}

% --- 6. Tabellen und Aufzählungen ---
\usepackage{enumitem}
\setlist[itemize]{left=0.5cm}

\newenvironment{romanenum}[1][]
  {%
    \ifx&#1&
    \else
      \textbf{#1}\quad
    \fi
    \begin{enumerate}[label=\roman*)]
  }
  {%
    \end{enumerate}%
  }

% --- 7. Inhaltsverzeichnis ---
\usepackage{tocloft}
\renewcommand{\cftsecfont}{\footnotesize}
\renewcommand{\cftsubsecfont}{\footnotesize}
\renewcommand{\cftsubsubsecfont}{\footnotesize}
\renewcommand{\cftsecpagefont}{\footnotesize}
\renewcommand{\cftsubsecpagefont}{\footnotesize}
\renewcommand{\cftsubsubsecpagefont}{\footnotesize}
\usepackage{etoc}

% --- 8. Abschnittsüberschriften ---
\usepackage{titlesec}
\titleformat{\section}{\normalfont\large\bfseries}{\thesection}{1em}{}
\titleformat{\subsection}{\normalfont\normalsize\bfseries}{\thesubsection}{0.5em}{}
\titleformat{\subsubsection}{\normalfont\normalsize\bfseries}{\thesubsubsection}{0.5em}{}
\setcounter{secnumdepth}{4}

% --- 9. Abstrakt-Umgebung ---
\usepackage{changepage}
\renewenvironment{abstract}
  {
    \begin{adjustwidth}{1.5cm}{1.5cm}
    \small
    \textsc{Abstract. –}%
  }
  {
    \end{adjustwidth}
  }

% --- 10. Todos/Notizen ---
\usepackage{todonotes}
% Hellblau definieren
\definecolor{hellblau}{rgb}{0.3,0.6,1}
\newenvironment{note}{\par\begingroup\color{hellblau}\itshape}{\par\endgroup}

% --- 11. Rahmen/Box-Umgebungen ---
\usepackage{mdframed}
\usepackage{tcolorbox}
\colorlet{shadecolor}{gray!25}

\newenvironment{customTheorem}
  {\vspace{10pt}%
   \begin{mdframed}[
     backgroundcolor=gray!20,
     linewidth=0pt,
     innertopmargin=10pt,
     innerbottommargin=10pt,
     skipabove=\dimexpr\topsep+\ht\strutbox\relax,
     skipbelow=\topsep,
   ]}
  {\end{mdframed}
   \vspace{10pt}%
  }

% --- 12. Python-Integration ---
% (Deaktiviert in dieser Version, aktiviere bei Bedarf)
% \usepackage{pythontex}
% \usepackage[makestderr]{pythontex}

% --- 13. Literaturverwaltung ---
\usepackage{csquotes}
\usepackage[backend=biber, style=alphabetic, citestyle=alphabetic]{biblatex}
\addbibresource{bibliography.bib}

% --- 14. Hyperlinks ---
\usepackage{hyperref}
\hypersetup{
  colorlinks   = true,
  urlcolor     = blue,
  linkcolor    = blue,
  citecolor    = blue,
  frenchlinks  = true
}

% --- 15. Absatzeinstellungen ---
\usepackage[parfill]{parskip}
\sloppy

% --- 16. Umgebungen ---
\usepackage{thmtools}

\newcommand{\CustomHeading}[3]{%
  \par\medskip\noindent%
  \textbf{#1 #2} \textnormal{(#3)}.\enskip%
}

\newenvironment{DEF}[2]{\begin{unitbox}\CustomHeading{Definition}{#1}{#2}}{\end{unitbox}}
\newenvironment{PROP}[2]{\begin{unitbox}\CustomHeading{Proposition}{#1}{#2}}{\end{unitbox}}
\newenvironment{THEO}[2]{\begin{unitbox}\CustomHeading{Theorem}{#1}{#2}}{\end{unitbox}}
\newenvironment{LEM}[2]{\begin{unitbox}\CustomHeading{Lemma}{#1}{#2}}{\end{unitbox}}
\newenvironment{KORO}[2]{\begin{unitbox}\CustomHeading{Corollar}{#1}{#2}}{\end{unitbox}}
\newenvironment{REM}[2]{\begin{unitbox}\CustomHeading{Remark}{#1}{#2}}{\end{unitbox}}
\newenvironment{EXA}[2]{\begin{unitbox}\CustomHeading{Example}{#1}{#2}}{\end{unitbox}}
\newenvironment{STUD}[2]{\begin{unitbox}\CustomHeading{Study}{#1}{#2}}{\end{unitbox}}
\newenvironment{CONC}[2]{\begin{unitbox}\CustomHeading{Concept}{#1}{#2}}{\end{unitbox}}
\newenvironment{OTH}[2]{\begin{unitbox}\CustomHeading{Other}{#1}{#2}}{\end{unitbox}}
\newenvironment{EXE}[2]{\begin{unitbox}\CustomHeading{Exercise}{#1}{#2}}{\end{unitbox}}
\newenvironment{MOT}[2]{\begin{unitbox}\CustomHeading{Motivation}{#1}{#2}}{\end{unitbox}}
\newenvironment{PROOF}[2]{\begin{unitbox}\CustomHeading{Proof}{#1}{#2}}{\end{unitbox}}

% --- Unit Umgebung für Source-Inhalte ---
\usepackage{mdframed}
\newmdenv[
  linewidth=1pt,
  topline=false,
  bottomline=false,
  rightline=false,
  leftmargin=0cm,
  rightmargin=0cm,
  skipabove=10pt,
  skipbelow=10pt,
  innertopmargin=0.5\baselineskip,
  innerbottommargin=0.5\baselineskip,
  backgroundcolor=gray!10,
  linecolor=gray
]{unitbox}

\newenvironment{unit}[1]
  {\begin{unitbox}\textbf{Unit #1}\par\smallskip}
  {\end{unitbox}}

% --- 17. ... ---

% --- 18. Extras ---
\usepackage{stmaryrd}
\usepackage{bbold}  % falls du \mathbb{1} nutzen willst

% --- 00. Titel und Autor ---
\title{Lineare Funktionalanalysis}
\author{Tim Jaschik}
\date{\today}

\begin{document}

\maketitle
\rule{\textwidth}{0.5pt}
\begin{abstract}
Kurze Beschreibung …
\end{abstract}
\rule{\textwidth}{0.5pt}
\vspace{0.5cm}

\tableofcontents

\pagebreak

\section*{§ 0 Einführung}
Schwerpunkt der Vorlesung ist die systematische Untersuchung unendlichdimensionaler normierter Räume und der stetigen linearen Abbildungen zwischen diesen Räumen. Die hierbei entwickelte Theorie ist Grundvoraussetzung für die Behandlung wichtiger Problemstellungen u.a. im Bereich der partiellen Differentialgleichungen, der mathematischen Physik, der Numerik und der Stochastik.\\
Ein wesentliches Merkmal der Funktionalanalysis ist das fruchtbare Zusammenspiel analytischer und linear algebraischer Zusammenhänge. Eine Ausgangsfrage der linearen Algebra ist die Frage der Lösbarkeit linearer Gleichungssysteme der Form

$$
T a=b
$$

Hier und $\operatorname{im}$ Folgenden sei $\mathbb{K}=\mathbb{R}$ oder $\mathbb{K}=\mathbb{C}, b \in \mathbb{K}^{m}$ und eine Matrix $T \in \mathbb{K}^{m \times n}$ seien gegeben, und die Menge der Lösungen $a \in \mathbb{K}^{n}$ ist gesucht. Lösbar ist dieses Gleichungssystem genau dann, wenn $b$ im Erzeugnis der Spalten von $T$ liegt, und dies ist leicht zu entscheiden. Im Falle der Lösbarkeit kann man den affinen Lösungsraum dann als Summe einer speziellen Lösung und des Kerns von $T$ schreiben, wobei man diesen Kern ebenfalls leicht durch Zeilenumformungen von $T$ bestimmen kann.\\
In der Funktionalanalysis betrachtet man u.a. ähnlich aussehende Probleme der Form

$$
T u=f .
$$

Hierbei ist $f \in F$ gegeben und $u \in E$ gesucht, wobei nun $E$ und $F$ unendlichdimensionale Vektorräume über $\mathbb{K}$ seien und eine $\mathbb{K}$-lineare Abbildung $T: E \rightarrow F$ gegeben ist. In Anwendungen sind dabei $u$ und $f$ oft Funktionen, d.h. Elemente von Funktionenräumen. Wir werden sehen, dass wir zur strukturellen Untersuchung solche unendlich-dimensionalen Probleme analytische Eigenschaften ins Spiel bringen müssen. Insbesondere werden wir vorraussetzen, dass $E$ und $F$ normierte Vektorräume sind und $T$ stetig ist.

\section*{Einfaches Beispiel}
Wir betrachten eine eingespannte Saite, welche durch den Graph einer Funktion $u:[0,1] \rightarrow \mathbb{R}$ mit $u(0)=u(1)=0$ beschrieben sei. Auf diese Saite möge nun die\\
vertikale Kraft $f(x)$ am Punkt $x \in[0,1]$ wirken.\\
Im Falle der Gewichtskraft wäre dies z.B. $f(x)=-c$ für $x \in[0,1]$, wobei $c$ als Produkt der Massendichte der Saite und der Schwerebeschleunigung $g=9,81 \mathrm{~m} / \mathrm{s}^{2}$ gegeben ist. Die Form der Saite wird dann (bis auf eine Konstante) durch die Gleichungen

$$
-u^{\prime \prime}(x)=f(x) \quad x \in(0,1), \quad u(0)=u(1)=0
$$

beschrieben. Diese (unendlich vielen) Gleichungen kann man zusammenfassend in der Form (1) schreiben, wenn man z.B. die Vektorräume

$$
E:=\left\{u \in C^{2}([0,1], \mathbb{K}): u(0)=u(1)=0\right\}, \quad F:=C([0,1], \mathbb{K})
$$

und die lineare Abbildung $T: E \rightarrow F, T u=-u^{\prime \prime}$ betrachtet. Glücklicherweise ist in diesem Fall das Problem eindeutig lösbar, und die Lösung ist gegeben durch

$$
u:=T^{-1} f \quad \text { d.h. } \quad u(x)=\left[T^{-1} f\right](x)=\int_{0}^{1} G(x, y) f(y) d y \quad \text { für } y \in[0,1]
$$

mit der Greenfunktion

$$
G(x, y)= \begin{cases}(1-y) x, & x \leq y \\ (1-x) y, & x>y\end{cases}
$$

(Beweis als Übungsaufgabe, Blatt 1). Man beachte, dass $T^{-1}: F \rightarrow E$ ein sogenannter Integraloperator ist, den man als Matrix mit kontinuierlich vielen Einträgen betrachten kann.\\
Damit wäre das einfache Problem (2) bereits gelöst, allerdings schließen sich z.B. folgende wichtige Fragen an, auf welche die Funktionalanalysis Antworten geben kann:\\
i). Welche Vektorräume muss man betrachten, wenn man unstetige vertikale Kräfte $f$ zulassen möchte?\\
ii). Wie sieht die Lösungsmenge des (ähnlich aussehenden) Problems

$$
-\left(p u^{\prime}\right)^{\prime}+q u=f \quad \text { in }(0,1), \quad u(0)=u(1)=0
$$

mit gegebenen Funktionen $p \in C^{1}([0,1]), q \in C([0,1])$ aus?\\
Hier ist $T: E \rightarrow F$ also gegeben durch $T u=-\left(p u^{\prime}\right)^{\prime}+q u$.\\
iii). Gibt es Eigenfunktionen für die Probleme (2) und (3), d.h. Funktionen $u \in E \subset F$ und $\lambda \in \mathbb{C}$ mit $T u=\lambda u$ ? Wie sieht die Menge der zugehörigen Eigenwerte aus?\\
iv). Wie behandelt man eine mehrdimensionale Version des Problems (Eingespannte Membran im vertikalen Kraftfeld)?

Wir werden auf Probleme der Form (1) in systematischer Weise zurückkommen. Zuvor müssen wir aber grundlegende Eigenschaften von unendlich-dimensionalen normierten Räumen und stetigen linearen Abbildungen verstehen.

\section*{§ 1 Normierte Räume und Prähilberträume}
Bezeichnungen: Stets sei $\mathbb{K}=\mathbb{R}$ oder $\mathbb{K}=\mathbb{C}$ und $E$ ein $\mathbb{K}$-Vektorraum.

\subsection*{Definition}
i) Eine Abbildung $\|\cdot\|: E \rightarrow \mathbb{R}$ heißt Halbnorm auf $E$, falls für alle $x, y \in E$ und $\lambda \in \mathbb{K}$ gilt:

$$
\begin{array}{ll}
(N 1) & \|\lambda x\|=|\lambda|\|x\| \quad \text { (Homogenität) } \\
(N 2) & \|x+y\| \leq\|x\|+\|y\| \quad \text { (Dreiecksungleichung) }
\end{array}
$$

Es folgt dann: $\left\|0_{E}\right\|=\left\|0_{\mathbb{K}} \cdot 0_{E}\right\| \stackrel{(N 1)}{=}\left|0_{\mathbb{K}}\right|\left\|0_{E}\right\|=0$ und damit $0=\|x+(-x)\| \stackrel{(N 2)}{\leq} \|x\|+\|-x\| \stackrel{(N 1)}{=} 2\|x\|$, also

$$
(N 3) \quad\|x\| \geq 0 \text { für alle } x \in E \text {. }
$$

ii) Eine Halbnorm $\|\cdot\|$ auf $E$ heißt Norm, falls zusätzlich für alle $x \in E$ gilt:

$$
\text { (N4) } \quad\|x\|=0 \Rightarrow x=0 \text { (Definitheit) }
$$

In diesem Fall heißt das Paar $(E,\|\cdot\|)$ normierter Raum.

\subsection*{Definition}
Zwei Halbnormen $\|\cdot\|_{1},\|\cdot\|_{2}$ auf $E$ heißen äquivalent, falls $a, b>0$ existieren mit

$$
a\|x\|_{1} \leq\|x\|_{2} \leq b\|x\|_{1} \text { für alle } x \in E \text {. }
$$

Wir schreiben dann: $\|\cdot\|_{1} \sim\|\cdot\|_{2}$.\\
Bemerkung: „\~{}" definiert Äquivalenzrelationen auf der Menge der Halbnormen und auf der Menge der Normen auf $E$.

\subsection*{Bemerkung und Beispiel}
i) $\operatorname{dim} E<\infty \Rightarrow$ Alle Normen auf $E$ sind äquivalent (bekannt aus dem Grundstudium).\\
ii) $E=\mathbb{K}^{N}$. Für $x \in \mathbb{K}^{N}$ und $p \in[1, \infty]$ sei

$$
|x|_{p}:= \begin{cases}\left(\left|x_{1}\right|^{p}+\ldots+\left|x_{N}\right|^{p}\right)^{\frac{1}{p}}, & p<\infty \\ \max _{1 \leq i \leq N}\left|x_{i}\right| & p=\infty\end{cases}
$$

$\underline{\text { Bekannt aus der Analysis II: }}|\cdot|_{p}: E \rightarrow \mathbb{R}$ ist eine Norm auf $E$. Dabei gilt die Höldersche Ungleichung:\\
$\sum_{i=1}^{N}\left|x_{i} y_{i}\right| \leq|x|_{p}|y|_{p^{\prime}}$ für alle $x, y \in \mathbb{K}^{N}$ mit $\quad p^{\prime}:=\left\{\begin{array}{lll}\frac{p}{p-1} & \text { falls } & 1<p<\infty \\ \infty & \text { falls } & p=1 \\ 1 & \text { falls } & p=\infty\end{array}\right.$\\
$p^{\prime}$ heißt konjugierter Exponent zu $p$. Beachte: $p, q \in(1, \infty)$ konjugiert $\Longleftrightarrow \frac{1}{p}+\frac{1}{q}=1$\\
iii) Sei $M$ eine beliebige Menge, und sei $\ell^{\infty}(M, \mathbb{K})$ der $\mathbb{K}$-Vektorraum der beschränkten Funktionen $u: M \rightarrow \mathbb{K}$ mit punktweisen linearen Operationen.\\
Dann ist $\ell^{\infty}(M, \mathbb{K})$ ein normierter $\mathbb{K}$-Vektorraum mit der Supremumsnorm $\|\cdot\|_{\infty}$ definiert durch $\|u\|_{\infty}:=\sup _{z \in M}|u(z)|$.\\
Der Beweis der Normeigenschaften ist sehr einfach.\\
Ist $M=\mathbb{N}$, so notiert man die Funktionswerte von $u \in \ell^{\infty}(\mathbb{N}, \mathbb{K})$ als Folge, d. h. man schreibt $u=\left(u_{k}\right)_{k}$ mit $u_{k}=u(k)$. Dann ist $\|u\|_{\infty}=\sup _{k \in \mathbb{N}}\left|u_{k}\right|$.

Kurzschreibweisen: $\ell^{\infty}(M)$ statt $\ell^{\infty}(M, \mathbb{K})$, $\ell^{\infty}$ statt $\ell^{\infty}(\mathbb{N})=\ell^{\infty}(\mathbb{N}, \mathbb{K})$.

Wir wollen im Folgenden weitere unendlich-dimensionale normierte Räume durch ein allgemeines Prinzip konstruieren. Dieses Prinzip liefert auch einen (alternativen) Beweis der Halbnormeigenschaften von $|\cdot|_{p}$ aus $1.3($ ii $)$.

\subsection*{Definition}
(a) Eine Teilmenge $A \subset E$ heißt\\
(i) konvex, wenn für alle $x, y \in A, \lambda \in[0,1]$ gilt: $(1-\lambda) x+\lambda y \in A$;\\
(ii) symmetrisch, wenn gilt: $x \in A \Longrightarrow \theta x \in A$ für alle $\theta \in \mathbb{K}$ mit $|\theta|=1$;\\
(iii) absolut konvex, wenn $A$ konvex und symmetrisch ist.\\
(iv) absorbierend, wenn für jedes $x \in E$ ein $s>0$ existiert mit $s x \in A$.\\
(b) Eine auf einer konvexen Teilmenge $A \subset E$ definierte Funktion

$$
f: A \rightarrow \mathbb{R} \cup\{\infty\}
$$

heißt konvex, wenn $f((1-\lambda) x+\lambda y) \leq(1-\lambda) f(x)+\lambda f(y)$ für alle $x, y \in A$, $\lambda \in[0,1]$ gilt.\\
Ist $A$ absolut konvex, so heißt $f$ absolut konvex, wenn $f$ konvex und gerade ist. Dabei heißt $f$ gerade, wenn $f(\theta x)=f(x)$ für alle $x \in A$ und $\theta \in \mathbb{K}$ mit $|\theta|=1$ gilt.

\subsection*{Bemerkung und Beispiel}
i). Ist $\|\cdot\|$ eine Halbnorm auf $E$ und $r>0$, so sind die Mengen $U_{r}(0):=\{x \in E$ : $\|x\|<r\}$ und $B_{r}(0):=\{x \in E:\|x\| \leq r\}$ absolut konvex und absorbierend.\\
ii). $A \subset E$ ist absolut konvex genau dann, wenn $\lambda x+\mu y \in A$ für alle $x, y \in A$, $\lambda, \mu \in \mathbb{K}$ mit $|\lambda|+|\mu| \leq 1$ gilt.\\
iii). Ist $A \subset E$ konvex, so gilt $\sum_{i=1}^{n} \lambda_{i} x_{i} \in A$ für alle $n \in \mathbb{N}, x_{1}, \ldots, x_{n} \in A$ und $\lambda_{1}, \ldots, \lambda_{n} \in[0, \infty)$ mit $\sum_{i=1}^{n} \lambda_{i}=1$.\\
iv). Ist $A \subset E$ absolut konvex, so gilt $\sum_{i=1}^{n} \lambda_{i} x_{i} \in A$ für alle $n \in \mathbb{N}, x_{1}, \ldots, x_{n} \in A$ und $\lambda_{1}, \ldots, \lambda_{n} \in \mathbb{K}$ mit $\sum_{i=1}^{n}\left|\lambda_{i}\right| \leq 1$.\\
v). Sei $I \subset \mathbb{R}$ ein Intervall. Dann ist $f \in C^{1}(I)$ genau dann konvex, wenn $f^{\prime}$ monoton wachsend ist.

Die Beweise von i)-v) sind allesamt recht einfach.

\subsection*{Satz}
Sei $A \subset E$ absolut konvex und absorbierend. Dann wird durch

$$
\|x\|_{A}:=\inf \left\{t>0: \frac{x}{t} \in A\right\}=\inf \{t>0: x \in t A\} \quad \text { für } x \in E
$$

eine Halbnorm auf $E$ definiert.

Beweis. Da $A$ absorbierend ist, erhält man $\|x\|_{A}<\infty$ für alle $x \in E$.\\
Zur Homogenität: Für $x \in E$ und $\lambda \in \mathbb{K}$ gilt

$$
\begin{aligned}
\|\lambda x\|_{A} & =\inf \{t>0: \lambda x \in t A\}=\inf \{t>0:|\lambda| x \in t A\} \\
& =\inf \{|\lambda| s: s>0, x \in s A\}=|\lambda|\|x\|_{A} .
\end{aligned}
$$

Zur Dreiecksungleichung: Seien $x, y \in E$, und seien $s, t>0$ mit $\frac{x}{s}, \frac{y}{t} \in A$. Mit $\lambda=\frac{t}{t+s}$ (also $1-\lambda=\frac{s}{t+s}$ ) ist dann

$$
\frac{x+y}{s+t}=(1-\lambda) \frac{x}{s}+\lambda \frac{y}{t} \in A
$$

aufgrund der Konvexität von $A$, und somit folgt $\|x+y\|_{A} \leq s+t$. Durch Infimumbildung in $s$ und $t$ folgt $\|x+y\|_{A} \leq\|x\|_{A}+\|y\|_{A}$. \(\square\)

\subsection*{Definition und Satz}
Für $p \in[1, \infty)$ sei $\ell^{p}(\mathbb{N}, \mathbb{K})$ die Menge aller Folgen $x=\left(x_{k}\right)_{k}$ in $\mathbb{K}$ mit $\sum_{k \in \mathbb{N}}\left|x_{k}\right|^{p}<\infty$. $\ell^{p}(\mathbb{N}, \mathbb{K}$ ) ist ein normierter $\mathbb{K}$-Vektorraum (bzgl. komponentenweiser Addition und Skalarmultiplikation) mit der Norm $\|\cdot\|_{p}$ definiert durch

$$
\|x\|_{p}=\left(\sum_{k \in \mathbb{N}}\left|x_{k}\right|^{p}\right)^{\frac{1}{p}} \quad \text { für } x=\left(x_{k}\right)_{k} \in \ell^{p}(\mathbb{N}, \mathbb{K}) .
$$

Beweis. Offensichtlich ist die Nullfolge $0 \in \ell^{p}(\mathbb{N}, \mathbb{K})$. Seien $x=\left(x_{k}\right)_{k}, y=\left(y_{k}\right)_{k} \in \ell^{p}(\mathbb{N}, \mathbb{K})$ und $\lambda \in \mathbb{K}$. Dann ist

$$
\sum_{k \in \mathbb{N}}\left|\lambda x_{k}\right|^{p}=|\lambda|^{p} \sum_{k}\left|x_{k}\right|^{p}<\infty,
$$

und somit $\lambda x=\left(\lambda x_{k}\right)_{k} \in \ell^{p}(\mathbb{N}, \mathbb{K})$. Da die Funktion $(0, \infty) \rightarrow \mathbb{R}, t \mapsto t^{p}$ gemäß 1.5 v ) konvex ist, gilt ferner

$$
\sum_{k \in \mathbb{N}}\left|x_{k}+y_{k}\right|^{p} \leq \sum_{k \in \mathbb{N}} 2^{p}\left(\frac{\left|x_{k}\right|}{2}+\frac{\left|y_{k}\right|}{2}\right)^{p} \leq 2^{p-1} \sum_{k \in \mathbb{N}}\left(\left|x_{k}\right|^{p}+\left|y_{k}\right|^{p}\right)<\infty
$$

und somit $x+y=\left(x_{k}+y_{k}\right)_{k} \in \ell^{p}(\mathbb{N}, \mathbb{K})$. Also ist $\ell^{p}(\mathbb{N}, \mathbb{K})$ ein $\mathbb{K}$-Vektorraum. Zudem ist

$$
A:=\left\{x \in \ell^{p}(\mathbb{N}, \mathbb{K}): \sum_{k \in \mathbb{N}}\left|x_{k}\right|^{p} \leq 1\right\}
$$

offensichtlich eine absorbierende Teilmenge von $\ell^{p}(\mathbb{N}, \mathbb{K})$, und aufgrund der Konvexität der Funktion $t \mapsto t^{p}$ ist $A$ auch absolut konvex. Somit definiert

$$
\|x\|_{A}=\inf \left\{t>0: \sum_{k \in \mathbb{N}}\left|\frac{x_{k}}{t}\right|^{p} \leq 1\right\}=\inf \left\{t>0: \sum_{k \in \mathbb{N}}\left|x_{k}\right|^{p} \leq t^{p}\right\}=\left(\sum_{k \in \mathbb{N}}\left|x_{k}\right|^{p}\right)^{\frac{1}{p}}=\|x\|_{p} .
$$

eine Halbnorm auf $\ell^{p}(\mathbb{N}, \mathbb{K})$ gemäß 1.6.\\
Zum Beweis der Definitheit (N4) beachte man:

$$
\|x\|_{p}=0 \Longrightarrow \sum_{k \in \mathbb{N}}\left|x_{k}\right|^{p}=0 \Longrightarrow x_{k}=0 \text { für alle } k \in \mathbb{N} \Longrightarrow x=0 \text {. }
$$ \(\square\)

\subsection*{Bemerkung}
i). Analog definiert man $\ell^{p}(I, \mathbb{K})$ für eine beliebige höchstens abzählbare Indexmenge $I$ anstelle von $\mathbb{N}$, z. B. $I=\mathbb{Z}$.\\
Kurzschreibweisen: $\ell^{p}(I)$ statt $\ell^{p}(I, \mathbb{K})$

$$
\ell^{p} \text { statt } \ell^{p}(\mathbb{N})=\ell^{p}(\mathbb{N}, \mathbb{K})
$$

Man beachte: Der Fall $I=\{1, \ldots, N\}$ führt wieder zurück auf $\mathbb{K}^{N}$.\\
ii). Sind $p, q \in[1, \infty]$ konjugierte Exponenten, so gilt die Höldersche Ungleichung

$$
\sum_{k \in I}\left|x_{k} y_{k}\right| \leq\|x\|_{p}\|y\|_{q} \quad \text { für } x=\left(x_{k}\right)_{k} \in l^{p}(I, \mathbb{K}), y=\left(y_{k}\right)_{k} \in l^{q}(I, \mathbb{K})
$$

Im Fall $p=\infty, q=1$ sieht man dies direkt:

$$
\sum_{k \in I}\left|x_{k} y_{k}\right| \leq\left(\sup _{k \in I}\left|x_{k}\right|\right)\left(\sum_{k \in I}\left|y_{k}\right|\right)=\|x\|_{1}\|y\|_{\infty}
$$

Im Fall $p, q \in(1, \infty)$ verwenden wir die Youngsche Ungleichung

$$
(Y) \quad a b \leq \frac{a^{p}}{p}+\frac{b^{q}}{q} \quad \text { für } a, b \geq 0 .
$$

Diese Ungleichung beweist man leicht durch Minimierung der Funktion $x \mapsto \frac{x^{p}}{p}-x b$ in $[0, \infty)$ für festes $b \geq 0$.\\
Unter Verwendung von (Y) folgt nun mit $s=\|x\|_{p}, t=\|y\|_{q}$ :

$$
\begin{aligned}
\sum_{k \in I}\left|x_{k} y_{k}\right|=s t \sum_{k \in I}\left|\frac{x_{k}}{s} \frac{y_{k}}{t}\right| & \leq s t \sum_{k \in I}\left(\frac{1}{p}\left|\frac{x_{k}}{s}\right|^{p}+\frac{1}{q}\left|\frac{y_{k}}{t}\right|^{q}\right) \\
& =\operatorname{st}\left(\frac{1}{p}\left\|\frac{x}{s}\right\|_{p}^{p}+\frac{1}{q}\left\|\frac{y}{t}\right\|_{q}^{q}\right)=\operatorname{st}\left(\frac{1}{p}+\frac{1}{q}\right)=\text { st. }
\end{aligned}
$$

\subsection*{Satz}
Sei $p \in[1, \infty)$.\\
i) Für $q \in(p, \infty]$ gilt $\ell^{p} \subset \ell^{q}, \ell^{p} \neq \ell^{q}$, und $\|x\|_{q} \leq\|x\|_{p}$ für alle $x \in \ell^{p}$, aber $\|\cdot\|_{q}$ und $\|\cdot\|_{p}$ sind nicht äquivalent auf $\ell^{p}$.\\
ii) $\lim _{\substack{q \rightarrow \infty \\ q \geq p}}\|x\|_{q}=\|x\|_{\infty}$ für alle $x \in \ell^{p}$.

Beweis. Übungsaufgabe, Blatt 1.

\subsection*{Definition}
i) Eine Abbildung $h: E \times E \rightarrow \mathbb{K}$ heißt hermitesche Form auf $E$, falls gilt:\\
(H1) $h(\cdot, y): E \rightarrow \mathbb{K}$ ist $\mathbb{K}$-linear für alle $y \in E$\\
(H2) $h(x, y)=\overline{h(y, x)}$ für alle $x, y \in E$\\
Insbesondere folgt: $h(x, x) \in \mathbb{R}$ für alle $x \in E$.\\
ii) Eine hermitesche Form $h$ heißt positiv, falls $h(x, x) \geq 0$ für alle $x \in E$.\\
iii) Eine positive hermitesche Form heißt Skalarprodukt, wenn für alle $x \in E$ die Implikation gilt: $\quad h(x, x)=0 \Longrightarrow x=0$.

In diesem Fall schreiben wir oft $\langle x, y\rangle$ statt $h(x, y)$. Das Paar $(E,\langle\cdot, \cdot\rangle)$ heißt dann Prähilbertraum.

\subsection*{Bemerkung}
Ist $h$ eine hermitesche Form auf $E$, so gilt

$$
h\left(x, \alpha y_{1}+\beta y_{2}\right)=\bar{\alpha} h\left(x, y_{1}\right)+\bar{\beta}\left(x, y_{2}\right) \quad \text { für alle } x, y_{1}, y_{2} \in E, \alpha, \beta \in \mathbb{K} .
$$

Ist $\mathbb{K}=\mathbb{R}$, so gilt:

$$
h \text { hermitesche Form } \Longleftrightarrow h \text { symmetrische Bilinearform }
$$

\subsection*{Satz}
Sei $h$ eine positive hermitesche Form auf $E$. Dann gilt:\\
i) $|h(x, y)|^{2} \leq h(x, x) h(y, y)$ für alle $x, y \in E \quad$ (Cauchy-Schwarz-Ungleichung)\\
ii) Durch $\|x\|:=\sqrt{h(x, x)}$ ist eine Halbnorm auf $E$ definiert. Dabei gilt:

$$
\|\cdot\| \text { Norm } \Longleftrightarrow h \text { Skalarprodukt. }
$$

Beweis. Übungsaufgabe auf Blatt 1.

\subsection*{Beispiel}
Sei $I$ eine höchstens abzählbare Indexmenge. Dann ist $\ell^{2}(I)$ ein Prähilbertraum mit Skalarprodukt $\langle\cdot, \cdot\rangle$ definiert durch

$$
\langle x, y\rangle=\sum_{k \in I} x_{k} \bar{y}_{k} \text { für } x=\left(x_{k}\right)_{k \in I}, y=\left(y_{k}\right)_{k \in I} \in \ell^{2}(I) \text {. }
$$

Die Konvergenz der obigen Reihe folgt aus dem Majorantenkriterium, denn es gilt

$$
\left|x_{k} \bar{y}_{k}\right|=\left|x_{k}\right|\left|y_{k}\right| \leq \frac{1}{2}\left(\left|x_{k}\right|^{2}+\left|y_{k}\right|^{2}\right) \quad \text { für alle } k \in \mathbb{N},
$$

und die Reihen $\sum_{k}\left|x_{k}\right|^{2}$ und $\sum_{k}\left|y_{k}\right|^{2}$ konvergieren.

\section*{§2 Topologie in metrischen und normierten Räumen}
\subsection*{Grundwissen über metrische Räume}
Sei $X$ eine Menge. Eine Metrik auf $X$ ist eine Abbildung $d: X \times X \rightarrow[0, \infty)$, so dass für $x, y, z \in X$ gilt:\\
(M1) $d(x, y)=0 \Longleftrightarrow x=y$\\
(M2) $d(x, y)=d(y, x)$\\
(M3) $d(x, z) \leq d(x, y)+d(y, z)$.\\
Das Paar $(X, d)$ heißt metrischer Raum.\\
(i) Für $A \subset X, a \in X$ und $\varepsilon>0$ setzen wir:

$$
\begin{aligned}
& \operatorname{diam} A:=\sup \{d(x, y): x, y \in A\} \quad(\text { Durchmesser von } A) \\
& \operatorname{dist}(a, A):=\inf \{d(a, y): y \in A\} \quad(\text { Abstand von } a \text { zu } A) \\
& \begin{array}{lll}
U_{\varepsilon}(a)=\{y \in X: d(y, a)<\varepsilon\} & \\
B_{\varepsilon}(a)=\{y \in X: d(y, a) \leq \varepsilon\} & \\
U_{\varepsilon}(A)=\{y \in X: \operatorname{dist}(y, A)<\varepsilon\} & \\
B_{\varepsilon}(A)=\{y \in X: \operatorname{dist}(y, A) \leq \varepsilon\} & \\
& \\
\bar{A}:=\{x \in X: \operatorname{dist}(x, A)=0\} & (\text { Abschluss von } A)
\end{array} \\
& \begin{array}{lll}
\mathring{\mathrm{A}}:=\{x \in A & : & \text { es existiert } \varepsilon>0 \\
\partial A:=\bar{A} \backslash \mathring{\mathrm{A}} & \text { mit } \left.U_{\varepsilon}(x) \subset A\right\} & \quad(\text { offener Kern von } A)
\end{array}
\end{aligned}
$$

A heißt

\begin{itemize}
  \item abgeschlossen, falls $A=\bar{A}$;
  \item offen, falls $A=\mathring{\mathrm{A}}$;
  \item Umgebung von $a$, falls $a \in \mathring{\mathrm{A}}$;
  \item dicht in $X$, falls $\bar{A}=X$;
  \item beschränkt, falls diam $A<\infty$.
\end{itemize}

Es gilt:

\begin{itemize}
  \item $U_{\varepsilon}(a), U_{\varepsilon}(A)$ offen, $B_{\varepsilon}(a), B_{\varepsilon}(A)$ abgeschlossen
  \item $A$ offen $\Longleftrightarrow X \backslash A$ abgeschlossen
  \item $X, \varnothing$ sind offen und abgeschlossen
  \item $\mathring{\mathrm{A}}$ offen, $\bar{A}$ abgeschlossen
  \item $\partial A$ abgeschlossen\\
(iii) Seien $A_{i} \subset X, i \in I$ ( $I$ beliebige Indexmenge). Dann gilt:
  \item Sind alle $A_{i}, i \in I$ offen, so auch $\bigcup_{i \in I} A_{i}$ und, falls $I$ endlich, auch $\bigcap_{i \in I} A_{i}$.
  \item Sind alle $A_{i}, i \in I$ abgeschlossen, so auch $\bigcap_{i \in I}$ und, falls $I$ endlich, auch $\bigcup_{i \in I} A_{i}$.\\
(iv) Sei $a \in X$ und $\left(x_{k}\right)_{k} \subset X$ eine Folge. Falls $d\left(x_{k}, a\right) \rightarrow 0$ für $k \rightarrow \infty$, so nennen wir a Grenzwert der Folge $\left(x_{k}\right)_{k}$ und schreiben $a=\lim _{k \rightarrow \infty} x_{k}$ bzw. $x_{k} \rightarrow a$ für $k \rightarrow \infty$.
\end{itemize}

Leicht zu sehen: Durch diese Eigenschaft ist $a$ eindeutig bestimmt.\\
(v) Ist $E$ ein $\mathbb{K}$-Vektorraum, so induziert jede Norm $\|\cdot\|$ auf $E$ genau eine Metrik $d$ gegeben durch $d(x, y)=\|x-y\|$ für alle $x, y \in E(*)$.\\
Ist ( $X, d$ ) ein metrischer Raum und $A \subset X$, so ist auch ( $A, d$ ) ein metrischer Raum. Insbesondere ist jede Teilmenge eines normierten Raumes ein metrischer Raum mit der durch (*) gegebenen Metrik.\\
(vi) Die Eigenschaften „offen" und „abgeschlossen" hängen vom umgebenden metrischen Raum ab. Ist ( $X, d$ ) ein metrischer Raum und $M \subset X$, so gilt für $A \subset M$ :

\begin{itemize}
  \item $A$ offen in $(M, d) \Longleftrightarrow$ es existiert $A^{\prime} \subset X$ offen mit $A=A^{\prime} \cap M$
  \item $A$ abgeschlossen in $(M, d) \Longleftrightarrow$ es existiert $A^{\prime} \subset X$ abgeschlossen mit $A=A^{\prime} \cap M$\\
(vii) Sei $E$ ein $\mathbb{K}$-Vektorraum. Zwei äquivalente Normen $\|\cdot\|_{1}$ und $\|\cdot\|_{2}$ auf $E$ erzeugen die gleiche Topologie, d. h. das gleiche System offener Teilmengen.
\end{itemize}

Für $A \subset E$ gilt also:\\
$A$ offen bzgl. $\|\cdot\|_{1} \Longleftrightarrow A$ offen bzgl. $\|\cdot\|_{2}$. Gleichermaßen invariant bleiben die Eigenschaften „abgeschlossen, beschränkt" und „dicht", die Definitionen $\bar{A}$ und $\stackrel{\circ}{A}$ und die Konvergenz von Folgen.

Insbesondere sind diese Begriffe im Vektorraum $\mathbb{K}^{N}$ unabhängig von der Wahl einer Norm wohldefiniert.

\subsection*{Definition}
Der metrische Raum ( $X, d$ ) heißt separabel, falls $X$ eine abzählbar dichte Teilmenge besitzt.

Hier und im Folgenden steht "abzählbar" für "endlich oder abzählbar unendlich".

\subsection*{Beispiel}
i). $\mathbb{K}^{N}$ ist separabel, denn

\begin{itemize}
  \item $\mathbb{Q}^{N}$ ist abzählbar und dicht in $\mathbb{R}^{N}$
  \item $(\mathbb{Q}+i \mathbb{Q})^{N}$ ist abzählbar und dicht in $\mathbb{C}^{N}$.\\
ii). $\mathbb{R} \backslash \mathbb{Q}$ ist separabel, denn $\sqrt{2}+\mathbb{Q}$ ist eine abzählbare und dichte Teilmenge von $\mathbb{R} \backslash \mathbb{Q}$.
\end{itemize}

\subsection*{Satz}
( $X, d$ ) separabel, $A \subset X \Longrightarrow(A, d)$ separabel.\\
Beweis. Nach Voraussetzung existiert eine abzählbare Menge $Y=\left\{y_{n}: n \in \mathbb{N}\right\} \subset X$, welche in $X$ dicht liegt. Wir setzen

$$
H:=\left\{(m, n) \in \mathbb{N} \times \mathbb{N}: B_{\frac{1}{m}}\left(y_{n}\right) \cap A \neq \varnothing\right\}
$$

und wählen $b_{m, n} \in B_{\frac{1}{m}}\left(y_{n}\right) \cap A$ für jedes $(m, n) \in H$. Da $H$ abzählbar ist, ist auch die Menge

$$
B:=\left\{b_{m, n}:(m, n) \in H\right\} \subset A
$$

abzählbar. Ferner ist $B$ dicht in $A$ (Übung!). Es folgt, dass $A$ separabel ist.

\subsection*{Definition}
Sei $(E,\|\cdot\|)$ ein normierter Raum. Eine Teilmenge $M \subset E$ heißt total, falls $\overline{\text { span } \mathrm{M}}= E$ gilt.

Hier und im Folgenden bezeichne span M das Vektorraumerzeugnis von M in E, welches als der Unterraum aller Linearkombinationen von Elementen aus $M$ bzw. äquivalent als der Schnitt aller Unterräume von $E$, welche $M$ enthalten, gegeben ist.

\subsection*{Satz}
Sei $(E,\|\cdot\|)$ ein normierter $\mathbb{K}$-Vektorraum. Dann: $E$ separabel $\Longleftrightarrow E$ besitzt eine abzählbare totale Teilmenge.

Beweis: „ $\Longrightarrow$ ": $E$ separabel $\Longrightarrow$ Es existiert $M \subset E$ abzählbar mit

$$
\bar{M}=E \Longrightarrow \overline{\operatorname{span} \mathrm{M}}=E, \quad \text { da } \bar{M} \subset \overline{\operatorname{span} \mathrm{M}}
$$

also ist $M$ total.\\
„ $\Longleftarrow$ ": Sei $M \subset E$ abzählbar und total. Wir setzen

$$
D:=\left\{\sum_{k=1}^{n} a_{k} v_{k}: n \in \mathbb{N}, a_{1}, \ldots, a_{n} \in R, v_{1}, \ldots, v_{n} \in M\right\}
$$

wobei

$$
R=\left\{\begin{array}{lll}
\mathbb{Q}, & \text { falls } & \mathbb{K}=\mathbb{R} \\
\mathbb{Q}+i \mathbb{Q}, & \text { falls } & \mathbb{K}=\mathbb{C}
\end{array}\right.
$$

sei. $D$ ist abzählbar, da $M$ und $R$ abzählbar sind.\\
Sei $x \in E, \varepsilon>0$. Da $M$ total ist, existiert $v_{1}, \ldots, v_{n} \in M, \lambda_{1}, \ldots, \lambda_{n} \in \mathbb{K}$ mit $\left\|x-\sum_{k=1}^{n} \lambda_{k} v_{k}\right\|<\frac{\varepsilon}{2} \quad\left(^{*}\right)$\\
Sei $C:=\sum_{k=1}^{n}\left\|v_{k}\right\|$. Da $R$ dicht in $\mathbb{K}$ ist, existieren $a_{1}, \ldots, a_{n} \in R$ mit

$$
\left|\lambda_{k}-a_{k}\right|<\frac{\varepsilon}{2 C} \quad \text { für } k=1, \ldots, n .
$$

Es folgt

$$
\left\|x-\sum_{k=1}^{n} a_{k} v_{k}\right\| \stackrel{(*)}{\leq} \frac{\varepsilon}{2}+\left\|\sum_{k=1}^{n}\left(\lambda_{k}-a_{k}\right) v_{k}\right\| \leq \frac{\varepsilon}{2}+\sum_{k=1}^{n}\left|\lambda_{k}-a_{k}\right|\left\|v_{k}\right\|<\frac{\varepsilon}{2}+\frac{\varepsilon}{2}=\varepsilon .
$$

Da $\varepsilon>0$ beliebig gewählt war, folgt $x \in \bar{D}$. Somit ist $\bar{D}=E$, und $E$ ist separabel.

\subsection*{Definition und Bemerkung}
i). Eine Folge $\left(x_{k}\right)_{k} \subset \mathbb{K}$ heißt finit, falls $x_{k} \neq 0$ für höchstens endlich viele $k \in \mathbb{N}$.

Sei $\mathcal{F}$ der $\mathbb{K}$-Vektorraum der finiten Folgen\\
Klar: $\mathcal{F} \subset \ell^{p}$ für alle $p \in[1, \infty]$.\\
Eine Basis von $\mathcal{F}$ ist gegeben durch die Einheitsvektoren $e_{k}:=\left(\delta_{k j}\right)_{j}= (0, \ldots, 0,1,0, \ldots), k \in \mathbb{N}$.

Sei $B:=\left\{e_{k} \mid k \in \mathbb{N}\right\}$.\\
Übungsaufgabe Blatt 2: Für $p \in[1, \infty)$ ist $\mathcal{F}$ dicht in $\ell^{p}$, also $B$ total in $\ell^{p}$. Es folgt mit 2.5. $\ell^{p}$ ist separabel für $p \in[1, \infty)$\\
ii). $\ell^{\infty}$ ist nicht separabel (Übungsaufgabe Blatt 2).

\section*{Stetige Abbildungen}
Seien im Folgenden ( $X, d_{X}$ ) und ( $Y, d_{Y}$ ) metrische Räume.

\subsection*{Definition und Bemerkung}
Sei $f: X \rightarrow Y$ eine Abbildung und $a \in X$.\\
i). $f$ heißt stetig in $a$, wenn $\lim _{n \rightarrow \infty} f\left(x_{n}\right)=f(a)$ für jede Folge $\left(x_{n}\right)_{n}$ in $X$ mit $\lim _{n \rightarrow \infty} x_{n}=a$ gilt. Dies ist bekanntermaßen äquivalent zu einer der folgenden Eigenschaften:

\begin{itemize}
  \item Für alle $\varepsilon>0$ existiert $\delta>0$ mit $f\left(U_{\delta}(a)\right) \subset U_{\varepsilon}(f(a))$.
  \item Für jede Umgebung $V \subset Y$ von $f(a)$ ist $f^{-1}(V)$ eine Umgebung von $a$.\\
ii). $f$ heißt stetig, wenn $f$ in jedem Punkt aus $X$ stetig ist. Dies ist äquivalent zu folgenden Eigenschaften:
  \item Für alle offenen Teilmengen $V \subset Y$ ist $f^{-1}(V) \subset X$ offen in $X$.
  \item Für alle abgeschlossenen Teilmengen $V \subset Y$ ist $f^{-1}(V) \subset X$ abgeschlossen in $X$.
  \item Für alle $A \subset X$ ist $f(\bar{A}) \subset \overline{f(A)}$.\\
iii). Wir setzen $\mathcal{C}(X, Y):=\{f: X \rightarrow Y: f$ stetig $\}$ und $\mathcal{C}_{b}(X, Y):=\{f \in \mathcal{C}(X, Y): f(X)$ beschränkt $\}.$\\
Ist speziell $Y=\mathbb{K}$ (mit $\mathbb{K}=\mathbb{R}$ oder $\mathbb{K}=\mathbb{C}$ ), so schreiben wir kurz $\mathcal{C}(X)$ anstelle von $\mathcal{C}(X, \mathbb{K})$ und $\mathcal{C}_{b}(X)$ anstelle von $\mathcal{C}_{b}(X, \mathbb{K})$.
\end{itemize}

\subsection*{Bemerkung}
i). Die Komposition stetiger Abbildungen ist stetig\\
ii). Sind $f_{n}: X \rightarrow Y$ stetig für $n \in \mathbb{N}$ und konvergiert die Funktionenfolge $\left(f_{n}\right)_{n}$ gleichmäßig gegen $f: X \rightarrow Y$ (d.h. $\sup _{x \in X} d_{Y}\left(f_{n}(x), f(x)\right) \rightarrow 0$ für $n \rightarrow \infty$ ), so ist auch $f$ stetig.\\
iii). Sind $f, g: X \rightarrow Y$ stetig und ist $\{x \in X: f(x)=g(x)\}$ dicht in $X$, so ist $f \equiv g$.\\
iv). Die Mengen $\mathcal{C}(X)$ und $\mathcal{C}_{b}(X)$ sind $\mathbb{K}$-Vektorräume. Genauer ist $\mathcal{C}_{b}(X)$ ein abgeschlossener Unterraum von $l^{\infty}(X)$ bzgl. der Norm $\|\cdot\|_{\infty}$. Dies folgt aus ii), da die Konvergenz bzgl. $\|\cdot\|_{\infty}$ genau der gleichmäßigen Konvergenz entspricht. Falls nicht explizit anders bemerkt, betrachten wir $\mathcal{C}_{b}(X)$ im Folgenden stets mit der Norm $\|\cdot\|_{\infty}$

\subsection*{Definition}
Sei $f: X \rightarrow Y$ eine Abbildung\\
i) $f$ heißt gleichmäßig stetig, falls für alle $\varepsilon>0$ ein $\delta>0$ existiert, so dass für alle $x_{1}, x_{2} \in X$ die Implikation gilt: $d_{X}\left(x_{1}, x_{2}\right)<\delta \Longrightarrow d_{Y}\left(f\left(x_{1}\right), f\left(x_{2}\right)\right)<\varepsilon$\\
ii) $f$ heißt Lipschitz-stetig, falls $L>0$ existiert mit

$$
d_{Y}\left(f\left(x_{1}\right), f\left(x_{2}\right)\right) \leq L d_{X}\left(x_{1}, x_{2}\right) \quad \text { für alle } x_{1}, x_{2} \in X
$$

iii) $f$ heißt Isometrie, falls $d_{Y}\left(f\left(x_{1}\right), f\left(x_{2}\right)\right)=d_{X}\left(x_{1}, x_{2}\right)$ für alle $x_{1}, x_{2} \in X$\\
iv) $f$ heißt Homöomorphismus, falls $f$ bijektiv ist und $f$ sowie $f^{-1}$ stetig sind. Existiert solch ein $f$, so heißen $X$ und $Y$ homöomorph. Schreibe: $X \approx Y$.

\subsection*{Bemerkung}
i) Für eine Abbildung $f: X \rightarrow Y$ gelten folgende Implikationen:\\
$f$ Isometrie $\Longrightarrow f$ Lipschitz-stetig $\Longrightarrow f$ gleichmäßig stetig $\Longrightarrow f$ stetig.\\
ii) Sei $\varnothing \neq A \subset X$. Dann ist die Abbildung $X \rightarrow \mathbb{R}, x \mapsto \operatorname{dist}(x, A)$ Lipschitz-stetig (mit $L=1$ ).

Seien im Folgenden $\left(E,\|\cdot\|_{E}\right),\left(F,\|\cdot\|_{F}\right)$ normierte $\mathbb{K}$-Vektorräume.

\subsection*{Satz}
Sei $T: E \rightarrow F \mathbb{K}$-linear.\\
i). Äquivalent sind:

\begin{itemize}
  \item $T$ stetig
  \item $T$ stetig in 0
  \item $T\left(B_{1}(0)\right) \subset F$ beschränkt (wobei hier $B_{1}(0):=\left\{x \in E:\|x\|_{E} \leq 1\right\}$ sei)
  \item $T$ Lipschitz-stetig\\
ii). Ist $\operatorname{dim} E<\infty$, so ist $T$ stetig.
\end{itemize}

Beweis. (i) bekannt und einfach.\\
(ii) Man sieht direkt, dass eine Norm $\|\cdot\|_{T}$ auf $E$ definiert ist durch $\|x\|_{T}:= \|x\|_{E}+\|T x\|_{F}$. Nach 1.3j) ist $\|\cdot\|_{T}$ äquivalent zu $\|\cdot\|_{E}$, da $E$ endlichdimensional ist. Somit existiert $C>0$ mit

$$
\|x\|_{T} \leq C\|x\|_{E} \quad \text { für alle } x \in E,
$$

also insbesondere

$$
\|T x\|_{F} \leq\|x\|_{T} \leq C\|x\|_{E} \leq C \quad \text { für alle } x \in B_{1}(0) .
$$

\subsection*{Definition}
i). Wir setzen $\mathcal{L}(E, F):=\{T: E \rightarrow F: T$ linear und stetig $\}$. Man rechnet leicht nach, dass $\mathcal{L}(E, F)$ ein normierter Raum ist mit

$$
\|T\|=\|T\|_{\mathcal{L}(E, F)}:=\sup _{x \in B_{1}(0)}\|T x\|_{F}=\sup _{x \in E \backslash\{0\}} \frac{\|T x\|_{F}}{\|x\|_{E}} \quad \text { für } T \in \mathcal{L}(E, F) .
$$

ii). $E^{\prime}:=\mathcal{L}(E, \mathbb{K})$ heißt topologischer Dualraum von $E$.\\
iii). Wir setzen $\mathcal{L}(E):=\mathcal{L}(E, E)$ (Raum der stetigen Endomorphismen von $E$ )\\
iv). Eine bijektive $\mathbb{K}$-lineare Abbildung $T: E \rightarrow F$ heißt topologischer Isomorphismus, falls $T$ und $T^{-1}$ stetig sind. Existiert eine solche Abbildung $T$, so heißen ( $E,\|\cdot\|_{E}$ ) und ( $F,\|\cdot\|_{F}$ ) topologisch isomorph. Wir setzen

$$
\text { Iso }(E, F):=\{T: E \rightarrow F: T \text { topologischer Isomorphismus }\}
$$

und schreiben kurz $\operatorname{Iso}(E)$ anstelle von $\operatorname{Iso}(E, E)$.

\subsection*{Bemerkung}
i). Die Elemente von $\mathcal{L}(E, F)$ nennt man auch beschränkte lineare Operatoren, und $\|\cdot\|=\|\cdot\|_{\mathcal{L}(E, F)}$ nennt man die Operatornorm.\\
ii). Ist $T \in \mathcal{L}(E, F)$, so ist $c=\|T\|$ die kleinste nichtnegative Zahl mit der Eigenschaft

$$
\|T x\|_{F} \leq c\|x\|_{E} \quad \text { für alle } x \in E
$$

iii). Ist ( $G,\|\cdot\|_{G}$ ) ein weiterer normierter Raum und sind $T \in \mathcal{L}(E, F), S \in \mathcal{L}(F, G)$, so gilt

$$
\|S T x\|_{G} \leq\|S\|\|T x\|_{F} \leq\|S\|\|T\|\|x\|_{E} \quad \text { für alle } x \in E,
$$

also\\
$S T \in \mathcal{L}(E, G) \quad$ und $\quad\|S T\| \leq\|S\|\|T\| \quad$ (Submultiplikativität der Norm) Insbesondere gilt $\left\|T^{n}\right\| \leq\|T\|^{n}$ für $T \in \mathcal{L}(E)$ und $n \in \mathbb{N}$.

\subsection*{Beispiel}
i). Sei der $\mathbb{K}$-Vektorraum $\mathcal{F}$ der finiten Folgen (siehe 2.6) versehen mit $\|\cdot\|_{\infty}$, d. h. $\|x\|_{\infty}=\sup _{k \in \mathbb{N}}\left|x_{k}\right|=\max _{k \in \mathbb{N}}\left|x_{k}\right|$ für $x \in \mathcal{F}$. Sei ferner $T: \mathcal{F} \rightarrow \mathcal{F}$ definiert durch $T\left(x_{k}\right)_{k}=\left(\frac{1}{k} x_{k}\right)_{k}$. Man sieht leicht, dass $T$ bijektiv ist. Da ferner $\|T x\|_{\infty} \leq\|x\|_{\infty}$ für alle $x \in \mathcal{F}$ gilt, ist $T$ auch stetig. Allerdings ist $T^{-1}$ nicht stetig, da

$$
\frac{1}{k} e_{k} \rightarrow 0 \text { in }\left(\mathcal{F},\|\cdot\|_{\infty}\right) \text { und }\left\|T^{-1} \frac{1}{k} e_{k}\right\|_{\infty}=\left\|e_{k}\right\|_{\infty}=1 \text { für alle } k .
$$

Hier sei $e_{k}:=\left(\delta_{n k}\right)_{n} \in \mathcal{F}$ wie in 2.6 der $k$-te Einheitsvektor in $\mathcal{F}$.\\
ii). Sei $1 \leq p<q \leq \infty$. Dann ist die identische Abbildung $T$ : $\left(\ell^{p},\|\cdot\|_{p}\right) \rightarrow\left(\ell^{p},\|\cdot\|_{q}\right)$ stetig (da $\|x\|_{q} \leq\|x\|_{p}$ für alle $x \in \ell^{p}$ ), aber $T^{-1}$ ist nicht stetig, da $\|\cdot\|_{q}$ und $\|\cdot\|_{p}$ auf $\ell^{p}$ nicht äquivalent sind.\\
iii). (Integraloperatoren) Seien $[a, b],[c, d] \subset \mathbb{R}$ kompakte Intervalle, und sei $k$ : $[c, d] \times[a, b] \rightarrow \mathbb{C}$ stetig. Wir definieren die lineare Abbildung $K: \mathcal{C}[a, b] \rightarrow \mathcal{C}[c, d]$ durch

$$
[K u](t):=\int_{a}^{b} k(t, s) u(s) d s \quad \text { für } t \in[c, d]
$$

Offensichtlich ist $K$ wohldefiniert und linear. $K$ ist stetig, denn für alle $u \in \mathcal{C}[a, b]$ und $t \in[c, d]$ gilt

$$
|[K u](t)| \leq \int_{a}^{b}|k(t, s)||u(s)| d s \leq M\|u\|_{\infty}
$$

mit

$$
M:=\max _{c \leq t \leq d} \int_{a}^{b}|k(t, s)| d s<\infty
$$

also $\|K u\|_{\infty} \leq M\|u\|_{\infty}$ für $u \in \mathcal{C}[a, b]$. Es folgt somit $K \in \mathcal{L}(\mathcal{C}[a, b], \mathcal{C}[c, d])$ mit $\|K\| \leq M$.\\
Wir werden nun sehen, dass sogar $\|K\|=M$ gilt. Sei dazu $t_{0} \in[c, d]$ ein Punkt, wo das Maximum in (*) angenommen wird. Für $\varepsilon>0$ sei ferner

$$
u_{\varepsilon} \in \mathcal{C}[a, b] \quad \text { definiert durch } \quad u_{\varepsilon}(s)=\frac{\overline{k\left(t_{0}, s\right)}}{\left|k\left(t_{0}, s\right)\right|+\varepsilon}
$$

Da $\left\|u_{\varepsilon}\right\|_{\infty} \leq 1$ unabhängig von $\varepsilon>0$, ist

$$
\|K\| \geq\left\|K u_{\varepsilon}\right\|_{\infty} \geq\left|\left[K u_{\varepsilon}\right]\left(t_{0}\right)\right|=\int_{a}^{b} \frac{\left|k\left(t_{0}, s\right)\right|^{2}}{\left|k\left(t_{0}, s\right)\right|+\varepsilon} d s
$$

wobei das letzte Integral für $\varepsilon \rightarrow 0$ nach dem Satz von der majorisierten Konvergenz gegen $\int_{a}^{b}\left|k\left(t_{0}, s\right)\right| d s=M$ konvergiert. Also ist auch $\|K\| \geq M$, und insgesamt folgt Gleichheit.\\
Spezielles Beispiel:\\
Sei $K \in \mathcal{L}(\mathcal{C}[0,1], \mathcal{C}[0,1])$ der Lösungsoperator zum Problem (2) aus Kapitel 0 , d.h. es gelte

$$
K f(t)=\int_{0}^{1} G(t, s) f(s) d s \quad \text { für } f \in \mathcal{C}[0,1], t \in[0,1]
$$

mit der Greenfunktion

$$
G(t, s)= \begin{cases}(1-t) s, & s \leq t \\ (1-s) t, & s \geq t\end{cases}
$$

Dann ist

$$
\|K\|=\max _{t \in[0,1]} h(t)
$$

mit

$$
\begin{aligned}
h(t) & =\int_{0}^{1}|G(t, s)| d s=(1-t) \int_{0}^{t} s d s+t \int_{t}^{1}(1-s) d s \\
& =(1-t) \int_{0}^{t} s d s+t \int_{0}^{1-t} s d s=\frac{1}{2}\left[(1-t) t^{2}+t(1-t)^{2}\right]=\frac{t(1-t)}{2}
\end{aligned}
$$

Es folgt $\|K\|=\frac{1}{8}$.

\section*{§ 3 Vollständigkeit und der Satz von Baire}
Sei stets $(X, d)$ ein metrischer Raum.

\subsection*{Definition}
i). Eine Folge $\left(x_{n}\right)_{n}$ in $X$ heißt Cauchyfolge, wenn $\lim _{n \rightarrow \infty} \sup _{m \geq n} d\left(x_{m}, x_{n}\right)=0$ gilt.\\
ii). ( $X, d$ ) heißt vollständig, wenn jede Cauchyfolge in $X$ konvergiert.\\
iii). Ein vollständiger normierter Raum ( $E,\|\cdot\|$ ) heißt Banachraum.\\
iv). Ein vollständiger Prähilbertraum ( $E,\langle\cdot, \cdot\rangle$ ) heißt Hilbertraum.

In (iii) und (iv) bezieht sich die Vollständigkeit dabei auf die induzierte Metrik.

\subsection*{Bemerkung}
Sei $\left(x_{n}\right)_{n} \subset X$ eine Folge.\\
i). Ist $\left(x_{n}\right)_{n}$ konvergent, so ist $\left(x_{n}\right)_{n}$ eine Cauchyfolge.\\
ii). Ist $\left(x_{n}\right)_{n}$ eine Cauchyfolge, so ist $\left(x_{n}\right)_{n}$ in $X$ beschränkt, d.h. die Menge $\left\{x_{n}: n \in \mathbb{N}\right\} \subset X$ ist beschränkt.\\
iii). Ist $\left(x_{n}\right)_{n}$ eine Cauchyfolge und besitzt $\left(x_{n}\right)_{n}$ eine gegen $a \in X$ konvergente Teilfolge, so konvergiert bereits die Folge $\left(x_{n}\right)_{n}$ selbst gegen $a$.\\
iv). Ist $\left(x_{n}\right)_{n}$ eine Cauchyfolge, $(Y, \tilde{d})$ ein weiterer metrischer Raum und $f: X \rightarrow Y$ gleichmäßig stetig, so ist auch $\left(f\left(x_{n}\right)\right)_{n}$ eine Cauchyfolge (in $Y$ ).

\subsection*{Satz}
Sei $A \subset X$.\\
i). Ist ( $X, d$ ) vollständig und $A \subset X$ abgeschlossen, so ist auch ( $A, d$ ) vollständig.\\
ii). Ist ( $A, d$ ) vollständig, so ist $A$ abgeschlossen in $X$.

Beweis. i) Sei $\left(x_{n}\right)_{n}$ eine Cauchyfolge in $A$. Da $X$ vollständig ist, existiert $x= \lim _{n \rightarrow \infty} x_{n} \in X$, wobei $x \in \bar{A}=A$ nach Voraussetzung. Also konvergiert $\left(x_{n}\right)_{n}$ in $A$. ii) Sei $x \in \bar{A}$. Dann ist $x=\lim _{n \rightarrow \infty} x_{n}$ für eine Folge $\left(x_{n}\right)_{n}$ in $A$. Nach 3.2 ) ist $\left(x_{n}\right)_{n}$ eine Cauchyfolge in $A$, und damit konvergiert $\left(x_{n}\right)_{n}$ nach Voraussetzung in $A$. Es folgt $x \in A$. Insgesamt folgt $A=\bar{A}$.

\subsection*{Bemerkung}
Sei $(Y, \tilde{d})$ ein weiterer metrischer Raum und $f: X \rightarrow Y$ eine Abbildung.\\
i). Ist $f$ eine Isometrie, so ist neben $f$ auch $f^{-1}: f(X) \rightarrow X$ gleichmäßig stetig, und aus 3.2 v ) folgt:

$$
(X, d) \text { vollständig } \Longleftrightarrow(f(X), \tilde{d}) \text { vollständig }
$$

Gilt dies, so ist $f(X)$ abgeschlossen in $Y$ nach 3.3,\\
ii). Ist $f$ ein Homöomorphismus, so muss $f$ nicht die Vollständigkeit erhalten; sei dazu speziell $X=\mathbb{R}$ und $Y=\left(-\frac{\pi}{2}, \frac{\pi}{2}\right)$, jeweils versehen mit der Betragsmetrik $d(x, y)=|x-y|$. Sei ferner der Homöomorphismus $f: X \rightarrow Y$ gegeben durch $f(x)=\arctan x$. Dann ist $X$ vollständig und $Y$ nicht, da $Y \subset X$ nicht abgeschlossen ist.

\subsection*{Bemerkung und Beispiele}
i). Sind $\|\cdot\|_{1},\|\cdot\|_{2}$ äquivalente Normen auf einem $\mathbb{K}$-Vektorraum $E$, so gilt:

$$
\left(E,\|\cdot\|_{1}\right) \text { Banachraum } \quad \Longleftrightarrow \quad\left(E,\|\cdot\|_{2}\right) \text { Banachraum }
$$

Allgemeiner gilt für zwei topologisch isomorphe $\mathbb{K}$-Vektorräume ( $E,\|\cdot\|_{E}$ ) und $\left(F,\|\cdot\|_{F}\right)$ :

$$
E \text { Banachraum } \Longleftrightarrow F \text { Banachraum }
$$

ii). $\mathbb{K}^{n}$ ist ein Banachraum (bzgl. jeder Norm). Ist ferner $E$ irgendein normierter Raum mit $n:=\operatorname{dim} E<\infty$, so ist $E$ topologisch isomorph zu $\mathbb{K}^{n}$ und damit ein Banachraum nach i).\\
iii). Jeder endlich-dimensionale Unterraum eines normierten Raums ist abgeschlossen nach ii) und 3.3 .\\
iv). $l^{p}$ ist ein Banachraum für $p \in[1, \infty]$. Wir beweisen dies zunächst im Fall $p<\infty$ : Sei $\left(x_{n}\right)_{n}$ eine Cauchyfolge in $l^{p}$, wobei wir $x_{n}=\left(\xi_{n k}\right)_{k}$ notieren. Sei $k \in \mathbb{N}$ fest; für $m, n \in \mathbb{N}$ ist dann

$$
\left|\xi_{m k}-\xi_{n k}\right|^{p} \leq \sum_{j \in \mathbb{N}}\left|\xi_{m j}-\xi_{n j}\right|^{p}=\left\|x_{m}-x_{n}\right\|_{p}^{p}
$$

und somit ist die Folge $\left(\xi_{n k}\right)_{n} \subset \mathbb{K}$ eine Cauchyfolge. Da $\mathbb{K}$ vollständig ist, existiert $\xi_{k}:=\lim _{n \rightarrow \infty} \xi_{n k} \in \mathbb{K}$. Setze $x:=\left(\xi_{k}\right)_{k}$. Für alle $l \in \mathbb{N}$ gilt dann

$$
\sum_{k=1}^{l}\left|\xi_{k}\right|=\lim _{n \rightarrow \infty} \sum_{k=1}^{l}\left|\xi_{n k}\right|^{p} \leq \limsup _{n \rightarrow \infty}\left\|x_{n}\right\|_{p}^{p} \stackrel{\text { (3.2) }}{\gtrless}(i)
$$

Also ist $\sum_{k \in \mathbb{N}}\left|\xi_{k}\right|^{p}<\infty$ und somit $x \in l^{p}$. Für alle $n, l \in \mathbb{N}$ gilt zudem

$$
\sum_{k=1}^{l}\left|\xi_{k}-\xi_{n k}\right|^{p}=\lim _{m \rightarrow \infty} \sum_{k=1}^{l}\left|\xi_{m k}-\xi_{n k}\right|^{p} \leq \sup _{m \geq n}\left\|x_{m}-x_{n}\right\|_{p}^{p}=: c_{n}
$$

und somit $\left\|x-x_{n}\right\|_{p}^{p} \leq c_{n}$. Nach Voraussetzung gilt zudem $\lim _{n \rightarrow \infty} c_{n}=0$ und somit $\lim _{n \rightarrow \infty} x_{n}=x$ in $l^{p}$; dies war zu zeigen.\\
Der Beweis für $p=\infty$ ist ähnlich, nur einfacher. Es gilt sogar:\\
v). Ist $M$ eine beliebige Menge, so ist $l^{\infty}(M)$ ein Banachraum (Übung, Blatt 2).\\
vi). Nach 2.8 v) ist der Unterraum $\mathcal{C}_{b}(X)$ abgeschlossen in $l^{\infty}(X)$ und damit ein Banachraum nach v) und 3.3.\\
vii). $\ell^{2}$ ist ein Hilbertraum mit dem Skalarprodukt $\langle\cdot, \cdot\rangle$ aus 1.13 .

\subsection*{Satz}
Seien $E, F$ normierte Räume. Ist $F$ ein Banachraum, so ist auch $\mathcal{L}(E, F)$ ein Banachraum.\\
Insbesondere ist der Dualraum $E^{\prime}$ eines normierten Raumes stets ein Banachraum.

Beweis. Sei $\left(T_{n}\right)_{n}$ eine Cauchyfolge in $\mathcal{L}(E, F)$ und $c_{n}:=\sup _{m \geq n}\left\|T_{n}-T\right\|$ für $n \in \mathbb{N}$, so dass $\lim _{n \rightarrow \infty} c_{n}=0$ gilt. Für alle $x \in E$ gilt dann

$$
\sup _{m \geq n}\left\|T_{m} x-T_{n} x\right\| \leq \sup _{m \geq n}\left\|T_{m}-T_{n}\right\|\|x\|=c_{n}\|x\| \rightarrow 0 \quad \text { für } n \rightarrow \infty,
$$

d.h. $\left(T_{n} x\right)_{n}$ ist eine Cauchyfolge in $F$. Aus der Voraussetzung folgt somit die Existenz von

$$
T x:=\lim _{n \rightarrow \infty} T_{n} x \in F \quad \text { für alle } x \in E .
$$

Aus der Linearität der Abbildungen $T_{n}$ folgt direkt, dass auch die Zuordnung $x \mapsto T x$ linear ist. Ferner ist

$$
\|T x\|=\lim _{n \rightarrow \infty}\left\|T_{n} x\right\| \leq\left(\sup _{n \in \mathbb{N}}\left\|T_{n}\right\|\right)\|x\| \quad \text { für alle } x \in E \text {, }
$$

d.h. $T$ ist stetig. Schließlich ist\\
$\left\|\left(T-T_{n}\right) x\right\|=\lim _{m \rightarrow \infty}\left\|\left(T_{m}-T_{n}\right) x\right\| \leq\left(\sup _{m \geq n}\left\|T_{m}-T_{n}\right\|\right)\|x\|=c_{n}\|x\| \quad$ für alle $x \in E$,\\
d.h. $\left\|T-T_{n}\right\| \leq c_{n} \rightarrow 0$ für $n \rightarrow \infty$. Dies zeigt die Vollständigkeit von $\mathcal{L}(E, F)$.

\subsection*{Satz (Cantor)}
Sei ( $X, d$ ) vollständig, und seien $A_{n} \subset X, n \in \mathbb{N}$ abgeschlossene, nichtleere Mengen mit $A_{n+1} \subset A_{n}$ für alle $n$ und $\lim _{n \rightarrow \infty} \operatorname{diam} A_{n}=0$. Dann ist $A:=\bigcap_{n \in \mathbb{N}} A_{n}$ nichtleer.

Beweis. Wähle $x_{n} \in A_{n}$ für $n \in \mathbb{N}$. Nach Voraussetzung ist dann

$$
x_{m} \in A_{n} \quad \text { für } m \geq n .
$$

Also folgt

$$
\sup _{m \geq n} d\left(x_{m}, x_{n}\right) \leq \operatorname{diam} A_{n} \rightarrow 0 \quad \text { für } n \rightarrow \infty,
$$

d.h. $\left(x_{n}\right)_{n}$ ist eine Cauchyfolge in $X$. Da ( $X, d$ ) vollständig ist, existiert $x:=\lim _{n \rightarrow \infty} x_{n}$. Wegen (*) ist $x \in \bar{A}_{n}=A_{n}$ für alle $n \in \mathbb{N}$; also $x \in A$. \(\square\)

Nachbemerkung: Es gilt sogar $A=\{x\}$, da $\operatorname{diam} A \leq \lim _{n \rightarrow \infty} A_{n}=0$ ist.

\subsection*{Satz von Baire}
Sei ( $X, d$ ) ein vollständiger metrischer Raum, und seien $D_{n} \subset X$ offen und dicht für $n \in \mathbb{N}$. Dann ist auch $D:=\bigcap_{n \in \mathbb{N}} D_{n}$ dicht in $X$.

Beweis. Es reicht, zu zeigen, dass $D \cap U \neq \varnothing$ für jede nichtleere offene Teilmenge $U$ von $X$ gilt. Sei also $U \subset X$ offen und nichtleer. Da $D_{1} \subset X$ offen und dicht ist, ist $D_{1} \cap U$ offen und nichtleer. Wähle $a_{1} \in D_{1} \cap U$ und $r_{1} \in\left(0, \frac{1}{2}\right)$ mit $A_{1}:=B_{r_{1}}\left(a_{1}\right) \subset D_{1} \cap U$. Dann gilt:

$$
\varnothing \neq \stackrel{\circ}{A}_{1} \subset A_{1}=\bar{A}_{1} \subset D_{1} \cap U \quad \text { und } \quad \operatorname{diam} A_{1} \leq 1 .
$$

Konstruiere nun sukzessive Mengen $A_{n} \subset X, n \geq 2$ mit

$$
\varnothing \neq \stackrel{\circ}{A}_{n} \subset A_{n}=\bar{A}_{n} \subset D_{n} \cap A_{n-1} \quad \text { und } \quad \operatorname{diam} A_{n} \leq \frac{1}{n} .
$$

Sei dazu $n \geq 2$, und seien $A_{1}, \ldots, A_{n-1}$ bereits konstruiert. Da $\varnothing \neq \stackrel{\circ}{A}_{n-1}$ offen in $X$ ist und $D_{n}$ offen und dicht ist, muss $D_{n} \cap \stackrel{\circ}{A}_{n-1}$ offen und nichtleer sein. Wähle $a_{n} \in D_{n} \cap \stackrel{\circ}{A}_{n-1}$ und $r_{n} \in\left(0, \frac{1}{2 n}\right)$ mit $A_{n}:=B_{r_{n}}\left(a_{n}\right) \subset D_{n} \cap \stackrel{\circ}{A}_{n-1}$. Mit dieser Wahl von $A_{n}$ gilt ( $* *$ ).\\
Anwendung von 3.7 auf die Mengen $A_{n}, n \in \mathbb{N}$ liefert schließlich $A:=\bigcap_{n \in \mathbb{N}} A_{n} \neq \varnothing$. Aus (*) und (**) folgt aber

$$
A_{n} \subset\left[D_{1} \cap \cdots \cap D_{n}\right] \cap U \quad \text { für alle } n \in \mathbb{N}, \quad \text { also } A \subset D \cap U .
$$

Es folgt $D \cap U \neq \varnothing$.

\subsection*{Korollar}
Sei ( $X, d$ ) vollständig, und sei $A_{n} \subset X$ abgeschlossen für $n \in \mathbb{N}$. Ferner sei $A:= \bigcup_{n \in \mathbb{N}} A_{n}$. Dann gilt:\\
i). Ist $\stackrel{\circ}{A}_{n}=\varnothing$ für alle $n \in \mathbb{N}$, so ist auch $\stackrel{\circ}{A}=\varnothing$.\\
ii). Enthält $A$ eine offene Kugel $U_{r}(a)$ mit $a \in X, r>0$, dann enthält mindestens eine der Mengen $A_{n}$ eine abgeschlossene Kugel der Form $B_{s}(b)$ mit $b \in X$, $s>0$.

Beweis. i) Nach Voraussetzung ist $D_{n}:=X \backslash A_{n}$ offen und dicht für alle $n \in \mathbb{N}$. Nach 3.8 ist somit $X \backslash A=\bigcap_{n \in \mathbb{N}} D_{n}$ dicht in $X$. Dies liefert $\stackrel{\circ}{A}=X \backslash \overline{X \backslash A}=\varnothing$. ii) folgt direkt aus i).

\subsection*{Definition}
Sei $A \subset X$.\\
i). $A$ heißt nirgends dicht (in $X$ ), wenn $\stackrel{\circ}{\bar{A}}=\varnothing$ gilt.\\
ii). $A$ heißt mager (oder von erster Kategorie), falls sich $A$ als abzählbare Vereinigung nirgends dichter Mengen schreiben lässt.

\subsection*{Bemerkung und Beispiele}
i). Die Vereinigung abzählbar vieler magerer Mengen ist mager\\
ii). Für $a \in X$ ist äquivalent:

$$
\{a\} \text { nirgends dicht in } X \Longleftrightarrow \operatorname{dist}(a, X \backslash\{a\})=0
$$

Gilt dies für alle $a \in X$, so ist jede abzählbare Teilmenge von $X$ mager. Dies ist insbesondere dann der Fall, wenn $\{0\} \neq X$ ein normierter Raum (mit induzierter Metrik) ist.\\
iii). Sei $(E,\|\cdot\|)$ ein normierter Raum und $F \subset E$ ein Unterrraum, welcher nicht dicht ist. Dann ist $F$ auch nirgends dicht, denn:\\
$F$ nicht dicht bedeutet, dass $v \in E \backslash \bar{F}$ existiert. Da $\bar{F}$ ein Unterraum von $E$ ist (siehe Übungsblatt 1), können wir dabei $0 \neq\|v\|<1$ annehmen. Für alle $x \in \bar{F}$ und $\delta>0$ ist dann $x+\delta v \in U_{\delta}(x) \backslash \bar{F}$. Es folgt $\bar{F}=\varnothing$.

\subsection*{Satz}
Sei ( $X, d$ ) vollständig und $A \subset X$ mager. Dann ist $X \backslash A$ dicht in $X$. Insbesondere ist $X$ nicht mager in sich, falls $X \neq \varnothing$ ist.

Beweis. Sei $A=\bigcup_{n \in \mathbb{N}} A_{n}$ mit $\stackrel{\circ}{A}_{n}=\varnothing$ für alle $n$. Nach 3.9 gilt dann $\stackrel{\circ}{A} \subset \stackrel{\circ}{B}=\varnothing$ für $B:=\bigcup_{n \in \mathbb{N}} \overline{A_{n}}$. Es folgt $\overline{X \backslash A}=X \backslash \stackrel{\circ}{A}=X$.

\subsection*{Korollar}
Sei $E$ ein unendlichdimensionaler Banachraum. Dann ist jede Basis von $E$ (im Sinne der linearen Algebra) überabzählbar.

Beweis. Angenommen, $E$ besäße eine abzählbare Basis $\left\{v_{1}, v_{2}, \ldots,\right\}$. Dann ist $E= \bigcup_{n \in \mathbb{N}} V_{n}$ mit $V_{n}:=\operatorname{span}\left\{v_{1}, \ldots, v_{n}\right\}$. Nach Voraussetzung und 3.5 iii) ist $V_{n}=\bar{V}_{n} \neq E$ für alle $n$, und somit ist $V_{n}$ nirgends dicht in $E$ nach 3.11.ii). Dann ist $E$ mager in sich, im Widerspruch zu 3.12

\subsection*{Bemerkung}
Wir werden bald sehen, dass der Satz von Baire wichtige Konsequenzen für die Struktur der Menge stetiger linearer Abbildungen hat (Satz von der gleichmäßigen Beschränktheit, Satz von der offenen Abbildung).

\section*{Vollständigkeit und Reihen}
\subsection*{Definition}
Sei $(E,\|\cdot\|)$ ein normierter Raum und $\left(x_{n}\right)_{n} \subset E$ eine Folge.\\
i). Wenn $S:=\lim _{n \rightarrow \infty} \sum_{k=1}^{n} x_{k} \in E$ existiert, so nennen wir die Reihe $\sum_{k} x_{k}$ konvergent und setzen $\sum_{k=1}^{\infty} x_{k}:=S$.\\
ii). Die Reihe $\sum_{k} x_{k}$ heißt absolut konvergent, wenn $\sum_{k=1}^{\infty}\left\|x_{k}\right\|<\infty$ ist.

\subsection*{Satz}
Sei $(E,\|\cdot\|)$ ein normierter Raum. Dann gilt:\\
$E$ ist genau dann ein Banachraum, wenn jede absolut konvergente Reihe in $E$ auch konvergent ist.

Beweis. $\Longrightarrow$ : Genau wie für $E=\mathbb{R}$ bzw. $\mathbb{C}$ ("Cauchykriterium"), man muss nur alle Beträge durch Normen ersetzen.\\
$\Longleftarrow$ : Sei $\left(x_{n}\right)_{n} \subset E$ eine Cauchyfolge. Dann existiert für alle $k \in \mathbb{N}$ ein $\tau_{k} \in \mathbb{N}$ mit

$$
\sup _{m \geq n}\left\|x_{m}-x_{n}\right\|<2^{-k} \quad \text { für } n \geq \tau_{k} \text {. }
$$

Wähle nun induktiv $n_{1} \geq \tau_{1}$ und $n_{k+1} \geq \max \left\{\tau_{k+1}, n_{k}+1\right\}$ für $k \in \mathbb{N}$. Dann gilt

$$
\left\|x_{n_{k+1}}-x_{n_{k}}\right\| \leq 2^{-k} \quad \text { für alle } k,
$$

also insbesondere

$$
\sum_{k=1}^{\infty}\left\|x_{n_{k+1}}-x_{n_{k}}\right\| \leq \sum_{k=1}^{\infty} 2^{-k}=1
$$

dh. die Reihe $\sum_{k}\left(x_{n_{k+1}}-x_{n_{k}}\right)$ konvergiert absolut. Nach Voraussetzung konvergiert sie damit auch, d.h. es existiert

$$
x:=\lim _{k \rightarrow \infty} x_{n_{k}}=x_{n_{1}}+\lim _{k \rightarrow \infty} \sum_{j=1}^{k}\left(x_{n_{j+1}}-x_{n_{j}}\right) \quad \in E
$$

Nach 3.2 v ) konvergiert die Folge damit insgesamt gegen $x$, was zu zeigen war.

\section*{§4 Kompaktheit}
Im Folgenden sei ( $X, d$ ) ein metrischer Raum.

\subsection*{Definition und Bemerkung}
i). ( $X, d$ ) heißt kompakt, falls gilt:

Ist $I$ eine beliebige Indexmenge und sind $U_{i} \subset X, i \in I$ offene Mengen mit $X=\bigcup_{i \in I} U_{i}$, so existiert eine endliche Teilmenge $J \subset I$ mit $X=\bigcup_{i \in J} U_{i}$. Mit anderen Worten: Jede offene Überdeckung von $X$ lässt sich auf eine endliche Teilüberdeckung reduzieren.\\
ii). ( $X, d$ ) heißt präkompakt, wenn für alle $\varepsilon>0$ eine endliche Teilmenge $M \subset X$ existiert mit $X=B_{\varepsilon}(M)$.\\
iii). $A \subset X$ heißt kompakt (bzw. präkompakt), wenn der metrische Teilraum ( $A, d$ ) kompakt (bzw. präkompakt) ist. Somit gilt:

\begin{itemize}
  \item $A$ ist kompakt genau dann, wenn für jedes Familie $U_{i} \subset X, i \in I$ offener Mengen mit $A \subset \bigcup_{i \in I} U_{i}$ eine endiche Teilmenge $J \subset I$ existiert mit $A \subset \bigcup_{i \in J} U_{i}$.
  \item $A$ ist präkompakt genau dann für jedes $\varepsilon>0$ eine endliche Teilmenge $M \subset A$ existiert mit $A \subset B_{\varepsilon}(M)$ (wobei $B_{\varepsilon}(M)$ hier die Umgebung von $M$ in $X$ bezeichne).\\
iv). $A \subset X$ heißt relativ kompakt, wenn $\bar{A}$ kompakt ist.
\end{itemize}

\subsection*{Satz}
Für einen metrischen Raum ( $X, d$ ) sind folgende Eigenschaften äquivalent.\\
i). $X$ ist kompakt\\
ii). Jede Folge $\left(x_{k}\right)_{k}$ in $X$ besitzt eine in $X$ konvergente Teilfolge.\\
iii). $X$ ist präkompakt und vollständig.

Beweis. "i) $\Longrightarrow$ ii) ": Sei $X$ kompakt, und sei $\left(x_{k}\right)_{k}$ eine Folge in $X$. Ohne Einschränkung sei $M:=\left\{x_{k}: k \in \mathbb{N}\right\}$ eine unendliche Menge. Angenommen, kein $a \in X$ ist Grenzwert einer Teilfolge von $\left(x_{k}\right)_{k}$. Dann existiert für jedes $a \in X$ ein $\varepsilon_{a}>0$ so, dass $U_{\varepsilon_{a}}(a) \cap M$ endlich ist. Da $X=\bigcup_{a \in X} U_{\varepsilon_{a}}(a)$ kompakt ist, existieren $a_{1}, \ldots, a_{n} \in X$ mit $X=\bigcup_{i=1}^{n} U_{\varepsilon_{a_{i}}}\left(a_{i}\right)$. Es folgt: $M=M \cap X=\bigcup_{i=1}^{n} M \cap U_{\varepsilon_{a_{i}}}\left(a_{i}\right)$ ist endlich. Widerspruch.\\
"ii) $\Longrightarrow$ iii)": Wir zeigen zunächst die Vollständigkeit von $X$. Sei $\left(x_{k}\right)_{k}$ eine Cauchyfolge in $X$; diese besitzt nach Voraussetzung eine in $X$ konvergente Teilfolge. Nach 3.2 ist dann aber schon $\left(x_{k}\right)_{k}$ selbst konvergent.

Wir zeigen nun die Präkompaktheit von $X$. Sei dazu $\varepsilon>0$. Angenommen, es gäbe keine endliche Menge $M \subset X$ mit $X=B_{\varepsilon}(M)$. Wir wählen dann $a_{1} \in X$ beliebig;

$$
\begin{aligned}
& \text { da } X \not \subset U_{\varepsilon}\left(a_{1}\right) \text {, existiert } a_{2} \in X \backslash U_{\varepsilon}\left(a_{1}\right) ; \\
& \text { da } X \not \subset U_{\varepsilon}\left(a_{1}\right) \cup U_{\varepsilon}\left(a_{2}\right), \text { existiert } a_{3} \in X \backslash\left[U_{\varepsilon}\left(a_{1}\right) \cap U_{\varepsilon}\left(a_{2}\right)\right] ; \\
& \quad \ldots \\
& \text { da } X \not \subset \bigcup_{j=1}^{n} U_{\varepsilon}\left(a_{j}\right), \text { existiert } a_{n+1} \in X \backslash \bigcup_{j=1}^{n} U_{\varepsilon}\left(a_{j}\right) .
\end{aligned}
$$

Diese Konstruktion definiert induktiv eine Folge $\left(a_{n}\right)_{n}$ in $X$ derart, dass $d\left(a_{n}, a_{m}\right) \geq \varepsilon$ für alle $m, n \in \mathbb{N}, m \neq n$. Nach Voraussetzung muss $\left(a_{n}\right)_{n}$ aber eine konvergente Teilfolge $\left(a_{n_{j}}\right)_{j}$ besitzen. Mit $a:=\lim _{j \rightarrow \infty} a_{n_{j}}$ hat man dann

$$
\varepsilon \leq d\left(a_{n_{j}}, a_{n_{j+1}}\right) \leq d\left(a_{n_{j}}, a\right)+d\left(a, a_{n_{j+1}}\right) \rightarrow 0 \quad \text { für } j \rightarrow \infty .
$$

Widerspruch. Es folgt die Präkompaktheit von $X$.\\
"iii) $\Longrightarrow$ i)": Seien $U_{i} \subset X, i \in I$ offen mit

$$
X=\bigcup_{i \in I} U_{i} .
$$

Da $X$ präkompakt ist, existiert für jedes $n \in \mathbb{N}$ eine endliche Menge $M_{n} \subset X$ mit

$$
X=B_{\frac{1}{n}}\left(M_{n}\right)=\bigcup_{x \in M_{n}} B_{\frac{1}{n}}(x)
$$

Sei nun angenommen, dass $X$ nicht von endlich vielen der Mengen $U_{i}$ überdeckt wird. Da $M_{1}$ endlich ist, existiert $x_{1} \in M_{1}$ so, dass $A_{1}:=B_{1}\left(x_{1}\right)$ nicht von endlich vielen $U_{i}$ überdeckt wird. Da ferner $M_{2}$ endlich ist und $A_{1}=\bigcup_{x \in M_{2}} A_{1} \cap B_{\frac{1}{2}}(x)$ gilt, existiert $x_{2} \in M_{2}$ so, dass $A_{2}:=A_{1} \cap B_{\frac{1}{2}}\left(x_{2}\right)$ nicht von endlich vielen $U_{i}$ überdeckt wird. Induktiv findet man:\\
(**) Es gibt $x_{n} \in M_{n}$ so, dass $A_{n}:=A_{n-1} \cap B_{\frac{1}{n}}\left(x_{n}\right)$ nicht von endl. vielen $U_{i}$ überdeckt wird.

Dies liefert eine Folge abgeschlossener Mengen mit $A_{n+1} \subset A_{n}$ und diam $A_{n} \leq \frac{2}{n}$ für alle $n \in \mathbb{N}$. Da $X$ vollständig ist, folgt mit 3.7, dass ein $a \in \bigcap_{n \in \mathbb{N}} A_{n}$ existiert. Ferner existiert $i \in I$ mit $a \in U_{i}$. Da $U_{i}$ offen ist, existiert auch ein $\delta>0$ mit $B_{\delta}(a) \subset U_{i}$. Für $n>\frac{2}{\delta}$ ist $\operatorname{diam} A_{n}<\delta$ und $a \in A_{n}$, also $A_{n} \subset B_{\delta}(a) \subset U_{i}$ im Widerspruch zu (**). Der Widerspruch zeigt, dass eine endliche Teilmenge $J \subset I$ existiert mit $X=\bigcup_{i \in J} U_{i}$. Es folgt die Kompaktheit von $X$.

\subsection*{Bemerkung}
i). Ist $X$ kompakt und $A \subset X$ abgeschlossen, so ist $A$ kompakt.\\
ii). Ist $X$ kompakt, $(Y, \tilde{d})$ ein weiterer metrischer Raum und $f: X \rightarrow Y$ stetig, so ist auch $f(X) \subset Y$ kompakt, und $f$ ist gleichmäßig stetig.\\
iii). (Satz von Heine-Borel) Sei $E$ ein endlich dimensionaler normierter Raum und $A \subset E$. Dann ist $A$ genau dann kompakt, wenn $A$ beschränkt und abgeschlossen ist.

Beweis. Bekannt aus der Analysis II.

\subsection*{Bemerkung}
i). Ist $X$ präkompakt, so ist $X$ offensichtlich beschränkt. Ferner ist $X$ dann auch separabel, denn für jedes $n \in \mathbb{N}$ existiert eine endliche Teilmenge $M_{n} \subset X$ mit $X=B_{\frac{1}{n}}\left(M_{n}\right)$. Somit ist die Menge $M:=\bigcup_{n \in \mathbb{N}} M_{n}$ abzählbar und dicht in $X$.\\
ii). Für $A \subset X$ sind äquivalent:

\begin{itemize}
  \item $A$ ist präkompakt
  \item Für jedes $\varepsilon>0$ existiert eine präkompakte Menge $C \subset X$ mit $A \subset B_{\varepsilon}(C)$.
  \item Für jedes $\varepsilon>0$ existieren endlich viele Teilmengen $A_{1}, \ldots, A_{n}$ von $A$ mit $A \subset \bigcup_{i=1}^{n} A_{i}$ und diam $A_{i} \leq \varepsilon$ für $i=1, \ldots, n$.
  \item $\bar{A}$ ist präkompakt.
\end{itemize}

Beweis als Übung (Blatt 3)\\
iii). Für $A \subset X$ gilt wegen 4.2 und ii).:

\begin{itemize}
  \item $A$ relativ kompakt $\Longrightarrow A$ präkompakt
  \item ( $X, d$ ) vollständig und $A$ präkompakt $\Longrightarrow A$ relativ kompakt.
\end{itemize}

\subsection*{Satz}
Sei $(E,\|\cdot\|)$ ein normierter Raum und $A \subset E$. Dann gilt: $A$ ist präkompakt genau dann, wenn $A$ beschränkt ist und für jedes $\varepsilon>0$ ein endlich dimensionaler Unterraum $F \subset E$ existiert mit $A \subset B_{\varepsilon}(F)$.

Beweis. " $\Longrightarrow$ ": Sei $A$ präkompakt. Wie bereits bemerkt, ist $A$ dann auch beschränkt. Sei ferner $\varepsilon>0$, und sei $M \subset A$ endlich mit $A \subset B_{\varepsilon}(M)$. Setze $F:=\operatorname{span} M$. Dann ist $\operatorname{dim} F<\infty$ und $A \subset B_{\varepsilon}(F)$.\\
" $\Longleftarrow$ ": Nach Voraussetzung ist $r:=\sup _{x \in A}\|x\|<\infty$. Sei $\varepsilon>0$, und sei $F \subset E$ ein endlich dimensionaler Unterraum mit $A \subset B_{\varepsilon}(F)$. Setze $C:=\{x \in F:\|x\| \leq r+\varepsilon\}$. Nach 4.3 iii) ist $C$ dann kompakt, also auch präkompakt. Ferner existiert für jedes $x \in A$ ein $y \in F$ mit $\|x-y\|<\varepsilon$, also insbesondere $\|y\| \leq\|x\|+\varepsilon \leq r+\varepsilon$ und somit $y \in C$. Also ist $A \subset B_{\varepsilon}(C)$, und mit 4.4 i) folgt die Behauptung.

\subsection*{Rieszsches Lemma}
Sei $(E,\|\cdot\|)$ ein normierter Raum und $F \subset E$ ein echter abgeschlossener Unterraum. Sei ferner $r>0$ und $\varepsilon \in(0,1)$. Dann existiert ein $x \in E$ mit $\|x\|=r$ und $\operatorname{dist}(x, F) \geq(1-\varepsilon) r$.

Beweis. Da $F=\bar{F} \neq E$ gilt, existiert $a \in E \backslash F$ mit $\operatorname{dist}(a, F):=\alpha>0$. Ferner existiert $y \in F$ mit $\|a-y\|<\frac{\alpha}{1-\varepsilon}$. Sei nun $\lambda=\frac{r}{\|a-y\|}$ und $x=\lambda(a-y)$. Dann ist $\|x\|=r$, und für alle $w \in F$ gilt:

$$
\|x-w\|=\lambda\|a-\underbrace{\left(y+\lambda^{-1} w\right)}_{\in F}\| \geq \lambda \alpha=\frac{r \alpha}{\|a-y\|}>(1-\varepsilon) r .
$$

Es folgt $\operatorname{dist}(x, F) \geq(1-\varepsilon) r$.

\subsection*{Satz}
Für einen normierten $\operatorname{Raum}(E,\|\cdot\|)$ sind äquivalent:\\
i). $\operatorname{dim} E<\infty$\\
ii). Jede abgeschlossene und beschränkte Teilmenge von $E$ ist kompakt\\
iii). $B_{1}(0) \subset E$ ist präkompakt.

Beweis. "i) $\Longrightarrow$ ii)": siehe 4.3 (ii).\\
"ii) $\Longrightarrow$ iii)": $B_{1}(0)$ ist abgeschlossen und beschränkt, also kompakt nach Voraussetzung und damit auch präkompakt.\\
"iii) $\Longrightarrow$ i)": Nach 4.5 existiert ein endlich dimensionaler (und damit abgeschlossener) Unterraum $F \subset E$ mit $(*) B_{1}(0) \subset B_{\frac{1}{3}}(F)$. Ist $F \neq E$, so existiert nach 4.6 zu $\varepsilon=\frac{1}{3}$ und $r=1$ ein $x \in E \backslash F$ mit $\|x\|=1$ und $\operatorname{dist}(x, F) \geq\left(1-\frac{1}{3}\right)=\frac{2}{3}>\frac{1}{3}$, also $x \notin B_{\frac{1}{3}}(F)$ im Widerspruch zu (*). Es folgt also $E=F$; somit ist $E$ endlich dimensional. \(\square\)

\section*{Präkompaktheit für Mengen stetiger Abbildungen}
Seien stets $X, Y$ metrische Räume und $E$ ein normierter Raum.

\subsection*{Bemerkung}
Ist $X$ kompakt, so ist $\mathcal{C}(X, Y)=\mathcal{C}_{b}(X, Y)$ nach 4.3j),iii). Insbesondere ist dann also $\mathcal{C}(X)=\mathcal{C}_{b}(X)$ ein Banachraum mit Norm $\|\cdot\|_{\infty}$ nach 3.5 vi .

\subsection*{Definition}
Eine Teilmenge $H \subset \mathcal{C}(X, Y)$ heißt\\
i). gleichgradig stetig in $a \in X$, wenn für jedes $\varepsilon>0$ ein $\delta>0$ existiert mit $f\left(U_{\delta}(a)\right) \subset U_{\varepsilon}(f(a))$ für alle $f \in H$.\\
ii). gleichgradig stetig, wenn i) für alle $a \in X$ gilt.

Eine Folge $\left(f_{n}\right)_{n}$ in $\mathcal{C}(X, Y)$ heißt gleichgradig stetig (in $a \in X$ ), wenn dies für die Menge $\left\{f_{n}: n \in \mathbb{N}\right\}$ gilt.

\subsection*{Bemerkung und Beispiel}
i). Gleichgradige Stetigkeit vererbt sich auf Teilmengen von Funktionen.\\
ii). Sei $L>0$ fest gewählt und

$$
H:=\{f \in \mathcal{C}(X, Y): d(f(x), f(y)) \leq L d(x, y) \quad \text { für alle } x, y \in X\} .
$$

Dann ist $H$ gleichgradig stetig.\\
iii). Sei $X=[0,1], Y=\mathbb{R}$ und sei die Funktionenfolge $\left(f_{n}\right)_{n}$ in $\mathcal{C}(X, Y)$ definiert durch $f_{n}(x)=x^{n}$ für $n \in \mathbb{N}$. Wie man leicht sieht, ist diese Folge ist gleichgradig stetig in jedem Punkt $x \in[0,1)$. Sie ist allerdings nicht gleichgradig stetig in $x=1$ : Für festes $\varepsilon \in\left(0, \frac{1}{2}\right)$ gilt nämlich

$$
\varepsilon^{\frac{1}{n}} \rightarrow 1 \quad \text { für } n \rightarrow \infty,
$$

aber

$$
\left|f_{n}\left(\varepsilon^{\frac{1}{n}}\right)-f_{n}(1)\right|=1-\varepsilon>\varepsilon \quad \text { für alle } n \in \mathbb{N} \text {. }
$$

Also existiert kein $\delta>0$ derart, dass für alle $n \in \mathbb{N}$ die Implikation

$$
|x-1|<\delta \quad \Longrightarrow \quad\left|f_{n}(x)-f_{n}(1)\right|<\varepsilon
$$

gilt.

\subsection*{Definition und Satz}
Sei $l^{\infty}(X, E):=\{f: X \rightarrow E: f(X) \subset E$ beschränkt $\}$. Dann gilt:\\
i). $l^{\infty}(X, E)$ ist ein normierter Raum mit Norm $\|\cdot\|_{\infty}$ definiert durch $\|f\|_{\infty}:= \sup _{x \in X}\|f(x)\|$ für $f \in l^{\infty}(X, E)$.\\
ii). $\mathcal{C}_{b}(X, E)$ ist ein abgeschlossener Teilraum von $l^{\infty}(X, E)$.\\
iii). Ist $X$ kompakt, so ist $\mathcal{C}(X, E)=\mathcal{C}_{b}(X, E)$.\\
iv). Es gelten folgende Äquivalenzen:\\
$E$ Banachraum $\Longleftrightarrow l^{\infty}(X, E)$ Banachraum $\Longleftrightarrow \mathcal{C}_{b}(X, E)$ Banachraum\\
Beweis. i)-iii) wie für $E=\mathbb{K}$. iv): leichte Übung. \(\square\)

\subsection*{Satz von Arzelà und Ascoli}
Ist $X$ kompakt, so gilt für $H \subset \mathcal{C}(X, E)$ :\\
$H$ präkompakt $\Longleftrightarrow\left\{\begin{array}{l}\bullet H \text { ist gleichgradig stetig; } \\ \bullet \text { Für alle } a \in X \text { ist } H(a):=\{f(a): f \in H\} \text { in } E \text { präkompakt }\end{array}\right.$\\
Beweis. " $\Longleftarrow$ ": Sei $\varepsilon>0$. Da $H$ gleichgradig stetig ist, existiert zu jedem $a \in X$ ein $\delta_{a}>0$ mit $f\left(U_{\delta_{a}}(a)\right) \subset U_{\frac{\varepsilon}{3}}(f(a))$ für alle $f \in H$. Da $X=\bigcup_{a \in X} U_{\delta_{a}}(a)$ kompakt ist, existieren $a_{1}, \ldots, a_{n}$ derart, dass

$$
X=\bigcup_{i=1}^{n} U_{i} \quad \text { mit } U_{i}:=U_{\delta_{a_{i}}\left(a_{i}\right)} \text { gilt. }
$$

Da die Mengen $H\left(a_{i}\right) \subset E, i=1, \ldots, n$ präkompakt sind, existiert eine endliche Menge $M \subset E$ mit $\bigcup_{i=1}^{n} H\left(a_{i}\right) \subset B_{\frac{\varepsilon}{6}}(M)$. Definiert man für $v=\left(v_{1}, \ldots, v_{n}\right) \in M^{n}$ also

$$
L_{v}:=\left\{f \in H:\left\|f\left(a_{i}\right)-v_{i}\right\| \leq \frac{\varepsilon}{6} \text { für } i=1, \ldots, n\right\},
$$

so folgt nach Konstruktion $H \subset \bigcup_{v \in M^{n}} L_{v}$. Da es sich hier um eine endliche Vereinigung handelt, reicht es gemäß 4.4ii) nun, zu zeigen:

$$
\operatorname{diam} L_{v} \leq \varepsilon \text { in } \mathcal{C}(X, E) \text { bzgl. }\|\cdot\|_{\infty} \text { für alle } v \in M^{n} \text {. }
$$

Seien dazu $f, g \in L_{v}$, und sei $x \in X$ beliebig. Dann ist $x \in U_{i}$ für ein $i \in\{1, \ldots, n\}$, und somit

$$
\|f(x)-g(x)\| \leq \underbrace{\left\|f(x)-f\left(a_{i}\right)\right\|}_{\leq \frac{\varepsilon}{3}}+\underbrace{\left\|f\left(a_{i}\right)-v_{i}\right\|}_{\leq \frac{\varepsilon}{6}}+\underbrace{\left\|v_{i}-g\left(a_{i}\right)\right\|}_{\leq \frac{\varepsilon}{6}}+\underbrace{\left\|g\left(a_{i}\right)-g(x)\right\|}_{\leq \frac{\varepsilon}{3}} \leq \varepsilon .
$$

Es folgt $\|f-g\|_{\infty} \leq \varepsilon$, und damit ist diam $L_{v} \leq \varepsilon$.\\
" $\Longrightarrow$ ": Sei $a \in X$ und $\varepsilon>0$. Nach Voraussetzung existiert eine endliche Teilmenge $M \subset H$ mit

$$
H \subset \bigcup_{f \in M} B_{\frac{\varepsilon}{4}}(f) \quad \text { in } \mathcal{C}(X, E) \text { bzgl. }\|\cdot\|_{\infty}
$$

Da $M$ endlich ist, existiert zudem $\delta>0$ so, dass $f\left(U_{\delta}(a)\right) \subset U_{\frac{\varepsilon}{3}}(f(a))$ für alle $f \in M$ gilt. Sei nun $g \in H$ beliebig. Dann existiert $f \in M$ mit $\|g-f\|_{\infty} \leq \frac{\varepsilon}{4}$, und somit gilt für $y \in U_{\delta}(a)$ :

$$
\begin{aligned}
\|g(y)-g(a)\| & \leq\|g(y)-f(y)\|+\|f(y)-f(a)\|+\|f(a)-g(a)\| \\
& \leq 2\|f-g\|_{\infty}+\|f(y)-f(a)\| \leq 2 \frac{\varepsilon}{4}+\frac{\varepsilon}{3}<\varepsilon .
\end{aligned}
$$

Es folgt $g\left(U_{\delta}(a)\right) \subset U_{\varepsilon}(g(a))$, und dies liefert die gleichgradige Stetigkeit von $H$ in a. Nach Definition von $\|\cdot\|_{\infty}$ liefert (**) aber auch direkt die Inklusion

$$
H(a) \subset \bigcup_{f \in M} B_{\frac{\varepsilon}{4}}(f(a)) \quad \text { in } E \text { bzgl. }\|\cdot\| ;
$$

also ist $H(a)$ in $E$ präkompakt. \(\square\)

\subsection*{Korollar}
Sei $X$ kompakt und $\left(f_{n}\right)_{n}$ eine beschränkte und gleichgradig stetige Folge in $\mathcal{C}(X)$. Dann hat $\left(f_{n}\right)_{n}$ eine gleichmäßig konvergente Teilfolge.

Beweis. Sei $H:=\left\{f_{n}: n \in \mathbb{N}\right\} \subset \mathcal{C}(X)$. Dann ist $H$ gleichgradig stetig und $H(a)=\{f(a): f \in H\} \subset \mathbb{K}$ beschränkt und damit auch präkompakt für alle $a \in X$. Mit 4.12 folgt, dass $H$ in $\mathcal{C}(X)$ präkompakt und somit auch relativ kompakt ist (nach 4.4 iii), da $\mathcal{C}(X)$ vollständig ist). Also hat $\left(f_{n}\right)_{n}$ eine in $\bar{H}$ konvergente Teilfolge. \(\square\)

\section*{Kompakte Operatoren}
Seien $E, F$ normierte $\mathbb{K}$-Vektorräume.

\subsection*{Definition}
Eine lineare Abbildung $T: E \rightarrow F$ heißt kompakt, falls für jede beschränkte Folge $\left(x_{n}\right)_{n} \subset E$ die Bildfolge $\left(T x_{n}\right)_{n}$ eine konvergente Teilfolge besitzt.\\
Wir setzen $\mathcal{K}(E, F):=\{T: E \rightarrow F: T$ kompakt $\}$ und $\mathcal{K}(E):=\mathcal{K}(E, E)$.

\subsection*{Beispiel}
Seien $[a, b],[c, d] \subset \mathbb{R}$ kompakte Intervalle und $k:[c, d] \times[a, b] \rightarrow \mathbb{C}$ eine stetige Funktion. Dann ist der Integraloperator $K \in \mathcal{L}(\mathcal{C}[a, b], \mathcal{C}[c, d])$ aus 2.14ii), gegeben durch

$$
[K u](t):=\int_{a}^{b} k(t, s) u(s) d s \quad \text { für } t \in[c, d]
$$

kompakt. Sei dazu $\left(u_{n}\right)_{n}$ eine Folge in $\mathcal{C}[a, b]$ mit $r:=\sup _{n \in \mathbb{N}}\left\|u_{n}\right\|_{\infty}<\infty$. Aufgrund der Stetigkeit von $K$ (siehe 2.14ii)) ist dann auch die Folge $\left(K u_{n}\right)_{n}$ beschränkt in $\mathcal{C}[c, d]$. Mit 4.13 reicht es also, zu zeigen, dass die Folge $\left(K u_{n}\right)_{n}$ gleichgradig stetig ist. Sei dazu $t_{0} \in[c, d]$ und $\varepsilon>0$. Da $k$ auf der kompakten Menge $[c, d] \times[a, b]$ gleichmäßig stetig ist, existiert $\delta>0$ mit

$$
\left|k(t, s)-k\left(t_{0}, s\right)\right| \leq \frac{\varepsilon}{2(b-a) r} \quad \text { für alle } t \in U_{\delta}\left(t_{0}\right) \text { und } s \in[a, b] \text {. }
$$

Es folgt dann für $t \in U_{\delta}\left(t_{0}\right)$ und $n \in \mathbb{N}$ :

$$
\left|K u_{n}(t)-K u_{n}\left(t_{0}\right)\right| \leq \int_{a}^{b}\left|k(t, s)-k\left(t_{0}, s\right)\right|\left|u_{n}(s)\right| d s \leq r \int_{a}^{b}\left|k(t, s)-k\left(t_{0}, s\right)\right| d s \leq \frac{\varepsilon}{2}<\varepsilon .
$$

Die liefert die benötigte gleichgradige Stetigkeit.

\subsection*{Bemerkung}
Sei $T: E \rightarrow F$ linear.\\
i). Äquivalent sind:

\begin{itemize}
  \item $T$ kompakt
  \item Ist $M \subset E$ beschränkt, so ist $T(M) \subset F$ relativ kompakt.
  \item $T\left(B_{1}(0)\right) \subset F$ ist relativ kompakt (hier sei wieder $B_{1}(0) \subset E$ die Einheitskugel).\\
ii). Ist $T$ kompakt, so ist $T$ auch stetig.\\
iii). Ist $T$ stetig und $\operatorname{dim}$ Bild $T<\infty$, so ist $T$ kompakt.\\
iv). Ist $\operatorname{dim} E<\infty$, so ist $T$ kompakt.
\end{itemize}

Ferner ist nach 4.7 die identische Abbildung id $: E \rightarrow E$ genau dann kompakt, wenn $\operatorname{dim} E<\infty$ ist.

\subsection*{Satz}
i). $\mathcal{K}(E, F)$ ist ein Unterraum von $\mathcal{L}(E, F)$.\\
ii). Ist $F$ vollständig, so ist $\mathcal{K}(E, F)$ in $\mathcal{L}(E, F)$ abgeschlossen.

Beweis. i) Offensichtlich ist die Nullabbildung in $\mathcal{K}(E, F)$. Seien $S, T \in \mathcal{K}(E, F)$ und $\alpha \in \mathbb{K}$. Sei ferner $\left(x_{n}\right)_{n} \subset E$ eine beschränkte Folge. Da $S$ kompakt ist, existiert eine Teilfolge $\left(x_{n_{k}}\right)_{k}$ und $y \in F$ mit $\lim _{k \rightarrow \infty} S x_{n_{k}}=y$. Da $T$ kompakt ist, existiert ferner eine Teilfolge $\left(x_{n_{k_{j}}}\right)_{j}$ und $z \in F$ mit $\lim _{j \rightarrow \infty} T x_{n_{k_{j}}}=z$. Folglich existiert auch

$$
\lim _{j \rightarrow \infty}(\alpha S+T) x_{n_{k_{j}}}=\alpha y+z \in F,
$$

und somit ist $\alpha S+T \in \mathcal{K}(E, F)$.\\
ii) Seien $T_{n} \in \mathcal{K}(E, F) n \in \mathbb{N}$ derart, dass $T:=\lim _{n \rightarrow \infty} T_{n}$ in $\mathcal{L}(E, F)$ existiert. Es reicht, zu zeigen, dass $T\left(B_{1}(0)\right) \subset F$ präkompakt ist; dann ist diese Menge aufgrund der Vollständigkeit von $F$ auch relativ kompakt. Sei dazu $\varepsilon>0$. Dann existiert $n \in \mathbb{N}$ mit $\left\|T_{n}-T\right\| \leq \varepsilon$, also $\left\|T_{n} x-T x\right\| \leq \varepsilon$ für alle $x \in B_{1}(0)$. Somit $T\left(B_{1}(0)\right) \subset B_{\varepsilon}\left(T_{n}\left(B_{1}(0)\right)\right)$, wobei $T_{n}\left(B_{1}(0)\right) \subset F$ nach Voraussetzung präkompakt ist. Mit 4.4i) folgt die Präkompaktheit von $T\left(B_{1}(0)\right)$.

\subsection*{Beispiel}
Sei $1 \leq p<\infty, E=\ell^{p},\left(\varepsilon_{n}\right)_{n}$ eine Nullfolge in $\mathbb{K}$ und $T \in \mathcal{L}(E)$ definiert durch

$$
T x=\left(\varepsilon_{n} x_{n}\right)_{n} \quad \text { für } x=\left(x_{n}\right)_{n} \in \ell^{p} .
$$

Dann ist $T \in \mathcal{K}(E)$, denn:\\
Seien $T_{k} \in \mathcal{L}(E), k \in \mathbb{N}$ definiert durch

$$
T_{k} x=\left(\varepsilon_{1} x_{1}, \ldots, \varepsilon_{k} x_{k}, 0,0, \ldots\right) \quad \text { für } x=\left(x_{n}\right)_{n} \in \ell^{p} .
$$

Da $\operatorname{dim}$ Bild $T_{k}<\infty$ ist, folgt $T_{k} \in \mathcal{K}(E)$ nach 4.16ii). Ferner ist

$$
\left\|\left(T-T_{k}\right) x\right\|_{p}^{p}=\sum_{n>k}\left|\varepsilon_{n} x_{n}\right|^{p} \leq \delta_{k}^{p}\|x\|_{p}^{p} \quad \text { für } x=\left(x_{n}\right)_{n} \in \ell^{p} \text { mit } \delta_{k}:=\sup _{n>k}\left|\varepsilon_{n}\right| .
$$

Es folgt $\left\|T-T_{k}\right\| \leq \delta_{k} \rightarrow 0$ für $k \rightarrow \infty$. Somit ist $T \in \mathcal{K}(E)$ nach 4.17i).

\subsection*{Satz}
Seien $D, G$ weitere normierte Räume, und sei $T \in \mathcal{K}(E, F)$ sowie $T_{1} \in \mathcal{L}(D, E)$ und $T_{2} \in \mathcal{L}(F, G)$. Dann gilt:

$$
T T_{1} \in \mathcal{K}(D, F) \quad \text { und } \quad T_{2} T \in \mathcal{K}(E, G) .
$$

Beweis. Sei zunächst $M \subset D$ beschränkt. Dann ist auch $T_{1}(M)$ in $E$ beschränkt, da $T_{1}$ stetig ist. Da ferner $T$ kompakt ist, ist $T\left(T_{1}(M)\right)$ relativ kompakt in $F$. Es folgt $T T_{1} \in \mathcal{K}(D, F)$.\\
Sei num $\left(x_{n}\right)_{n}$ eine beschränkte Folge in $E$. Da $T$ kompakt ist, existiert eine Teilfolge $\left(x_{n_{k}}\right)_{k}$ und $y \in F$ mit $\lim _{k \rightarrow \infty} T x_{n_{k}}=y$. Da $T_{2}$ stetig ist, existiert $\lim _{k \rightarrow \infty} T 2 T x_{n_{k}}=T_{2} y \in G$. Es folgt $T_{2} T \in \mathcal{K}(\stackrel{k \rightarrow \infty}{E}, \stackrel{-}{G})$.

\subsection*{Definition}
Man sagt, der normierte Raum $E$ sei stetig in $F$ eingebettet, wenn es eine injektive Abbildung $i \in \mathcal{L}(E, F)$ gibt. Die Abbildung $i$ nennt man dann Einbettung von $E$ in $F$ und schreibt $E \stackrel{i}{\hookrightarrow} F$.\\
Ist ferner $i$ kompakt, so sagt man: $E$ ist in $F$ kompakt eingebettet.

\subsection*{Beispiel}
Seien $a, b \in \mathbb{R}, a<b$. Wir betrachten den Banachraum $\mathcal{C}^{1}[a, b]$ mit Norm $\|\cdot\|_{C^{1}}$ definiert durch $\|u\|_{C^{1}}:=\|u\|_{\infty}+\left\|u^{\prime}\right\|_{\infty}$, vgl. Übungsblatt 3 .

Dann ist $\mathcal{C}^{1}[a, b]$ kompakt in $\mathcal{C}[a, b]$ eingebettet, denn:\\
Für $u \in B_{1}(0) \subset \mathcal{C}^{1}[a, b]$ und $x, y \in[a, b]$ ist nach dem Mittelwertsatz

$$
|u(x)-u(y)| \leq\left\|u^{\prime}\right\|_{\infty}|x-y| \leq|x-y| .
$$

Gemäß 4.10ii) folgt daher die gleichgradige Stetigkeit von $B_{1}(0) \subset \mathcal{C}^{1}[a, b]$. Ferner ist $\|u\|_{\infty} \leq\|u\|_{C^{1}} \leq 1$ für alle $u \in B_{1}(0) \subset \mathcal{C}^{1}[a, b]$. Mit 4.12 (Satz von Arzelà und Ascoli) folgt daher, dass die Einheitskugel $B_{1}(0)$ von $\mathcal{C}^{1}[a, b]$ im Raum $\mathcal{C}[a, b]$ präkompakt und damit auch relativ kompakt ist.

\section*{Der Satz von Stone-Weierstraß}
Im Folgenden sei ( $X, d$ ) stets ein kompakter metrischer Raum.

\subsection*{Bemerkung}
Im Folgenden wollen wir ein hinreichendes Kriterium dafür herleiten, dass eine Teilmenge im normierten Raum ( $\mathcal{C}(X),\|\cdot\|_{\infty}$ ) dicht liegt. Dabei nutzen wir die Tatsache, dass dieser Raum auch eine kommutative $\mathbb{K}$-Algebra mit Eins ist, d.h.:

\begin{itemize}
  \item $\mathcal{C}(X)$ ist abgeschlossen unter der (punktweisen) Multiplikation - von Funktionen, und diese ist assoziativ, distributiv (bzgl. der Addition in $\mathcal{C}(X)$ ) und kommutativ.
  \item Die konstante Einsfunktion $1=1_{X}$ ist ein Einselement für die Multiplikation.
  \item Die Skalarmultiplikation verträgt sich mit der Multiplikation von Funktionen, d.h. es gilt $\lambda(f \cdot g)=(\lambda f) \cdot g=f \cdot(\lambda g)$ für $f, g \in \mathcal{C}(X)$ und $\lambda \in \mathbb{K}$.
\end{itemize}

Ein Unterraum $A \subset \mathcal{C}(X)$ heißt Teilalgebra von $\mathcal{C}(X)$ (mit 1), wenn gilt:

\begin{itemize}
  \item $1 \in A$
  \item Sind $f, g \in A$, so ist auch $f \cdot g \in A$.
\end{itemize}

\subsection*{Satz von Stone-Weierstraß}
Sei $A$ eine Teilalgebra von $\mathcal{C}(X)=\mathcal{C}(X, \mathbb{K})$ mit folgenden Eigenschaften.\\
(i) $A$ trennt die Punkte von $X$, d.h. für je zwei Punkte $x, y \in X$ mit $x \neq y$ existiert eine Funktion $f \in A$ mit $f(x) \neq f(y)$.\\
(ii) Im Fall $\mathbb{K}=\mathbb{C}$ ist $A$ abgeschlossen unter komplexer Konjugation, d.h. für $f \in A$ ist auch $\bar{f} \in A$.

Dann liegt $A$ dicht in $\mathcal{C}(X)$.\\
Um dies zu beweisen, benötigen wir zunächst zwei Hilfssätze

\subsection*{Hilfssatz (Spezialfall der binomischen Reihe)}
Für $s \in[-1,1]$ gilt

$$
\sqrt{1-s}=\sum_{n=0}^{\infty} a_{n} s^{n}
$$

mit

$$
a_{n}=(-1)^{n}\binom{\frac{1}{2}}{n}=(-1)^{n} \prod_{j=1}^{n} \frac{\frac{3}{2}-j}{j} \quad \text { für } n \in \mathbb{N}_{0}
$$

und die Reihe konvergiert absolut.

Beweis. Die aus der Analysis bekannte binomische Reihenentwicklung zur Potenz $\frac{1}{2}$ liefert (4.1) zunächst für $s \in(-1,1)$. Wir zeigen nun die Abschätzung

$$
\left|a_{n}\right| \leq n^{-\frac{3}{2}} \quad \text { für } n \geq 1 \text {. }
$$

Dies ist klar für $a_{1}=-\frac{1}{2}$. Für $n \geq 2$ gilt zudem

$$
\frac{\left|a_{n}\right|}{\left|a_{n-1}\right|}=\left|\frac{\frac{3}{2}-n}{n}\right|=1-\frac{3}{2 n} \leq\left(1-\frac{1}{n}\right)^{\frac{3}{2}}=\left(\frac{n}{n-1}\right)^{-\frac{3}{2}} .
$$

Hier haben wir die Konvexität der Funktion $(0, \infty) \rightarrow \mathbb{R}, t \mapsto t^{\frac{3}{2}}$ verwendet (diese liefert die Ungleichung $(1+t)^{\frac{3}{2}} \geq 1+\frac{3}{2} t$ für $t \in(-1, \infty)$ ). Per Induktion folgt also für $n \geq 2$ :

$$
\left|a_{n}\right| \leq\left(\frac{n}{n-1}\right)^{-\frac{3}{2}}\left|a_{n-1}\right| \leq\left(\frac{n}{n-1}\right)^{-\frac{3}{2}}(n-1)^{-\frac{3}{2}}=n^{-\frac{3}{2}} .
$$

Insbesondere konvergiert die Reihe in (4.1) für alle $s \in[-1,1]$ absolut, und mit dem Abelschen Grenzwertsatz folgt, dass 4.24 auch für $s= \pm 1$ gilt.\\
(Randbemerkung: Es gilt $a_{n}<0$ für $n \geq 1$ )

\subsection*{Hilfssatz}
Sei $A$ eine Teilalgebra von $\mathcal{C}(X, \mathbb{R})$. Dann gilt:\\
(a) Der Abschluss $\bar{A}$ von $A$ ist auch eine Teilalgebra von $\mathcal{C}(X, \mathbb{R})$.\\
(b) Ist $A$ abgeschlossen in $\mathcal{C}(X, \mathbb{R})$, so gilt:\\
(i) Ist $f \in A$ mit $f \geq 0$, so ist auch $\sqrt{f} \in A$.\\
(ii) Sind $f, g \in A$, so sind auch die Funktionen $\min \{f, g\}, \max \{f, g\} \in A$.

Beweis. (a) Sei $f \in \mathcal{C}(X)$ fest. Dann ist die lineare Abbildung $\mathcal{C}(X) \rightarrow \mathcal{C}(X)$, $g \mapsto f g$ stetig wegen der Ungleichung $\|f g\|_{\infty} \leq\|f\|_{\infty}\|g\|_{\infty}$. Sukzessive folgt daraus:

\begin{itemize}
  \item $f \in A, g \in \bar{A} \quad \Longrightarrow \quad f \cdot g \in \bar{A}$;
  \item $f, g \in \bar{A} \quad \Longrightarrow \quad f \cdot g \in \bar{A}$.
\end{itemize}

Da ferner $1 \in A \subset \bar{A}$ ist, folgt die Behauptung.\\
(b)(i) Ohne Einschränkung sei $0 \leq f \leq 1$. Dann ist $g=1-f \in A$ und $0 \leq g \leq 1$. Gemäß 4.24 gilt

$$
\sqrt{f(x)}=\sqrt{1-g(x)}=\sum_{n=0}^{\infty} a_{n} g^{n}(x) \quad \text { für } x \in X
$$

mit den Binomialkoeffizienten aus (4.2). Da $A$ abgeschlossen ist, reicht es, zu zeigen, dass die Reihe in (4.4) bzgl. der Norm $\|\cdot\|_{\infty}$ von $\mathcal{C}(X, \mathbb{R})$ konvergiert; dann ist $\sqrt{f} \in A$. Wegen $\left\|g^{n}\right\|_{\infty} \leq 1$ für alle $n \in \mathbb{N}_{0}$ ist aber

$$
\sum_{n=0}^{\infty}\left\|a_{n} g^{n}\right\|_{\infty} \leq \sum_{n=0}^{\infty}\left|a_{n}\right|\|g\|_{\infty}^{n} \leq \sum_{n=0}^{\infty}\left|a_{n}\right|<\infty
$$

gemäß 4.24. d.h. die Reihe $\sum_{n=0}^{\infty} a_{n} g^{n}$ konvergiert absolut im Banachraum $\mathcal{C}(X, \mathbb{R})$. Somit konvergiert sie auch in $\mathcal{C}(X, \mathbb{R})$ nach 3.16.\\
(b)(ii) Mit $f$ ist auch $f^{2} \in A$ und somit $|f|=\sqrt{f^{2}} \in A$ wegen (b)(i). Mit $f, g \in A$ sind also auch

$$
\min \{f, g\}=\frac{1}{2}(f+g-|f-g|) \quad \text { und } \quad \max \{f, g\}=-\min \{-f,-g\}
$$

Funktionen in $A$.

Beweis von 4.23. Wir betrachten zunächst den Fall $\mathbb{K}=\mathbb{R}$. Sei $f \in \mathcal{C}(X, \mathbb{R})$ und $\varepsilon>0$ gegeben. Es reicht, zu zeigen, dass $g \in \bar{A}$ existiert mit $\|f-g\|_{\infty} \leq \varepsilon$. Zur Konstruktion der Funktion $g$ betrachten wir zunächst $x, y \in X$ beliebig mit $x \neq y$. Nach Voraussetzung existiert eine Funktion $h=h_{x, y} \in A$ mit $h(x) \neq h(y)$. Mit dieser Wahl definieren wir num $g_{x, y} \in A$ durch

$$
g_{x, y}(z)=f(y)+(f(x)-f(y)) \frac{h(z)-h(y)}{h(x)-h(y)} \quad \text { für } z \in X
$$

dann gilt $g_{x, y}(x)=f(x)$ und $g_{x, y}(y)=f(y)$. Ferner definieren wir die konstante Funktion $g_{x, x} \in A, g_{x x} \equiv f(x)$ für $x \in X$.\\
Sei nun $x \in X$ fest gewählt. Für $y \in X$ betrachten wir dann die offene Umgebung

$$
U_{y}:=\left\{z \in X: g_{x, y}(z)<f(z)+\varepsilon\right\}
$$

von $y$. Aufgrund der Kompaktheit von $X$ existieren $y_{1}, \ldots, y_{n} \in X$ mit $X=\bigcup_{j=1}^{n} U_{y_{j}}$. Ferner ist wegen 4.25 die Funktion $g_{x}:=\min \left\{g_{x, y_{i}}: i=1, \ldots, n\right\} \in \bar{A}$. Dabei gilt $g_{x}(x)=f(x)$ und $g_{x}(z)<f(z)+\varepsilon$ für alle $z \in X$.\\
Daher ist $V_{x}:=\left\{z \in X: g_{x}(z)>f(z)-\varepsilon\right\}$ eine offene Umgebung von $x$. Wiederum aufgrund der Kompaktheit von $X$ existieren $x_{1}, \ldots, x_{m} \in X$ mit $X=\bigcup_{j=1}^{m} V_{x_{j}}$. Ferner ist auch $g:=\max \left\{g_{x_{j}}: j=1, \ldots, m\right\} \in \bar{A}$, und es gilt

$$
f(z)-\varepsilon<g<f(z)+\varepsilon \quad \text { für alle } z \in X .
$$

Es folgt $\|f-g\|_{\infty} \leq \varepsilon$, wie benötigt.\\
Wir betrachten schließlich den Fall $\mathbb{K}=\mathbb{C}$. Für $A_{0}:=\{\operatorname{Re} f: f \in A\}$ gilt dann:\\
1.) Für $f \in A$ ist auch $\operatorname{Re} f=\frac{1}{2}(f+\bar{f}) \in A$, und somit folgt $A_{0} \subset A$.\\
2.) Für $f \in A$ ist $\operatorname{Im} f=\operatorname{Re}(-i f) \in A_{0}$.

Aus 1.) und 2.) folgt die Mengengleichheit

$$
\text { (*) } \quad A=A_{0}+i A_{0},
$$

und diese impliziert insbesondere, dass $A_{0}$ die Punkte von $X$ trennt. Ferner ist $A_{0}$ eine Teilalgebra von $\mathcal{C}(X, \mathbb{R})$, denn:\\
Für $\operatorname{Re} f, \operatorname{Re} g \in A_{0}$ mit $f, g \in A$ ist

$$
\operatorname{Re} f \cdot \operatorname{Re} g=\operatorname{Re}[f \operatorname{Re} g] \in A_{0}
$$

da nach 1.) auch $\operatorname{Re} g$ und somit $f \operatorname{Re} g$ in $A$ liegt. Ferner ist offensichtlich $1 \in A_{0}$. Wie bereits gezeigt, liegt $A_{0}$ also dicht in $\mathcal{C}(X, \mathbb{R})$, und somit liegt $A$ wegen (*) auch dicht in $\mathcal{C}(X, \mathbb{C})$.

\subsection*{Korollar}
i) Sei $X \subset \mathbb{R}^{N}$ kompakt. Dann liegen die Polynomfunktionen in $x_{1}, \ldots, x_{N}$ (mit Koeffizienten in $\mathbb{K}$ ) dicht in $\mathcal{C}(X, \mathbb{K})$.\\
ii) Sei $X \subset \mathbb{C}^{N}$ kompakt. Dann liegen die Polynomfunktionen in $z_{1}, \ldots, z_{N}$, $\bar{z}_{1}, \ldots, \bar{z}_{N}$ mit komplexen Koeffizienten dicht in $\mathcal{C}(X, \mathbb{C})$.

Beweis. i) Offensichtlich bilden die genannten Funktionen eine Teilalgebra von $\mathcal{C}(X, \mathbb{K})$, welche die Punkte von $X$ trennt, da zu je zwei verschiedenen Punkten $x, y \in \mathbb{R}^{N}$ eine affin lineare Funktion $g: \mathbb{R}^{N} \rightarrow \mathbb{R}$ existiert mit $g(x) \neq g(y)$ (z.B. die Funktion $\left.t \mapsto g(t)=(t-y) \cdot(x-y), t \in \mathbb{R}^{N}\right)$. Also folgt die Behauptung aus den Satz von Stone-Weierstraß.\\
ii) Die Behauptung folgt wie i) aus 4.23, wenn man zusätzlich beachtet, dass die hier betrachteten Funktionen eine Teilalgebra bilden, welche auch unter der komplexen Konjugation abgeschlossen ist.

\subsection*{Definition und Korollar}
Sei $\mathcal{C}_{2 \pi}(\mathbb{R}, \mathbb{C})$ der Vektorraum der stetigen $2 \pi$-periodischen Funktionen $\mathbb{R} \rightarrow \mathbb{C}$. Dann liegt der Teilraum $T$ der trigonometrischen Polynome dicht in $\mathcal{C}_{2 \pi}(\mathbb{R}, \mathbb{C})$ bzgl. $\|\cdot\|_{\infty}$. Dabei heißt eine Funktion $f \in \mathcal{C}_{2 \pi}(\mathbb{R}, \mathbb{C})$ trigonometrischen Polynom, wenn es $n \in \mathbb{N}$ und $a_{k} \in \mathbb{C}, k=-n, \ldots, n$ gibt mit $f(t)=\sum_{k=-n}^{n} a_{k} e^{i k t}$ für $t \in \mathbb{R}$.

Beweis. Sei $S^{1}:=\{z \in \mathbb{C}:|z|=1\} \subset \mathbb{C}$ der Einheitskreis. Dann ist die Abbildung

$$
I: C\left(S^{1}, \mathbb{C}\right) \rightarrow \mathcal{C}_{2 \pi}(\mathbb{R}, \mathbb{C}), \quad[I f](t):=f\left(e^{i t}\right)
$$

ist ein isometrischer Isomorphismus (bzgl. $\|\cdot\|_{\infty}$ ), welcher Polynome in $z$ und $\bar{z}$ genau auf trigonometrische Polynome abbildet.\\
Die Behauptung folgt also durch Anwendung von 4.26 (ii) auf $X:=S^{1} \subset \mathbb{C}$.

\subsection*{Korollar}
Ist ( $X, d$ ) ein kompakter metrischer Raum, so ist $\mathcal{C}(X)$ separabel.\\
Beweis. Da $X$ separabel ist (siehe 4.4), existieren $x_{n} \in X, n \in \mathbb{N}$ derart, dass $\left\{x_{n}: n \in \mathbb{N}\right\}$ in $X$ dicht liegt. Für $n \in \mathbb{N}$ sei nun $\varphi_{n} \in \mathcal{C}(X)$ definiert durch $\varphi_{n}(x):= d\left(x, x_{n}\right)$. Sei ferner $A \subset \mathcal{C}(X)$ die kleinste Teilalgebra, welche die Funktionen $\varphi_{n}$, $n \in \mathbb{N}$ enthält. Als Vektorraum ist $A=\operatorname{span} M$, wobei $M$ die Menge der endlichen

Produkte $\varphi_{n_{1}} \cdots \varphi_{n_{m}}$ mit $m \in \mathbb{N}_{0}$ und $n_{l} \in \mathbb{N}, l=1, \ldots, m$ ist (leeres Produkt ist die Einsfunktion). Offensichtlich ist $M$ abzählbar. Ferner trennt die Algebra $A$ die Punkte von $X$ : Sind nämlich $x, y \in X$ mit $x \neq y$, so existiert $n \in \mathbb{N}$ mit $\varphi_{n}(x)=d\left(x, x_{n}\right)<\frac{d(x, y)}{2}$. Es folgt dann

$$
\varphi_{n}(y)=d\left(y, x_{n}\right)>d(y, x)-d\left(x, x_{n}\right)>\frac{d(x, y)}{2}>\varphi_{n}(x) .
$$

$\operatorname{Im}$ Fall $\mathbb{K}=\mathbb{C}$ ist $A$ zudem abgeschlossen unter komplexer Konjugation, da die Funktionen $\varphi_{n}, n \in \mathbb{N}$ reellwertig sind. Also folgt mit dem Satz von Stone-Weierstraß, dass $A$ in $\mathcal{C}(X)$ dicht liegt. Somit ist $M$ eine abzählbare totale Teilmenge von $\mathcal{C}(X)$, und damit ist $\mathcal{C}(X)$ separabel nach 2.5. \(\square\)

\section*{§5 Struktursätze über stetige lineare Abbildungen}
Stets seien $E, F$ normierte $\mathbb{K}$-Vektorräume

\subsection*{Satz}
i). Seien $X, Y$ metrische Räume. Ist $Y$ vollständig, $D \subset X$ dicht und $f_{0}: D \rightarrow Y$ gleichmäßig stetig, so hat $f_{0}$ genau eine stetige Fortsetzung $f: X \rightarrow Y$. Ist ferner $f_{0}$ eine Isometrie, so ist $f$ ebenfalls eine Isometrie.\\
ii). Sei $F$ ein Banachraum, $D \subset E$ ein dichter linearer Teilraum und $T_{0} \in \mathcal{L}(D, F)$ gegeben. Dann existiert genau eine Fortsetzung $T \in \mathcal{L}(E, F)$ von $T_{0}$, d.h. es existiert genau ein $T \in \mathcal{L}(E, F)$ mit $\left.T\right|_{D}=T_{0}$. Dabei gilt $\|T\|=\left\|T_{0}\right\|$.

Beweis. (i) Ist $x \in X$ und $\left(x_{n}\right)_{n}$ eine Folge in $X$ mit $x_{n} \rightarrow x$ für $n \rightarrow \infty$, so ist $\left(f_{0}\left(x_{n}\right)\right)_{n}$ eine Cauchyfolge in $Y$, da $f$ gleichmäßig stetig ist. Da $Y$ vollständig ist, existiert somit $f(x):=\lim _{n \rightarrow \infty} f_{0}\left(x_{n}\right) \in Y$. Diese Definition ist auch unabhängig von der Wahl der Folge $\left(x_{n}\right)_{n}$, wie ein Folgenmischverfahren zeigt.\\
Ferner ist die so definierte Funktion $f: X \rightarrow Y$ offensichtlich die einzig stetige Funktion mit $\left.f\right|_{D}=f_{0}$ (da $D \subset X$ dicht liegt). Man sieht leicht, dass $f$ eine Isometrie ist, falls dies für $f_{0}$ gilt.\\
(ii) Übungsaufgabe (man verwende (i)).

\subsection*{Hauptsatz (von der gleichmäßigen Beschränktheit)}
Sei $E$ ein Banachraum, und sei $H \subset \mathcal{L}(E, F)$ punktweise beschränkt, d. h.,

$$
\text { für jedes } x \in E \text { existiere } C_{x}>0 \text { mit }\|T x\|_{F} \leq C_{x} \text { für alle } T \in H \text {. }
$$

Dann ist $H$ beschränkt in $\mathcal{L}(E, F)$, d. h. es existiert $C>0$ mit $\|T\| \leq C$ für alle $T \in H$.

Beweis. Für $n \in \mathbb{N}$ sei

$$
A_{n}:=\left\{x \in E:\|T x\|_{F} \leq n \text { für alle } T \in H\right\}=\bigcap_{T \in H} T^{-1}(\underbrace{B_{n}(0)}_{\subset F}) \text {. }
$$

Dann ist $A_{n}$ abgeschlossen in $E$ für alle $n \in \mathbb{N}$, und nach Voraussetzung gilt $E= \bigcup_{n \in \mathbb{N}} A_{n}$. Da $E$ vollständig ist, existiert nach 3.9 i) ein $n \in \mathbb{N}, b \in E$ und $s>0$ mit $B_{s}(b) \subset A_{n}$, d.h.

$$
\|T x\|_{F} \leq n \quad \text { für alle } x \in B_{s}(b), T \in H .
$$

Sei num $v \in E \backslash\{0\}$ und $T \in H$ beliebig. Dann ist $x_{v}:=b+s \frac{v}{\|v\|_{E}} \in B_{s}(b)$ und $v=\frac{\|v\|_{E}}{s}\left(x_{v}-b\right)$. Es folgt

$$
\|T v\|_{F} \leq \frac{\|v\|_{E}}{s}\left\|T x_{v}-T b\right\|_{F} \leq \frac{\|v\|_{E}}{s}\left(\left\|T x_{v}\right\|_{F}+\|T b\|_{F}\right) \stackrel{(*)}{\leq} \frac{2 n}{s}\|v\|_{E} .
$$

Insgesamt folgt $\sup _{T \in H}\|T\| \leq \frac{2 n}{s}<\infty$. \(\square\)

\subsection*{Korollar}
Sei $E$ ein Banachraum, und sei $\left(T_{n}\right)_{n}$ eine Folge in $\mathcal{L}(E, F)$.\\
(i) Gilt $\limsup _{n \rightarrow \infty}\left\|T_{n}\right\|=\infty$, so existiert $x \in E$ derart, dass die Folge $\left(T_{n} x\right)_{n}$ in $F$ nicht konvergiert.\\
(ii) Existiert $T x:=\lim _{n \rightarrow \infty} T_{n} x \in F$ für alle $x \in E$, so gilt $T \in \mathcal{L}(E, F)$ für die so definierte (offensichtlich lineare) Abbildung $x \mapsto T x$.

Beweis. Wenn $\lim _{n \rightarrow \infty} T_{n} x \in F$ für alle $x \in E$ existiert, so ist die Menge $H:=\left\{T_{n}:\right. n \in \mathbb{N}\} \subset \mathcal{L}(E, F)$ punktweise beschränkt im Sinne von 5.2. Der Satz liefert dann ein $c>0$ mit $\left\|T_{n}\right\| \leq c$ für alle $n \in \mathbb{N}$. Damit folgt (i). Für $T$ wie in (ii) folgt ferner

$$
\|T x\|=\lim _{n \rightarrow \infty}\left\|T_{n} x\right\| \leq c\|x\| \quad \text { für alle } x \in E \text {, }
$$

und somit ist $T$ (als offensichtlich lineare Abbildung) stetig. \(\square\)

\subsection*{Definition und Bemerkung}
Seien $X, Y$ metrische Räume. Eine Abbildung $f: X \rightarrow Y$ heißt offen, falls gilt:

$$
A \subset X \text { offen } \quad \Longrightarrow \quad f(A) \subset Y \text { offen }
$$

Man beachte: Ist $f: X \rightarrow Y$ offen und bijektiv, so ist $f^{-1}$ stetig (vgl. 2.7i))

\subsection*{Hauptsatz (von der offenen Abbildung)}
Sei $E$ ein Banachraum, und sei $T \in \mathcal{L}(E, F)$ derart, dass Bild $T$ nicht mager in $F$ ist. Dann ist $T$ surjektiv und offen.

Beweis. Folgt später. \(\square\)

\subsection*{Beispiel}
Sei $1 \leq p<q \leq \infty$. Nach 1.9 ist die Inklusion $i:\left(\ell^{p},\|\cdot\|_{p}\right) \hookrightarrow\left(\ell^{q},\|\cdot\|_{q}\right)$ stetig. Da $\ell^{p}$ ein Banachraum und $i$ nicht surjektiv ist, folgt mit 5.5, dass $\ell^{p}$ in ( $\ell^{q}\|\cdot\|_{q}$ ) mager ist.

\subsection*{Korollar}
Seien $E, F$ Banachräume und $T \in \mathcal{L}(E, F)$. Dann gilt:\\
i) Ist $T$ surjektiv, so ist $T$ offen.\\
ii) Ist $T$ bijektiv, so ist $T$ ein topologischer Isomorphismus.

Beweis. i) Nach Voraussetzung und 3.12 ist Bild $T=F$ ist nicht mager in $F$. Also ist $T$ offen nach 5.5.\\
ii) Aus i) und 5.4 folgt die Stetigkeit von $T^{-1}$.

\subsection*{Korollar}
Seien $\|\cdot\|_{1},\|\cdot\|_{2}$ Normen auf einem $\mathbb{K}$-Vektorraum $V$ derart, dass $\left(V,\|\cdot\|_{1}\right),\left(V,\|\cdot\|_{2}\right)$ Banachräume sind. Ferner möge $C>0$ existieren mit $\|x\|_{1} \leq C\|x\|_{2}$ für alle $x \in V$. Dann sind $\|\cdot\|_{1}$ und $\|\cdot\|_{2}$ äquivalent.

Beweis. Nach Voraussetzung ist id : $\left(V,\|\cdot\|_{2}\right) \rightarrow\left(V,\|\cdot\|_{1}\right)$ stetig, also ein topologischer Isomorphismus nach 5.7i). Insbesondere ist id : $\left(V,\|\cdot\|_{1}\right) \rightarrow\left(V,\|\cdot\|_{2}\right)$ stetig, d.h. es existiert $D>0$ mit $\|x\|_{2} \leq D\|x\|_{1}$ für alle $x \in V$. Die Behauptung folgt.

\subsection*{Definition und Bemerkung}
Auf dem $\mathbb{K}$-Vektorraum $E \times F$ sei die Norm $\|\cdot\|$ definiert durch $\|(x, y)\|:=\|x\|_{E}+ \|y\|_{F}$. Dann gilt:

$$
(E \times F,\|\cdot\|) \text { ist Banachraum } \quad \Longleftrightarrow \quad E \text { und } F \text { sind Banachräume }
$$

\subsection*{Satz (vom abgeschlossenen Graphen)}
Seien $E, F$ Banachräume, und sei $T: E \rightarrow F$ eine lineare Abbildung, deren Graph $\Gamma:=\{(x, y) \in E \times F: y=T x\}$ in $E \times F$ abgeschlossen ist. Dann ist $T \in \mathcal{L}(E, F)$, also stetig.

Beweis. Nach Voraussetzung und 5.9 ist ( $E \times F,\|\cdot\|$ ) ein Banachraum und somit auch $(\Gamma,\|\cdot\|)$, da $\Gamma \subset E \times F$ abgeschlossen ist. Sei $G: \Gamma \rightarrow E$ definiert durch $G(x, y)=x$. Da $\|G(x, y)\|_{E}=\|x\|_{E} \leq\|(x, y)\|$ für $(x, y) \in \Gamma$, ist $G \in \mathcal{L}(\Gamma, E)$. Ferner ist $G$ bijektiv mit $G^{-1}(x)=(x, T x)$ für $x \in E$. Gemäß 5.7 ii) ist nun $G^{-1} \in \mathcal{L}(E, \Gamma)$; also existiert $C>0$ mit

$$
\|x\|_{E}+\|T x\|_{F}=\left\|G^{-1}(x)\right\| \leq C\|x\|_{E} \quad \text { für alle } x \in E,
$$

und somit $\|T x\|_{F} \leq C\|x\|_{E}$ für alle $x \in E$. Es folgt die Stetigkeit von $T$.

\subsection*{Bemerkung}
Seien $E, F$ normierte Räume über $\mathbb{K}$, und sei $T: E \rightarrow F \mathbb{K}$-linear. Offensichtlich sind dann folgende Eigenschaften äquivalent:

\begin{itemize}
  \item Der Graph von $T$ ist abgeschlossen.
  \item Für jede Folge $\left(x_{n}\right)_{n}$ in $E$, für die Grenzwerte
\end{itemize}

$$
x:=\lim _{n \rightarrow \infty} x_{n} \in E \quad \text { und } \quad y:=\lim _{n \rightarrow \infty} T x_{n} \in F
$$

existieren, gilt $y=T x$.

\begin{itemize}
  \item Für jede Nullfolge $\left(x_{n}\right)_{n}$ in $E$, für die der Grenzwert $y:=\lim _{n \rightarrow \infty} T x_{n} \in F$ existiert, gilt $y=0$.
\end{itemize}

\section*{Stetige Projektionen und topologische Zerlegungen}
Sei stets $(E,\|\cdot\|)$ ein normierter Raum

\subsection*{Definition und Bemerkung}
i). Seien $V, W$ Teilräume von $E$ mit $E=V \oplus W$. Aus der linearen Algebra ist bekannt, dass genau eine Projektion $P: E \rightarrow E$ (d. h. $P$ linear und $P^{2}=P$ ) mit $V=\operatorname{Kern} P$ und $W=$ Bild $P$ existiert. Wir nennen $P$ die Projektion auf $W$ längs $V$. Dann ist id $-P: E \rightarrow E$ die Projektion auf $V$ längs $W$. Offensichtlich sind äquivalent:

$$
P \in \mathcal{L}(E) \Longleftrightarrow \mathrm{id}-P \in \mathcal{L}(E)
$$

Gilt dies, so schreiben wir $E=V \bigoplus_{\text {top }} W$ (topologische direkte Summe) und nennen $W$ ein topologisches Komplement von $V$ (bzw. umgekehrt). Insbesondere sind in diesem Fall $V$ und $W$ abgeschlossen in $E$.\\
ii). Ein Unterraum $V \subset E$ heißt stetig projiziert, falls es eine Projektion $P \in \mathcal{L}(E)$ mit $V=$ Bild $P$ gibt, d. h. falls $V$ ein topologisches Komplement in $E$ besitzt.

\subsection*{Beispiel}
i). Sei $E=\left(\mathbb{R}^{2},|\cdot|_{\infty}\right), V=\operatorname{span}\{(1,0)\}$ und $W_{n}=\operatorname{span}\left\{\left(1, \frac{1}{n}\right)\right\}$ für $n \in \mathbb{N}$. Dann ist $E=V \oplus W_{n}$. Sei $P_{n}$ die Projektion mit Bild $P=W_{n}$ und $\operatorname{Kern} P=V$. Da $(0,1)=-(n, 0)+(n, 1) \in V \oplus W_{n}$, ist $P_{n}(0,1)=(n, 1)$, also $\left|P_{n}(0,1)\right|_{\infty}= n=n|(0,1)|_{\infty}$. Es folgt $\left\|P_{n}\right\| \geq n$.\\
ii). Sei nun $E=\mathcal{F}$ der Raum der finiten Folgen mit Norm $\|\cdot\|_{\infty}$. Dann ist $E=V \oplus W$ mit\\
$V=\left\{x=\left(x_{j}\right)_{j} \in \mathcal{F}: x_{j}=0\right.$ für $j$ gerade $\} \quad$ und $\quad W=\operatorname{span}\left\{d_{n} \mid n \in \mathbb{N}\right\}$, wobei $d_{n}=\left(0, \ldots, 0,1, \frac{1}{n}, 0, \ldots\right)=e_{2 n-1}+\frac{1}{n} e_{2 n}$ für $n \in \mathbb{N}$ sei. Sei $P$ die Projektion auf $W$ längs $V$. Wie in i) sieht man

$$
e_{2 n}=-n e_{2 n-1}+n d_{n} \in V \oplus W \quad \text { für alle } n \in \mathbb{N},
$$

also $\left\|P e_{2 n}\right\|_{\infty}=\left\|-n d_{n}\right\|_{\infty}=n=n\left\|e_{2 n}\right\|_{\infty}$. Es folgt, dass $P$ nicht stetig ist.

\subsection*{Satz}
Seien $V, W$ Teilräume von $E$ mit $E=V \oplus W$. Dann gilt:\\
i). Ist $V$ endlichdimensional und $W$ in $E$ abgeschlossen, so ist $E=V \oplus_{\text {top }} W$.\\
ii). Ist $E$ ein Banachraum und sind $V$ und $W$ abgeschlossen in $E$, so ist $E= V \oplus_{\text {top }} W$.

Beweis. Sei $P: E \rightarrow E$ die Projektion mit Bild $P=V$ und $\operatorname{Kern} P=W$.\\
Zu i): Angenommen, $P$ ist nicht stetig. Dann existiert eine Folge $\left(x_{n}\right)_{n} \subset E \backslash\{0\}$ mit $\left\|P x_{n}\right\| \geq n\left\|x_{n}\right\|$ für alle $n \in \mathbb{N}$. Setze $y_{n}:=\frac{x_{n}}{\left\|P x_{n}\right\|}$. Dann gilt $y_{n} \rightarrow 0$ für $n \rightarrow \infty$, aber $P y_{n} \in S_{V}:=\{z \in V \mid\|z\|=1\}$ für alle $n$. Da $\operatorname{dim} V<\infty$, ist $S_{V}$ kompakt. Also existiert $z \in S_{V}$ und eine Teilfolge $\left(y_{n_{\ell}}\right)_{\ell}$ mit $\lim _{\ell \rightarrow \infty} P y_{n_{\ell}}=z$. Somit ist

$$
-z=\lim _{\ell \rightarrow \infty} \underbrace{(\mathrm{id}-P) y_{n_{\ell}}}_{\in W} \in W,
$$

da $W$ abgeschlossen ist. Es folgt $z \in S_{V} \cap W$, im Widerspruch zu $V \cap W=\{0\}$. Der Widerspruch zeigt die Stetigkeit von $P$.\\
ii) Nach 5.10 reicht es, zu zeigen, dass $P$ einen abgeschlossenen Graphen besitzt. Zum Beweis verwenden wir 5.11. Sei also $\left(x_{n}\right)_{n}$ eine Nullfolge in $E$ derart, dass $y:=\lim _{n \rightarrow \infty} P x_{n} \in E$ existiert. Da $V=$ Bild $P$ in $E$ abgeschlossen ist, folgt $y \in V$. Da ferner $W=\operatorname{Bild}(\mathrm{id}-P)$ in $E$ abgeschlossen ist, folgt $y=-\lim _{n \rightarrow \infty}(\mathrm{id}-P) x_{n} \in W$. Somit ist $y \in V \cap W=\{0\}$, d.h. $y=0$. Gemäß 5.11 folgt die Behauptung.

Zum Abschluss des Kapitels beweisen wir nun den Satz von der offenen Abbildung 5.5. Dazu vereinbaren wir zunächst ein paar praktische Bezeichnungen.

\subsection*{Definition und Bemerkung}
Für $M \subset E, \alpha \in \mathbb{K}$ und $v \in E$ sei

$$
\alpha M:=\{\alpha x: x \in M\} \subset E \quad \text { und } \quad v+M:=\{v+x: x \in M\} \subset E .
$$

Da für $\alpha \neq 0$ die Abbildungen

$$
E \rightarrow E, \quad x \mapsto \alpha x \quad \text { und } \quad x \mapsto v+x
$$

Homöomorphismen sind, gilt:

\begin{itemize}
  \item $\overline{\alpha M}=\alpha \bar{M}$ und $(\alpha \stackrel{\circ}{M})=\alpha \stackrel{\circ}{M}$, falls $\alpha \neq 0$;
  \item $\overline{v+M}=v+\bar{M}$ und $(v+\stackrel{\circ}{M})=v+\stackrel{\circ}{M}$.
\end{itemize}

\section*{Beweis von 5.5 .}
Sei $U:=U_{1}(0) \subset E$ und $V:=U_{1}(0) \subset F$. Dann ist $E=\bigcup_{n \in \mathbb{N}} n U$, also

$$
\text { Bild } T=\bigcup_{n \in \mathbb{N}} T(n U)=\bigcup_{n \in \mathbb{N}} n T(U)
$$

Da Bild $T$ nicht mager ist, existiert $n \in \mathbb{N}$ derart, dass die Menge $\frac{{ }^{\circ}}{n T(U)} \stackrel{\text { 5.15 }}{=} n \stackrel{\circ}{T(U)}$ nichtleer ist. Also existiert $y \in F$ und $\rho>0$ mit $U_{\rho}(y) \subset \overline{T(U)}$.\\
Beh. 1: $\rho V \subset \overline{T(U)}$.\\
Sei dazu $z \in \rho V$, d.h. $z \in F$ und $\|z\|<\rho$. Dann ist $y \pm z \in U_{\rho}(y) \subset \overline{T(U)}$, also $y \pm z=\lim _{n \rightarrow \infty} T x_{n}^{ \pm}$für geeignete Folgen $\left(x_{n}^{+}\right)_{n},\left(x_{n}^{-}\right)_{n}$ in $U$. Dies liefert $\frac{1}{2}\left(x_{n}^{+}-x_{n}^{-}\right) \in U$ für alle $n$, da $\left\|\frac{1}{2}\left(x_{n}^{+}-x_{n}^{-}\right)\right\| \leq \frac{1}{2}\left(\left\|x_{n}^{+}\right\|+\left\|x_{n}^{-}\right\|\right)<1$. Es folgt

$$
z=\frac{1}{2}((y+z)-(y-z))=\lim _{n \rightarrow \infty} T\left(\frac{1}{2}\left[x_{n}^{+}-x_{n}^{-}\right]\right) \in \overline{T(U)} .
$$

Beh. 2: Für alle $z \in F \backslash\{0\}$ und $\varepsilon>0$ existiert ein $x \in E$ mit $\|T x-z\|<\varepsilon$ und $\|x\|<\frac{\|z\|}{\rho}$.\\
Ist nämlich $z \in F \backslash\{0\}$, so ist $\frac{\rho}{\|z\|} z \in \overline{\rho V} \subset \overline{T(U)}$. Für jedes $\varepsilon>0$ existiert also ein $\tilde{x} \in U$ mit $\left\|T \tilde{x}-\frac{\rho}{\|z\|} z\right\|<\frac{\rho}{\|z\|} \varepsilon$. Für $x:=\frac{\|z\|}{\rho} \tilde{x}$ gilt also $\|T x-z\|<\varepsilon$ und $\|x\|<\frac{\|z\|}{\rho}$.\\
Beh. 3: $\rho V \subset T(U)$.\\
Sei dazu $z \in \rho V \backslash\{0\}$, und sei $\left(\varepsilon_{n}\right)_{n} \subset(0, \infty)$ eine Folge mit $\sum_{n \in \mathbb{N}} \varepsilon_{n}=\rho-\|z\|$.

Nach Beh. 2 existiert $x_{1} \in E$ mit $\left\|T x_{1}-z\right\|<\varepsilon_{1}$ und $\left\|x_{1}\right\|<\frac{\|z\|}{\rho}$. Wir konstruieren nun rekursiv eine Folge $\left(x_{n}\right)_{n \geq 2} \subset E$ mit

$$
\left\|\sum_{k=1}^{n} T x_{k}-z\right\|<\varepsilon_{n} \quad \text { und } \quad\left\|x_{n}\right\| \leq \frac{1}{\rho}\left\|\sum_{k=1}^{n-1} T x_{k}-z\right\|
$$

Seien dazu $x_{1}, \ldots, x_{n-1}$ bereits konstruiert, und sei $z_{n}:=z-\sum_{k=1}^{n-1} T x_{k}$.

\begin{enumerate}
  \item Fall: $z_{n}=0$. Dann setzen wir $x_{n}=0$ und erhalten ( $*$ ).
  \item Fall: $z_{n} \neq 0$. Dann existiert nach Behauptung 2 ein $x_{n} \in E$ mit $\left\|x_{n}\right\|<\frac{\left\|z_{n}\right\|}{\rho}$ und
\end{enumerate}

$$
\varepsilon_{n}>\left\|T x_{n}-z_{n}\right\|=\left\|\sum_{k=1}^{n} T x_{k}-z\right\|
$$

Wiederum gilt dann (*).\\
Nach Konstruktion ist nun

$$
\sum_{n \in \mathbb{N}}\left\|x_{n}\right\|<\frac{\|z\|}{\rho}+\sum_{n \geq 2} \frac{1}{\rho} \overbrace{\left\|\sum_{k=1}^{n-1} T x_{k}-z\right\|}^{<\varepsilon_{n}-1}<\frac{1}{\rho}\left(\|z\|+\sum_{n \geq 1} \varepsilon_{n}\right)=\frac{1}{\rho}(\|z\|+\rho-\|z\|)=1 .
$$

Da $E$ nach Voraussetzung ein Banachraum ist, existiert gemäß 3.16 nun $x:= \sum_{k=1}^{\infty} x_{k} \in E$, wobei

$$
\|x\|=\lim _{n \rightarrow \infty}\left\|\sum_{k=1}^{n} x_{k}\right\| \leq \lim _{n \rightarrow \infty} \sum_{k=1}^{n}\left\|x_{k}\right\|=\sum_{k \in \mathbb{N}}\left\|x_{k}\right\|<1
$$

und somit $x \in U$ ist. Ferner gilt

$$
\|T x-z\|=\lim _{n \rightarrow \infty}\left\|T\left(\sum_{k=1}^{n} x_{k}\right)-z\right\| \stackrel{(*)}{\leq} \lim _{n \rightarrow \infty} \varepsilon_{n}=0
$$

und somit $z=T x \in T(U)$. Es folgt Beh. 3.\\
Aus Beh. 3 folgt

$$
F=\bigcup_{n=1}^{\infty} n \rho V \subset \bigcup_{n=1}^{\infty} n T(U)=\bigcup_{n=1}^{\infty} T(n U)=\text { Bild } T
$$

und somit ist $T$ surjektiv. Wir zeigen schließlich die Offenheit der Abbildung $T$. Sei dazu $A \subset E$ offen und $y \in T(A)$, also $y=T x$ mit $x \in A$. Dann existiert $\delta>0$ mit $U_{\delta}(x) \subset A$, also

$$
T(A) \supset T\left(U_{\delta}(x)\right)=T(x+\delta U)=T x+\delta T(U) \supset y+\delta \rho V=U_{\delta \rho}(y)
$$

Es folgt, dass $T(A) \subset F$ offen ist.

\section*{§6 $L^{p}$-Räume}
Stets sei im Folgenden ( $\Omega, \mathcal{M}, \mu$ ) ein Maßraum. Ferner verwenden wir die Vereinbarung $\frac{r}{\infty}=0$ für $r \in \mathbb{R}$. Wir wiederholen zunächst zentrale Begriffe der Maß- und Integrationstheorie und Konvergenzsätze, welche größtenteils aus der Integrationstheorie bekannt sind.

\section*{Vorbereitender Exkurs über Konvergenzsätze der Integrationstheorie}
\subsubsection*{Definition}
Sei $(Y, d)$ ein metrischer Raum und $f: \Omega \rightarrow Y$ eine Funktion.

\begin{itemize}
  \item $f$ heißt $\mu$-messbar (oder einfach messbar), wenn für jede offene Teilmenge $U \subset Y$ die Menge $f^{-1}(U) \subset \Omega \mu$-messbar ist.
  \item Ist $f \mu$-messbar und nimmt $f$ nur endlich viele Werte an, so nennt man $f$ eine Stufenfunktion (oder Elementarfunktion).
\end{itemize}

Wir schreiben:

$$
\begin{aligned}
\mathcal{M}(\Omega, Y) & :=\{f: \Omega \rightarrow Y: f \mu \text {-messbar }\} \\
\mathcal{E}(\Omega, Y) & :=\{f: \Omega \rightarrow Y: f \text { Stufenfunktion }\} .
\end{aligned}
$$

$\operatorname{Im}$ Fall $Y=\mathbb{K}(\mathbb{K}=\mathbb{R}$ oder $\mathbb{K}=\mathbb{C})$ schreiben wir auch $\mathcal{M}(\Omega)$ bzw. $\mathcal{E}(\Omega)$. Im Fall $Y=\mathbb{R}$ gilt insbesondere

$$
f \in \mathcal{M}(\Omega, \mathbb{R}) \Longleftrightarrow\{x \in \Omega: f(x)>c\} \text { messbar für alle } c \in \mathbb{R} .
$$

Wir werden im Folgenden (nicht-endliche) Funktionen $f: \Omega \rightarrow \overline{\mathbb{R}} \mu$-messbar nennen, wenn sie die Bedingung auf der rechten Seite von (6.1) erfüllen. Auch für solche Funktionen schreiben wir dann $f \in \mathcal{M}(\Omega, \mathbb{R})$ bzw. $f \in \mathcal{M}(\Omega)$.

\subsubsection*{Bemerkung}
Für $f \in \mathcal{E}(\Omega)$ gilt $f=\sum_{y \in \mathbb{K}} y 1_{\{f=y\}}$, wobei wir hier und im Folgenden die Kurzschreibweise $\{f=y\}:=f^{-1}(y)=\{x \in \Omega: f(x)=y\} \subset \Omega$ verwenden.

\subsubsection*{Definition und Bemerkung}
(i) Für eine $\mu$-messbare Funktion $f: \Omega \rightarrow \overline{\mathbb{R}}$ ist offenbar äquivalent:

$$
\int_{\Omega} f^{ \pm} d \mu<\infty \quad \Longleftrightarrow \quad \int_{\Omega}|f| d \mu<\infty
$$

Hier seien $f^{+}:=\max \{f, 0\}$ bzw. $f^{-}:=-\min \{f, 0\}$ der Positiv- bzw. Negativteil von $f$. Die Funktion $f$ heißt $\mu$-integrierbar, falls die Bedingungen in (6.2) gelten. In diesem Fall setzt man

$$
\int_{\Omega} f d \mu:=\int_{\Omega} f^{+} d \mu-\int_{\Omega} f^{-} d \mu \quad \in \mathbb{R}
$$

(ii) Eine $\mu$-messbare Funktion $f: \Omega \rightarrow \mathbb{C}$ heißt $\mu$-integrierbar, wenn $\operatorname{Re} f$ und $\operatorname{Im} f \mu$-integrierbar sind. In diesem Fall setzt man

$$
\int_{\Omega} f d \mu:=\int_{\Omega} \operatorname{Re} f d \mu+i \int_{\Omega} \operatorname{Im} f d \mu
$$

(iii) Die Menge der $\mu$-integrierbaren $\mathbb{K}$-wertigen Funktionen bezeichnen wir mit $\mathfrak{L}^{1}(\Omega)$ oder $\mathfrak{L}^{1}(\Omega, d \mu)$.

\subsubsection*{Satz}
(i) Ist $f: \Omega \rightarrow[0, \infty]$ messbar, so existiert eine Folge $\left(f_{k}\right)_{k}$ in $\mathcal{E}(\Omega)$ mit

$$
0 \leq f_{k}(x) \leq f_{k+1}(x) \leq f(x) \quad \text { für alle } x \in \Omega, k \in \mathbb{N}
$$

und $\lim _{k \rightarrow \infty} f_{k}(x)=f(x)$ für alle $x \in \Omega$.\\
(ii) Ist $f \in \mathcal{M}(\Omega, \mathbb{C})$, so existiert eine Folge in $\mathcal{E}(\Omega, \mathbb{C})$ mit $\left|f_{k}\right| \leq|f|$ in $\Omega$ für alle $k$ und $f_{k} \rightarrow f$ für $k \rightarrow \infty$ punktweise in $\Omega$.

Bemerkung: Ist $\Omega=\bigcup_{k \in \mathbb{N}} M_{k}$ mit $\mu$-messbaren Teilmengen $M_{k}$ und $\mu\left(M_{k}\right)<\infty$ für alle $k \in \mathbb{N}$, so können die Folgen $\left(f_{k}\right)_{k}$ in (i) und (ii) derart gewählt werden, dass gilt:

$$
\mu\left(\left\{f_{k} \neq 0\right\}\right)<\infty \quad \text { für alle } k \in \mathbb{N} .
$$

Dies gilt insbesondere im Fall $\Omega=\mathbb{R}^{N}$ und $\mu=\lambda$ (Lebesguemaß).

\subsubsection*{Satz von der monotonen Konvergenz}
i). Seien $f_{k} \in \mathcal{M}(\Omega, \mathbb{R}), k \in \mathbb{N}$ nichtnegative Funktionen mit $f_{k} \leq f_{k+1}$ f. ü. auf $\Omega$ für alle $k$, und sei $f \in \mathcal{M}(\Omega, \mathbb{R})$ f. ü. definiert durch $f(x):=\lim _{k \rightarrow \infty} f_{k}(x)$. Dann gilt: $\int_{\Omega} f d \mu=\lim _{k \rightarrow \infty} \int_{\Omega} f_{k} d \mu$\\
ii). Seien $f_{k} \in \mathfrak{L}^{1}(\Omega) \mathbb{R}$-wertige Funktionen für $k \in \mathbb{N}$ mit $f_{k} \leq f_{k+1}$ f. ü. auf $\Omega$ für alle $k$, und sei $f \in \mathcal{M}(\Omega, \mathbb{R})$ f. ü. definiert durch $f(x):=\lim _{k \rightarrow \infty} f_{k}(x)$. Dann gilt: Ist die Folge der Integrale $\int_{\Omega} f_{k} d \mu$ beschränkt, so ist $f$ integrierbar und $\int_{\Omega} f d \mu=\lim _{k \rightarrow \infty} \int_{\Omega} f_{k} d \mu$.

\subsubsection*{Lemma von Fatou}
Sei $\left(f_{k}\right)_{k} \subset \mathfrak{L}^{1}(\Omega)$ eine Folge $\mathbb{R}$-wertiger Funktionen.\\
a) Ist $\liminf _{k \rightarrow \infty} \int_{\Omega} f_{k} d \mu<\infty$ und existiert $g \in \mathfrak{L}^{1}(\Omega, \mathbb{R})$ mit $f_{k} \geq g$ f. ü. auf $\Omega$ für alle $k \in \mathbb{N}$, so ist $f:=\liminf _{k \rightarrow \infty} f_{k}$ integrierbar und $\int_{\Omega} f d \mu \leq \liminf _{k \rightarrow \infty} \int_{\Omega} f_{k} d \mu$.\\
b) Ist $\limsup _{k \rightarrow \infty} \int_{\Omega} f_{k} d \mu>-\infty$ und existiert $g \in \mathfrak{L}^{1}(\Omega, \mathbb{R})$ mit $f_{k} \leq g$ f. ü. auf $\Omega$ für alle $k \in \mathbb{N}$, so ist $f:=\limsup _{k \rightarrow \infty} f_{k}$ integrierbar und $\int_{\Omega} f d \mu \geq \limsup _{k \rightarrow \infty} \int_{\Omega} f_{k} d \mu$.

\subsubsection*{Satz von der majorisierten Konvergenz (Satz von Lebesgue)}
Seien $f_{k} \in \mathfrak{L}^{1}(\Omega), k \in \mathbb{N}$ derart, dass der punktweise Grenzwert $f:=\lim _{k \rightarrow \infty} f_{k}$ punktweise auf $\Omega$ existiert. Ferner existiere eine $\mathbb{R}$-wertige Funktion $g \in \mathfrak{L}^{1}(\Omega)$ derart, dass für alle $k \in \mathbb{N}$ die Ungleichung $\left|f_{k}\right| \leq g$ f. ü. auf $\Omega$ gilt.\\
Dann ist $f \in \mathfrak{L}^{1}(\Omega)$, und es gilt $\int_{\Omega} f d \mu=\lim _{k \rightarrow \infty} \int_{\Omega} f_{k} d \mu$.

\section*{Ende des Exkurses}
\subsection*{Definition und Bemerkung}
Für eine $\mu$-messbare Funktion $u: \Omega \rightarrow \overline{\mathbb{R}}$ sei

$$
\underset{\Omega}{\operatorname{esssup}} u:=\inf \{t \in \mathbb{R}: u \leq t \text { f.ü. auf } \Omega\} \quad \text { (wesentliches Supremum von } u \text { ) }
$$

Man sieht leicht, dass stets $u \leq \operatorname{esssup} u$ f.ü. auf $\Omega$ gilt.

\subsection*{Satz und Definition}
Sei $1 \leq p \leq \infty$. Die Menge $\mathfrak{L}^{p}(\Omega)$ aller messbaren Funktionen $u: \Omega \rightarrow \mathbb{K}$ mit

$$
\begin{cases}\int_{\Omega}|u|^{p} d \mu<\infty & \text { im Fall } p<\infty \\ \underset{\Omega}{\operatorname{essup}}|u|<\infty & \text { im Fall } p=\infty\end{cases}
$$

ist ein $\mathbb{K}$-Vektorraum, und durch

$$
\|u\|_{p}:= \begin{cases}\left(\int_{\Omega}|u|^{p} d \mu\right)^{\frac{1}{p}}, & \text { im Fall } p<\infty \\ \underset{\Omega}{\operatorname{esssup}}, & \text { im Fall } p=\infty\end{cases}
$$

wird eine Halbnorm auf $\mathfrak{L}^{p}(\Omega)$ definiert.\\
Beweis. Die Vektorraumeigenschaften sieht man sehr ähnlich wie in 1.7. Sei nun

$$
A_{p}:= \begin{cases}\left\{u \in \mathfrak{L}^{p}(\Omega): \int_{\Omega}|u|^{p} d \mu \leq 1\right\} & \text { im Fall } p<\infty \\ \left\{u \in \mathfrak{L}^{p}(\Omega): \operatorname{esssup}_{\Omega}|u| \leq 1\right\} & \text { im Fall } p=\infty\end{cases}
$$

Offensichtlich ist $A_{p}$ dann absolut konvex und absorbierend in $\mathfrak{L}^{p}(\Omega)$, und gemäß 1.6 definiert

$$
u \mapsto \inf \left\{t>0: \frac{u}{t} \in A_{p}\right\}=\|u\|_{p} .
$$

eine Halbnorm auf $\mathfrak{L}^{p}(\Omega)$.

\subsection*{Satz (Höldersche Ungleichung)}
Seien $p, q \in[1, \infty]$ konjugierte Exponenten, d.h. es gilt $\frac{1}{p}+\frac{1}{q}=1$. Seien ferner $u \in \mathfrak{L}^{p}(\Omega)$ und $v \in \mathfrak{L}^{q}(\Omega)$. Dann gilt

$$
\int_{\Omega}|u v| d \mu \leq\|u\|_{p}\|v\|_{q}
$$

Beweis. Die Ungleichung ist klar im Fall $p \in\{1, \infty\}$. Sei also nun $1<p<\infty$, also auch $1<q<\infty$. Ohne Einschränkung können wir $\|u\|_{p},\|v\|_{q}>0$ annehmen, denn andenfalls ist $u v=0$ f.ü. in $\Omega$ und die Ungleichung ist trivial erfüllt. Mit der Youngschen Ungleichung (vgl. 1.8(ii)) folgt dann mit $s=\|u\|_{p}, t=\|v\|_{q}$ :

$$
\begin{aligned}
\int_{\Omega}|u v| d \mu & =s t \int_{\Omega}\left|\frac{u}{s} \| \frac{v}{t}\right| d \mu \leq s t \int_{\Omega}\left(\frac{1}{p}\left|\frac{u}{s}\right|^{p}+\frac{1}{q}\left|\frac{v}{t}\right|^{q}\right) d \mu \\
& =s t\left(\frac{1}{p}\left\|\frac{u}{s}\right\|_{p}^{p}+\frac{1}{q}\left\|\frac{v}{t}\right\|_{q}^{q}\right)=s t\left(\frac{1}{p}+\frac{1}{q}\right)=s t
\end{aligned}
$$

wie behauptet.

\subsection*{Bemerkung (Verallgemeinerte Höldersche Ungleichung)}
Seien $p_{1}, \ldots, p_{n} \in[1, \infty]$ mit $\sum_{i=1}^{n} \frac{1}{p_{i}}=\frac{1}{p}$ für ein $p \in[1, \infty]$. Seien ferner $u_{i} \in \mathfrak{L}^{p_{i}}(\Omega)$, $i=1, \ldots, n$ gegeben. Dann ist

$$
u:=u_{1} \cdots \cdots u_{n} \in \mathfrak{L}^{p}(\Omega) \quad \text { und } \quad\|u\|_{p} \leq\left\|u_{1}\right\|_{p_{1}} \cdots \cdots\left\|u_{n}\right\|_{p_{n}} .
$$

Dies folgt aus 6.3 per Induktion (Übung, Blatt 5).

\subsection*{Bemerkung und Beispiel}
Sei $1 \leq p<q \leq \infty$.\\
i). Ist $u \in \mathfrak{L}^{p}(\Omega) \cap \mathfrak{L}^{q}(\Omega)$, so ist auch $u \in \mathfrak{L}^{s}(\Omega)$ für $s \in(p, q)$, wobei gilt:

$$
\|u\|_{s} \leq\|u\|_{p}^{\alpha}\|u\|_{q}^{1-\alpha} \quad \text { mit } \alpha:=\frac{p\left(1-\frac{s}{q}\right)}{s\left(1-\frac{p}{q}\right)}
$$

Merkregel: $\frac{1}{s}=\frac{\alpha}{p}+\frac{1-\alpha}{q}$.\\
ii). Ist $\mu(\Omega)<\infty$, so ist $\mathfrak{L}^{q}(\Omega) \subset \mathfrak{L}^{p}(\Omega)$ und

$$
\|u\|_{p} \leq \mu(\Omega)^{\frac{1}{p}-\frac{1}{q}}\|u\|_{q} \quad \text { für alle } u \in \mathfrak{L}^{q}(\Omega)
$$

iii). Ist $\mu=\lambda$ das Lebesguemaß auf $\mathbb{R}^{N}$, so gilt

$$
\mathfrak{L}^{p}\left(\mathbb{R}^{N}\right) \backslash \mathfrak{L}^{q}\left(\mathbb{R}^{N}\right) \neq \varnothing \neq \mathfrak{L}^{q}\left(\mathbb{R}^{N}\right) \backslash \mathfrak{L}^{p}\left(\mathbb{R}^{N}\right)
$$

Beweis. (i) Nach Voraussetzung ist $|u|^{\alpha} \in \mathfrak{L}^{p / \alpha}(\Omega)$ und $|u|^{1-\alpha} \in \mathfrak{L}^{q /(1-\alpha)}(\Omega)$. Ferner gilt

$$
\frac{1}{p / \alpha}+\frac{1}{q /(1-\alpha)}=\frac{\alpha}{p}+\frac{1-\alpha}{q}=\frac{1}{s}
$$

nach Voraussetzung. Gemäß 6.4 ist also $|u|=|u|^{\alpha}|u|^{1-\alpha} \in \mathfrak{L}^{s}(\Omega)$, also

$$
u \in \mathfrak{L}^{s}(\Omega) \quad \text { mit } \quad\|u\|_{s} \leq\left\||u|^{\alpha}\right\|_{\frac{p}{\alpha}}\left\||u|^{1-\alpha}\right\|_{\frac{q}{1-\alpha}}=\|u\|_{p}^{\alpha}\|u\|_{q}^{1-\alpha}
$$

(ii) Sei $u \in \mathfrak{L}^{q}(\Omega)$. Wir schreiben $u=u \cdot 1_{\Omega}$. Mit $\beta=\left(\frac{1}{p}-\frac{1}{q}\right)^{-1}$ ist dann $1_{\Omega} \in \mathfrak{L}^{\beta}$ mit

$$
\left\|1_{\Omega}\right\|_{\beta}=\left(\int_{\Omega} 1 d \mu\right)^{1 / \beta}=[\mu(\Omega)]^{\frac{1}{p}-\frac{1}{q}}
$$

Ferner ist $\frac{1}{q}+\frac{1}{\beta}=\frac{1}{p}$, und somit folgt mit 6.4:

$$
u \in \mathfrak{L}^{p}(\Omega) \quad \text { mit } \quad\|u\|_{p} \leq\left\|1_{\Omega}\right\|_{\beta}\|u\|_{q}=[\mu(\Omega)]^{\frac{1}{p}-\frac{1}{q}}\|u\|_{q}
$$

(iii) Übung, Blatt 5.

\subsection*{Definition und Satz ( $L^{p}$-Raum)}
Sei $Z:=\{u \in \mathcal{M}(\Omega, \mathbb{K}): u=0$ f.ü. auf $\Omega\}$, und sei $1 \leq p \leq \infty$. Dann ist $Z \subset \mathfrak{L}^{p}(\Omega)$ ein Unterraum, und der Faktorraum $L^{p}(\Omega)=L^{p}(\Omega, \mathbb{K}):=\mathfrak{L}^{p}(\Omega) / Z$ ist ein normierter Raum mit Norm $\|\cdot\|_{p}$ gegeben durch

$$
\|u+Z\|_{p}:=\|u\|_{p} \quad \text { für } u \in \mathfrak{L}^{p}(\Omega)
$$

Beweis. 1. Wohldefiniertheit: Ist $u+Z=v+Z$ für $u, v \in \mathfrak{L}^{p}(\Omega)$, so ist $u=v$ f.ü. auf $\Omega$ und somit $\|u\|_{p}=\|v\|_{p}$.\\
2. Definitheit: Ist $0=\|u+Z\|_{p}=\|u\|_{p}$ für ein $u \in \mathfrak{L}^{p}(\Omega)$, so ist $u=0$ f.ü. auf $\Omega$ (dies ist bekannt aus der Integrationstheorie), also $u \in Z$ und damit $u+Z=Z=0_{L^{p}(\Omega)}$. Die restlichen Normeigenschaften von $\|\cdot\|_{p}$ folgen direkt aus 6.2.

\subsection*{Bemerkung (u.a. zur Notation)}
i). Man schreibt anstelle von $u+Z$ oder $[u]$ (Äquivalenzklasse von $u$ ) oft nur $u \in L^{p}(\Omega)$. Man betrachtet die Elemente von $L^{p}(\Omega)$ also als Funktionen, die nur bis auf Abänderung auf Nullmengen eindeutig definiert sind.\\
ii). Angenommen, $\Omega$ ist ein metrischer Raum mit der Eigenschaft

Ist $V \subset \Omega$ eine Nullmenge, so ist $\Omega \backslash V$ dicht in $\Omega$\\
(z.B. $\Omega \subset \mathbb{R}^{N}$ offen mit Lebesgue-Maß). Sind dann $u, v \in \mathfrak{L}^{p}(\Omega)$ stetige Funktionen mit $u=v$ f.ü. auf $\Omega$, so gilt $u=v$ auf ganz $\Omega$. Es folgt also, dass dann die Abbildung

$$
\mathfrak{L}^{p}(\Omega) \cap \mathcal{C}(\Omega) \rightarrow L^{p}(\Omega), \quad u \mapsto u+Z
$$

injektiv ist.

\subsection*{Satz}
Sei $1 \leq p \leq \infty$.\\
i). $L^{p}(\Omega)$ ist ein Banachraum.\\
ii). Seien $u, u_{k} \in L^{p}(\Omega), k \in \mathbb{N}$ gegeben mit $\left\|u_{k}-u\right\|_{p} \rightarrow 0$ für $k \rightarrow \infty$. Dann existiert eine Teilfolge $\left(u_{k_{j}}\right)_{j}$, welche punktweise f.ü. auf $\Omega$ gegen $u$ konvergiert.

Beweis. Gemäß 3.16 reicht es, zu zeigen, dass jede absolut konvergente Reihe in $L^{p}(\Omega)$ auch konvergiert. Seien also $u_{k} \in L^{p}(\Omega), k \in \mathbb{N}$ gegeben mit

$$
C_{m}:=\sum_{k=m}^{\infty}\left\|u_{k}\right\|_{p} \rightarrow 0 \quad \text { für } m \rightarrow \infty .
$$

Wir betrachten im Folgenden die Hilfsfunktionen

$$
h_{m}:=\sum_{k=m}^{\infty}\left|u_{k}\right|: \Omega \rightarrow[0, \infty], \quad m \in \mathbb{N}
$$

Wir zeigen zunächst:

$$
h_{m} \in L^{p}(\Omega) \quad \text { und } \quad\left\|h_{m}\right\|_{p} \leq C_{m} \quad \text { für } m \in \mathbb{N} \text {. }
$$

Für festes $m \in \mathbb{N}$ und $j>m$ liefert die Dreiecksungleichung

$$
\left\|\sum_{k=m}^{j}\left|u_{k}\right|\right\|_{p} \leq \sum_{k=m}^{j}\left\|u_{k}\right\|_{p} \leq C_{m}
$$

Im Fall $p<\infty$ folgt dann aus dem Satz von der monotonen Konvergenz

$$
\int_{\Omega}\left|h_{m}\right|^{p} d \mu=\lim _{j \rightarrow \infty} \int_{\Omega}\left(\sum_{k=m}^{j}\left|u_{k}\right|\right)^{p} d \mu=\lim _{j \rightarrow \infty}\left\|\sum_{k=m}^{j}\left|u_{k}\right|\right\|_{p}^{p} \leq C_{m}^{p}
$$

und somit (6.3). Im Fall $p=\infty$ folgt ebenfalls

$$
h_{m}=\lim _{j \rightarrow \infty} \sum_{k=m}^{j}\left|u_{k}\right| \leq C_{m} \quad \text { f.ü. in } \Omega
$$

und somit wiederum (6.3). Aus (6.3) folgt insbesondere:

$$
h_{1}=\sum_{k=1}^{\infty}\left|u_{k}\right|<\infty \quad \text { f.ü. auf } \Omega .
$$

Somit existiert $u(x):=\sum_{k=1}^{\infty} u_{k}(x) \in \mathbb{K}$ für fast alle $x \in \Omega$. Sei $u$ trivial auf ganz $\Omega$ fortgesetzt; dann ist $u \mu$-messbar und $|u| \leq h_{1}$ f.ü. auf $\Omega$. Gemäß (6.3) ist also $u \in L^{p}(\Omega)$ mit $\|u\|_{p} \leq C_{1}$.\\
Sei schließlich $g_{m}:=u-\sum_{k=1}^{m-1} u_{k}$ für $m \geq 2$. Dann gilt $\left|g_{m}\right| \leq h_{m}$ f. ü. in $\Omega$, und wegen (6.3) folgt

$$
\left\|g_{m}\right\|_{p} \leq C_{m} \rightarrow 0 \quad \text { für } m \rightarrow \infty .
$$

Somit konvergiert die Reihe $\sum_{k=1}^{\infty} u_{k}$ bzgl. $\|\cdot\|_{p}$ gegen $u$.\\
Zu ii): Wähle sukzessive $k_{j} \in \mathbb{N}$ mit $k_{j+1}>k_{j}$ und derart, dass für alle $j \in \mathbb{N}$ gilt:

$$
\left\|v_{j}\right\|_{p} \leq 2^{-j} \quad \text { für } v_{j}:=u_{k_{j+1}}-u_{k_{j}} \in L^{p}(\Omega)
$$

Insbesondere konvergiert die Reihe $\sum_{j \in \mathbb{N}} v_{j}$ absolut in $L^{p}(\Omega)$. Wie im Beweis von i) folgt daher, dass sie sowohl bzgl. $\|\cdot\|_{p}$ als auch punktweise f.ü. auf $\Omega$ gegen eine Funktion $v \in L^{p}(\Omega)$ konvergiert. Nach Voraussetzung gilt aber auch

$$
\left\|u-u_{k_{1}}-\sum_{j=1}^{m} v_{j}\right\|_{p}=\left\|u-u_{k_{m+1}}\right\|_{p} \rightarrow 0 \quad \text { für } m \rightarrow \infty,
$$

und somit folgt $v=u-u_{k_{1}}$. Es gilt also

$$
u_{k_{m}}=\sum_{j=1}^{m-1} v_{j}+u_{k_{1}} \rightarrow v+u_{k_{1}}=u \quad \text { für } m \rightarrow \infty \text { punktweise fast überall auf } \Omega,
$$

d.h. die Teilfolge ( $u_{k_{m}}$ ) hat die gewünschte Eigenschaft.

\subsection*{Beispiel (Wandernder Buckel)}
Der Übergang zu einer Teilfolge in 6.8(ii) ist i.A. notwendig. Um dies einzusehen, betrachten wir $1 \leq p<\infty$ und $f_{n}:[0,1] \rightarrow \mathbb{R}, n \in \mathbb{N}$ definiert durch

$$
f_{n}(t):= \begin{cases}1 & \text { für } t \in\left[k 2^{-\nu},(k+1) 2^{-\nu}\right] \\ 0 & \text { sonst. }\end{cases}
$$

Hier seien $k, \nu \in \mathbb{N}_{0}$ die jeweils eindeutig bestimmten Zahlen mit $n=2^{\nu}+k$ und $k<2^{\nu}$. Dann ist $f_{n} \in L^{p}([0,1])$ für alle $n$ und $\lim _{n \rightarrow \infty}\left\|f_{n}\right\|_{p}=0$; aber in keinem Punkt $x \in[0,1]$ konvergiert die Folge $\left(f_{n}(x)\right)_{n}$ gegen 0.

\subsection*{Definition und Satz}
Auf $L^{2}(\Omega)$ ist ein Skalarprodukt $\langle\cdot, \cdot\rangle$ definiert durch

$$
\langle f, g\rangle=\int_{\Omega} f \bar{g} d \mu \quad \text { für } f, g \in L^{2}(\Omega)
$$

(Im Fall $\mathbb{K}=\mathbb{R}$ erübrigt sich hierbei natürlich die komplexe Konjugation).\\
Dabei ist $\|\cdot\|_{2}$ die von $\langle\cdot, \cdot\rangle$ induzierte Norm. Nach 6.8 ist $L^{2}(\Omega)$ also ein Hilbertraum.

Beweis. Die Skalarprodukteigenschaften rechnet man problemlos nach.

\subsection*{Satz}
Sei $1 \leq p<\infty$, und sei $u \in L^{p}(\Omega)$. Dann existiert eine Folge $\left(u_{k}\right)_{k} \subset \mathcal{E}(\Omega)$ mit $\left|u_{k}\right| \leq|u|$ in $\Omega$ und $\mu\left(\left\{u_{k} \neq 0\right\}\right)<\infty$ für alle $k \in \mathbb{N}$ sowie $\lim _{k \rightarrow \infty}\left\|u_{k}-u\right\|_{p}=0$.

Beweis. Nach 6.0.4 existieren $u_{k} \in \mathcal{E}(\Omega), k \in \mathbb{N}$ mit $\left|u_{k}\right| \leq|u|$ für alle $k$ und $u_{k} \rightarrow u$ punktweise. Für festes $k$ existiert dabei $c=c(k)>0$ mit $\left|u_{k}\right| \geq c$ auf der Menge $\left\{u_{k} \neq 0\right\}$, und somit folgt

$$
\mu\left(\left\{u_{k} \neq 0\right\}\right) c^{p} \leq \int_{\Omega}\left|u_{k}\right|^{p} d \mu \leq \int_{\Omega}|u|^{p} d \mu<\infty
$$

Dies zeigt $\mu\left(\left\{u_{k} \neq 0\right\}\right)<\infty$. Sei nun $g_{k}:=\left|u_{k}-u\right|^{p}$ für $k \in \mathbb{N}$. Dann konvergiert die Folge $\left(g_{k}\right)_{k}$ punktweise gegen 0 , wobei

$$
\left|g_{k}\right| \leq\left(\left|u_{k}\right|+|u|\right)^{p} \leq 2^{p}|u|^{p} \quad \text { auf } \Omega \text { für alle } k \in \mathbb{N} \text { gilt. }
$$

Mit dem Satz von Lebesgue (siehe 6.0.7) folgt daher

$$
\lim _{k \rightarrow \infty} \int_{\Omega} g_{k} d \mu=0
$$

und dies liefert $\lim _{k \rightarrow \infty}\left\|u_{k}-u\right\|_{p}=0$. \(\square\)

Bisher haben wir noch nicht geklärt, unter welchem Umständen die Räume $L^{p}(\Omega)$ separabel sind. Ist darüber hinaus $\Omega$ ein metrischer Raum, so ist es in Anwendung oft sehr nützlich, zu wissen, inwieweit stetige Funktionen in $L^{p}(\Omega)$ dicht liegen. Diese Fragen wollen wir im Folgenden erörtern. Dazu ist der folgende Satz sehr nützlich; des Weiteren benötigen wir einige Sachverhalte über Borel- und Radonmaße.

\subsection*{Satz von Egorov (und Severini)}
Sei $Y$ ein separabler metrischer Raum, und seien $u_{n}: \Omega \rightarrow Y$ messbare Funktionen für $n \in \mathbb{N}$, welche punktweise f.ü. gegen eine Funktion $u: \Omega \rightarrow Y$ konvergieren. Ist $\mu(\Omega)<\infty$, so existiert zu jedem $\varepsilon>0$ eine messbare Teilmenge $A \subset \Omega$ mit $\mu(\Omega \backslash A)<\varepsilon$ und derart, dass $\left(u_{n}\right)_{n}$ auf $A$ gleichmäßig gegen $u$ konvergiert.

Beweis. Sei $\varepsilon>0$, und für $n, k \in \mathbb{N}$ sei

$$
E_{n, k}:=\bigcup_{m=n}^{\infty}\left\{x \in \Omega: d\left(u_{m}(x), u(x)\right) \geq \frac{1}{k}\right\} .
$$

Dann gilt $E_{n+1, k} \subset E_{n, k}$ für alle $n, k \in \mathbb{N}$. Da ferner $\left(u_{n}\right)_{n}$ punktweise f.ü. gegen $u$ konvergiert und $\mu\left(E_{1, k}\right) \leq \mu(\Omega)<\infty$ gilt, folgt

$$
\lim _{n \rightarrow \infty} \mu\left(E_{n, k}\right)=\mu\left(\bigcap_{n=1}^{\infty} E_{n, k}\right)=0 \quad \text { für alle } k \in \mathbb{N} .
$$

Somit existiert zu jedem $k \in \mathbb{N}$ ein $n_{k} \in \mathbb{N}$ mit $\mu\left(E_{n_{k}, k}\right)<\varepsilon 2^{-k}$. Sei nun

$$
A:=\Omega \backslash \bigcup_{k \in \mathbb{N}} E_{n_{k}, k} .
$$

Für $x \in A$ und $k \in \mathbb{N}$ gilt dann $x \notin E_{n_{k}, k}$, also

$$
d\left(u_{m}(x), u(x)\right)<\frac{1}{k} \quad \text { für } m \geq n_{k} .
$$

Es folgt, dass die Folge $\left(u_{n}\right)_{n}$ auf $A$ gleichmäßig gegen $u$ konvergiert. Ferner gilt

$$
\mu(\Omega \backslash A) \leq \sum_{k=1}^{\infty} \mu\left(E_{n_{k}, k}\right)<\varepsilon .
$$ \(\square\)

Bemerkung: Die Separabilität von $Y$ wird im obigen Beweis benötigt, um die Messbarkeit der Mengen $E_{n, k}$ sicherzustellen.

\section*{Exkurs über Borel- und Radonmaße}
Wir benötigen zunächst einige ergänzende Begriffsbildungen zur Kompaktheit. Sei im Folgenden ( $\Omega, d$ ) ein metrischer Raum.

\subsection*{Definition}
Sei $E$ ein normierter Raum und $U \subset E$ eine Teilmenge mit $0 \in U$. Mit $\mathcal{C}_{c}(\Omega, U)$ bezeichnen wir die Menge der stetigen Funktionen $u: \Omega \rightarrow U$ derart, dass der Träger

$$
\operatorname{supp} u:=\overline{\{x \in \Omega: u(x) \neq 0\}} \subset \Omega
$$

kompakt ist. Im Fall $U=E$ ist diese Menge ein (i.A. nicht abgeschlossener !) Unterraum des normierten Raums $\mathcal{C}_{b}(\Omega, E)$.\\
Kurzschreibweise: $\mathcal{C}_{c}(\Omega)$ anstelle von $\mathcal{C}_{c}(\Omega, \mathbb{K})$.

\subsection*{Definition}
$\Omega$ heißt\\
i). lokal kompakt, wenn jeder Punkt $x \in \Omega$ eine kompakte Umgebung besitzt.\\
ii). $\sigma$-kompakt, wenn kompakte Mengen $K_{j} \subset \Omega, j \in \mathbb{N}$ existieren mit $\Omega=\bigcup_{j=1}^{\infty} K_{j}$.

\subsection*{Bemerkungen und Beispiel}
(i) Ist $\Omega$ lokal kompakt und $K \subset \Omega$ kompakt, so besitzt $K$ eine kompakte Umgebung. Ferner existiert in diesem Fall eine Funktion $u \in \mathcal{C}_{c}(\Omega,[0,1])$ mit $u \equiv 1$ in einer offenen Umgebung von $K$.\\
(ii) Ist $\Omega \sigma$-kompakt, so ist $\Omega$ separabel: Sind nämlich $K_{j} \subset \Omega, j \in \mathbb{N}$ kompakt mit $\Omega:=\bigcup_{j=1}^{\infty} K_{j}$, so sind die Mengen $K_{j}$ separabel nach 4.4 und somit auch $\Omega$ (als abzählbare Vereinigung separabler Mengen).\\
(iii) Ein normierter Raum $E$ ist lokal kompakt genau dann, wenn $E$ endlich dimensional ist (vgl. 4.7).\\
(iv) Jeder endlich dimensionale normierte Raum ist $\sigma$-kompakt.\\
(v) Der Raum $\mathcal{F} \subset l^{\infty}$ der finiten $\mathbb{K}$-wertigen Folgen (versehen mit der Norm $\|\cdot\|_{\infty}$ ) ist $\sigma$-kompakt, denn es ist $\mathcal{F}=\bigcup_{j=1}^{\infty} K_{j}$, wobei $K_{j} \subset \mathcal{F}$ die kompakte Teilmenge der Folgen $x=\left(x_{\ell}\right)_{\ell} \in \mathcal{F}$ mit $\|x\|_{\infty} \leq j$ und $x_{\ell}=0$ für $\ell \geq j$ sei. $\mathcal{F}$ ist aber nicht lokal kompakt, da $\operatorname{dim} \mathcal{F}=\infty$ gilt.

\subsection*{Satz}
Sei $\Omega$ lokal kompakt und $\sigma$-kompakt. Dann gilt:\\
(i) Es existiert eine Folge kompakter Teilmengen $M_{j}, j \in \mathbb{N}$ von $\Omega$ derart, dass jede kompakte Menge $K \subset \Omega$ in einer der Mengen $M_{j}$ enthalten ist.\\
(ii) Der Raum ( $\mathcal{C}_{c}(\Omega),\|\cdot\|_{\infty}$ ) ist separabel.

Beweis. Zu (i): Seien $K_{j} \subset \Omega, j \in \mathbb{N}$ kompakt mit $\Omega:=\bigcup_{j=1}^{\infty} K_{j}$. Gemäß 6.15 (i) existiert für jedes $j \in \mathbb{N}$ eine kompakte Umgebung $L_{j}$ von $K_{j}$. Setze $M_{j}:=\bigcup_{i=1}^{j} L_{i}$ für $j \in \mathbb{N}$. Dann bilden die Mengen $U_{j}:=\stackrel{\circ}{M}_{j}, j \in \mathbb{N}$ eine offene Überdeckung von $\Omega$ mit $U_{j} \subset U_{j+1}$ für $j \in \mathbb{N}$. Für jede kompakte Teilmenge $K \subset \Omega$ existiert somit ein $j \in \mathbb{N}$ mit $K \subset U_{j} \subset M_{j}$, wie behauptet.\\
Zu (ii): Sei $M_{j}, j \in \mathbb{N}$ wie in (i), und sei

$$
\mathcal{C}_{j}:=\left\{f \in \mathcal{C}_{c}(\Omega): \operatorname{supp} f \subset M_{j}\right\} \quad \text { für } j \in \mathbb{N} .
$$

Dabei ist $\mathcal{C}_{j}$ also isometrisch isomorph zu einem Unterraum von $\mathcal{C}\left(M_{j}\right)$ vermöge der isometrischen Einbettung $\mathcal{C}_{j} \rightarrow \mathcal{C}\left(M_{j}\right),\left.f \mapsto f\right|_{M_{j}}$. Gemäß 4.28 ist dabei $\mathcal{C}\left(M_{j}\right)$ und somit auch $\mathcal{C}_{j}$ separabel bzgl. $\|\cdot\|_{\infty}$. Ferner ist $\mathcal{C}_{c}(\Omega)=\bigcup_{j \in \mathbb{N}} \mathcal{C}_{j}$ wegen (i), und somit ist $\mathcal{C}_{c}(\Omega)$ ebenfalls separabel.

\subsection*{Definition}
(i) Für ein Mengenteilsystem $\mathcal{O} \subset \mathcal{P}(\Omega)$ definiert man $\langle\mathcal{O}\rangle^{\sigma}$ als den Schnitt aller $\sigma$-Algebren, welche $\mathcal{O}$ enthalten. Man nennt $\langle\mathcal{O}\rangle^{\sigma}$ dann die von $\mathcal{O}$ erzeugte $\sigma$-Algebra.\\
(ii) Ist $\mathcal{O} \subset \mathcal{P}(\Omega)$ das System aller offenen Teilmengen von $\Omega$, so nennt man $\mathcal{B}(\Omega):=\langle\mathcal{O}\rangle^{\sigma}$ die Borelalgebra von $\Omega$, und die Elemente von $\mathcal{B}(\Omega)$ heißen Borelmengen.\\
(iii) Ist $\mu$ ein Maß auf einer $\sigma$-Algebra $\mathcal{M}_{\mu}$ auf $\Omega$ mit $\mathcal{B}(\Omega) \subset \mathcal{M}_{\mu}$, so nennt man $\mu$ ein Borelmaß auf $\Omega$.\\
Vorsicht: In Teilen der Literatur wird bei einem Borelmaß $\mu$ noch zusätzlich die Eigenschaft verlangt, dass $\mu$ lokal endlich ist, d.h. das jeder Punkt in $\Omega$ eine offene Umgebung $U$ besitzt mit $\mu(U)<\infty$ (vgl. Def. 6.19(ii)).

\subsection*{Beispiele und Bemerkung zur Notation}
(i) Bekanntermaßen ist (nach Konstruktion, s. Vorlesung Integrationstheorie) das Lebesguemaß auf $\mathbb{R}^{N}$ ein Borelmaß. Ferner ist z.B. das Diracmaß $\delta_{x}$ zu $x \in \Omega$, gegeben durch

$$
\delta_{x}=\left\{\begin{array}{ll}
1, & x \in A, \\
0, & x \notin A
\end{array} \quad \text { für } A \subset \Omega\right.
$$

ein Borelmaß.\\
(ii) Ist im Folgenden von einem Borelmaß $\mu$ auf $\Omega$ die Rede, so betrachten wir $\mu$ stets auf einer gegebenen, zugehörigen $\sigma$-Algebra $\mathcal{M}_{\mu} \supset \mathcal{B}(\Omega)$, deren Elemente wir im Folgenden $\mu$-messbar nennen.

\subsection*{Definition}
Sei $\mu$ ein Borelmaß auf $\Omega$.\\
(i) $\mu$ heißt von außen regulär, wenn für jede $\mu$-messbare Teilmenge $A \subset \Omega$ gilt:

$$
\mu(A)=\inf \{\mu(U): U \text { offen, } A \subset U\}
$$

(ii) $\mu$ heißt von innen regulär, wenn für jede $\mu$-messbare Teilmenge $A \subset \Omega$ gilt:

$$
\mu(A)=\sup \{\mu(K): K \text { kompakt, } K \subset A\}
$$

(iii) $\mu$ heißt lokal endlich, wenn jeder Punkt $x \in \Omega$ eine offene Umgebung $U$ besitzt mit $\mu(U)<\infty$.\\
(iv) $\mu$ heißt Radon-Maß auf $\Omega$, wenn $\mu$ von innen regulär und lokal endlich ist.

Bemerkung: Ist $\Omega$ lokal kompakt, so ist (iii) offensichtlich äquivalent zur Eigenschaft

$$
\mu(K)<\infty \quad \text { für jede kompakte Teilmenge } K \subset \Omega .
$$

\subsection*{Satz}
Sei $\Omega$ lokal kompakt und $\sigma$-kompakt, und sei $\mu$ ein von außen reguläres und lokal endliches Borelmaß auf $\Omega$. Dann gilt:\\
(i) $\mu$ ist auch von innen regulär und somit ein Radon-Maß.\\
(ii) Für jede $\mu$-messbare Teilmenge $A \subset \Omega$ und jedes $\varepsilon>0$ existiert eine offene Menge $U \subset \Omega$ und eine abgeschlossene Menge $V \subset \Omega$ mit $V \subset A \subset U$ und $\mu(U \backslash V)<\varepsilon$.

Beweis. Sei $\Omega=\bigcup_{j \in \mathbb{N}} K_{j}$ mit kompakten Teilmengen $K_{j} \subset \Omega$. Ohne Einschränkung können wir dabei $K_{j} \subset K_{j+1}$ für $j \in \mathbb{N}$ annehmen. Sei $A \subset \Omega \mu$-messbar. Wir zeigen zunächst (ii) für $A$.

\begin{enumerate}
  \item Behauptung: $Z u$ jedem $\varepsilon>0$ existiert $U \subset \Omega$ offen mit $A \subset U$ und $\mu(U \backslash A)<\frac{\varepsilon}{2}$. Zum Beweis sei $A_{j}:=A \cap K_{j}$ für $j \in \mathbb{N}$. Dann ist $A_{j} \subset K_{j}$ und somit $\mu\left(A_{j}\right) \leq \mu\left(K_{j}\right)<\infty$ gemäß (6.4). Da $\mu$ von außen regulär ist, existieren nun offene Mengen $U_{j} \subset \Omega$ mit $A_{j} \subset U_{j}$ und $\mu\left(U_{j}\right)<\mu\left(A_{j}\right)+\varepsilon 2^{-j-1}$, somit also
\end{enumerate}

$$
\mu\left(U_{j} \backslash A_{j}\right)<\varepsilon 2^{-j-1} \quad \text { für } j \in \mathbb{N} \text {. }
$$

Mit $U:=\bigcup_{j \in \mathbb{N}} U_{j}$ folgt $U \backslash A \subset \bigcup_{j \in \mathbb{N}}\left(U_{j} \backslash A_{j}\right)$ und

$$
\mu(U \backslash A) \leq \sum_{j \in \mathbb{N}} \mu\left(U_{j} \backslash A_{j}\right)<\frac{\varepsilon}{2},
$$

wie behauptet.\\
Anwendung der 1. Behauptung auf die $\mu$-messbare Menge $\Omega \backslash A$ liefert nun auch $W \subset \Omega$ offen mit $\Omega \backslash A \subset W$ und $\mu(W \backslash(\Omega \backslash A))<\frac{\varepsilon}{2}$. Sei $V:=\Omega \backslash W$. Dann ist $V$ abgeschlossen in $\Omega$ mit $V \subset A$; ferner gilt $W \backslash(\Omega \backslash A)=W \cap A=A \backslash V$ und somit

$$
\mu(U \backslash V)=\mu(U \backslash A)+\mu(A \backslash V)<\frac{\varepsilon}{2}+\frac{\varepsilon}{2}=\varepsilon
$$

Somit ist (ii) gezeigt. Nun zu (i). Wir müssen zeigen:

$$
\mu(A) \leq \sup \{\mu(K): K \text { kompakt, } K \subset A\} ;
$$

die umgekehrte Ungleichung ist offensichtlich. Für die oben definierten Mengen $A_{j}:= A \cap K_{j}$ gilt $A_{j} \subset A_{j+1}$ für $j \in \mathbb{N}$ und $A=\bigcup_{j=1}^{\infty} A_{j}$, somit also

$$
\mu(A)=\lim _{j \rightarrow \infty} \mu\left(A_{j}\right)=\sup _{j \in \mathbb{N}} \mu\left(A_{j}\right)
$$

Sei nun $\varepsilon>0$. Anwendung von (ii) liefert abgeschlossene Mengen $V_{j} \subset \Omega, j \in \mathbb{N}$ mit $V_{j} \subset A_{j}$ und $\mu\left(A_{j} \backslash V_{j}\right)<\varepsilon$, d.h $\mu\left(A_{j}\right) \leq \mu\left(V_{j}\right)+\varepsilon$. Ferner ist $V_{j}$ kompakt, da $V_{j} \subset K_{j}$. Es folgt

$$
\mu(A) \leq \sup _{j \in \mathbb{N}} \mu\left(V_{j}\right)+\varepsilon \leq \sup \{\mu(K): K \text { kompakt, } K \subset A\}+\varepsilon
$$

Da $\varepsilon>0$ beliebig gewählt war, folgt die Behauptung.

\subsection*{Beispiel}
Bekannterweise ist das Lebesguemaß $\lambda$ auf $\mathbb{R}^{N}$ ein von außen reguläres und lokal endliches Borelmaß. Somit ist $\lambda$ ein Radonmaß.

\subsection*{Hilfssatz}
Sei $\Omega$ lokal kompakt und $\mu$ ein lokal endliches Borelmaß auf $\Omega$. Dann existiert zu jeder kompakten Teilmenge $K \subset \Omega$ und zu jedem $\varepsilon>0$ eine kompakte Umgebung $L$ von $K$ mit $\mu(L \backslash K)<\varepsilon$.

Beweis. Übungsaufgabe, Blatt 6. Man betrachte die Mengen $K_{n}:=\{x \in \Omega$ : $\left.\operatorname{dist}(x, K) \leq \frac{1}{n}\right\}$ für $n \in \mathbb{N}$.

\subsection*{Satz von Lusin}
Sei $\mu$ ein Radon-Maß auf $\Omega$. Sei ferner $u: \Omega \rightarrow \mathbb{K} \mu$-messbar. Dann gilt:\\
i). Ist $A \subset \Omega \mu$-messbar mit $\mu(A)<\infty$, so existiert zu jedem $\varepsilon>0$ eine kompakte Teilmenge $K \subset A$ mit $\mu(A \backslash K)<\varepsilon$ und derart, dass $\left.u\right|_{K}$ stetig ist.\\
ii). Ist $\Omega$ lokal kompakt und $\mu(\{x \in \Omega: u(x) \neq 0\})<\infty$, so existiert zu jedem $\varepsilon>0$ eine Funktion $u_{\varepsilon} \in \mathcal{C}_{c}(\Omega)$ mit

$$
\mu\left(\left\{x \in \Omega: u(x) \neq u_{\varepsilon}(x)\right\}\right)<\varepsilon \quad \text { und } \quad \sup _{x \in \Omega}\left|u_{\varepsilon}(x)\right| \leq \sup _{x \in \Omega}|u(x)| .
$$

Beweis. (i) Ohne Einschränkung sei $u \equiv 0$ in $\Omega \backslash A$ (sonst betrachte $u 1_{A}$ anstelle von $u$ ).

\begin{enumerate}
  \item Spezialfall: $u$ ist eine Elementarfunktion, d.h. $u=\sum_{k=1}^{n} a_{k} 1_{B_{k}}$ mit $\mu$-messbaren disjunkten Teilmengen $B_{k} \subset A$ und $a_{k} \in \mathbb{K} \backslash\{0\}, k=1, \ldots, n$. Sei ferner $B_{0}:= \{x \in A: u(x)=0\}$. Wegen $\mu\left(B_{k}\right) \leq \mu(A)<\infty$ und der inneren Regularität von $\mu$ existieren dann kompakte Teilmengen $C_{k} \subset B_{k}$ mit $\mu\left(B_{k} \backslash C_{k}\right)<\frac{\varepsilon}{n+1}$ für $k=0, \ldots, n$. Dabei gilt $\min _{k \neq j} \operatorname{dist}\left(C_{k}, C_{j}\right)>0$, da es sich um disjunkte kompakte Mengen handelt. Sei num $K:=\bigcup_{k=0}^{n} C_{k} \subset A$. Dann ist $K$ kompakt. Ferner ist $u$ lokal konstant und somit stetig auf $K$. Schließlich ist $A \backslash K$ die disjunkte Vereinigung der Mengen $B_{k} \backslash C_{k}$ und somit
\end{enumerate}

$$
\mu(A \backslash K)=\sum_{k=0}^{n} \mu\left(B_{k} \backslash C_{k}\right)<\varepsilon
$$

\begin{enumerate}
  \setcounter{enumi}{1}
  \item Allgemeiner Fall: Gemäß 6.0.4 existiert eine Folge von Elementarfunktionen $u_{n}$ mit $\left|u_{n}\right| \leq|u|$ für alle $n$ und $u_{n} \rightarrow u$ punktweise in $\Omega$ für $n \rightarrow \infty$. Nach dem Spezialfall existieren ferner kompakte Teilmengen $K_{n} \subset A$ mit $\mu\left(A \backslash K_{n}\right)<\varepsilon 2^{-n-1}$ und derart, dass $\left.u_{n}\right|_{K_{n}}$ stetig ist. Ferner existiert gemäß 6.12 eine $\mu$-messbare Teilmenge $L \subset A$ mit $\mu(A \backslash L)<\frac{\varepsilon}{4}$ und $u_{n} \rightarrow u$ gleichmäßig auf $L$ für $n \rightarrow \infty$. Schließlich existiert aufgrund der inneren Regularität von $\mu$ eine kompakte Teilmenge $K_{0} \subset L$ mit $\mu\left(L \backslash K_{0}\right)<\frac{\varepsilon}{4}$, also
\end{enumerate}

$$
\mu\left(A \backslash K_{0}\right) \leq \mu(A \backslash L)+\mu\left(L \backslash K_{0}\right)<\frac{\varepsilon}{2}
$$

Sei nun $K:=\bigcap_{n=0}^{\infty} K_{n}$. Dann ist $K \subset A$ kompakt und

$$
\mu(A \backslash K) \leq \sum_{n=0}^{\infty} \mu\left(A \backslash K_{n}\right)<\varepsilon \sum_{n=0}^{\infty} 2^{-n-1}=\varepsilon .
$$

In $K$ sind ferner alle Funktionen $u_{n}, n \in \mathbb{N}$ stetig; ferner gilt $u_{n} \rightarrow u$ gleichmäßig auf $K$. Es folgt, dass $\left.u\right|_{K}$ stetig ist.\\
(ii) Anwendung von (i) auf $A:=\{x \in \Omega: u(x) \neq 0\}$ liefert eine kompakte Teilmenge $K \subset A$ mit $\mu(A \backslash K)<\frac{\varepsilon}{2}$ und derart, dass $\left.u\right|_{K}$ stetig ist. Der Fortsetzungssatz von Tietze (aus der mengentheoretischen Topologie) besagt, dass sich $\left.u\right|_{K}$ zu einer stetigen Funktion $\tilde{u}: \Omega \rightarrow \mathbb{K}$ fortsetzen lässt mit $\|\tilde{u}\|_{\infty} \leq \sup _{x \in K}|u(x)|$. Gemäß 6.22 existiert ferner eine kompakte Umgebung $L$ von $K$ mit $\mu(L \backslash K)<\frac{\varepsilon}{2}$. Sei nun

$$
\xi \in \mathcal{C}_{c}(\Omega,[0,1]) \quad \text { definiert durch } \quad \xi(x)=\frac{\operatorname{dist}(x, \Omega \backslash L)}{\operatorname{dist}(x, K)+\operatorname{dist}(x, \Omega \backslash L)},
$$

und sei $u_{\varepsilon}:=\xi \tilde{u} \in \mathcal{C}_{c}(\Omega,[0,1])$. Dann ist $\xi \equiv 1$ auf $K$ und $\xi \equiv 0$ auf $\Omega \backslash L$. Nach Konstruktion ist also

$$
\left\{x \in \Omega: u(x) \neq u_{\varepsilon}(x)\right\} \subset[A \cup L] \backslash K=[A \backslash K] \cup[L \backslash K]
$$

und somit

$$
\mu\left(\left\{x \in \Omega: u(x) \neq u_{\varepsilon}(x)\right\}\right) \leq \frac{\varepsilon}{2}+\frac{\varepsilon}{2}=\varepsilon .
$$

Ferner ist

$$
\sup _{x \in \Omega}\left|u_{\varepsilon}(x)\right| \leq\|\tilde{u}\|_{\infty} \leq \sup _{x \in K}|u(x)| \leq \sup _{x \in \Omega}|u(x)|,
$$

wie gewünscht. \(\square\)

\subsection*{Beispiel}
Sei $\Omega=[0,1] \subset \mathbb{R}$, versehen mit der Betragsmetrik. Das Lebesguemaß $\lambda$ auf $\Omega$ erfüllt dann die Voraussetzungen des Satzes von Lusin. Ist insbesondere

$$
u:[0,1] \rightarrow \mathbb{R}, \quad u(x)= \begin{cases}1, & x \in \mathbb{Q}, \\ 0, & x \notin \mathbb{Q}\end{cases}
$$

die Dirichletfunktion und $\varepsilon>0$ gegeben, so kann man wie folgt eine kompakte Menge $K \subset[0,1]$ mit $\lambda([0,1] \backslash K)<\varepsilon$ und derart konstruieren, dass $\left.u\right|_{K}$ stetig ist: Sei $\left\{q_{n}: n \in \mathbb{N}\right\}$ eine Abzählung von $\mathbb{Q} \cap[0,1]$ und $U_{n}:=U_{\varepsilon 2^{-n-2}}\left(q_{n}\right)$ für $n \in \mathbb{N}$. Sei ferner

$$
K:=[0,1] \backslash \bigcup_{n \in \mathbb{N}} U_{n} .
$$

Dann ist $K$ kompakt, $u \equiv 0$ auf $K$ (und somit dort stetig) sowie

$$
\lambda([0,1] \backslash K) \leq \sum_{n \in \mathbb{N}} \lambda\left(U_{n}\right)=\varepsilon \sum_{n \in \mathbb{N}} 2^{-n-1}=\frac{\varepsilon}{2}<\varepsilon .
$$

\section*{Ende des Exkurses}
\subsection*{Satz}
Sei $\Omega$ ein lokal kompakter metrischer Raum und $\mu$ ein Radon-Maß auf $\Omega$. Sei ferner $1 \leq p<\infty$. Dann gilt:\\
i). $\mathcal{C}_{c}(\Omega) \subset L^{p}(\Omega)$ ist dicht.\\
ii). Ist $\Omega$ ferner $\sigma$-kompakt, so ist $L^{p}(\Omega)$ separabel.

Beweis. (i) Sei $u \in L^{p}(\Omega)$ und $\varepsilon>0$ gegeben. Nach 6.11 existiert $v \in \mathcal{E}(\Omega)$ mit $|v| \leq|u|$ in $\Omega, \mu(\{v \neq 0\})<\infty$ und $\|u-v\|_{p}<\frac{\varepsilon}{2}$. Ohne Einschränkung sei dabei $v \neq 0$, d.h. $\|v\|_{\infty}>0$. Nach dem Satz von Lusin 6.23(ii) existiert nun $w \in \mathcal{C}_{c}(\Omega)$ mit $\|w\|_{\infty} \leq\|v\|_{\infty}$ und $\mu(\{w \neq v\}) \leq\left(\frac{\varepsilon}{4\|v\|_{\infty}}\right)^{p}$, also

$$
\|w-v\|_{p}^{p}=\int_{\Omega}|w-v|^{p} d \mu \leq \mu(\{w \neq v\}) 2^{p}\|v\|_{\infty}^{p} \leq\left(\frac{\varepsilon}{2}\right)^{p}
$$

Es folgt $\|v-w\|_{p} \leq \frac{\varepsilon}{2}$ und somit $\|u-w\|_{p} \leq\|u-v\|_{p}+\|v-w\|_{p}<\varepsilon$.\\
(ii) Gemäß 6.16(i) existieren kompakte Mengen $M_{j} \subset \Omega, j \in \mathbb{N}$ mit $\Omega=\bigcup_{j=1}^{\infty} M_{j}$ und derart, dass jede kompakte Menge $K \subset \Omega$ in einer der Mengen $M_{j}$ enthalten ist. Wie im Beweis von 6.16(ii) setzen wir

$$
\mathcal{C}_{j}:=\left\{f \in \mathcal{C}_{c}(\Omega): \operatorname{supp} f \subset M_{j}\right\} \subset \mathcal{C}\left(M_{j}\right) \quad \text { für } j \in \mathbb{N} \text {. }
$$

Gemäß 4.28 ist $\mathcal{C}\left(M_{j}\right)$ und somit auch $\mathcal{C}_{j}$ separabel bzgl. $\|\cdot\|_{\infty}$, daher existieren abzähbare und dichte Teilmengen $A_{j} \subset \mathcal{C}_{j}$ für $j \in \mathbb{N}$.\\
Wir zeigen: Die abzählbare Menge $A:=\bigcup_{j \in \mathbb{N}} A_{j}$ liegt dicht in $L^{p}(\Omega)$. Sei dazu $u \in L^{p}(\Omega)$ und $\varepsilon>0$ beliebig. Gemäß (i) existiert $v \in \mathcal{C}_{c}(\Omega)$ mit $\|u-v\|_{p}<\frac{\varepsilon}{2}$; dabei ist $v \in \mathcal{C}_{j}$ für ein $j \in \mathbb{N}$. Somit existiert $w \in A_{j}$ mit $\mu\left(M_{j}\right)^{\frac{1}{p}}\|v-w\|_{\infty} \leq \frac{\varepsilon}{2}$. Dann ist

$$
\|w-v\|_{p}^{p}=\int_{\Omega}|w-v|^{p} d \mu \leq \mu\left(M_{j}\right)\|v-w\|_{\infty}^{p} \leq\left(\frac{\varepsilon}{2}\right)^{p}
$$

und somit $\|w-v\|_{p} \leq \frac{\varepsilon}{2}$. Es folgt $\|u-w\|_{p} \leq\|u-v\|_{p}+\|v-w\|_{p}<\varepsilon$. Also ist $A$ dicht in $L^{p}(\Omega)$, wie behauptet.

\subsection*{Bemerkung}
Ist $\Omega \subset \mathbb{R}^{N}$ Lebesgue-messbar, so ist $L^{p}(\Omega)$ separabel, da man $L^{p}(\Omega)$ als eine Teilmenge des gemäß 6.25 separablen Raums $L^{p}\left(\mathbb{R}^{N}\right)$ auffassen kann (vermöge trivialer Fortsetzung).

\subsection*{Satz}
Ist $\Omega \subset \mathbb{R}^{N}$ Lebesgue-messbar mit $\lambda(\Omega)>0$, so ist $L^{\infty}(\Omega)$ (definiert bzgl. des Lebesgue-Maßes) nicht separabel.

Beweisskizze. Für $r \geq 0$ sei $\Omega_{r}:=B_{r}(0) \cap \Omega$. Mit Satz von Lebesgue (siehe 6.0.7) sieht man, dass die monoton wachsende Funktion

$$
g:[0, \infty) \rightarrow[0, \infty), \quad g(r)=\lambda\left(\Omega_{r}\right)
$$

stetig ist. Ferner ist $g(0)=0$ und $\lim _{r \rightarrow \infty} g(r)=\lambda(\Omega)$ nach dem Satz von der monotonen Konvergenz (siehe 6.0.5). Für $s \in(0, \lambda(\Omega))$ existiert nach dem Zwischenwertsatz also ein $r_{s} \in[0, \infty)$ mit $g\left(r_{s}\right)=s$. Betrachte nun die Menge

$$
M:=\left\{1_{\Omega_{r_{s}}}: 0<s<\lambda(\Omega)\right\} \subset L^{\infty}(\Omega) .
$$

Für $0<s<t<\lambda(\Omega)$ ist dann $r_{s}<r_{t}$ und somit $\Omega_{r_{s}} \subset \Omega_{r_{t}}$ sowie $\lambda\left(\Omega_{r_{t}} \backslash \Omega_{r_{s}}\right)= t-s>0$; somit also $\left\|1_{\Omega_{r_{s}}}-1_{\Omega_{r_{t}}}\right\|_{\infty}=1$. Somit ist $M$ überabzählbar und diskret in $L^{\infty}(\Omega)$. Mit 2.6) folgt, dass $L^{\infty}(\Omega)$ nicht separabel ist. \(\square\)

\subsection*{Definition und Satz}
Sei $\Omega \subset \mathbb{R}^{N}$ offen und

$$
\mathcal{C}_{c}^{\infty}(\Omega):=\left\{f \in C^{\infty}(\Omega, \mathbb{K}): \operatorname{supp} f \subset \Omega \text { kompakt }\right\} .
$$

Dann liegt $C_{c}^{\infty}(\Omega)$ dicht in $L^{p}(\Omega)$ für $1 \leq p<\infty$, aber (offensichtlich) nicht dicht in $L^{\infty}(\Omega)$. Hierbei sei das Lebesguemaß auf $\Omega$ zugrunde gelegt.

Beweisskizze. Sei $u \in L^{p}(\Omega), \varepsilon>0$. Nach 6.25 existiert $v \in \mathcal{C}_{c}(\Omega)$ mit $\|u-v\|_{p}<\frac{\varepsilon}{2}$. Wir setzen $v$ trivial auf $\mathbb{R}^{N}$ fort. Wir betrachten ferner den Glättungskern

$$
\rho \in C_{c}^{\infty}\left(\mathbb{R}^{N}\right), \quad \rho(x)= \begin{cases}c \exp \left(\frac{1}{|x|^{2}-1}\right) & ,|x|<1 \\ 0 & ,|x| \geq 1,\end{cases}
$$

wobei $c>0$ so gewählt ist, dass $\int_{\mathbb{R}^{N}} \rho d x=1$ ist. Für $h>0$ sei ferner

$$
\rho_{h} \in C_{c}^{\infty}\left(\mathbb{R}^{N}\right), \quad \rho_{h}(x)=\frac{1}{h^{N}} \rho\left(\frac{x}{h}\right),
$$

dann ist $\operatorname{supp} \rho_{h}=B_{h}(0)$ und es gilt ebenfalls $\int_{\mathbb{R}^{N}} \rho_{h} d x=1$ für $h>0$. Nun sei

$$
v_{h}: \mathbb{R}^{N} \rightarrow \mathbb{R}, \quad v_{h}(x)=\int_{\mathbb{R}^{N}} \rho_{h}(x-y) v(y) d y
$$

für $0<h<\operatorname{dist}(\operatorname{supp} v, \partial \Omega)$. Dann ist $v_{h} \in C_{c}^{\infty}\left(\mathbb{R}^{N}\right)$ mit $\operatorname{supp} v_{h} \subset B_{h}(\operatorname{supp} v) \subset \Omega$; ferner gilt mit

$$
\varepsilon_{h}:=\sup \{|v(x)-v(y)|: x, y \in \Omega,|x-y| \leq h\}
$$

die Abschätzung

$$
\left|v_{h}(x)-v(x)\right| \leq \int_{\mathbb{R}^{N}} \rho_{h}(x-y)|v(y)-v(x)| d y \leq \varepsilon_{h} \int_{\mathbb{R}^{N}} \rho_{h}(x-y) d y=\varepsilon_{h} \quad \text { für } x \in \Omega .
$$

Es folgt $v_{h} \rightarrow v$ bzgl. $\|\cdot\|_{\infty}$ in $\Omega$ und somit auch $\left\|v_{h}-v\right\|_{L^{p}(\Omega)} \rightarrow 0$, da $\operatorname{supp} v_{h}$ für $h \leq$ 1 in der kompakten Menge $B_{1}(\operatorname{supp} v)$ enthalten ist (vgl. 6.5(ii)). Die Einschränkung $w_{h}:=\left.v_{h}\right|_{\Omega} \in \mathcal{C}_{c}^{\infty}(\Omega)$ erfüllt also $\left\|w_{h}-u\right\|_{p} \leq\left\|w_{h}-v\right\|_{p}+\|v-u\|_{p}<\frac{\varepsilon}{2}+\frac{\varepsilon}{2}=\varepsilon$ für $0<h<\operatorname{dist}(\operatorname{supp} v, \partial \Omega)$ genügend klein gewählt.

\section*{§7 Hilberträume und ihre Geometrie}
Stets sei $H$ ein Hilbertraum über $\mathbb{K}$ mit Skalarprodukt $\langle\cdot, \cdot\rangle$ und induzierter Norm $\|\cdot\|$.

\subsection*{Satz und Definition}
Sei $C \subset H$ konvex, abgeschlossen und nichtleer. Dann existiert zu jedem $x \in H$ genau ein $y \in C$ mit $\|x-y\|=\operatorname{dist}(x, C)$. Wir setzen $y=P(x)$ und nennen $P: H \rightarrow H$ die (abstandsminimierende) Projektion auf $C$.

Beweis. Zur Existenz: Sei $x \in H$ und $\alpha:=\operatorname{dist}(x, C)$. Sei ferner $\left(y_{n}\right)_{n} \subset C$ eine Folge mit

$$
\left\|x-y_{n}\right\| \rightarrow \alpha \quad \text { für } n \rightarrow \infty .
$$

Für $n, m \in \mathbb{N}$ gilt $\frac{y_{n}+y_{m}}{2} \in C$ aufgrund der Konvexität von $C$, also

$$
\begin{aligned}
4 \alpha^{2} \leq 4\left\|x-\frac{y_{n}+y_{m}}{2}\right\|^{2}=\left\|x-y_{m}+x-y_{n}\right\|^{2} \\
\stackrel{\text { Ü.-Bl. }}{=} 2\left(\left\|x-y_{m}\right\|^{2}+\left\|x-y_{n}\right\|^{2}\right)-\underbrace{\left\|x-y_{m}-\left(x-y_{n}\right)\right\|^{2}}_{\left\|y_{n}-y_{m}\right\|^{2}}
\end{aligned}
$$

und somit

$$
\sup _{m \geq n}\left\|y_{m}-y_{n}\right\|^{2} \leq 2 \sup _{m \geq n}\left(\left\|x-y_{m}\right\|^{2}+\left\|x-y_{n}\right\|^{2}\right)-4 \alpha^{2} \xrightarrow{(*)} 0 \quad \text { für } n \rightarrow \infty
$$

Es folgt, dass $\left(y_{n}\right)_{n} \subset C$ eine Cauchyfolge ist. Ferner ist $C$ als abgeschlossene Teilmenge des Hilbertraums $H$ nach 3.3 vollständig, also existiert $y:=\lim _{n} y_{n} \in C$. Ferner gilt $\|x-y\|=\lim _{n}\left\|x-y_{n}\right\|=\alpha$.

Zur Eindeutigkeit: Sei \begin{tabular}{l}
$n$ \\
$y^{\prime} \in C$ mit $\left\|x-y^{\prime}\right\|=\alpha=\|x-y\|$. Dann erfüllt die Folge \\
\hline
\end{tabular} ( $y, y^{\prime}, y, y^{\prime}, \ldots$ ) Bedingung ( $*$ ) und ist daher nach obigem Argument eine Cauchyfolge. Dies liefert $y=y^{\prime}$.

\subsection*{Satz}
Sind $C \subset H$ und $P: H \rightarrow H$ wie in 7.1 und $Q:=\operatorname{id}-P: H \rightarrow H$, so gilt\\
i). $\left.P\right|_{C}=\mathrm{id}_{C}$\\
ii). Für alle $x \in H, y \in C$ ist $\operatorname{Re}\langle Q(x), P(x)-y\rangle \geq 0$\\
iii). $\left\|P\left(x_{1}\right)-P\left(x_{2}\right)\right\| \leq\left\|x_{1}-x_{2}\right\|$ für alle $x_{1}, x_{2} \in H$, insbesondere ist $P$ Lipschitzstetig.

Beweis. i) ist klar nach Definition.\\
ii) Für $\varepsilon \in(0,1)$ ist $z_{\varepsilon}:=P(x)+\varepsilon(y-P(x)) \in C$ (da $C$ konvex), also

$$
\begin{aligned}
0 & \leq\left\|x-z_{\varepsilon}\right\|^{2}-\operatorname{dist}(x, C)^{2}=\|Q(x)+\varepsilon(P(x)-y)\|^{2}-\|Q(x)\|^{2} \\
& =2 \varepsilon \operatorname{Re}\langle Q(x), P(x)-y\rangle+\varepsilon^{2}\|P(x)-y\|^{2}=: t(\varepsilon)
\end{aligned}
$$

und somit

$$
0 \leq \lim _{\varepsilon \rightarrow 0^{+}} \frac{t(\varepsilon)}{2 \varepsilon}=\operatorname{Re}\langle Q(x), P(x)-y\rangle
$$

iii) Es ist

$$
\begin{aligned}
& \left\|x_{1}-x_{2}\right\|^{2}=\left\|Q\left(x_{1}\right)-Q\left(x_{2}\right)+P\left(x_{1}\right)-P\left(x_{2}\right)\right\|^{2} \\
& =\left\|Q\left(x_{1}\right)-Q\left(x_{2}\right)\right\|^{2}+\left\|P\left(x_{1}\right)-P\left(x_{2}\right)\right\|^{2}+2 \underbrace{\operatorname{Re}\left\langle Q\left(x_{1}\right)-Q\left(x_{2}\right), P\left(x_{1}\right)-P\left(x_{2}\right)\right\rangle}_{\geq 0 \text { nach ii) }} \\
& \geq\left\|P\left(x_{1}\right)-P\left(x_{2}\right)\right\|^{2}
\end{aligned}
$$

\subsection*{Bemerkung}
Die Sätze 7.1 und 7.2 gelten auch, wenn ( $H,\langle\cdot, \cdot\rangle$ ) ein Prähilbertraum und $C \subset H$ nichtleer, konvex und vollständig ist, insbesondere also dann, wenn die abgeschlossene, nichtleere und konvexe Teilmenge $C$ in einem endlichdimensionalen Teilraum des Prähilbertraums $H$ enthalten ist.

\subsection*{Beispiel}
Sei $\Omega \subset \mathbb{R}^{N}$ Lebesgue-messbar, $H=L^{2}(\Omega, \mathbb{R})$ sowie $C:=\left\{u \in L^{2}(\Omega)\right.$ : $u \geq 0$ f.ü. auf $\Omega\}$. Man sieht leicht, dass $C$ eine nichtleere, abgeschlossene und konvexe Teilmenge von $H$ ist. Wir zeigen nun, dass die Projektion $P$ von $H$ auf $C$ gegeben ist durch $P(u)=u^{+}=\max \{0, u\}$ für $u \in H$. Tatsächlich ist $u^{+} \in C$ für alle $u \in H$, und für $v \in C$ gilt

$$
\|v-u\|_{2}^{2}=\int_{\Omega}(v-u)^{2} d x \geq \int_{\{u<0\}}(v-u)^{2} d x \geq \int_{\{u<0\}} u^{2} d x=\int_{\Omega}\left(u^{+}-u\right)^{2} d x=\left\|u^{+}-u\right\|_{2}^{2}
$$

Also ist $u^{+}$das zu $u$ abstandsminimierende Element in $C$, d.h. $P(u)=u^{+}$gemäß 7.1.

\subsection*{Definition und Satz}
Für $M \subset H$ sei $M^{\perp}:=\{x \in H:\langle x, y\rangle=0$ für alle $y \in M\}$. Dann gilt:\\
i) $M^{\perp}$ ist ein abgeschlossener Unterraum von $H$\\
ii) $M^{\perp}=(\overline{\operatorname{span} M})^{\perp}$

Beweis. i) folgt direkt aus der Linearität und Stetigkeit des Skalarprodukts in der ersten Komponente.\\
Zu ii): Offensichtlich gilt für Teilmengen $A \subset B \subset H$ die Inklusionsumkehrung $B^{\perp} \subset A^{\perp}$. Wegen $M \subset \overline{\operatorname{span} M}$ folgt also $(\overline{\operatorname{span} M})^{\perp} \subset M^{\perp}$. Sei umgekehrt $x \in M^{\perp}$, und sei zunächst $y=\alpha_{1} y_{1}+\cdots+\alpha_{n} y_{n} \in \operatorname{span} M$ mit $y_{i} \in M, \alpha_{i} \in K$. Dann ist

$$
\langle x, y\rangle=\sum_{i=1}^{n} \bar{\alpha}_{i}\left\langle x, y_{i}\right\rangle=0
$$

Ist schließlich $y \in \overline{\operatorname{span} M}$ und $\left(y^{n}\right)_{n}$ eine Folge in $\operatorname{span} M$ mit $\lim _{n \rightarrow \infty} y^{n}=y$, so folgt $\langle x, y\rangle=\lim _{n \rightarrow \infty}\left\langle x, y^{n}\right\rangle=0$. Somit ist $x \in(\overline{\operatorname{span} M})^{\perp}$, und insegesamt folgt $M^{\perp}= (\overline{\operatorname{span} M})^{\perp}$.

\subsection*{Satz und Definition}
Sind $V, W \subset H$ Unterräume mit $H=V+W$ und $\langle v, w\rangle=0$ für alle $v \in V, w \in W$, so gilt:\\
i). $\|v+w\|^{2}=\|v\|^{2}+\|w\|^{2}$ für alle $v \in V, w \in W$ („Satz des Pythagoras")\\
ii). $W=V^{\perp}$ und $V=W^{\perp}$, insbesondere sind $V, W \subset H$ abgeschlossen nach 7.5.\\
iii). $H=V \bigoplus_{\text {top }} W$

Wir schreiben dann $H=V \bigoplus_{\perp} W$ und nennen diese Zerlegung orthogonal.\\
Beweis. i) Für alle $v \in V, w \in W$ ist

$$
\|v+w\|^{2}=\|v\|^{2}+\|w\|^{2}+2 \operatorname{Re} \underbrace{\langle v, w\rangle}_{=0}=\|v\|^{2}+\|w\|^{2} .
$$

ii) Nach Voraussetzung ist $W \subset V^{\perp}$ und $V \subset W^{\perp}$. Sei nun $w \in V^{\perp}$. Da $H=V+W$ ist, existiert $v \in V$ und $\tilde{w} \in W$ mit $w=v+\tilde{w}$, also

$$
\|v\|^{2}=\langle v, v\rangle=\langle\underbrace{w-\tilde{w}}_{\in V^{\perp}}, v\rangle=0 .
$$

Es folgt $w=\tilde{w} \in W$. Dies zeigt $W=V^{\perp}$. Genauso folgt $V=W^{\perp}$.\\
iii) Es gilt $V \cap W=\{0\}$, da $\|x\|^{2}=\langle x, x\rangle=0$ für alle $x \in V \cap W$. Also gilt: $H=V \oplus W$. Da $V, W \subset H$ gemäß ii) abgeschlossen sind, folgt $H=V \bigoplus_{\text {top }} W$ nach 5.14ii).

\subsection*{Satz und Definition}
Sei $V$ ein abgeschlossener Unterraum von $H$. Dann gilt $H=V \bigoplus_{\perp} V^{\perp}$, und die (abstandsminimierende) Projektion $P$ auf $V$ aus 7.1 erfüllt:\\
i) $P \in \mathcal{L}(H), P^{2}=P$ und Bild $P=V$, $\operatorname{Kern} P=V^{\perp}$\\
ii) $\|P\|=1$ falls $V \neq\{0\}$

Man nennt $P$ die orthogonale Projektion auf $V$.\\
Beweis. Wir zeigen zunächst: Für $x \in H$ ist $Q(x):=x-P(x) \in V^{\perp} . \quad(*)$ Nach 7.2 i) gilt für alle $y \in V, \alpha \in \mathbb{K}$ :

$$
0 \leq \operatorname{Re}\langle Q(x), P(x)-(\underbrace{-\bar{\alpha} y+P(x)}_{\in V})\rangle=\operatorname{Re}\langle Q(x), \bar{\alpha} y\rangle=\operatorname{Re}[\alpha\langle Q(x), y\rangle] .
$$

Durch Wahl von $\alpha=1$ und $\alpha=-1$ folgt $\operatorname{Re}\langle Q(x), y\rangle=0$ für alle $y \in V$. Ist $\mathbb{K}=\mathbb{C}$, so folgt durch Wahl von $\alpha= \pm i$ auch

$$
0 \leq \operatorname{Re}[ \pm i\langle Q(x), y\rangle]=\mp \operatorname{Im}\langle Q(x), y\rangle, \quad \text { also } \quad \operatorname{Im}\langle Q(x), y\rangle=0 \quad \text { für alle } y \in V .
$$

Insgesamt folgt $\langle Q(x), y\rangle=0$ für alle $y \in V$, und somit ist $Q(x) \in V^{\perp}$, wie behauptet.\\
Aus (*) folgt nun $x=P(x)+Q(x) \in V+V^{\perp}$ für alle $x \in H$, also $H=V+V^{\perp}$. Nach 7.6 ist also $H=V \bigoplus_{\perp} V^{\perp}$.\\
Sei nun $\widetilde{P} \in \mathcal{L}(H)$ die Projektion auf $V$ längs $V^{\perp}$ im Sinne von 5.12. Dann gilt für alle $x \in H$ :

$$
\begin{aligned}
\operatorname{dist}(x, V)^{2} & =\|x-P(x)\|^{2}=\|\underbrace{x-\widetilde{P}(x)}_{\in V^{\perp}}+\underbrace{\widetilde{P}(x)-P(x)}_{\in V}\|^{2} \\
& \stackrel{\text { (7.6) }}{=}\|x-\widetilde{P}(x)\|^{2}+\|\widetilde{P}(x)-P(x)\|^{2} \geq\|x-\widetilde{P}(x)\|^{2} \geq \operatorname{dist}(x, V)^{2} .
\end{aligned}
$$

Dies erzwingt $\|\widetilde{P}(x)-P(x)\|^{2}=0$, also $\widetilde{P}(x)=P(x)$ für alle $x \in H$.\\
Insbesondere ist $P \in \mathcal{L}(H), P^{2}=P$ sowie Bild $P=V$ und Kern $P=V^{\perp}$.\\
Weiterhin ist $\|P x\| \leq\|x\|$ für $x \in H$ nach 7.2iii) und $\|P x\|=\|x\|$ für $x \in V$. Dies liefert $\|P\|=1$, falls $V \neq\{0\}$.

\subsection*{Korollar}
Für $M \subset H$ gilt\\
i) $M^{\perp \perp}=\overline{\operatorname{span} M}$\\
ii) $M$ total in $H \Longleftrightarrow M^{\perp}=\{0\}$

Beweis. i) Setze $V=\overline{\operatorname{span} M}$. Dann ist $H=V \bigoplus_{\perp} V^{\perp}$ nach 7.7, also $V=V^{\perp \perp} ={\text { 7.5 }} M^{\perp \perp}$.\\
ii) Ist $M$ total, so ist $M^{\perp} \stackrel{\text { 7.5 }}{=} \overline{\operatorname{span}}^{\perp}=H^{\perp}=\{0\}$.

Ist umgekehrt $M^{\perp}=\{0\}$, so ist $H=\{0\}^{\perp}=M^{\perp \perp} \stackrel{i}{=} \overline{\operatorname{span} M}$. Also ist $M$ total.

\subsection*{Definition und Bemerkung}
Für $x \in H$ sei

$$
\varphi_{x} \in H^{\prime}=\mathcal{L}(H, \mathbb{K}) \quad \text { definiert durch } \quad \varphi_{x}(y)=\langle y, x\rangle \quad \text { für } y \in H .
$$

Die Linearität von $\varphi_{x}$ ist offensichtlich. Da ferner

$$
\left|\varphi_{x}(y)\right|=|\langle y, x\rangle| \leq\|y\|\|x\| \quad \text { für alle } y \in H \quad \text { und } \quad\left|\varphi_{x}(x)\right|=\|x\|^{2},
$$

ist $\varphi_{x}$ in der Tat auch stetig mit

$$
\left\|\varphi_{x}\right\|=\|x\|
$$

Es folgt, dass die Abbildung $J: H \rightarrow H^{\prime}, x \mapsto J(x):=\varphi_{x}$ eine Isometrie ist und damit insbesondere injektiv.\\
Man beachte ferner: $J$ ist antilinear, d. h. es gilt\\
$J\left(x_{1}+x_{2}\right)=J\left(x_{1}\right)+J\left(x_{2}\right) \quad$ und $\quad J(\lambda x)=\bar{\lambda} J(x) \quad$ für alle $x, x_{1}, x_{2} \in H, \lambda \in \mathbb{K}$.

\subsection*{Rieszscher Darstellungssatz}
i) Für alle $\varphi \in H^{\prime}$ existiert $x \in H$ mit $\varphi=\varphi_{x}$. Die Abbildung $J: H \rightarrow H^{\prime}$ aus 7.9 ist also bijektiv.\\
ii) Die Operatornorm auf $H^{\prime}$ wird induziert von einem Skalarprodukt $\langle\cdot, \cdot\rangle_{*}$ gegeben durch

$$
\langle\varphi, \psi\rangle_{*}:=\left\langle J^{-1}(\psi), J^{-1}(\varphi)\right\rangle \quad \text { für alle } \varphi, \psi \in H^{\prime} .
$$

Insbesondere ist ( $H^{\prime},\langle\cdot, \cdot\rangle_{*}$ ) ein Hilbertraum.

Beweis. i) Sei $\varphi \in H^{\prime} \backslash\{0\}$. Dann ist $V:=\operatorname{Kern} \varphi$ ein abgeschlossener Unterraum von $H$. Nach 7.7 ist also $H=V \bigoplus_{\perp} V^{\perp}$ mit $V^{\perp}=\operatorname{span}\left\{x_{0}\right\}$ für ein $x_{0} \in V^{\perp} \backslash\{0\}$. Setze $x:=\frac{\overline{\varphi\left(x_{0}\right)}}{\left\|x_{0}\right\|^{2}} x_{0}$. Dann gilt für alle $y=v+\lambda x \in H$ mit $v \in V$ :

$$
\langle y, x\rangle=\langle v, x\rangle+\lambda\|x\|^{2}=\lambda\|x\|^{2} \quad \text { und } \quad \varphi(y)=\lambda \varphi(x)=\lambda \frac{\left|\varphi\left(x_{0}\right)\right|^{2}}{\left\|x_{0}\right\|^{2}}=\lambda\|x\|^{2} .
$$

Es folgt $\varphi=\varphi_{x}$.\\
ii) Die Skalarprodukt-Eigenschaften von $\langle\cdot, \cdot\rangle_{*}$ folgen direkt aus denen von $\langle\cdot, \cdot\rangle$. Ferner gilt für $\varphi \in H^{\prime}$ mit $x=J^{-1}(\varphi)$ :

$$
\langle\varphi, \varphi\rangle_{*} \stackrel{\text { Def. }}{=}\langle x, x\rangle=\|x\|^{2} \stackrel{\boxed{7.1}}{=}\|\varphi\|^{2}
$$

Somit induziert $\langle\cdot, \cdot\rangle_{*}$ also die Operatornorm auf $H^{\prime}$.

\section*{Exkurs: Anwendung auf ein Neumann-Randwertproblem}
Wir betrachten stets $\mathbb{K}=\mathbb{R}$ und $\Omega=(-1,1) \subset \mathbb{R}$. Gegeben seien Funktionen $a \in L^{\infty}(\Omega), f \in L^{2}(\Omega)$, wobei $\underset{\Omega}{\operatorname{essinf}} a>0$ sei. Wir suchen eine Lösung des NeumannRandwertproblems

$$
(N P) \quad\left\{\begin{aligned}
-u^{\prime \prime}+a(x) u & =f(x), \quad x \in \Omega \\
u^{\prime}(1)=u^{\prime}(-1) & =0
\end{aligned}\right.
$$

Dazu benötigen wir einige Vorbereitungen.

\subsection*{Definition}
i). Sei $u \in L^{2}(\Omega)$. Wir nennen $v \in L^{2}(\Omega)$ eine schwache Ableitung von $u$, wenn gilt:

$$
\int_{\Omega} v w d x=-\int_{\Omega} u w^{\prime} d x \quad \text { für alle } w \in C_{c}^{1}(\Omega)
$$

Dadurch ist $v$ eindeutig bestimmt, denn gilt (*) auch für $\tilde{v} \in L^{2}(\Omega)$, so ist

$$
\langle v-\tilde{v}, w\rangle_{L^{2}(\Omega)}=-\int_{\Omega}(v-\tilde{v}) w d x=\int_{\Omega}(u-u) w^{\prime} d x=0 \quad \text { für alle } w \in C_{c}^{1}(\Omega)
$$

Da $C_{c}^{1}(\Omega)$ in $L^{2}(\Omega)$ dicht liegt, folgt $v-\tilde{v} \in C_{c}^{1}(\Omega)^{\perp}=\{0\}$.\\
Wir schreiben $u^{\prime}$ anstelle von $v$.\\
ii). Wir setzen

$$
H^{1}(\Omega):=\left\{u \in L^{2}(\Omega): u \text { besitzt eine schwache Ableitung } u^{\prime} \in L^{2}(\Omega)\right\}
$$

und versehen $H^{1}(\Omega)$ mit dem Skalarprodukt

$$
\left\langle u_{1}, u_{2}\right\rangle_{H^{1}}:=\left\langle u_{1}, u_{2}\right\rangle_{L^{2}(\Omega)}+\left\langle u_{1}^{\prime}, u_{2}^{\prime}\right\rangle_{L^{2}(\Omega)} .
$$

Die induzierte Norm sei mit $\|\cdot\|_{H^{1}}$ bezeichnet. $H^{1}$ gehört zu der Klasse der sogenannten Sobolevräume.

\subsection*{Bemerkung und Beispiel}
i). Ist $u \in C^{1}(\Omega) \cap L^{2}(\Omega)$ mit klassischer Ableitung $u^{\prime} \in L^{2}(\Omega)$, so ist $u^{\prime}$ auch die schwache Ableitung von $u$. Dies folgt aus der gewöhnlichen Regel zur partiellen Integration. Dies gilt insbesondere, wenn $u \in C^{1}(\bar{\Omega})$ ist.\\
ii). Sei $u \in L^{2}(\Omega)$ definiert durch $u(x)=|x|^{\alpha}$ für ein $\alpha>\frac{1}{2}$. Dann ist $u \in H^{1}(\Omega)$ mit schwacher Ableitung

$$
v=u^{\prime} \in L^{2}(\Omega), \quad v(x)=\alpha|x|^{\alpha-2} x \quad \text { für } x \neq 0 .
$$

In der Tat ist $v \in L^{2}(\Omega)$, denn wegen $2(\alpha-1)>-1$ ist

$$
\int_{\Omega}|v(x)|^{2} d x=\alpha^{2} \int_{\Omega}|x|^{2(\alpha-1)}<\infty
$$

Ferner ist $v(x)=u^{\prime}(x)$ im klassischen Sinne klassisch für $x \neq 0$, und somit gilt für $w \in \mathcal{C}_{c}^{1}(\Omega)$ :

$$
\begin{aligned}
\int_{\Omega} v w d x & =\lim _{\varepsilon \rightarrow 0^{+}}\left(\int_{-1}^{-\varepsilon} u^{\prime}(x) w(x) d x+\int_{\varepsilon}^{1} u^{\prime}(x) w(x) d x\right) \\
& =\lim _{\varepsilon \rightarrow 0^{+}}\left([u w](-\varepsilon)-\int_{-1}^{-\varepsilon} u(x) w^{\prime}(x) d x-[u w](\varepsilon)-\int_{\varepsilon}^{1} u(x) w^{\prime}(x) d x\right) \\
& =-\lim _{\varepsilon \rightarrow 0^{+}}\left(\int_{-1}^{-\varepsilon} u(x) w^{\prime}(x) d x+\int_{\varepsilon}^{1} u(x) w^{\prime}(x) d x\right)=\int_{\Omega} u w^{\prime} d x
\end{aligned}
$$

Dies zeigt $u^{\prime}=v$ im schwachen Sinne.\\
Wir werden weiter unten in 7.14 sehen, dass die oben definierte Funktion im Fall $\alpha<\frac{1}{2}$ nicht mehr in $H^{1}(\Omega)$ liegt.

\subsection*{Satz}
$H^{1}(\Omega)$ ist ein Hilbertraum.

Beweis. Die Skalarprodukteigenschaften sind klar. Sei $\left(u_{k}\right)_{k}$ eine Cauchyfolge in $H^{1}(\Omega)$. Nach Definition des Skalarprodukts ist $\left(u_{k}\right)_{k}$ dann eine Cauchyfolge in $L^{2}(\Omega)$ und $\left(u_{k}^{\prime}\right)_{k}$ ist eine Cauchyfolge in $L^{2}(\Omega)$. Somit gilt

$$
u_{k} \rightarrow u \in L^{2}(\Omega) \quad \text { und } \quad u_{k}^{\prime} \rightarrow v \in L^{2}(\Omega) \quad \text { für } k \rightarrow \infty .
$$

Ferner gilt für alle $w \in C_{c}^{1}(\Omega)$ :

$$
\int_{\Omega} v w d x=\lim _{k \rightarrow \infty} \int_{\Omega} u_{k}^{\prime} w d x=-\lim _{k \rightarrow \infty} \int_{\Omega} u_{k} w^{\prime} d x=-\int_{\Omega} u w^{\prime} d x .
$$

Somit ist also $u \in H^{1}(\Omega)$ mit $u^{\prime}=v$, und

$$
\left\|u_{k}-u\right\|_{H^{1}}^{2}=\left\|u_{k}-u\right\|_{2}^{2}+\left\|u_{k}^{\prime}-v\right\|_{2}^{2} \rightarrow 0 \quad \text { für } n \rightarrow \infty .
$$

Dies zeigt die Vollständigkeit von $H^{1}(\Omega)$.

\subsection*{Satz}
Sei $u \in H^{1}(\Omega)$. Dann gilt:\\
i). Ist $u^{\prime} \equiv 0$, so besitzt $u$ einen konstanten Repräsentanten.\\
ii). Die Funktion $u$ lässt sich (nach Wahl eines geeigneten Repräsentanten) zu einer stetigen Funktion auf $\bar{\Omega}$ fortsetzen, welche

$$
u(x)-u(y)=\int_{x}^{y} u^{\prime}(t) d t \quad \text { für alle } x, y \in \bar{\Omega}
$$

erfüllt.\\
iii). Es gilt $|u(x)-u(y)| \leq\left\|u^{\prime}\right\|_{2}|x-y|^{\frac{1}{2}}$ für $x, y \in \bar{\Omega}$ und $\|u\|_{\infty} \leq \sqrt{2}\|u\|_{H^{1}}$.

Beweis. Zu (i): Sei $\psi \in C_{c}(\Omega)$ mit $\int_{\Omega} \psi d x=1$ fest gewählt. Sei ferner $w \in C_{c}(\Omega)$ beliebig, und sei

$$
f_{w}:=w-\left(\int_{\Omega} w d x\right) \psi \in C_{c}(\Omega)
$$

Diese Funktion besitzt eine Stammfunktion gegeben durch

$$
F_{w} \in C_{c}^{1}(\Omega), \quad F_{w}(x)=\int_{-1}^{x} f_{w}(t) d t
$$

und somit gilt (nach Voraussetzung und Definition der schwachen Ableitung)

$$
\begin{aligned}
0 & =-\int_{\Omega} u^{\prime} F_{w} d x=\int_{\Omega} u f_{w} d x=\int_{\Omega} u w d x-\left(\int_{\Omega} w d x\right)\left(\int_{\Omega} u \psi d y\right) \\
& =\int_{\Omega}\left[u-c_{u}\right] w d x \quad \text { mit } c_{u}:=\int_{\Omega} u \psi d y
\end{aligned}
$$

Da $w \in C_{c}(\Omega)$ beliebig gewählt war, folgt $u-c_{u} \in\left[C_{c}(\Omega)\right]^{\perp}$ bzgl. des $L^{2}$ Skalarprodukts, also $u-c_{u}=0$ in $L^{2}(\Omega)$ aufgrund der Dichtheit von $C_{c}(\Omega)$ in $L^{2}(\Omega)$. Somit gilt $u \equiv c_{u}$ f.ü. in $\Omega$.\\
Zu (ii): Sei $u \in H^{1}(\Omega)$, und sei $v: \Omega \rightarrow \mathbb{R}$ definiert durch $v(x)=\int_{0}^{x} u^{\prime}(t) d t$. Für $x, y \in \Omega$ gilt dann

$$
|v(x)-v(y)|=\left|\int_{y}^{x} u^{\prime}(t) d t\right| \leq|x-y|^{\frac{1}{2}}\left(\int_{\min \{x, y\}}^{\max \{x, y\}}\left[u^{\prime}(t)\right]^{2} d t\right)^{\frac{1}{2}} \leq|x-y|^{\frac{1}{2}}\left\|u^{\prime}\right\|_{2} .
$$

Insbesondere ist $v$ gleichmäßig stetig und lässt sich somit zu einer stetigen Funktion auf $\bar{\Omega}$ fortsetzen, welche dann (7.3) sogar für $x, y \in \bar{\Omega}$ erfüllt. Insbesondere ist $v \in L^{2}(\Omega)$. Wir zeigen nun, dass $v$ die schwache Ableitung $u^{\prime}$ besitzt. Sei dazu $\varphi \in C_{c}^{1}(\Omega)$. Dann gilt mit dem Satz von Fubini

$$
\begin{aligned}
& \int_{\Omega} \varphi^{\prime} v d x=\int_{-1}^{1} \varphi^{\prime}(x) \int_{0}^{x} u^{\prime}(t) d t d x=-\int_{-1}^{0} \varphi^{\prime}(x) \int_{x}^{0} u^{\prime}(t) d t d x+\int_{0}^{1} \varphi^{\prime}(x) \int_{0}^{x} u^{\prime}(t) d t d x \\
& =-\int_{1}^{0} u^{\prime}(t) \int_{-1}^{t} \varphi^{\prime}(x) d x d t+\int_{0}^{1} u^{\prime}(t) \int_{t}^{1} \varphi^{\prime}(x) d x d t=-\int_{1}^{0} u^{\prime}(t) \varphi(t) d t-\int_{0}^{1} u^{\prime}(t) \varphi(t) d t \\
& =-\int_{-1}^{1} u^{\prime} \varphi d t
\end{aligned}
$$

Dies zeigt $v^{\prime}=u^{\prime}$ im schwachen Sinne. Somit ist die Funktion $u-v$ gemäß i) nach Übergang zu einem geeigneten Repräsentanten konstant, und dies liefert die Behauptung; insbesondere ist

$$
u(x)-u(y)=v(x)-v(y)=\int_{x}^{y} u^{\prime}(t) d t \quad \text { für alle } x, y \in \bar{\Omega}
$$

Zu (iii): Die erste Ungleichung folgt direkt aus (7.3). Ist zudem $y \in \bar{\Omega}$ mit $|u(y)|= \min _{\bar{\Omega}}|u|$ gewählt, so folgt

$$
|u(x)| \leq|u(x)-u(y)|+|u(y)| \leq|x-y|^{\frac{1}{2}}\left\|u^{\prime}\right\|_{2} \leq \sqrt{2}\left\|u^{\prime}\right\|_{2}+\min _{\bar{\Omega}}|u|,
$$

wobei

$$
\left(\min _{\bar{\Omega}}|u|\right)^{2} \leq \frac{1}{2} \int_{\Omega}|u|^{2} d y=\frac{\|u\|_{2}^{2}}{2}
$$

gilt. Es folgt $\|u\|_{\infty} \leq \sqrt{2}\left\|u^{\prime}\right\|_{2}+\frac{1}{\sqrt{2}}\|u\|_{2} \leq \sqrt{2}\|u\|_{H^{1}}$, wie behauptet.

\subsection*{Bemerkung}
i). Aus 7.14 (iii) folgt, dass $H^{1}(\Omega)$ stetig in $C(\bar{\Omega})$ eingebettet ist. Die Einbettung ist sogar kompakt, da sie in der Form

$$
H^{1}(\Omega) \hookrightarrow C^{\frac{1}{2}}(\bar{\Omega}) \hookrightarrow C(\bar{\Omega})
$$

faktorisiert, wobei hier $C^{\frac{1}{2}}(\bar{\Omega})$ den Raum der $\frac{1}{2}$-Hölder-stetigen Funktionen auf $\bar{\Omega}$ bezeichne und die zweite Einbettung gemäß Übungsblatt 3, Aufg. 4, kompakt ist.\\
ii). Besitzt $u \in H^{1}(\Omega)$ eine schwache Ableitung $u^{\prime} \in C(\bar{\Omega})$, so ist $u \in C^{1}(\bar{\Omega})$ (nach Wahl eines geeigneten Repräsentanten). Dies folgt unmittelbar aus (7.2) und dem klassischen Hauptsatz der Integral- und Differentialrechnung.

\subsection*{Satz und Definition}
Seien Funktionen $a \in L^{\infty}(\Omega)$, $f \in L^{2}(\Omega)$ gegeben.\\
i). Wir nennen $u \in H^{1}(\Omega)$ eine schwache Lösung des Randwertproblems (NP), wenn

$$
\langle u, w\rangle_{*}=\int_{\Omega} f w d x \quad \text { für alle } w \in H^{1}(\Omega)
$$

gilt, wobei

$$
\langle u, w\rangle_{*}:=\int_{\Omega}\left(u^{\prime} w^{\prime}+a(x) u w\right) d x
$$

sei.\\
ii). Ist $a_{0}:=\underset{\Omega}{\operatorname{essinf}} a>0$, so ist durch

$$
(u, w) \mapsto\langle u, w\rangle_{*}
$$

ein Skalarprodukt definiert mit der Eigenschaft, dass die induzierte Norm $\|\cdot\|_{*}$ äquivalent zur Norm $\|\cdot\|_{H^{1}}$ ist. Insbesondere ist ( $H^{1}(\Omega),\langle\cdot, \cdot\rangle_{*}$ ) ebenfalls ein Hilbertraum.

Beweis von (ii). Offensichtlich ist $\langle\cdot, \cdot\rangle_{*}$ eine symmetrische Bilinearform auf $H^{1}(\Omega)$, und für $u \in H^{1}(\Omega)$ gilt

$$
a_{0}\|u\|_{2}^{2} \leq \int_{\Omega} a(x) u^{2} d x \leq\|a\|_{\infty}\|u\|_{2}^{2}
$$

also

$$
\min \left\{a_{0}, 1\right\}\|u\|_{H^{1}}^{2} \leq\langle u, u\rangle_{*} \leq \max \left\{\|a\|_{\infty}, 1\right\}\|u\|_{H^{1}}^{2} .
$$

Somit ist $\langle\cdot, \cdot\rangle_{*}$ ein Skalarprodukt derart, dass die induzierte Norm $\|\cdot\|_{*}$ äquivalent zur Norm $\|\cdot\|_{H^{1}}$ ist.

\subsection*{Satz}
Sind $a \in L^{\infty}(\Omega), f \in L^{2}(\Omega)$ mit $a_{0}:=\operatorname{essinf} a>0$ gegeben, so besitzt das Randwertproblem ( $N P$ ) genau eine schwache Lösung.

Beweis. Sei $H:=H^{1}(\Omega)$, versehen mit dem Skalarprodukt $\langle\cdot, \cdot\rangle_{*}$, und sei $\varphi \in H^{\prime}$ definiert durch $\varphi(w)=\int_{\Omega} f w d x$. Die Stetigkeit von $\varphi$ folgt aus der Abschätzung

$$
|\varphi(w)|=\left|\langle f, w\rangle_{L^{2}(\Omega)}\right| \leq\|f\|_{2}\|w\|_{2} \leq\|f\|_{2}\left(\frac{1}{a_{0}} \int_{\Omega} a(x) w^{2} d x\right)^{\frac{1}{2}} \leq \frac{1}{\sqrt{a_{0}}}\|f\|_{2}\|w\|_{*}
$$

Nach dem Rieszschen Darstellungssatz 7.10 existiert nun genau ein $u \in H^{1}(\Omega)$ mit

$$
\langle u, w\rangle_{*}=\varphi(w)=\int_{\Omega} f w d x \quad \text { für alle } w \in H^{1}(\Omega)
$$

Also ist $u \in H^{1}(\Omega)$ die gesuchte eindeutige schwache Lösung von (NP).

\subsection*{Bemerkung}
Seien Funktionen $a \in L^{\infty}(\Omega)$, $f \in L^{2}(\Omega)$ gegeben.\\
i). Ist $u$ eine schwache Lösung von ( $N P$ ), so folgt

$$
\int_{\Omega} u^{\prime} w^{\prime}=-\int_{\Omega}(a u-f) w d x \quad \text { für alle } w \in C^{1}(\bar{\Omega}) \subset H^{1}(\Omega)
$$

also insbesondere für $w \in C_{c}^{1}(\bar{\Omega})$. Somit ist

$$
u^{\prime} \in H^{1}(\Omega) \quad \text { mit schwacher Ableitung } \quad u^{\prime \prime}=a u-f .
$$

Insbesondere folgt mit 7.15(i) (nach Wahl eines geeigneten Repräsentanten) $u^{\prime} \in C(\bar{\Omega})$, und mit 7.15 (ii) folgt $u \in C^{1}(\bar{\Omega})$.\\
ii). Sind sogar $a, f \in C(\bar{\Omega})$, so folgt $u^{\prime \prime}=a u-f \in C(\bar{\Omega})$ für die schwache Ableitung von $u^{\prime}$, d.h. $u^{\prime} \in C^{1}(\bar{\Omega})$ gemäß 7.15 (ii). In diesem Fall ist also $u \in C^{2}(\bar{\Omega})$, und die Differentialgleichung $u^{\prime \prime}=a u-f$ ist im klassischen Sinne in $\Omega$ erfüllt. Aus $(*)$ folgt dann durch partielle Integration für alle $w \in C^{1}(\bar{\Omega})$ die Gleichung

$$
u^{\prime}(1) w(1)-u^{\prime}(-1) w(-1)=\int_{\Omega}\left(u^{\prime} w\right)^{\prime} d x=\int_{\Omega}\left(u^{\prime} w^{\prime}+u^{\prime \prime} w\right) d x \stackrel{(*)}{=} 0
$$

und dies liefert die Neumann-Randbedingungen

$$
(N R) \quad u^{\prime}(1)=u^{\prime}(-1)=0
$$

iii). Die Bedingungen ( $N R$ ) sind auch ohne die Zusatzvoraussetzung $a, f \in C(\bar{\Omega})$ erfüllt. Tatsächlich gilt (s. Übungsblatt 8):\\
Ist $\tilde{u} \in L^{2}(\Omega)$ und existiert $C>0$ mit

$$
\left|\int_{\Omega} \tilde{u} w^{\prime} d x\right| \leq C\|w\|_{2} \quad \text { für alle } w \in C^{1}(\bar{\Omega})
$$

so ist $\tilde{u} \in H^{1}(\Omega)$, und für den Repräsentanten $\tilde{u} \in C(\bar{\Omega})$ gilt $\tilde{u}(1)=\tilde{u}(-1)=0$.

\section*{Ende des Exkurses}
\section*{Orthonormalsysteme und abstrakte Fourierreihen}
Nun sei wieder $(H,\langle\cdot, \cdot\rangle)$ ein beliebiger Hilbertraum.

\subsection*{Definition}
i). Eine Menge $A \subset H$ heißt Orthonormalsystem (kurz: ONS), falls

$$
\langle x, y\rangle=\left\{\begin{array}{ll}
0 & x \neq y \\
1 & x=y
\end{array} \quad \text { für alle } x, y \in A .\right.
$$

ii). Ein totales Orthonormalsystem $A$ heißt Orthonormalbasis (kurz: ONB). Man beachte, dass dies nicht bedeutet, dass A eine Basis im Sinne der Linearen Algebra ist.

\subsection*{Satz}
Sei $A=\left\{e_{1}, e_{2}, \ldots\right\}$ ein abzählbares ONS in $H$. Dann gilt\\
i). Sind $c_{n} \in \mathbb{K}, n \in \mathbb{N}$ derart, dass $x=\sum_{n=1}^{\infty} c_{n} e_{n} \in H$ existiert, so gilt $c_{n}=\left\langle x, e_{n}\right\rangle$ für alle $n \in \mathbb{N}$.\\
ii). Ist $\left(c_{n}\right)_{n} \subset \ell^{2}$ gegeben, so konvergiert $\sum_{n=1}^{\infty} c_{n} e_{n}$ in $H$.

Beweis. i) $\operatorname{Zu} n \in \mathbb{N}$ sei $\varphi_{n} \in H^{\prime}$ definiert durch $\varphi_{n}(y)=\left\langle y, e_{n}\right\rangle$. Dann ist\\
$\left\langle x, e_{n}\right\rangle=\varphi_{n}(x)=\varphi_{n}\left(\lim _{k \rightarrow \infty} \sum_{j=1}^{k} c_{j} e_{j}\right)=\lim _{k \rightarrow \infty} \varphi_{n}\left(\sum_{j=1}^{k} c_{j} e_{j}\right)=\lim _{k \rightarrow \infty} \sum_{j=1}^{k} c_{j} \underbrace{\varphi_{n}\left(e_{j}\right)}_{=\left\langle e_{j}, e_{n}\right\rangle=\delta_{j_{n}}}=c_{n}$\\
ii) Sei $s_{n}:=\sum_{k=1}^{n} c_{k} e_{k}$. Dann gilt\\
$\left\|s_{m}-s_{n}\right\|^{2}=\left\|\sum_{k=n+1}^{m} c_{k} e_{k}\right\|^{2}=\sum_{i, j=n+1}^{m} c_{i} \overline{c_{j}}\langle\underbrace{e_{i}, e_{j}}_{\delta_{i j}}\rangle=\sum_{k=n+1}^{m}\left|c_{k}\right|^{2} \quad$ für alle $n, m \in \mathbb{N}, m>n$.\\
Da ferner $\left(c_{k}\right)_{k} \in \ell^{2}$ ist, folgt

$$
\sup _{m \geq n}\left\|s_{m}-s_{n}\right\|^{2} \leq \sum_{k=n+1}^{\infty}\left|c_{k}\right|^{2} \rightarrow 0 \quad \text { für } n \rightarrow \infty .
$$

Also ist $\left(s_{n}\right)_{n}$ eine Cauchyfolge in $H$, und somit folgt die Behauptung.

\subsection*{Satz}
Sei $A=\left\{e_{1}, \ldots, e_{n}\right\}$ ein endliches ONS, $V=\operatorname{span} A$ und $P \in \mathcal{L}(H)$ die orthogonale Projektion auf $V$. Dann gilt für alle $x, y \in H$ :\\
i) $P x=\sum_{k=1}^{n}\left\langle x, e_{k}\right\rangle e_{k}$.\\
ii) $\langle P x, P y\rangle=\sum_{k=1}^{n}\left\langle x, e_{k}\right\rangle \overline{\left\langle y, e_{k}\right\rangle}$.

Beweis. i) Sei $z:=\sum_{k=1}^{n}\left\langle x, e_{k}\right\rangle e_{k}$. Dann ist\\
$\left\langle x-z, e_{j}\right\rangle=\left\langle x, e_{j}\right\rangle-\sum_{k=1}^{n}\left\langle x, e_{k}\right\rangle\left\langle e_{k}, e_{j}\right\rangle=\left\langle x, e_{j}\right\rangle-\left\langle x, e_{j}\right\rangle=0 \quad$ für alle $j \in\{1, \ldots, n\}$.

Also ist $x-z \in A^{\perp}=V^{\perp}=\operatorname{Kern} P$. Es folgt $P x=P z+P(x-z)=P z \stackrel{z \in V}{=} z$.\\
ii) Es ist\\
$\langle P x, P y\rangle=\left\langle\sum_{k=1}^{n}\left\langle x, e_{k}\right\rangle e_{k}, \sum_{j=1}^{n}\left\langle y, e_{j}\right\rangle e_{j}\right\rangle=\sum_{k, j=1}^{n}\left\langle x, e_{k}\right\rangle \overline{\left\langle y, e_{j}\right\rangle}\left\langle e_{k}, e_{j}\right\rangle=\sum_{k=1}^{n}\left\langle x, e_{k}\right\rangle \overline{\left\langle y, e_{k}\right\rangle}$. \(\square\)

\subsection*{Definition}
Für ein fest gewähltes, abzählbares ONS $A=\left\{e_{1}, e_{2}, \ldots\right\}$ in $H$ und $x \in H$ nennen wir\\
i) $\hat{x}(n):=\left\langle x, e_{n}\right\rangle$ den $n$-ten Fourierkoeffizient von $x$ (bzgl. $A$ ).\\
ii) $\sum_{n} \hat{x}(n) e_{n}$ die (abstrakte) Fourierreihe von $x$ (bzgl. $A$ ).

\subsection*{Satz}
Sei $A=\left\{e_{1}, e_{2}, \ldots\right\}$ ein abzählbar unendliches ONS in $H$. Sei ferner $V=\overline{\operatorname{span} A}$ und $P \in \mathcal{L}(H)$ die orthogonale Projektion auf $V$. Dann gilt für alle $x, y \in H$ :\\
i) $P x=\sum_{n=1}^{\infty} \hat{x}(n) e_{n}$\\
ii) $\langle P x, P y\rangle=\sum_{n=1}^{\infty} \hat{x}(n) \overline{\hat{y}(n)}$\\
iii) $\sum_{n=1}^{\infty}|\hat{x}(n)|^{2} \leq\|x\|^{2} \quad$ (Besselsche Ungleichung)

Beweis. Für $n \in \mathbb{N}$ sei $P_{n} \in \mathcal{L}(H)$ die orthogonale Projektion auf $\operatorname{span}\left\{e_{1}, \ldots, e_{n}\right\}$. Dann gilt

$$
\sum_{k=1}^{n}|\hat{x}(k)|^{2} \stackrel{\text { [7.21] }}{=}\left\|P_{n} x\right\|^{2} \stackrel{\text { [.7] } i i)}{=}\|x\|^{2} \quad \text { für alle } n \in \mathbb{N} \text {. }
$$

Es folgt iii), also ist $(\hat{x}(n))_{n} \in \ell^{2}$.\\
Gemäß 7.20ii) existiert somit $z:=\sum_{n=1}^{\infty} \hat{x}(n) e_{n} \in H$ (d.h. die Reihe konvergiert). Aus der Stetigkeit des Skalarprodukts folgt

$$
\left\langle x-z, e_{j}\right\rangle=\hat{x}(j)-\sum_{n=1}^{\infty} \hat{x}(n)\left\langle e_{n}, e_{j}\right\rangle=\hat{x}(j)-\hat{x}(j)=0 \quad \text { für alle } j \in \mathbb{N},
$$

also $x-z \in A^{\perp}=V^{\perp}$ und somit $P x=P z+P(x-z)=P z=z$. Dies liefert i).\\
Zu ii):

$$
\langle P x, P y\rangle=\left\langle\sum_{n=1}^{\infty} \hat{x}(n) e_{n}, \sum_{k=1}^{\infty} \hat{y}(k) e_{k}\right\rangle=\sum_{n, k=1}^{\infty} \hat{x}(n) \overline{\hat{y}(k)}\left\langle e_{n}, e_{k}\right\rangle=\sum_{n=1}^{\infty} \hat{x}(n) \overline{\hat{y}(n)}
$$

Hier haben wir wieder die Stetigkeit des Skalarprodukts verwendet.

\subsection*{Korollar}
Sei $A=\left\{e_{1}, e_{2}, \ldots\right\}$ ein abzählbares ONS in $H$. Dann sind äquivalent:\\
i) $A$ ist eine ONB .\\
ii) Für alle $x \in H$ gilt $x=\sum_{k=1}^{\infty} \hat{x}(k) e_{k}$.\\
iii) Für alle $x \in H$ gilt $\|x\|^{2}=\sum_{k=1}^{\infty}|\hat{x}(k)|^{2}$\\
iv) Für alle $x, y \in H$ gilt $\langle x, y\rangle=\sum_{k=1}^{\infty} \hat{x}(k) \overline{\hat{y}(k)}$ (Parsevalsche Gleichung)

Beweis. Sei $V=\overline{\operatorname{span} A}$ und $P \in \mathcal{L}(H)$ die orthogonale Projektion auf $V$. Dann gilt:

$$
A \text { total } \Longleftrightarrow V=H \quad \Longleftrightarrow \quad P=\mathrm{id}
$$

Somit erhalten wir die Implikationen:

$$
i) \stackrel{\text { 7.23 }}{\Longrightarrow} i i) \stackrel{\text { Stet.d.SKP }}{\Longrightarrow} i v) \stackrel{\text { trivial }}{\Longrightarrow} i i i) \text {. }
$$

Gilt iii), so folgt für alle $x \in H$ :

$$
\|P x\|^{2}+\|x-P x\|^{2} \stackrel{\text { (7.6) }}{=}\|x\|^{2}=\sum_{k=1}^{\infty}|\hat{x}(k)|^{2} \stackrel{\text { (7.23) }}{=}\|P x\|^{2}
$$

also $x=P x \in V$. Also ist $H=V$ und somit $A$ total, d.h. i) gilt.

\subsection*{Bemerkung}
Ist $A$ ein abzählbares ONS und existiert eine totale Teilmenge $M \subset H$ mit

$$
\|x\|^{2}=\sum_{k=1}^{\infty}|\hat{x}(k)|^{2} \quad \text { für alle } x \in M
$$

so ist $A$ bereits eine ONB.

Beweis. Sei wieder $V=\overline{\operatorname{span} A}$. Wie im Beweisschritt iii) auf i) in 7.24 erhält man mit (*) die Implikation: $x \in M \Longrightarrow x \in V$.\\
Somit $M \subset V$ und damit $H=\overline{\operatorname{span} M} \subset V$. Es folgt, dass $A$ eine ONB ist.

\subsection*{Bemerkung}
Sei $A \subset H$ eine abzählbare ONB. Dann ist die Abbildung $\mathcal{F}: H \rightarrow \ell^{2}, \mathcal{F} x:=(\hat{x}(k))_{k}$ ein isometrischer Isomorphismus, d. h. $\mathcal{F}$ ist bijektiv und es gilt

$$
\|\mathcal{F} x\|_{2}=\|x\|_{H} \quad \text { für alle } x \in H . \quad(*)
$$

Beweis. Die Linearität von $\mathcal{F}$ ist klar nach Definition, und aus 7.24ii) folgt (*). Insbesondere ist $\mathcal{F}$ injektiv. Die Surjektivität folgt aus 7.20i).

\subsection*{Satz}
Äquivalent sind:\\
i) $H$ ist unendlich-dimensional und separabel\\
ii) $H$ besitzt eine abzählbar unendliche ONB\\
iii) $H$ ist isometrisch isomorph $\mathrm{zu} \ell^{2}$

Man beachte: Die Aussagen gelten insbesondere im Fall $H=L^{2}(\Omega)$, wobei $\Omega \subset \mathbb{R}^{N}$ eine Lebesgue-messbare Menge mit $\lambda(\Omega)>0$ sei (vgl. 6.26).

Beweis. Aus ii) folgt iii) mit 7.26. Aus iii) folgt i), da $\ell^{2}$ separabel ist. Noch zu zeigen: i) $\Longrightarrow$ ii).\\
Sei dazu $M=\left\{a_{1}, a_{2}, \ldots\right\} \subset H$ eine abzählbare totale Menge. Ohne Einschränkung gelte $a_{1} \neq 0$ und $a_{k+1} \notin W_{k}:=\operatorname{span}\left\{a_{1}, \ldots, a_{k}\right\}$ für $k \in \mathbb{N}$. Sei $P_{k} \in \mathcal{L}(H)$ die orthogonale Projektion auf $W_{k}$ für $k \in \mathbb{N}$. Wir setzen $e_{1}:=\frac{a_{1}}{\left\|a_{1}\right\|}$ und

$$
e_{k}:=\frac{a_{k}-P_{k-1} a_{k}}{\left\|a_{k}-P_{k-1} a_{k}\right\|} \in W_{k} \cap W_{k-1}^{\perp} \quad \text { für } k \geq 2 .
$$

Nach Konstruktion ist $A:=\left\{e_{1}, e_{2}, \ldots\right\}$ ein ONS in $H$. Per Induktion sieht man ferner: $W_{k}=\operatorname{span}\left\{e_{1}, \ldots, e_{k}\right\}$ für alle $k \in \mathbb{N}$. Es folgt

$$
\operatorname{span} M=\bigcup_{k \in \mathbb{N}} W_{k}=\operatorname{span} A
$$

also ist $A$ total und somit eine ONB. Es folgt ii).

\section*{Klassische Fourierreihen}
Im Folgenden betrachten wir den Hilbertraum $H:=L^{2}[0,2 \pi]=L^{2}([0,2 \pi], \mathbb{C})$ (bezüglich des Lebesguemaßes) mit Skalarprodukt $\langle\cdot, \cdot\rangle$ definiert durch

$$
\langle f, g\rangle=\int_{0}^{2 \pi} f(t) \overline{g(t)} d t \quad \text { für } f, g \in H
$$

\subsection*{Definition und Satz}
Für $k \in \mathbb{Z}$ sei $e_{k} \in L^{2}[0,2 \pi]$ definiert durch $e_{k}(t)=\frac{1}{\sqrt{2 \pi}} e^{i k t}$.\\
Dann ist $A=\left\{e_{k} \mid k \in \mathbb{Z}\right\}$ eine ONB von $L^{2}[0,2 \pi]$.\\
Die Fourierkoeffizienten von $f \in L^{2}[0,2 \pi]$ bzgl. $A$ sind gegeben durch

$$
\hat{f}(k)=\frac{1}{\sqrt{2 \pi}} \int_{0}^{2 \pi} f(t) e^{-i k t} d t, \quad \text { für } k \in \mathbb{Z}
$$

Beweis. Es ist

$$
\left\langle e_{k}, e_{k}\right\rangle=\frac{1}{2 \pi} \int_{0}^{2 \pi} \underbrace{e^{i k t} e^{-i k t}}_{=1} d t=1 \quad \text { für alle } k \in \mathbb{Z}
$$

und

$$
\left\langle e_{k}, e_{j}\right\rangle=\frac{1}{2 \pi} \int_{0}^{2 \pi} e^{i(k-j) t} d t=\left.\frac{1}{2 \pi} \frac{1}{i(k-j)} e^{i(k-j) t}\right|_{0} ^{2 \pi}=0 \quad \text { für } k \neq j .
$$

Es folgt, dass $A$ ein ONS in $H=L^{2}([0,2 \pi])$ ist. Ferner ist $A$ total in $H$, da die Menge der trigonometrischen Polynome dicht in $H$ liegt. Die letztere Aussage wurde in Übungsaufgabe 1(ii), Blatt 7 gezeigt; es handelt sich um eine leichte Folgerung aus 4.27 und 6.25.

\subsection*{Bemerkung}
Aus 7.28 folgt, dass für jedes $f \in L^{2}[0,2 \pi]$ die Fourierreihe $\sum_{k \in \mathbb{Z}} \hat{f}(k) e_{k}$ bzgl. $\|\cdot\|_{2}$ gegen $f$ konvergiert. Weiterhin gilt:\\
i) Die Fourierreihe konvergiert f.ü. punktweise gegen $f$ (Satz von Carleson 1966).\\
ii) Ist $(\hat{f}(k))_{k} \in \ell^{1}(\mathbb{Z})$, so konvergiert die Fourierreihe in ( $\mathcal{C}_{2 \pi}(\mathbb{R}, \mathbb{C}),\|\cdot\|_{\infty}$ ), also gleichmäßig. Es folgt dann, dass $f$ (nach Wahl eines geeigneten Repräsentanten) zu einem Element in $\mathcal{C}_{2 \pi}(\mathbb{R}, \mathbb{C})$ fortgesetzt werden kann.\\
iii) Ist $f \in \mathcal{C}_{2 \pi}(\mathbb{R}, \mathbb{C})$, so muss die Fourierreihe i.A. trotzdem nicht punktweise gegen $f$ konvergieren.

Beweisskizze hierzu. Sei $t \in[0,2 \pi]$, und sei im Folgenden $\mathcal{C}_{2 \pi}:=\mathcal{C}_{2 \pi}(\mathbb{R}, \mathbb{C})$. Wir zeigen, dass ein $f \in \mathcal{C}_{2 \pi}$ existiert derart, dass die Fourierreihe von $f$ in\\
$t$ nicht gegen $f$ konvergiert. Wir verwenden den Satz von der gleichmäßigen Beschränktheit. Dazu definieren wir die Operatoren

$$
S_{n} \in \mathcal{L}\left(\mathcal{C}_{2 \pi}\right), \quad S_{n} f=\sum_{k=-n}^{n}\left\langle f, e_{k}\right\rangle e_{k}
$$

mit $e_{k}$ wie in 7.28. Dann gilt\\
$\left[S_{n} f\right](t)=\frac{1}{2 \pi} \int_{0}^{2 \pi} f(s)\left(\sum_{k=-n}^{n} e^{i k(t-s)}\right) d s=\frac{1}{2 \pi} \int_{0}^{2 \pi} d_{n}(t-s) f(s) d s \quad$ für $t \in[0,2 \pi]$\\
mit dem $n$-ten Dirichletkern

$$
\tau \mapsto d_{n}(\tau)=\sum_{k=-n}^{n} e^{i k \tau}= \begin{cases}\frac{\sin \left(\left(n+\frac{1}{2}\right) \tau\right)}{\sin \frac{\tau}{2}}, & \tau \in \mathbb{R} \backslash\{k 2 \pi: k \in \mathbb{Z}\} \\ 2 n+1, & \tau \in\{k 2 \pi: k \in \mathbb{Z}\}\end{cases}
$$

Die letztere Darstellung des Kerns kann man dabei z.B. mit der geometrischen Summenformel herleiten (Übungsaufgabe, Blatt 9). Wir betrachten nun die Linearformen

$$
T_{n} \in \mathcal{C}_{2 \pi}^{\prime}, \quad T_{n}(f)=\left[S_{n} f\right](t)=\frac{1}{2 \pi} \int_{0}^{2 \pi} d_{n}(t-s) f(s) d s
$$

Ähnlich wie in 2.14(iii) sieht man, dass für die zugehörige Operatornorm dann

$$
\left\|T_{n}\right\|=\int_{0}^{2 \pi}\left|d_{n}(\tau)\right| d \tau
$$

gilt (Übungsaufgabe, Blatt 9). Allerdings erhält man mit (7.4) auch

$$
\left\|T_{n}\right\|=\int_{0}^{2 \pi}\left|d_{n}(\tau)\right| d \tau \rightarrow \infty \quad \text { für } n \rightarrow \infty
$$

(Übungaufgabe, Blatt 9). Somit ist die Folge der Linearformen $T_{n}, n \in \mathbb{N}$ nicht beschränkt. Nach 5.2 ist sie damit auch nicht punktweise beschränkt, d.h. es existiert ein $f \in \mathcal{C}_{2 \pi}$ derart, dass die im Punkt $t$ ausgewerteten Partialsummen

$$
T_{n} f=\left[S_{n} f\right](t), \quad n \in \mathbb{N}
$$

der Fourierreihe von $f$ eine unbeschränkte Folge in $\mathbb{C}$ bilden und somit nicht konvergieren können.

\section*{§8 Der Satz von Lax-Milgram und der adjungierte Operator}
Stets sei $(H,\langle\cdot, \cdot\rangle)$ ein Hilbertraum über $\mathbb{K}$.

\subsection*{Definition}
Eine Abbildung $a: H \times H \rightarrow \mathbb{K}$ heißt sesquilinear, falls gilt:\\
1.) $a(\cdot, y): H \rightarrow \mathbb{K}$ ist linear für alle $y \in H$.\\
2.) $a(x, \cdot): H \rightarrow \mathbb{K}$ ist antilinear für alle $x \in H$, d. h. $a(x, \lambda y+z)=\bar{\lambda} a(x, y)+ a(x, z)$ für alle $x, y, z \in H, \lambda \in \mathbb{K}$.

Wir setzen $S_{q} \ell(H)=\{a: H \times H \rightarrow \mathbb{K}: a$ sesquilinear $\}$.

\subsection*{Satz und Definition}
Zu $T \in \mathcal{L}(H)$ sei $a_{T} \in S_{q} \ell(H)$ definiert durch $a_{T}(x, y)=\langle x, T y\rangle$ für alle $x, y \in H$. Dann gilt:\\
i). $\left|a_{T}(x, y)\right| \leq\|T\|\|x\|\|y\|$ für alle $x, y \in H$.\\
ii). Existiert $c_{1}>0$ mit $\left|a_{T}(x, x)\right| \geq c_{1}\|x\|^{2}$ für alle $x \in H$, so ist $T$ ein topologischer Isomorphismus.

Beweis. i) $\left|a_{T}(x, y)\right|=|\langle x, T y\rangle| \leq\|x\|\|T y\| \leq\|x\|\|T\|\|y\|$ für alle $x, y \in H$.\\
ii) Für alle $x \in H$ ist $c_{1}\|x\|^{2} \leq|\langle x, T x\rangle| \leq\|T x\|\|x\|$, also

$$
\|x\| \leq \frac{1}{c_{1}}\|T x\|
$$

Somit ist $T$ injektiv. Sei $V:=$ Bild $T$.\\
Beh. 1: $V$ ist abgeschlossen.\\
Sei dazu $\left(x_{n}\right)_{n} \subset H$ und $y \in H$ mit $T x_{n} \rightarrow y$ für $n \rightarrow \infty$. Dann gilt

$$
\sup _{m \geq n}\left\|x_{m}-x_{n}\right\| \stackrel{(*)}{\leq} \frac{1}{c_{1}} \sup _{m \geq n}\left\|T x_{m}-T x_{n}\right\| \rightarrow 0 \quad \text { für } n \rightarrow \infty
$$

Also ist $\left(x_{n}\right)_{n}$ eine Cauchyfolge, und somit existiert $x:=\lim _{n \rightarrow \infty} x_{n} \in H$. Da $T$ stetig ist, gilt $y=T x \in V$. Somit folgt Beh. 1.\\
Für $y \in V^{\perp}$ gilt nun $c_{1}\|y\|^{2} \leq|\langle y, T y\rangle|=0$, also $V^{\perp}=\{0\}$. Es folgt $V=\bar{V}= V^{\perp \perp}=\{0\}^{\perp}=H$, d. h. $T$ ist surjektiv.\\
Schließlich gilt $\left\|T^{-1} y\right\| \stackrel{(*)}{\leq} \frac{1}{c_{1}}\left\|T\left(T^{-1} y\right)\right\|=\frac{1}{c_{1}}\|y\|$ für alle $y \in H \Longrightarrow T^{-1} \in \mathcal{L}(H)$, $\left\|T^{-1}\right\| \leq \frac{1}{c_{1}}$. \(\square\)

\subsection*{Satz (von Lax-Milgram)}
Sei $a \in S_{q} \ell(H)$ und $c_{0}>0$ derart, das\\
(i) $|a(x, y)| \leq c_{0}\|x\|\|y\|$ für alle $x, y \in H$.

Dann gilt:\\
(a) Es existiert genau ein $T \in \mathcal{L}(H)$ mit $a=a_{T}$, d.h. mit $a(x, y)=\langle x, T y\rangle$ für alle $x, y \in H$.\\
(b) Gilt ferner\\
(ii) $|a(x, x)| \geq c_{1}\|x\|^{2}$ für alle $x \in H$ mit einem festen $c_{1}>0$,\\
so ist $T$ ein topologischer Isomorphismus.\\
(c) Sind (i) und (ii) erfüllt, so existiert zu jedem $f \in H$ genau ein $x \in H$ mit

$$
\text { (*) } \quad a(y, x)=\langle y, f\rangle \quad \text { für alle } y \in H .
$$

Beweis. (a) Für $x \in H$ sei $\psi_{x}: H \rightarrow \mathbb{K}$ definiert durch $\psi_{x}(y)=a(y, x)$. Dann ist $\psi_{x}$ linear und stetig mit $\left\|\psi_{x}\right\| \leq c_{0}\|x\|$, denn es gilt $\left|\psi_{x}(y)\right| \stackrel{(i)}{\leq} c_{0}\|y\|\|x\|$ für alle $y \in H$. Gilt $a=a_{T}$ für ein $T \in \mathcal{L}(H)$, so folgt also

$$
\langle y, T x\rangle=a_{T}(y, x)=\psi_{x}(y) \quad \text { für alle } x, y \in H,
$$

also

$$
\text { (**) } \quad T x=J^{-1} \psi_{x} \quad \text { für alle } x \in H \quad \text { (vgl. 7.10) }
$$

Insbesondere ist $T$ eindeutig bestimmt. Ist $T$ durch ( $* *$ ) definiert, so gilt für $x_{1}, x_{2} \in H, \lambda \in \mathbb{K}$ :\\
$T\left(\lambda x_{1}+x_{2}\right)=J^{-1} \psi_{\lambda x_{1}+x_{2}} \stackrel{a \in S_{\underline{q}} \ell(H)}{=} J^{-1}\left(\bar{\lambda} \psi_{x_{1}}+\psi_{x_{2}}\right) \stackrel{7.9}{=} \lambda J^{-1} \psi_{x_{1}}+J^{-1} \psi_{x_{2}}=\lambda T x_{1}+T x_{2}$.

Also ist $T$ ist linear. Ferner gilt

$$
\|T x\|=\left\|J^{-1} \psi_{x}\right\| \stackrel{\sqrt{7.10}}{=}\left\|\psi_{x}\right\| \leq c_{0}\|x\| \quad \text { für alle } x \in H
$$

Somit ist $T$ stetig, also $T \in \mathcal{L}(H)$.\\
(b) Dies folgt direkt aus (a) und 8.2 ii).\\
(c) Ist $T \in \mathcal{L}(H)$ gemäß (a) gewählt, so ist (*) äquivalent zu

$$
(* * *) \quad\langle y, T x\rangle=\langle y, f\rangle \quad \text { für alle } y \in H
$$

wobei $T$ gemäß (b) ein topologischer Isomorphismus ist. Also existiert genau ein $x \in H$, welches ( $* * *$ ) erfüllt, und dieses ist gegeben durch $x:=T^{-1} f$.

\subsection*{Bemerkungen und Beispiele}
(i) Unter den Voraussetzung (i) und (ii) von 8.3 existiert zu jedem $\varphi \in H^{\prime}$ genau ein $x \in H$ mit

$$
a(y, x)=\varphi(y) \quad \text { für alle } y \in H .
$$

Dies folgt direkt aus 8.3 und dem Rieszschen Darstellungssatz 7.10. Der Vorteil dieses Resultats gegenüber dem Rieszschen Darstellungssatz liegt darin, dass hier keine Symmetrie der Sesquilinearform $a$ verlangt wird.\\
(ii) Sei $\mathbb{K}=\mathbb{R}$ und $\Omega=(-1,1) \subset \mathbb{R}$. Gegeben seien Funktionen $h, b \in L^{\infty}(\Omega)$, $f \in L^{2}(\Omega)$, wobei

$$
\underset{\Omega}{\operatorname{essinf}} h \geq 1>\frac{\|b\|_{\infty}}{2}
$$

sei. Wir suchen eine Lösung des Neumann-Randwertproblems

$$
(N P)_{1} \quad\left\{\begin{aligned}
-u^{\prime \prime}+b(x) u^{\prime}+h(x) u & =f(x), \quad x \in \Omega \\
u^{\prime}(1)=u^{\prime}(-1) & =0
\end{aligned}\right.
$$

Ähnlich wie in 7.16 nennen wir $u \in H^{1}(\Omega)$ eine schwache Lösung von $(N P)_{1}$, wenn gilt:

$$
a(w, u):=\int_{\Omega}\left(u^{\prime} w^{\prime}+b u^{\prime} w+h u w\right) d x=\int_{\Omega} f w \quad \text { für alle } w \in H^{1}(\Omega)
$$

Dabei ist $a: H^{1}(\Omega) \times H^{1}(\Omega) \rightarrow \mathbb{R}$ eine Bilinearform (wegen $\mathbb{K}=\mathbb{R}$ also auch eine Sesquilinearform) mit\\
$|a(u, w)| \leq\left\|u^{\prime}\right\|_{2}\|w\|_{2}+\|b\|_{\infty}\left\|u^{\prime}\right\|_{2}\|w\|_{2}+\|h\|_{\infty}\|u\|_{2}\|w\|_{2} \leq\left(1+\|b\|_{\infty}+\|h\|_{\infty}\right)\|u\|_{H^{1}}\|w\|_{H^{1}}$\\
für alle $u, w \in H^{1}(\Omega)$; ferner gilt aufgrund der Voraussetzung ( $V$ )\\
$|a(u, u)| \geq\left\|u^{\prime}\right\|_{2}^{2}-\|b\|_{\infty}\left\|u^{\prime}\right\|\|u\|+\left(\operatorname{essinf}_{\Omega} h\right)\|u\|_{2}^{2} \geq\left(1-\frac{\|b\|_{\infty}}{2}\right)\left(\left\|u^{\prime}\right\|_{2}^{2}+\|u\|_{2}^{2}\right)=c_{1}\|u\|_{H^{1}}^{2}$\\
für $u \in H^{1}(\Omega)$ mit $c_{1}:=\left(1-\frac{\|b\|_{\infty}}{2}\right)>0$. Somit erfüllt die Bilinearform $a$ auf $H^{1}(\Omega)$ die Voraussetzungen des Satzes von Lax-Milgram. Ähnlich wie im Beweis von 7.17 betrachten wir nun die Linearform

$$
\varphi \in H^{1}(\Omega)^{\prime}, \quad \varphi(w)=\int_{\Omega} f w d x
$$

Wie in (i) bemerkt, existiert genau ein $u \in H^{1}(\Omega)$ mit $a(w, u)=\varphi(w)$ für alle $w \in H^{1}(\Omega)$, d.h. genau eine schwache Lösung $u \in H^{1}(\Omega)$ von $(N P)_{1}$. Für diese schwache Lösung kann man analoge Regularitätsaussagen wie in 7.18 folgern.\\
Man beachte: Die hier definierte Bilinearform a ist nicht symmetrisch; man kann die (eindeutige) Existenz schwacher Lösungen daher nicht allein aus dem Rieszschen Darstellungssatz folgern.\\
(iii) Selbst das einfachere Problem

$$
(N P) \quad\left\{\begin{aligned}
-u^{\prime \prime}+h(x) u & =f(x), \quad x \in \Omega \\
u^{\prime}(1)=u^{\prime}(-1) & =0
\end{aligned}\right.
$$

(vgl. Kapitel 7) ohne den Term erster Ordnung führt auf eine nicht symmetrische Bilinearform, wenn man eine komplexwertige Funktion $x \mapsto h(x)$ betrachtet. Analog wie in 7.17 kann man dann die Existenz genau einer schwachen Lösung aus dem Satz von Lax-Milgram folgern, wenn $\underset{\Omega}{\operatorname{essinf}} \operatorname{Re} h>0$ gilt. Hierzu betrachtet man $L^{2}(\Omega)=L^{2}(\Omega, \mathbb{C})$ und einen analog definierten komplexwertigen Sobolevraum $H^{1}(\Omega)=H^{1}(\Omega, \mathbb{C})$.\\
(iv) Das folgende Beispiel wurde nicht in der Vorlesung behandelt und ist somit nicht Teil des Prüfungsstoffs!\\
Seien $h, f \in C([0,1], \mathbb{C})$ Funktionen mit $\operatorname{Re} h>0$ auf $[0,1]$. Wir zeigen, dass das Dirichletproblem

$$
-u^{\prime \prime}+h u=f \quad \text { in }(0,1), \quad u(0)=u(1)=0
$$

eine Lösung $u \in C^{2}([0,1], \mathbb{C})$ besitzt. Dazu schreiben wir das Problem um in der Form

$$
u=K(f-h u)
$$

mit dem Integraloperator

$$
K: H \rightarrow H, \quad[K v](t)=\int_{0}^{1} G(t, s) v(s) d y \quad \text { für } t \in[0,1]
$$

zur Greenfunktion

$$
G \in C\left([0,1]^{2}\right), \quad G(t, s)= \begin{cases}(1-s) t, & t \leq s \\ (1-t) s, & t>s\end{cases}
$$

aus Kapitel 10 . Ferner setzen wir $H:=L^{2}([0,1], \mathbb{C})$ und definieren $a \in S_{q} \ell(H)$ durch

$$
a(w, v)=\int_{0}^{1} w \overline{\left(K v+\frac{v}{h}\right)} d t \quad \text { für } v, w \in H
$$

Aufgrund der Beschränktheit der Funktionen $G$ und $\frac{1}{h}$ sieht man leicht, dass dann die Voraussetzung (i) aus 8.3 erfüllt ist. Betrachtet man ferner $v \in C_{c}(0,1)$ und setzt $w=K v$, so ist $w \in C^{2}[0,1]$ eine Lösung von

$$
-w^{\prime \prime}=v \quad \text { in }(0,1), \quad w(0)=w(1)=0
$$

(vgl. Übungsblatt 1) und somit folgt mit partieller Integration

$$
\int_{0}^{1} v \overline{[K v]} d t=-\int_{0}^{1} w^{\prime \prime} \bar{w} d t=\int_{0}^{1} w^{\prime} \bar{w}^{\prime} d t=\int_{0}^{1}\left|w^{\prime}\right|^{2} d t \geq 0
$$

also

$$
\operatorname{Re} a(v, v) \geq \operatorname{Re} \int_{0}^{1} v \overline{\left(\frac{v}{h}\right)} d t \geq h_{0} \int_{0}^{1}|v|^{2} d x=h_{0}\|v\|_{2}^{2}
$$

mit

$$
h_{0}:=\min _{t \in[0,1]} \operatorname{Re} \frac{1}{h(t)}=\min _{t \in[0,1]} \operatorname{Re} \frac{\overline{h(t)}}{|h(t)|^{2}}=\min _{t \in[0,1]} \frac{\operatorname{Re} h(t)}{|h(t)|^{2}}>0
$$

nach Voraussetzung. Da $C_{c}(0,1)$ gemäß 6.25 in $H$ dicht liegt, erhält man durch Approximation auch

$$
|a(v, v)| \geq \operatorname{Re} a(v, v) \geq h_{0}\|v\|_{2}^{2} \quad \text { für alle } v \in H .
$$

Also erfüllt $a$ die Voraussetzungen (i) und (ii) von 8.3, und somit existiert gemäß 8.3(c) zu $\tilde{f}=K f \in H$ genau ein $\tilde{u} \in H$ mit

$$
a(w, \tilde{u})=\langle w, \tilde{f}\rangle \quad \text { für alle } w \in H,
$$

d.h.

$$
\int_{0}^{1} w \overline{\left.w\left(K \tilde{u}+\frac{\tilde{u}}{h}-K f\right]\right)} d t \quad \text { für } w \in H
$$

Dies liefert $K \tilde{u}+\frac{\tilde{u}}{h}-K f \in H^{\perp}=\{0\}$, d.h. $u:=\frac{\tilde{u}}{h}$ erfüllt (8.2). Aus Eigenschaften der Greenfunktion $G$ (vgl. Übungsblatt 1) erhält man dann, dass $u \in C^{2}([0,1])$ eine Lösung von (8.1) ist.

\subsection*{Satz und Definition}
Zu jedem $T \in \mathcal{L}(H)$ existiert genau ein $T^{*} \in \mathcal{L}(H)$ mit $\langle T x, y\rangle=\left\langle x, T^{*} y\right\rangle$ für alle $x, y \in H$. Man nennt $T^{*}$ den $z u T$ adjungierten Operator.

Beweis. Definiere $a \in S_{q} \ell(H)$ durch $a(x, y)=\langle T x, y\rangle$. Dann ist

$$
|a(x, y)| \leq\|T x\|\|y\| \leq\|T\|\|x\|\|y\| \quad \text { für } x, y \in H,
$$

also gilt 8.3 (i) mit $c_{0}=\|T\|$. Somit existiert genau ein $T^{*} \in \mathcal{L}(H)$ mit

$$
\langle T x, y\rangle=a(x, y)=\left\langle x, T^{*} y\right\rangle \quad \text { für alle } x, y \in H .
$$

\subsection*{Satz}
Für alle $S, T \in \mathcal{L}(H)$ gilt:\\
i) $T^{* *}=T$\\
ii) $(\lambda S+\mu T)^{*}=\bar{\lambda} S^{*}+\bar{\mu} T^{*}$ für $\lambda, \mu \in \mathbb{K}$.\\
iii) $(S T)^{*}=T^{*} S^{*}$\\
iv) Ist $T$ ein topologischer Isomorphismus, so auch $T^{*}$; und dann gilt $\left(T^{*}\right)^{-1}= \left(T^{-1}\right)^{*}$.\\
v) $\left\|T^{*}\right\|=\|T\|$\\
vi) $\left\|T^{*} T\right\|=\|T\|^{2}$

Beweis. i) - iii) Leichte Übung, z.B. iii): Für alle $x, y \in H$ ist $\langle S T x, y\rangle=\left\langle T x, S^{*} y\right\rangle= \left\langle x, T^{*} S^{*} y\right\rangle$, und mit 8.5 (Eindeutigkeit) folgt $(S T)^{*}=T^{*} S^{*}$.\\
Zu iv): Gemäß (iii) ist $\left(T^{-1}\right)^{*} T^{*}=\left(T T^{-1}\right)^{*}=\operatorname{id}_{H}^{*}=\operatorname{id}_{H}$ und analog $T^{*}\left(T^{-1}\right)^{*}= \mathrm{id}_{H}$. Somit folgt die Behauptung.\\
Zu v ) und vi): Für alle $x \in H$ ist

$$
\|T x\|^{2}=\left\langle x, T^{*} T x\right\rangle \leq\left\|T^{*} T x\right\|\|x\| \leq\left\|T^{*} T\right\|\|x\|^{2},
$$

also

$$
\|T\|^{2}=\left(\sup _{\|x\| \leq 1}\|T x\|\right)^{2}=\sup _{\|x\| \leq 1}\|T x\|^{2} \leq\left\|T^{*} T\right\| \leq\left\|T^{*}\right\|\|T\| .
$$

Es folgt $\|T\| \leq\left\|T^{*}\right\|$, und dann auch $\left\|T^{*}\right\| \leq\left\|T^{* *}\right\|=\|T\|$. Somit ist $\|T\|=\left\|T^{*}\right\|$, und aus (*) folgt nun $\left\|T^{*} T\right\|=\|T\|^{2}$.

\subsection*{Definition}
$T \in \mathcal{L}(H)$ heißt

\begin{itemize}
  \item normal, falls $T T^{*}=T^{*} T$ gilt.
  \item hermitesch oder selbstadjungiert, falls $T=T^{*}$ gilt.
  \item unitär, falls $T T^{*}=\mathrm{id}=T^{*} T$, also $T^{*}=T^{-1}$ gilt.
\end{itemize}

\subsection*{Beispiel (Multiplikationsoperatoren)}
Sei $\Omega \subset \mathbb{R}^{N}$ (Lebesgue-)messbar, $H=L^{2}(\Omega), q \in L^{\infty}(\Omega)$ und $T \in \mathcal{L}(H)$ definiert durch $[T f](x)=q(x) f(x)$ für $f \in H$ und $x \in \Omega$. Wegen

$$
\|T f\|_{2}^{2}=\int_{\Omega}|q(x) f(x)|^{2} \leq\|q\|_{\infty}^{2}\|f\|_{2}^{2} \quad \text { für alle } f \in H
$$

ist $T$ tatsächlich stetig mit $\|T\| \leq\|q\|_{\infty}$. Dabei gilt:

$$
\left\langle f, T^{*} g\right\rangle=\langle T f, g\rangle=\int_{\Omega} q(x) f(x) \overline{g(x)} d x=\int_{\Omega} f(x) \overline{\overline{q(x)} g(x)} d x=\langle f, \bar{q} g\rangle
$$

also ist $T^{*} \in \mathcal{L}(H)$ gegeben durch $T^{*} g=\bar{q} g$. Beachte: $T$ ist normal, denn $T T^{*} g= T^{*} T g=|q|^{2} g$ für alle $g \in H$.\\
Ferner gilt:

\begin{itemize}
  \item $T$ unitär $\Longleftrightarrow|q|=1$ f.ü. auf $\Omega$
  \item $T$ selbstadjungiert $\Longleftrightarrow q$ f.ü. auf $\Omega$ reellwertig.
\end{itemize}

Randbemerkung: $T$ definiert ebenfalls einen stetigen Operator $L^{p}(\Omega) \rightarrow L^{p}(\Omega)$ für $1 \leq p \leq \infty$, wobei auch dann $\|T\| \leq\|q\|_{\infty}$ gilt.

\subsection*{Beispiel (Integraloperatoren)}
Sei $\Omega \subset \mathbb{R}^{N}$ (Lebesgue-)messbar und $H=L^{2}(\Omega)$. Sei ferner $k \in L^{2}(\Omega \times \Omega)$. Der Integraloperator $K \in \mathcal{L}(H)$ zur Kernfunktion $k$ ist definiert durch

$$
[K f](x)=\int_{\Omega} k(x, y) f(y) d y \quad \text { für } f \in H \text { und } x \in \Omega \text { (f.ü.). }
$$

Beachte:

\begin{itemize}
  \item $K$ ist wohldefiniert, denn: Da $\int_{\Omega \times \Omega}|k(x, y)|^{2} d(x, y)<\infty$ ist, gilt nach dem Satz von Fubini:\\
i) $\int_{\Omega}|k(x, y)|^{2} d y<\infty$ für fast alle $x \in \Omega$\\
ii) $x \mapsto \int_{\Omega}|k(x, y)|^{2} d y$ ist integrierbar
\end{itemize}

Aufgrund der Cauchy-Schwarz-Ungleichung ist somit $K f(x)$ für $f \in H$ und fast alle $x \in \Omega$ wohldefiniert, und es gilt

$$
|K f(x)|^{2}=\left|\int_{\Omega} k(x, y) f(y) d y\right|^{2} \leq\|f\|_{2}^{2} \int_{\Omega}|k(x, y)|^{2} d y .
$$

Es folgt

$$
\|K f\|_{2}^{2}=\int_{\Omega}|K f(x)|^{2} d x \leq\|f\|_{2}^{2} \int_{\Omega} \int_{\Omega}|k(x, y)|^{2} d y d x=\|f\|_{2}^{2}\|k\|_{L^{2}(\Omega \times \Omega)}^{2},
$$

also $K \in \mathcal{L}(H)$ mit $\|K\| \leq\|k\|_{L^{2}(\Omega \times \Omega)}$.

\begin{itemize}
  \item $K^{*} \in \mathcal{L}(H)$ ist der Integraloperator zur Kernfunktion
\end{itemize}

$$
k^{*} \in L^{2}(\Omega \times \Omega), \quad k^{*}(x, y)=\overline{k(y, x)},
$$

denn für $f, g \in H$ gilt

$$
\begin{aligned}
\langle K f, g\rangle=\int_{\Omega}\left(\int_{\Omega} k(x, y) f(y) d y\right) \overline{g(x)} d x & =\int_{\Omega} f(y) \overline{\left(\int_{\Omega} \overline{k(x, y)} g(x) d x\right)} d y \\
& =\int_{\Omega} f(x) \overline{\left(\int_{\Omega} k^{*}(y, x) g(y) d y\right)} d x=\left\langle f, K^{*} g\right\rangle
\end{aligned}
$$

\subsection*{Satz}
Für alle $T \in \mathcal{L}(H)$ gilt $\operatorname{Kern} T^{*}=(\operatorname{Bild} T)^{\perp}$ und $\overline{\operatorname{Bild} T^{*}}=(\operatorname{Kern} T)^{\perp}$.\\
Beweis. Für $y \in H$ gilt

$$
\begin{aligned}
y \in \operatorname{Kern} T^{*} & \Longleftrightarrow\left\langle x, T^{*} y\right\rangle=0 \quad \text { für alle } x \in H \\
& \Longleftrightarrow\langle T x, y\rangle=0 \quad \text { für alle } x \in H \quad \Longleftrightarrow \quad y \in(\operatorname{Bild} T)^{\perp} .
\end{aligned}
$$

Also ist $\operatorname{Kern} T^{*}=(\operatorname{Bild} T)^{\perp}$. Weiterhin ist\\
$\operatorname{Kern} T=\operatorname{Kern} T^{* *}=\left(\operatorname{Bild} T^{*}\right)^{\perp}, \quad$ also $\quad(\operatorname{Kern} T)^{\perp}=\left(\operatorname{Bild} T^{*}\right)^{\perp \perp}=\overline{\operatorname{Bild} T^{*}}$. \(\square\)

\subsection*{Satz}
$P \in \mathcal{L}(H)$ ist eine orthogonale Projektion genau dann, wenn $P^{2}=P=P^{*}$\\
Beweis. „ ": Da $P^{2}=P$, ist $P$ eine Projektion. Ferner ist Bild $P=\operatorname{Kern}(\mathrm{id}-P)$ abgeschlossen, da $P$ stetig ist. Da schließlich $P^{*}=P$ gilt, ist $\operatorname{Kern} P=\operatorname{Kern} P^{*} \stackrel{\text { 8.10 }}{=}$ (Bild $P)^{\perp}$. Somit ist $P$ eine orthogonale Projektion.\\
„" Nach 7.7 gilt $P^{2}=P$, und für $Q=$ id $-P$ gilt Bild $Q=(\text { Bild } P)^{\perp}$. Also gilt für alle $x, y \in H$ :

$$
\langle P x, y\rangle=\langle P x, P y+Q y\rangle=\langle P x, P y\rangle=\langle P x+Q x, P y\rangle=\langle x, P y\rangle
$$

Mit 8.5 folgt $P=P^{*}$.

\subsection*{Bemerkung}
Für $T \in \mathcal{L}(H)$ gilt: $T$ unitär $\Longleftrightarrow T$ bijektive Isometrie\\
Beweis. Es gilt\\
$T$ bijektive Isometrie $\Longleftrightarrow$\\
$T$ bijektiv, und $\|T x\|=\|x\|$ für alle $x \in H \quad \Longleftrightarrow$\\
$T$ bijektiv, und $\langle T x, T y\rangle=\langle x, y\rangle$ für alle $x, y \in H \quad \Longleftrightarrow$\\
$T$ bijektiv, und $T^{*} T=\mathrm{id}$, also $T^{-1}=T^{*} \Longleftrightarrow$\\
$T$ unitär.\\
Bei der zweiten Äquivalenz haben wir die Polarisationsgleichungen

$$
\begin{aligned}
& \langle x, y\rangle=\frac{1}{2}\left(\|x+y\|^{2}-\|x\|^{2}-\|y\|^{2}\right) \quad \text { falls } \mathbb{K}=\mathbb{R}, \\
& \langle x, y\rangle=\frac{1}{2}\left(\|x+y\|^{2}-\|x\|^{2}-\|y\|^{2}\right)+\frac{i}{2}\left(\|x+i y\|^{2}-\|x\|^{2}-\|y\|^{2}\right) \quad \text { falls } \mathbb{K}=\mathbb{C} .
\end{aligned}
$$

(siehe Übungsblatt 1) verwendet.

\subsection*{Bemerkung}
Ist $\mathbb{K}=\mathbb{C}$ und $T \in \mathcal{L}(H)$, so ist $T^{*}$ bereits durch

$$
\langle T x, x\rangle=\left\langle x, T^{*} x\right\rangle \quad \text { für alle } x \in H
$$

eindeutig festgelegt.

Beweis. Seien $x, y \in H$. Anwendung von ( $*$ ) auf $x+i y$ und $x+y$ liefert die Gleichung $\langle T x, y\rangle=\left\langle x, T^{*} y\right\rangle$ (Übung!).

\section*{§9 Der Satz von Hahn-Banach}
\subsection*{Motivation}
Sei $(E,\|\cdot\|)$ ein normierter Raum, $F \subset E$ ein Unterraum und $\varphi \in F^{\prime}$ eine stetige Linearform auf $F$ bzgl. $\|\cdot\|$. Ein Ziel dieses Kapitels ist es, die Existenz einer stetigen Fortsetzung $\psi \in E^{\prime}$ von $\varphi$ auf $E$ mit $\|\psi\|=\|\varphi\|$ nachzuweisen. Unter den folgenden Zusatzvoraussetzungen können wir dies bereits:\\
i). $F \subset E$ ist dicht. Dann existiert genau eine solche Fortsetzung $\psi$ gemäß 5.1.\\
ii). $E$ ist ein Hilbertraum. Dann kann man mit 5.1 zunächst $\varphi$ auf $\bar{F} \subset E$ fortsetzen und dann $\psi:=\varphi \circ P \in E^{\prime}$ betrachten, wobei $P$ die orthogonale Projektion auf $\bar{F}$ sei. Dann hat $\psi$ die gewünschten Eigenschaften.

Im Folgenden betrachten wir nun einen allgemeinen Rahmen für dieses Problem.

\subsection*{Satz von Hahn-Banach (1. Version)}
Sei $V$ ein $\mathbb{R}$-Vektorraum, $U \subset V$ ein linearer Teilraum, $p: V \rightarrow \mathbb{R}$ konvex und $\varphi: U \rightarrow \mathbb{R}$ eine $\mathbb{R}$-lineare Abbildung mit $\varphi(x) \leq p(x)$ für alle $x \in U$. Dann existiert eine $\mathbb{R}$-lineare Fortsetzung $\psi: V \rightarrow \mathbb{R}$ von $\varphi$ mit $\psi(x) \leq p(x)$ für alle $x \in V$.

Beweis. 1. Fall: $V=U \oplus \operatorname{span}\{z\}$ für ein $z \in V \backslash U$.\\
Dann ist jede lineare Fortsetzung $\psi$ von $\varphi$ auf $V$ durch $\psi(z)$ bereits festgelegt, da dann

$$
(*) \quad \psi(x+\mu z)=\varphi(x)+\mu \psi(z) \quad \text { für alle } x \in U \text { und } \mu \in \mathbb{R} \text { gilt. }
$$

Seien $x, y \in U$ und $\alpha, \beta>0$, und sei $\lambda=\frac{\beta}{\beta+\alpha}$. Dann ist $0<\lambda<1$ und $1-\lambda=\frac{\alpha}{\beta+\alpha}$, also

$$
\begin{aligned}
\beta \varphi(x)+\alpha \varphi(y) & =\varphi(\beta x+\alpha y)=(\alpha+\beta) \varphi(\lambda x+[1-\lambda] y) \\
& \leq(\alpha+\beta) p(\lambda(x-\alpha z)+[1-\lambda](y+\beta z)) \\
& \leq \beta p(x-\alpha z)+\alpha p(y+\beta z)
\end{aligned}
$$

und somit

$$
\frac{1}{\alpha}[\varphi(x)-p(x-\alpha z)] \leq \frac{1}{\beta}[p(y+\beta z)-\varphi(y)]
$$

Also existiert $c \in \mathbb{R}$ mit

$$
\sup _{\substack{x \in U \\ \alpha>0}} \frac{1}{\alpha}[\varphi(x)-p(x-\alpha z)] \leq c \leq \inf _{\substack{x \in U \\ \alpha>0}} \frac{1}{\alpha}[p(x+\alpha z)-\varphi(x)]
$$

Definiere $\psi$ nun durch ( $*$ ) mit $\psi(z):=c$. Dann gilt

$$
\psi(x \pm \alpha z)=\varphi(x) \pm \alpha c \leq p(x \pm \alpha z) \quad \text { für } x \in U \text { und } \alpha>0,
$$

und für $\alpha=0$ gilt dies nach Voraussetzung. Also hat $\psi$ die gewünschten Eigenschaften.\\
2. Allgemeiner Fall: Sei $\mathcal{Z}$ die Menge aller Paare ( $W, \psi$ ), wobei $W \subset V$ ein Unterraum mit $U \subset W$ und $\psi: W \rightarrow \mathbb{R}$ eine lineare Abbildung ist mit

$$
\left.\psi\right|_{U}=\varphi \quad \text { und } \quad \psi(x) \leq p(x) \quad \text { für alle } x \in W .
$$

Dann ist eine Halbordnung auf $\prec$ auf $\mathcal{Z}$ definiert durch

$$
\left(W_{1}, \psi_{1}\right) \prec\left(W_{2}, \psi_{2}\right) \quad \Longleftrightarrow \quad W_{1} \subset W_{2} \quad \text { und }\left.\quad \psi_{2}\right|_{W_{1}}=\psi_{1} .
$$

Dabei besitzt jede Kette (d.h. totalgeordnete Teilmenge) $\mathcal{N} \subset \mathcal{Z}$ eine obere Schranke gegeben durch $\left(W_{0}, \psi_{0}\right)$ mit

$$
W_{0}=\{x \in V: \text { es gibt }(W, \psi) \in \mathcal{N} \text { mit } x \in W\}
$$

und $\psi_{0}: W_{0} \rightarrow \mathbb{R}$ definiert durch $\psi_{0}(x):=\psi(x)$ falls $(W, \psi) \in \mathcal{N}$ und $x \in W$. Aus der Ketteneigenschaft erhält man, dass $W_{0} \subset V$ ein Unterraum und $\psi_{0}$ eine wohldefinierte lineare Abbildung mit $\left.\psi\right|_{U}=\varphi$ und $\psi(x) \leq p(x)$ für alle $x \in W$ ist. Das Zornsche Lemma besagt nun, dass die Menge $\mathcal{Z}$ ein maximales Element $(W, \psi)$ bezüglich der Halbordnung $\prec$ besitzt. Dabei muss $W=V$ sein, denn sonst könnte man $z \in V \backslash W$ wählen und $\psi$ entsprechend dem ersten Fall unter Erhalt der Abschätzung $\psi \leq p$ auf $W \oplus \operatorname{span}\{z\}$ fortsetzen, was der Maximalität von $(W, \psi)$ wiederspricht.\\
Es folgt, dass $\psi: V \rightarrow \mathbb{R}$ die gewünschten Eigenschaften hat. \(\square\)

\subsection*{Satz von Hahn-Banach (2. Version)}
Sei $V$ ein $\mathbb{K}$-Vektorraum, $U \subset V$ ein linearer Teilraum, $p: V \rightarrow \mathbb{R}$ absolut konvex und $\varphi: U \rightarrow \mathbb{K}$ eine $\mathbb{K}$-lineare Abbildung mit $|\varphi(x)| \leq p(x)$ für alle $x \in U$. Dann existiert eine $\mathbb{K}$-lineare Fortsetzung $\psi: V \rightarrow \mathbb{K}$ von $\varphi$ mit $|\psi(x)| \leq p(x)$ für alle $x \in V$.

Beweis. Für $\mathbb{K}=\mathbb{R}$ folgt dies aus 9.2. Hat man nämlich eine $\mathbb{R}$-lineare Fortsetzung $\psi: V \rightarrow \mathbb{R}$ gefunden mit $\psi(x) \leq p(x)$ für alle $x \in V$, so gilt auch

$$
-\psi(x)=\psi(-x) \leq p(-x)=p(x) \quad \text { für alle } x \in V
$$

und damit $|\psi| \leq p$.\\
Sei daher im Folgenden $\mathbb{K}=\mathbb{C}$ angenommen. Wir betrachten $\gamma=\operatorname{Re} \varphi: U \rightarrow \mathbb{R}$. Dann ist $\gamma \mathbb{R}$-linear, und es gilt

$$
\gamma(i x)=\operatorname{Re}(i \varphi(x))=-\operatorname{Im} \varphi(x), \quad \text { also } \quad \varphi(x)=\gamma(x)-i \gamma(i x) \quad \text { für } x \in U .
$$

Ferner gilt $\gamma(x) \leq|\varphi(x)| \leq p(x)$ für alle $x \in U$, und nach 9.2 besitzt $\gamma$ eine $\mathbb{R}$-lineare Fortsetzung $\Gamma: V \rightarrow \mathbb{R}$ mit $\Gamma(x) \leq p(x)$ für alle $x \in V$. Sei nun $\psi: V \rightarrow \mathbb{C}$ definiert durch $\psi(x)=\Gamma(x)-i \Gamma(i x)$. Dann ist $\psi$ eine $\mathbb{R}$-lineare Fortsetzung von $\varphi$. Ferner gilt $\psi(i x)=\Gamma(i x)-i \Gamma(-x)=i(\Gamma(x)-i \Gamma(i x))=i \psi(x)$, d.h. $\psi$ ist auch $\mathbb{C}$-linear. Schließlich existiert für jedes $x \in V$ ein $a \in \mathbb{C},|a|=1$ mit $|\psi(x)|=a \psi(x)=\psi(a x) \in \mathbb{R}$ und somit

$$
|\psi(x)|=\Gamma(a x) \leq p(a x)=p(x) .
$$

Damit ist alles gezeigt.

Sei im Folgenden $(E,\|\cdot\|)$ ein normierter $\mathbb{K}$-Vektorraum.

\subsection*{Satz von Hahn-Banach (3. Version)}
Ist $F \subset E$ ein linearer Teilraum und $\varphi \in F^{\prime}$ gegeben, so existiert eine Fortsetzung $\psi \in E^{\prime}$ von $\varphi$ auf $E$ mit $\|\psi\|=\|\varphi\|$.

Beweis. Sei $p: E \rightarrow \mathbb{R}$ definiert durch $p(x)=\|\varphi\|\|x\|$. Dann ist $p$ absolut konvex und $|\varphi| \leq p$ auf $F$. Nach 9.3 existiert eine $\mathbb{K}$-lineare Fortsetzung $\psi: E \rightarrow \mathbb{K}$ von $\varphi$ mit $|\psi(x)| \leq p(x)=\|\varphi\|\|x\|$ für alle $x \in E$, also $\psi \in E^{\prime}$ und $\|\psi\| \leq\|\varphi\|$. Da ferner $\left.\psi\right|_{F}=\varphi$ gilt, folgt $\|\psi\|=\|\varphi\|$.

Bemerkung: Ist E separabel, so kann man 9.4 auch ohne Verwendung des Zornschen Lemmas beweisen.

\subsection*{Korollar}
$\mathrm{Zu} x \in E \backslash\{0\}$ existiert ein $\psi \in E^{\prime}$ mit $\psi(x)=\|x\|$ und $\|\psi\|=1$.\\
Beweis. Sei $F:=\operatorname{span}\{x\} \subset E$ und $\varphi \in F^{\prime}$ definiert durch $\varphi(\lambda x)=\lambda\|x\|$ für $\lambda \in \mathbb{K}$. Nach 9.4 existiert $\psi \in E^{\prime}$ mit $\left.\psi\right|_{F}=\varphi$ und $\|\psi\|=\|\varphi\|=1$.

\subsection*{Definition}
Wir setzen $E^{\prime \prime}:=\left(E^{\prime}\right)^{\prime}=\mathcal{L}\left(E^{\prime}, \mathbb{K}\right)$ und definieren die Abbildung $i_{E}: E \rightarrow E^{\prime \prime}$ durch

$$
i_{E}(x)(\varphi)=\varphi(x) \quad \text { für } x \in E \text { und } \varphi \in E^{\prime} .
$$

\subsection*{Satz}
$i_{E}: E \rightarrow E^{\prime \prime}$ ist eine lineare Isometrie.\\
Beweis. Die Linearität von $i_{E}$ ist einfach zu sehen. Ferner gilt für $x \in E$ und $\varphi \in E^{\prime}$ :

$$
\left|i_{E}(x)(\varphi)\right|=|\varphi(x)| \leq\|\varphi\|\|x\|, \quad \text { also } \quad\left\|i_{E}(x)\right\| \leq\|x\| .
$$

Sei nun $x \in E \backslash\{0\}$. Nach 9.5 existiert $\psi \in E^{\prime}$ mit $\|\psi\|=1$ und $\psi(x)=\|x\|$, also $\left|i_{E}(x)(\psi)\right|=|\psi(x)|=\|x\|$. Es folgt $\left\|i_{E}(x)\right\| \geq\|x\|$, und insgesamt folgt Gleichheit. Also ist $i_{E}$ eine Isometrie.

\subsection*{Korollar}
Jeder normierte Raum $(E,\|\cdot\|)$ ist isometrisch und dicht in einem Banachraum $\hat{E}$ eingebettet, welcher bis auf isometrische Isomorphie eindeutig bestimmt ist. Man nennt $\hat{E}$ die Vervollständigung von $E$.

Beweis. Wähle $\hat{E}:=\overline{\operatorname{Bild} i_{E}} \subset E^{\prime \prime}$. Da $E^{\prime \prime}$ ein Banachraum ist, ist $\hat{E}$ dies auch. Ferner ist $i_{E}: E \rightarrow E^{\prime \prime}$ eine Isometrie und Bild $i_{E}$ dicht in $\hat{E}$, wie gewünscht. Die Eindeutigkeit von $\hat{E}$ bis auf Isomorphie folgt leicht aus 5.1.

Im Folgenden sei $(F,\|\cdot\|)$ stets ein weiterer normierter $\mathbb{K}$-Vektorraum.

\subsection*{Definition und Satz}
Für $T \in \mathcal{L}(E, F)$ ist der duale Operator $T^{\prime} \in \mathcal{L}\left(F^{\prime}, E^{\prime}\right)$ wohldefiniert durch

$$
\left[T^{\prime} \varphi\right](x)=\varphi(T x) \quad \text { für } \varphi \in F^{\prime} \text { und } x \in E \text {. }
$$

Ferner sei $T^{\prime \prime}:=\left(T^{\prime}\right)^{\prime} \in \mathcal{L}\left(E^{\prime \prime}, F^{\prime \prime}\right)$\\
Beweis (Wohldefiniertheit und Stetigkeit von $T^{\prime}$ ). Für $\varphi \in F^{\prime}$ und $x \in E$ ist

$$
\left|\left[T^{\prime} \varphi\right](x)\right|=|\varphi(T x)| \leq\|\varphi\|\|T\|\|x\|,
$$

also $T^{\prime} \varphi \in E^{\prime}$ mit $\left\|T^{\prime} \varphi\right\| \leq\|T\|\|\varphi\|$.\\
Da $T^{\prime}$ offensichtlich linear ist, folgt $T^{\prime} \in \mathcal{L}\left(F^{\prime}, E^{\prime}\right)$ mit $\left\|T^{\prime}\right\| \leq\|T\|$.

\subsection*{Satz}
Für $T \in \mathcal{L}(E, F)$ gilt\\
i). $\|T\|=\left\|T^{\prime}\right\|$\\
\includegraphics[max width=\textwidth, center]{2025_10_06_e8fc9b2ca9d05628f654g-100}

Beweis. Zuerst zu (ii): Für $x \in E$ und $\varphi \in F^{\prime}$ gilt

$$
i_{F}(T x)(\varphi)=\varphi(T x)=\left[T^{\prime} \varphi\right](x)=i_{E}(x)\left(T^{\prime} \varphi\right)=\left[T^{\prime \prime} i_{E}(x)\right](\varphi)
$$

also $\left(i_{F} \circ T\right)(x)=\left(T^{\prime \prime} \circ i_{E}\right)(x)$. Dies zeigt (ii).\\
Nun zu (i): Der Beweis von 9.9 zeigt direkt: $\|T\| \geq\left\|T^{\prime}\right\| \geq\left\|T^{\prime \prime}\right\|$.\\
Da $i_{E}$ und $i_{F}$ Isometrien sind, ist andererseits

$$
\left\|T^{\prime \prime}\right\| \geq\left\|T^{\prime \prime} \circ i_{E}\right\| \stackrel{(i i)}{=}\left\|i_{F} \circ T\right\|=\|T\|
$$

Es folgt: $\left\|T^{\prime \prime}\right\|=\|T\|=\left\|T^{\prime}\right\|$.

\subsection*{Satz}
Sei $T \in \mathcal{L}(E, F)$.\\
i). Ist ( $G,\|\cdot\|$ ) ein weiterer normierter $\mathbb{K}$-Vektorraum und $S \in \mathcal{L}(F, G)$, so gilt $(S T)^{\prime}=T^{\prime} S^{\prime} \in \mathcal{L}\left(G^{\prime}, E^{\prime}\right)$.\\
ii). Ist $T$ ein topologischer Isomorphismus, so auch $T^{\prime}$, und es gilt $\left(T^{\prime}\right)^{-1}= \left(T^{-1}\right)^{\prime} \in \mathcal{L}\left(E^{\prime}, F^{\prime}\right)$.

Beweis. (i) Für $\varphi \in G^{\prime}$ und $x \in E$ ist

$$
\left[(S T)^{\prime}(\varphi)\right](x)=\varphi(S T x)=\left[S^{\prime} \varphi\right](T x)=\left[T^{\prime} S^{\prime} \varphi\right](x)
$$

also $(S T)^{\prime}(\varphi)=T^{\prime} S^{\prime} \varphi$ für $\varphi \in G^{\prime}$ und somit $(S T)^{\prime}=T^{\prime} S^{\prime} \in \mathcal{L}\left(G^{\prime}, E^{\prime}\right)$.\\
(ii) folgt direkt aus (i).

\subsection*{Definition}
Für $M \subset E$ und $N \subset E^{\prime}$ sei

$$
\begin{aligned}
M^{\perp} & :=\left\{\varphi \in E^{\prime}: \varphi(x)=0 \text { für alle } x \in M\right\}, \\
N^{\perp} & :=\{x \in E: \varphi(x)=0 \text { für alle } \varphi \in N\} .
\end{aligned}
$$

Vorsicht: Diese Notation kann für Teilmengen von $E^{\prime}$ doppeldeutig sein; es muss aus dem Zusammenhang klar sein, dass man das duale Paar ( $E, E^{\prime}$ ) und nicht das Paar ( $E^{\prime}, E^{\prime \prime}$ ) betrachtet.\\
Ist $E$ ein Hilbertraum und identifiziert man $E^{\prime}$ mit $E$ mittels des Rieszschen Darstellungssatzes, d.h. vermöge der Abbildung $E \rightarrow E^{\prime}, x \mapsto\langle\cdot, x\rangle$, so stimmt die obige Definition der "Orthogonalkomplemente" mit der in Kapitel 7 überein.

\subsection*{Satz}
Für $M \subset E$ gilt:\\
i). $M^{\perp}$ ist ein abgeschlossener Unterraum von $E^{\prime}$.\\
ii). $M^{\perp}=(\overline{\operatorname{span} M})^{\perp}$.

Für $N \subset E^{\prime}$ gelten analoge Aussagen.\\
Beweis. Leichte Übung.

\subsection*{Satz}
i). Für $M \subset E$ gilt $M^{\perp \perp}=\overline{\operatorname{span} M}$.\\
ii). Für $N \subset E^{\prime}$ gilt $\overline{\operatorname{span} N} \subset N^{\perp \perp}$.\\
iii). Ist $T \in \mathcal{L}(E, F)$, so gilt $\operatorname{Kern} T^{\prime}=(\operatorname{Bild} T)^{\perp}$ und $\operatorname{Kern} T=\left(\operatorname{Bild} T^{\prime}\right)^{\perp}$.

Bem: In (ii) gilt im Allgemeinen keine Gleichheit, siehe Beispiel 9.15 (iii). Wählt man in (iii) insbesondere $E=F$ und $T=\mathrm{id}$, so erhält man $\left\{0_{E^{\prime}}\right\}=E^{\perp}$ und $\left\{0_{E}\right\}=\left(E^{\prime}\right)^{\perp}$.

Beweis. (i) " $\supset$ ": Wegen 9.13 ist $M^{\perp}=(\overline{\operatorname{span} M})^{\perp}$; für $x \in \overline{\operatorname{span} M}$ und $\varphi \in M^{\perp}$ gilt also $\varphi(x)=0$. Es folgt $\overline{\operatorname{span} M} \subset M^{\perp \perp}$.\\
" $\subset$ ": Sei $x \notin F:=\overline{\operatorname{span} M}$. Dann ist die Projektion $P: F \oplus \operatorname{span}\{x\} \rightarrow \operatorname{span}\{x\}$ mit Kern $P=F$ stetig nach 5.14(i). Sei nun $\varphi \in(F \oplus \operatorname{span}\{x\})^{\prime}$ definiert durch $\varphi=h \circ P$,\\
wobei die lineare Abbildung $h: \operatorname{span}\{x\} \rightarrow \mathbb{K}$ definiert sei durch $h(x)=1$. Nach 9.4 existiert eine Fortsetzung $\psi \in E^{\prime}$ von $\varphi$, d.h. insbesondere ist $\left.\psi\right|_{M}=\left.\varphi\right|_{M}=0$ und $\psi(x)=\varphi(x)=1$. Also ist $\psi \in M^{\perp}$ und $x \notin M^{\perp \perp}$.\\
(ii) Wegen 9.13 ist $N^{\perp}=(\overline{\operatorname{span} N})^{\perp}$; für $\varphi \in \overline{\operatorname{span} N}$ und $x \in N^{\perp}$ gilt also $\varphi(x)=0$. Es folgt $\overline{\operatorname{span} N} \subset N^{\perp \perp}$.\\
(iii) Es ist

$$
\text { Kern } T^{\prime}=\left\{\varphi \in F^{\prime}: 0=\left[T^{\prime} \varphi\right](x)=\varphi(T x) \text { für alle } x \in E\right\}=(\text { Bild } T)^{\perp} .
$$

Ist ferner $x \in \operatorname{Kern} T$, d.h. $T x=0$, so gilt $0=\varphi(T x)=\left[T^{\prime} \varphi\right](x)$ für alle $\varphi \in F^{\prime}$, also $x \in\left(\text { Bild } T^{\prime}\right)^{\perp}$.\\
Ist anderseits $x \notin \operatorname{Kern} T$, d.h. $T x \neq 0$, so existiert nach 9.5 ein $\psi \in F^{\prime}$ mit $0 \neq \psi(T x)=\left[T^{\prime} \psi\right](x)$, also $x \notin\left(\text { Bild } T^{\prime}\right)^{\perp}$.

\subsection*{Beispiele}
i). Sei $U \subset E$ ein Unterraum und $j: U \rightarrow E$ die Inklusion. Dann ist $R:=j^{\prime}$ : $E^{\prime} \rightarrow U^{\prime}$ die Restriktion definiert durch $R \varphi=\left.\varphi\right|_{U}$ für $\varphi \in E^{\prime}$, da

$$
[R \varphi](x)=\varphi(j(x)) \quad \text { für } x \in U, \varphi \in E^{\prime} .
$$

Dabei ist $R$ surjektiv nach 9.4. Ferner gilt

$$
\text { Kern } R=(\text { Bild } j)^{\perp}=U^{\perp}
$$

ii). Sei $E=\ell^{1}$. Für $x=\left(x_{k}\right)_{k} \in \ell^{\infty}$ ist dann eine Linearform $\varphi_{x} \in\left(\ell^{1}\right)^{\prime}$ definiert durch

$$
\varphi_{x}(y)=\sum_{k=1}^{\infty} x_{k} y_{k} \quad \text { für } y=\left(y_{k}\right)_{k} \in \ell^{1} .
$$

Wir zeigen, dass die Abbildung

$$
\ell^{\infty} \rightarrow\left(\ell^{1}\right)^{\prime}, \quad x \mapsto \varphi_{x}
$$

ein isometrische Isomorphismus ist: Sei dazu zunächst $x \in \ell^{\infty}$. Wegen $\left|\varphi_{x}(y)\right| \leq \|x\|_{\infty}\|y\|_{1}$ für $y \in \ell^{1}$ ist dann $\left\|\varphi_{x}\right\| \leq\|x\|_{\infty}$. Betrachtet man ferner $k \in \mathbb{N}$ und $y:=\bar{x}_{k} e_{k}$, so folgt

$$
\left|x_{k}\right|^{2}=\varphi_{x}(y) \leq\left\|\varphi_{x}\right\|\|y\|_{1}=\left\|\varphi_{x}\right\|\left|x_{k}\right|, \quad \text { also }\left\|\varphi_{x}\right\| \geq\left|x_{k}\right|
$$

Somit folgt $\left\|\varphi_{x}\right\| \geq \sup _{k \in \mathbb{N}}\left|x_{k}\right|=\|x\|_{\infty}$, und insesamt folgt $\left\|\varphi_{x}\right\|=\|x\|_{\infty}$.\\
Es bleibt noch die Surjektivität der Abbildung zu zeigen. Sei dazu $\varphi \in\left(\ell^{1}\right)^{\prime}$ beliebig und $x_{k}:=\varphi\left(e_{k}\right)$ für $k \in \mathbb{N}$. Wie oben folgt dann

$$
\left|x_{k}\right|^{2}=\varphi\left(\bar{x}_{k} e_{k}\right) \leq\|\varphi\|\left|x_{k}\right|, \quad \text { also }\left|x_{k}\right| \leq\|\varphi\|
$$

und somit $\|x\|_{\infty} \leq\|\varphi\|$ für $x=\left(x_{k}\right)_{k}$. Also ist $x \in \ell^{\infty}$, und für $y \in \ell^{1}$ gilt

$$
\varphi_{x}(y)=\sum_{k=1}^{\infty} x_{k} y_{k}=\sum_{k=1}^{\infty} y_{k} \varphi\left(e_{k}\right)=\varphi\left(\sum_{k=1}^{\infty} y_{k} e_{k}\right)=\varphi(y),
$$

d.h. $\varphi=\varphi_{x}$.\\
iii). Die Abbildung $x \mapsto \varphi_{x}$ aus ii) definiert im Fall $p \in(1, \infty)$ auch einen isometrischen Isorphismus $\ell^{p^{\prime}} \rightarrow\left(\ell^{p}\right)^{\prime}$ (Übung, Blatt 10). Hier ist $p^{\prime}=\frac{p}{p-1}$ wie üblich der konjugierte Exponent zu $p$. Im Fall $p=\infty$ ist die Abbildung

$$
\ell^{1} \rightarrow\left(\ell^{\infty}\right)^{\prime}, \quad x \mapsto \varphi_{x}
$$

zwar immer noch eine Isometrie, aber nicht mehr surjektiv (Übung, Blatt 11).\\
iv). Wir zeigen, das für $N \subset E^{\prime}$ im Allgemeinen nicht $\overline{\operatorname{span} N}=N^{\perp \perp}$ gilt. Sei dazu $E=\ell^{1}$ und $\ell^{\infty} \rightarrow\left(\ell^{1}\right)^{\prime}, x \mapsto \varphi_{x}$ der isometrische Isomorphismus aus (ii). Sei ferner

$$
N:=\left\{\varphi_{e_{k}}: k \in \mathbb{N}\right\} \subset\left(\ell^{1}\right)^{\prime}
$$

Für $y=\left(y_{k}\right)_{k} \in N^{\perp} \subset \ell^{1}$ gilt dann

$$
y_{k}=\varphi_{e_{k}}(y)=0 \quad \text { für alle } k
$$

also $y=0$. Somit ist $N^{\perp}=\{0\}$, also $N^{\perp \perp}=\left(\ell^{1}\right)^{\prime}$. Andererseits ist

$$
\operatorname{span} N=\left\{\varphi_{x}: x \in \mathcal{F}\right\}
$$

wobei $\mathcal{F}$ wie bisher den Raum der finiten Folgen bezeichne. Dies liefert wegen (ii) dann\\
$\overline{\operatorname{span} N}=\left\{\varphi_{x}: x \in \overline{\mathcal{F}}^{\ell^{\infty}}\right\}=\left\{\varphi_{x}: x\right.$ ist Nullfolge $\} \subsetneq\left\{\varphi_{x}: x \in \ell^{\infty}\right\}=\left(\ell^{1}\right)^{\prime}$. Somit folgt $\overline{\operatorname{span} N} \subsetneq N^{\perp \perp}$.

\section*{§ 10 Trennungssätze}
Sei im Folgenden $E$ ein normierter $\mathbb{R}$-Vektorraum. Für Teilmengen $B, C \subset E$ und $\lambda \in \mathbb{R}$ setzen wir

$$
\lambda B:=\{\lambda x: x \in B\} \quad \text { und } \quad B+C:=\{x+y: x \in B, y \in C\}
$$

Ist ferner $a \in E$, so setzen wir $a+B:=\{a\}+B=\{a+x: x \in B\}$.

\subsection*{Definition und Satz}
Sei $C \subset E$ konvex mit $0 \in \stackrel{\circ}{C}$, und sei $\gamma_{C}: E \rightarrow \mathbb{R}$ die Eichfunktion definiert durch

$$
\gamma_{C}(x)=\inf \{t>0: x \in t C\}=\inf \left\{t>0: \frac{x}{t} \in C\right\} \text { für } x \in E \text {. }
$$

Dann ist $\gamma_{C}$ sublinear, d.h. für $x, y \in E$ und $\lambda \geq 0$ gilt

$$
\gamma_{C}(\lambda x)=\lambda \gamma_{C}(x) \quad \text { und } \quad \gamma_{C}(x+y) \leq \gamma_{C}(x)+\gamma_{C}(y) \text {. }
$$

Insbesondere ist $\gamma_{C}$ also konvex.\\
Beweis. Komplett analog zum Beweis von 1.6 (ii).

\subsection*{Trennungssatz von Mazur}
Sei $C \subset E$ konvex und offen, und sei $W \subset E$ ein Unterraum. Sei ferner $a \in E$ mit $C \cap[a+W]=\varnothing$. Dann existiert $\psi \in E^{\prime}$ und $\alpha \in \mathbb{R}$ mit

$$
\psi(x)<\alpha \quad \text { für } x \in C \quad \text { und }\left.\quad \psi\right|_{a+W} \equiv \alpha .
$$

Beweis. 1. Fall: $0 \in C$. Dann ist $0 \notin a+W$ und somit $a \notin W$. Setze $U:=W+ \operatorname{span}\{a\}$ und definiere die lineare Abbildung $\varphi: U \rightarrow \mathbb{R}$ durch $\varphi(w+\lambda a)=\lambda$ für $w \in W, \lambda \in \mathbb{R}$. Für $w \in W$ und $\lambda>0$ gilt nun $\frac{1}{\lambda} w+a \notin C$, also $w+\lambda a \notin \lambda C$ und somit (da $0 \in C$ )

$$
w+\lambda a \notin t C \quad \text { für } t \leq \lambda \text {. }
$$

Es folgt

$$
\gamma_{C}(w+\lambda a)=\inf \{t>0: w+\lambda a \in t C\} \geq \lambda=\varphi(w+\lambda a)
$$

Dies gilt offensichtlich auch für $\lambda \leq 0$, und somit ist $\varphi \leq \gamma_{C}$ auf ganz $U$. Nach 9.2 und 10.1 existiert also eine lineare Fortsetzung $\psi: E \rightarrow \mathbb{R}$ mit $\psi \leq \gamma_{C}$ auf $E$. Insbesondere ist dann $\psi \equiv 1$ auf $a+W$ und $\psi \leq \gamma_{C}<1$ auf $C$ nach Definition von $\gamma_{C}$, da $C$ offen ist. Ferner ist $\psi$ stetig, denn nach Voraussetzung existiert $\varepsilon>0$ mit $B_{\varepsilon}(0) \subset C$, und somit ist

$$
\psi \leq \gamma_{C} \leq 1 \quad \text { auf } B_{\varepsilon}(0)
$$

Aufgrund der Linearität folgt dann $|\psi| \leq 1$ auf $B_{\varepsilon}(0)$ und somit $\sup _{B_{1}(0)}|\psi| \leq \frac{1}{\varepsilon}$; und dies liefert die Stetigkeit von $\psi$. Insgesamt folgt also die Behauptung mit $\psi$ und $\alpha=1$.\\
2. Fall: $0 \notin C$. Wähle dann $x \in C$ und betrachte $C^{\prime}:=C-x$ sowie $a^{\prime}:=a-x$. Der erste Fall liefert dann $\psi \in E^{\prime}$ mit $\left.\psi\right|_{C^{\prime}}<1$ und $\left.\psi\right|_{a^{\prime}+W} \equiv 1$. Mit $\alpha:=1+\psi(x)$ folgt also

$$
\left.\psi\right|_{C}<\alpha \quad \text { und }\left.\quad \psi\right|_{a+W} \equiv \alpha
$$

wie gewünscht.

\subsection*{Korollar}
Seien $C_{1}, C_{2} \subset E$ nichtleer und konvex mit $C_{1} \cap C_{2}=\varnothing$.\\
i). Ist $C_{1}$ offen, so existiert $\psi \in E^{\prime}$ und $\alpha \in \mathbb{R}$ mit

$$
\psi(x)<\alpha \text { für } x \in C_{1} \quad \text { und } \quad \psi(x) \geq \alpha \text { für } x \in C_{2} \text {. }
$$

Ist zusätzlich $C_{2}$ offen, so ist die zweite Ungleichung ebenfalls strikt.\\
ii). Ist $C_{1}$ kompakt und $C_{2}$ abgeschlossen, so existiert $\psi \in E^{\prime}$ und $\alpha \in \mathbb{R}$ mit

$$
\psi(x)<\alpha<\psi(y) \quad \text { für } x \in C_{1}, y \in C_{2} \text {. }
$$

Beweis. (i) Wir setzen $C:=C_{1}-C_{2}$. Dann ist $C$ offen, konvex und $0 \notin C$, da $C_{1} \cap C_{2}=\varnothing$ ist. Anwendung von 10.2 auf $C, W=\{0\}$ und $a=0$ liefert $\psi \in E^{\prime}$ mit

$$
\psi(z)<\psi(0)=0 \quad \text { für alle } z \in C
$$

also

$$
\psi(x)<\psi(y) \quad \text { für alle } x \in C_{1}, y \in C_{2} \text {. }
$$

Für alle $x \in C_{1}$ ist somit $\psi(x) \leq \alpha:=\inf _{C_{2}} \psi$. Da $C_{1}$ offen ist, folgt gemäß Übungsaufgabe 3, Blatt 10 also $\psi(x)<\alpha$ (dieser Beweisschritt wird später ergänzt). Ist $C_{2}$\\
ebenfalls offen, so folgt genauso auch $\psi(x)>\alpha$ für alle $x \in C_{2}$.\\
(ii) Nach Voraussetzung existiert $\varepsilon>0$ mit

$$
U_{\varepsilon}\left(C_{1}\right) \cap U_{\varepsilon}\left(C_{2}\right)=\varnothing
$$

Anwendung von (i) auf die konvexen offenen Mengen $U_{\varepsilon}\left(C_{1}\right)$ und $U_{\varepsilon}\left(C_{2}\right)$ liefert also die Behauptung.

Nachbemerkung: Auf die Voraussetzung der Offenheit kann man im Allgemeinen nicht verzichten, siehe Übungsblatt 11.

\section*{§ 11 Schwache Konvergenz und schwach*Konvergenz}
Seien im Folgenden stets $E, F$ normierte Räume über $\mathbb{K}$.

\subsection*{Definition}
Eine Folge $\left(x_{n}\right)_{n} \subset E$ heißt schwach konvergent gegen $x \in E$, wenn gilt:

$$
\lim _{n \rightarrow \infty} \psi\left(x_{n}\right)=\psi(x) \quad \text { für alle } \psi \in E^{\prime} .
$$

Wir schreiben dann $x_{n} \rightharpoonup x$.

\subsection*{Bemerkungen und Beispiele}
i). Der schwache Grenzwert $x$ einer schwach konvergenten Folge $\left(x_{n}\right)_{n} \subset E$ ist eindeutig bestimmt, denn gilt auch $x_{n} \rightharpoonup y \in E$, so folgt

$$
\psi(y-x)=\lim _{n \rightarrow \infty} \psi\left(x_{n}-x_{n}\right)=0 \quad \text { für alle } \psi \in E^{\prime}
$$

und somit $y-x=0$ gemäß 9.5 .\\
ii). Konvergenz in $E$ impliziert schwache Konvergenz.\\
iii). Die schwach konvergenten Folgen bilden einen Unterraum des $\mathbb{K}$-Vektorraums aller Folgen in $E$, und die schwache Grenzwertbildung ist $\mathbb{K}$-linear.\\
iv). Ist $\mathbb{K}=\mathbb{C}$ und ist $\left(x_{n}\right)_{n} \subset E$ eine Folge mit $x_{n} \rightharpoonup x \in E$, so folgt auch\\
(*) $\quad \varphi\left(x_{n}\right) \rightarrow \varphi(x) \quad$ für alle $\mathbb{R}$-linearen (!) stetigen Abbildungen $\varphi: E \rightarrow \mathbb{R}$.\\
Ist nämlich eine solche Abbildung $\varphi$ gegeben, so ist

$$
\psi: E \rightarrow \mathbb{C}, \quad \psi(y)=\varphi(y)-i \varphi(i y)
$$

$\mathbb{C}$-linear und stetig und erfüllt daher $\psi\left(x_{n}\right) \rightarrow \psi(x)$. Da $\varphi=\operatorname{Re} \psi$ ist, folgt (*).\\
v). Sei $E=l^{p}$ mit $p \in[1, \infty)$, und sei $q=\frac{p}{p-1}$. Gemäß Übungsaufgabe 4, Blatt 10 ist $\left(l^{p}\right)^{\prime}$ isometrisch isomorph zu $l^{q}$ vermöge der Abbildung $l^{q} \rightarrow\left(l^{p}\right)^{\prime}, x \mapsto \varphi_{x}$ mit

$$
\varphi_{x}(y)=\sum_{n \in \mathbb{N}} y_{n} x_{n} \quad \text { für } x=\left(x_{n}\right)_{n} \in l^{q}, y=\left(y_{n}\right)_{n} \in l^{p} .
$$

Sei nun $e_{n}=\left(\delta_{k n}\right)_{k} \in l^{p}$ für $n \in \mathbb{N}$. Ist $p>1$, so gilt gilt $e_{n} \rightharpoonup 0$, denn für $x \in l^{q}$ gilt

$$
\varphi_{x}\left(e_{n}\right)=x_{n} \rightarrow 0 \quad \text { für } n \rightarrow \infty .
$$

vi). Ist $\left(x_{n}\right)_{n} \subset l^{1}$ eine Folge mit $x_{n} \rightharpoonup x \in l^{1}$ so folgt $x_{n} \rightarrow x$ (Für einen Beweis siehe Schröder, Funktionalanalysis) .\\
vii). Sei $H$ ein Hilbertraum. Dann ist $H^{\prime}$ isometrisch (antilinear) isomorph zu $H$ vermöge der Abbildung $H \rightarrow H^{\prime}, y \mapsto\langle\cdot, y\rangle$. Somit gilt für eine Folge $\left(x_{n}\right)_{n} \subset H$ und $x \in H$ :

$$
x_{n} \rightharpoonup x \quad \Longleftrightarrow \quad\left\langle x_{n}, y\right\rangle \rightarrow\langle x, y\rangle \quad \text { für alle } y \in H .
$$

Ist $\left\{e_{n}: n \in \mathbb{N}\right\}$ ein Orthonormalsystem in $H$, so gilt $e_{n} \rightharpoonup 0$, denn für alle $y \in H$ bilden die Fourierkoeffizienten $\overline{\left\langle e_{n}, y\right\rangle}=\left\langle y, e_{n}\right\rangle$ gemäß der Besselschen Ungleichung 7.23 eine Folge in $l^{2}$, also eine Nullfolge.\\
viii). Sei $H=L^{2}(\mathbb{R})$ (bzgl. des Lebesgue-Maßes), sei $u \in H \backslash\{0\}$ und $u_{n} \in H$ definiert durch $u_{n}(x)=u(x-n)$. Dann gilt $\left\|u_{n}\right\|=\|u\|>0$ für alle $n \in \mathbb{N}$, und somit $u_{n} \nrightarrow 0$ für $n \rightarrow \infty$.\\
Behauptung: $u_{n} \rightharpoonup 0$ für $n \rightarrow \infty$.\\
Zum Beweis verwenden wir (11.1). Sei dazu $v \in H$ beliebig. Dann ist

$$
\left\langle u_{n}, v\right\rangle=\int_{\mathbb{R}} u(x-n) \overline{v(x)} d x \quad \text { für alle } n \in \mathbb{N}
$$

Sind $u, v \in \mathcal{C}_{c}(\mathbb{R})$, so folgt offensichtlich, dass die rechte Seite für $n \rightarrow \infty$ gegen Null konvergiert. Zu beliebigen $\varepsilon>0$ existieren nun aber $\tilde{u}, \tilde{v} \in \mathcal{C}_{c}(\mathbb{R})$ mit $\|v-\tilde{v}\|<\varepsilon$ und $\|u-\tilde{u}\|<\varepsilon$, also auch $\left\|u_{n}-\tilde{u}_{n}\right\|<\varepsilon$ für $\tilde{u}_{n} \in H$ definiert $\operatorname{durch} \tilde{u}_{n}(x)=\tilde{u}(x-n), n \in \mathbb{N}$. Somit folgt

$$
\left|\left\langle u_{n}, v\right\rangle-\left\langle\tilde{u}_{n}, \tilde{v}\right\rangle\right| \leq\left\|u_{n}\right\|\|v-\tilde{v}\|+\left\|u_{n}-\tilde{u}_{n}\right\|\|\tilde{v}\| \leq \varepsilon(\|u\|+\|v\|+\varepsilon) .
$$

für alle $n \in \mathbb{N}$. Die Behauptung folgt.

\subsection*{Satz}
Sei $\left(x_{n}\right)_{n} \subset E$ eine Folge und $x \in E$. Dann gilt:\\
i). $x_{n} \rightharpoonup x \Longrightarrow\left(x_{n}\right)_{n}$ beschränkt in $E$.\\
ii). Ist $\left(x_{n}\right)_{n} \subset E$ beschränkt und $M \subset E^{\prime}$ eine totale Menge mit

$$
\lim _{n \rightarrow \infty} \psi\left(x_{n}\right)=\psi(x) \quad \text { für alle } \psi \in M
$$

so gilt $x_{n} \rightharpoonup x$.\\
Beweis. (i) Sei $\tilde{x}_{n}=i_{E} x_{n} \in E^{\prime \prime}$ für $n \in \mathbb{N}$. Dann ist für jedes $\varphi \in E^{\prime}$ die Menge $\left\{\tilde{x}_{n}(\varphi)=\varphi\left(x_{n}\right): n \in \mathbb{N}\right\} \subset \mathbb{K}$ beschränkt. Nach dem Satz von der gleichmäßigen Beschränktheit 5.2 ist also auch die Menge $\left\{\tilde{x}_{n}: n \in \mathbb{N}\right\} \subset E^{\prime \prime}$ beschränkt und somit auch die Menge $\left\{x_{n}: n \in \mathbb{N}\right\} \subset E$, da $i_{E}$ eine Isometrie ist.\\
(ii) Offensichtlich ist

$$
\varphi(x)=\lim _{n \rightarrow \infty} \varphi\left(x_{n}\right) \quad \text { für alle } \varphi \in G:=\operatorname{span} M \subset E^{\prime} \text {. }
$$

Nach Voraussetzung ist ferner $c:=\sup \left\{\left\|x_{n}\right\|,\|x\|: n \in \mathbb{N}\right\}<\infty$. Sei nun $\psi \in E^{\prime}$ und $\varepsilon>0$, und sei $\varphi \in G$ mit $\|\psi-\varphi\|<\frac{\varepsilon}{3 c}$. Da $\varphi\left(x_{n}\right) \rightarrow \varphi(x)$ für $n \rightarrow \infty$, existiert $n_{0} \in \mathbb{N}$ mit $\left|\varphi\left(x_{n}\right)-\varphi(x)\right|<\frac{\varepsilon}{3}$ für $n \geq n_{0}$. Für diese $n$ gilt dann auch\\
$\left|\psi\left(x_{n}\right)-\psi(x)\right| \leq\left|\psi\left(x_{n}\right)-\varphi\left(x_{n}\right)\right|+\left|\varphi\left(x_{n}\right)-\varphi(x)\right|+|\varphi(x)-\psi(x)|<\frac{\varepsilon}{3 c}\left\|x_{n}\right\|+\frac{\varepsilon}{3}+\frac{\varepsilon}{3 c}\left\|x_{n}\right\| \leq \varepsilon$.\\
Dies zeigt $\psi\left(x_{n}\right) \rightarrow \psi(x)$.

\subsection*{Satz}
Sei $C \subset E$ abgeschlossen und konvex, und sei $\left(x_{n}\right)_{n}$ eine Folge in $C$ mit $x_{n} \rightharpoonup x \in E$. Dann ist $x \in C$.

Beweis. Wäre $x \notin C$, so würde nach 10.3 eine stetige $\mathbb{R}$-lineare (!) Abbildung $\varphi$ : $E \rightarrow \mathbb{R}$ und $\alpha \in \mathbb{R}$ existieren mit $\left.\varphi\right|_{C}<\alpha<\varphi(x)$, also

$$
\varphi(x) \stackrel{\text { 11.2 } v}{=} \lim _{n \rightarrow \infty} \varphi\left(x_{n}\right) \leq \alpha<\varphi(x)
$$

\subsection*{Korollar}
Sei $\left(x_{n}\right)_{n} \subset E$ eine Folge mit $x_{n} \rightharpoonup x \in E$. Dann gilt $\|x\| \leq \liminf _{n \in \mathbb{N}}\left\|x_{n}\right\|$.\\
Beweis. Sei $r:=\liminf _{n \rightarrow \infty}\left\|x_{n}\right\|$. Nach Übergang zu einer Teilfolge können wir annehmen, dass $\lim _{n \rightarrow \infty}\left\|x_{n}\right\|=r$ gilt.\\
Wir geben nun zwei alternative Argumente zum Abschluss des Beweises:

\begin{enumerate}
  \item Sei $\varepsilon>0$ und $C_{\varepsilon}:=B_{r+\varepsilon}(0) \subset E$. Dann ist $C_{\varepsilon} \subset E$ abgeschlossen und konvex, und es existiert $n_{0} \in \mathbb{N}$ mit $x_{n} \in C_{\varepsilon}$ für $n \geq n_{0}$. Mit 11.4 folgt $x \in C_{\varepsilon}$, also $\|x\| \leq r+\varepsilon$. Es folgt $\|x\| \leq r$, da $\varepsilon>0$ beliebig gewählt war.\\
2 . Für alle $\varphi \in E^{\prime}$ gilt
\end{enumerate}

$$
\left|\left[i_{E} x\right](\varphi)\right|=|\varphi(x)|=\lim _{n \rightarrow \infty}\left|\varphi\left(x_{n}\right)\right| \leq\|\varphi\| \lim _{n \rightarrow \infty}\left\|x_{n}\right\|=r\|\varphi\| \text {, }
$$

also $\left\|i_{E} x\right\| \leq r$. Da $i_{E}$ eine Isometrie ist, folgt $\|x\| \leq r$.

\subsection*{Satz}
Sei $\left(x_{n}\right)_{n} \subset E$ eine Folge mit $x_{n} \rightharpoonup x \in E$.\\
i). Ist $B \subset E^{\prime}$ präkompakt, so ist die Konvergenz $\varphi\left(x_{n}\right) \rightarrow \varphi(x)$ gleichmäßig in $\varphi \in B$, d.h.

$$
\sigma_{n}:=\sup _{\varphi \in B}\left|\varphi\left(x_{n}\right)-\varphi(x)\right| \rightarrow 0 \quad \text { für } n \rightarrow \infty .
$$

ii). Ist $\left(\varphi_{n}\right)_{n} \subset E^{\prime}$ eine Folge mit $\varphi_{n} \rightarrow \varphi$, so folgt $\varphi_{n}\left(x_{n}\right) \rightarrow \varphi(x)$.

Beweis. (i) Sei $y_{n}:=x_{n}-x$ für $n \in \mathbb{N}$, dann gilt $y_{n} \rightharpoonup 0$. Sei ferner $\varepsilon>0$. Wegen 11.3 ist $M:=\sup _{n \in \mathbb{N}}\left\|y_{n}\right\|<\infty$. Da $B$ präkompakt ist, existieren $\varphi_{1}, \ldots, \varphi_{k} \in B$ mit $B \subset \bigcup_{i=1}^{k} B_{\frac{\varepsilon}{2 M}}\left(\varphi_{i}\right)$. Ferner existiert $n_{0} \in \mathbb{N}$ mit $\left|\varphi_{i}\left(y_{n}\right)\right| \leq \frac{\varepsilon}{2}$ für $i=1, \ldots, k$ und $n \geq n_{0}$. Zu beliebigen $\varphi \in B$ existiert nun $i \in\{1, \ldots, k\}$ mit $\left\|\varphi-\varphi_{i}\right\| \leq \frac{\varepsilon}{2 M}$. Für $n \geq n_{0}$ ist also

$$
\left|\varphi\left(y_{n}\right)\right| \leq\left|\left[\varphi-\varphi_{i}\right]\left(y_{n}\right)\right|+\left|\varphi_{i}\left(y_{n}\right)\right| \leq \frac{\varepsilon}{2 M}\left\|y_{n}\right\|+\frac{\varepsilon}{2} \leq \varepsilon
$$

und somit $\sigma_{n}=\sup _{\varphi \in B}\left|\varphi\left(y_{n}\right)\right| \leq \varepsilon$.\\
(ii) Nach Voraussetzung ist $B:=\left\{\varphi_{k}: k \in \mathbb{N}\right\} \subset E^{\prime}$ relativ kompakt, also präkompakt. Ferner gilt

$$
\sigma_{n}=\sup _{k \in \mathbb{N}}\left|\varphi_{k}\left(x_{n}\right)-\varphi_{k}(x)\right| \rightarrow 0 \quad \text { für } n \rightarrow \infty
$$

also

$$
\left|\varphi_{n}\left(x_{n}\right)-\varphi(x)\right| \leq \underbrace{\left|\varphi_{n}\left(x_{n}\right)-\varphi_{n}(x)\right|}_{\leq \sigma_{n}}+\left|\varphi_{n}(x)-\varphi(x)\right| \rightarrow 0 \quad \text { für } n \rightarrow \infty .
$$ \(\square\)

\subsection*{Bemerkung}
Sei $T \in \mathcal{L}(E, F)$. Dann ist $T$ auch "schwach stetig" im folgenden Sinne: Ist $\left(x_{n}\right)_{n} \subset E$ eine Folge mit $x_{n} \rightharpoonup x \in E$, so folgt $T x_{n} \rightharpoonup T x$ in $F$. Für alle $\psi \in F^{\prime}$ gilt nämlich $\operatorname{mit} \varphi=\psi \circ T \in E^{\prime}$ :

$$
\lim _{n \rightarrow \infty} \psi\left(T x_{n}\right)=\lim _{n \rightarrow \infty} \varphi\left(x_{n}\right)=\varphi(x)=\psi(T x) .
$$

\subsection*{Satz}
Sei $T \in \mathcal{K}(E, F)$. Dann gilt:\\
i). $T^{\prime} \in \mathcal{K}\left(F^{\prime}, E^{\prime}\right)$.\\
ii). Ist $\left(x_{n}\right)_{n} \subset E$ eine Folge mit $x_{n} \rightharpoonup x \in E$, so folgt $T x_{n} \rightarrow T x$ in $F$.

Beweis. (i) Sei $\left(\psi_{n}\right)_{n} \subset F^{\prime}$ eine beschränkte Folge und $M:=\sup _{n \in \mathbb{N}}\left\|\psi_{n}\right\|$. Nach Voraussetzung ist $K:=\overline{T\left(B_{1}(0)\right)} \subset F$ kompakt. Sei $f_{n}:=\left.\psi_{n}\right|_{K} \in \mathcal{C}(K)$. Dann ist

$$
\left\|f_{n}\right\|_{\infty} \leq\left\|\psi_{n}\right\| C \leq M C \quad \text { für alle } n \in \mathbb{N} \text { mit } C:=\max _{y \in K}\|y\|
$$

und\\
$\left|f_{n}(y)-f_{n}(x)\right|=\left|\psi_{n}(y-x)\right| \leq\left\|\psi_{n}\right\|\|y-x\| \leq M\|y-x\| \quad$ für alle $x, y \in K, n \in \mathbb{N}$. Somit ist die Folge $\left(f_{n}\right)_{n}$ gleichgradig stetig und beschränkt in $\mathcal{C}(K)$. Nach Übergang zu einer Teilfolge ist $\left(f_{n}\right)_{n}$ gemäß 4.13 (Satz von Arzelà-Ascoli) also gleichmäßig konvergent und somit insbesondere eine Cauchyfolge in $\mathcal{C}(K)$. Da ferner\\
$\left\|T^{\prime} \psi_{m}-T^{\prime} \psi_{n}\right\|_{E^{\prime}}=\sup _{x \in B_{1}(0)}\left|\psi_{m}(T x)-\psi_{n}(T x)\right| \leq \sup _{y \in K}\left|f_{m}(y)-f_{n}(y)\right|=\left\|f_{m}-f_{n}\right\|_{\infty}$\\
für $m, n \in \mathbb{N}$, ist auch $\left(T^{\prime} \psi_{n}\right)_{n}$ eine Cauchyfolge in $E^{\prime}$ und somit konvergent, da $E^{\prime}$ vollständig ist. Dies zeigt die Kompaktheit von $T^{\prime}$.\\
(ii) Nach Voraussetzung gilt $y_{n}:=x_{n}-x \rightharpoonup 0$. Sei $\psi_{n} \in F^{\prime}$ gemäß 9.5 gewählt mit $\left\|\psi_{n}\right\|=1$ und $\left|\psi_{n}\left(T y_{n}\right)\right|=\left\|T y_{n}\right\|$. Nach (i) ist die Menge $\left\{T^{\prime} \psi_{n}: n \in \mathbb{N}\right\}$ präkompakt in $E^{\prime}$, und mit 11.6 (i) folgt

$$
\left\|T y_{n}\right\|=\left|\psi_{n}\left(T y_{n}\right)\right|=\left|\left[T^{\prime} \psi_{n}\right]\left(y_{n}\right)\right| \rightarrow 0 \quad \text { für } n \rightarrow \infty,
$$

also $T x_{n} \rightarrow T x$ in $F$. \(\square\)

\subsection*{Definition}
Eine Folge $\left(\psi_{n}\right)_{n} \subset E^{\prime}$ heißt schwach*-konvergent gegen $\psi \in E^{\prime}$, wenn sie auf $E$ punktweise gegen $\psi$ konvergiert, d.h. wenn gilt:

$$
\lim _{n \rightarrow \infty} \psi_{n}(x)=\psi(x) \quad \text { für alle } x \in E .
$$

Wir schreiben dann $\psi_{n} \xrightarrow{w^{*}} \psi$.

\subsection*{Beispiel und Bemerkung}
i). Vermöge der Abbildung $\ell^{\infty} \rightarrow\left(\ell^{1}\right)^{\prime}, x \mapsto \varphi_{x}$ aus Beispiel 9.15 können wir $\ell^{\infty}$ mit $\left(\ell^{1}\right)^{\prime}$ identifizieren. Sei $x^{n}=\left(x_{k}^{n}\right)_{k} \in l^{\infty}$ für $n \in \mathbb{N}$ definiert durch $x_{k}^{n}=0$ für $k \leq n$ und $x_{k}^{n}=1$ für $k>n$. Dann gilt $x^{n} \xrightarrow{w^{*}} 0$, d.h.

$$
\varphi_{x^{n}}(y)=\sum_{k=n+1}^{\infty} y_{k} \rightarrow 0 \quad \text { für alle } y \in \ell^{1} .
$$

Allerdings gilt $x^{n} \nrightarrow 0$ in $\ell^{\infty}$, vgl. Übungsblatt 11.\\
ii). Satz 11.4 gilt i.A. nicht für die schwach*-Konvergenz in $E^{\prime}$. Um dies zu sehen, betrachten wir den abgeschlossenen Teilraum $E \subset l^{\infty}$ aller konvergenten Folgen in $\mathbb{K}$ (in der Literatur oft mit (c) bezeichnet). Für $n \in \mathbb{N}$ sei ferner $x^{n} \in E$ wie in i) definiert. Schließlich sei $C \subset E^{\prime}$ die Menge aller $\varphi \in E^{\prime}$ mit $\lim _{n \rightarrow \infty} \varphi\left(x^{n}\right)=0$ für $n \rightarrow \infty$. Man sieht leicht, dass $C$ ein abgeschlossener Teilraum von $E^{\prime}$ (also insbesondere konvex) ist. Sei num $\psi_{m} \in C, m \in \mathbb{N}$ definiert durch $\psi_{m}(y)=y_{m}$ für $y=\left(y_{k}\right)_{k} \in E$, und sei $\psi \in E^{\prime}$ die Grenzwertabbildung definiert durch

$$
\psi(y)=\lim _{k \rightarrow \infty} y_{k} \quad \text { für } y=\left(y_{k}\right)_{k} \in E .
$$

Offensichtlich gilt

$$
\lim _{m \rightarrow \infty} \psi_{m}(y)=\psi(y) \quad \text { für alle } y \in E, \text { also } \quad \psi_{m} \stackrel{w^{*}}{\psi} \psi \in E^{\prime} .
$$

Allerdings ist $\psi\left(x^{m}\right)=1$ für alle $m \in \mathbb{N}$ und somit $\psi \notin C$.\\
iii). Korollar 11.5 gilt analog auch für die schwach*-Konvergenz: Sei $\left(\psi_{n}\right)_{n} \subset E^{\prime}$ eine Folge mit $\psi_{n} \xrightarrow{w^{*}} \psi \in E^{\prime}$. Dann gilt

$$
|\psi(x)|=\lim _{n \rightarrow \infty}\left|\psi_{n}(x)\right| \leq \liminf _{n \rightarrow \infty}\left\|\psi_{n}\right\|\|x\| \quad \text { für } x \in E
$$

und somit $\|\psi\| \leq \liminf _{n \rightarrow \infty}\left\|\psi_{n}\right\|$.

\subsection*{Hauptsatz}
Ist $E$ separabel, so hat jede beschränkte Folge in $E^{\prime}$ eine schwach*-konvergente Teilfolge.

Beweis. Sei $\left(x_{k}\right)_{k} \subset E$ eine dichte Folge in $E$, und sei $\left(\varphi_{n}\right)_{n}$ eine beschränkte Folge in $E^{\prime}$. Dann sind auch die Folgen $\left[\varphi_{n}\left(x_{k}\right)\right]_{n} \subset \mathbb{K}, k \in \mathbb{N}$ beschränkt. Insbesondere existiert eine Teilfolge $\left(\varphi_{n}^{1}\right)_{n}$ von $\left(\varphi_{n}\right)_{n}$ und $\lambda_{1} \in \mathbb{K}$ mit

$$
\lim _{n \rightarrow \infty} \varphi_{n}^{1}\left(x_{1}\right)=\lambda_{1} .
$$

Ferner existiert eine Teilfolge $\left(\varphi_{n}^{2}\right)_{n}$ von $\left(\varphi_{n}^{1}\right)_{n}$ und $\lambda_{2} \in \mathbb{K}$ mit

$$
\lim _{n \rightarrow \infty} \varphi_{n}^{2}\left(x_{2}\right)=\lambda_{2}
$$

Sukzessive findet man für alle $k \in \mathbb{N}, k \geq 2$ eine Teilfolge $\left(\varphi_{n}^{k}\right)_{n}$ von $\left(\varphi_{n}^{k-1}\right)_{n}$ und $\lambda_{k} \in \mathbb{K}$ mit

$$
\lim _{n \rightarrow \infty} \varphi_{n}^{k}\left(x_{k}\right)=\lambda_{k}
$$

Wir betrachten nun die Diagonalfolge $\left(\psi_{n}\right)_{n}$ gegeben durch $\psi_{n}:=\varphi_{n}^{n}$. Für diese Folge gilt offensichtlich

$$
(*) \quad \lim _{n \rightarrow \infty} \psi_{n}\left(x_{k}\right)=\lambda_{k} \quad \text { für alle } k \in \mathbb{N} \text {. }
$$

Sei nun $x \in E$ beliebig. Wir zeigen: $\left(\psi_{n}(x)\right)_{n}$ ist Cauchyfolge in $\mathbb{K}$, also konvergent. Sei dazu $\varepsilon>0$ und $k \in \mathbb{N}$ mit $\left\|x_{k}-x\right\|<\frac{\varepsilon}{3 M}$, wobei $M:=\sup _{n \in \mathbb{N}}\left\|\psi_{n}\right\|$ sei. Nach Voraussetzung ist $M<\infty$. Wegen (*) existiert $n_{\varepsilon} \in \mathbb{N}$ mit $\left|\psi_{n}\left(x_{k}\right)-\psi_{m}\left(x_{k}\right)\right|<\frac{\varepsilon}{3}$ für $n, m \geq n_{\varepsilon}$. Dann gilt für $n, m \geq n_{\varepsilon}$ :\\
$\left|\psi_{n}(x)-\psi_{m}(x)\right| \leq\left|\psi_{n}(x)-\psi_{n}\left(x_{k}\right)\right|+\left|\psi_{n}\left(x_{k}\right)-\psi_{m}\left(x_{k}\right)\right|+\left|\psi_{m}\left(x_{k}\right)-\psi_{m}(x)\right|<2 M\left\|x_{k}-x\right\|+\frac{\varepsilon}{3} \leq \varepsilon$.\\
Es folgt, dass $\psi(x):=\lim _{n \rightarrow \infty} \psi_{n}(x)$ für alle $x \in E$ existiert, und $\psi: E \rightarrow \mathbb{K}$ ist linear. Ferner ist

$$
|\psi(x)|=\lim _{n \rightarrow \infty}\left|\psi_{n}(x)\right| \leq M\|x\| \quad \text { für alle } x \in E,
$$

d.h. $\psi \in E^{\prime}$. Nach Konstruktion gilt schließlich $\psi_{n} \stackrel{w^{*}}{\rightharpoonup} \psi$.

\section*{§ 12 Reflexivität und gleichmäßige Konvexität}
Sei stets $(E,\|\cdot\|)$ ein normierter Raum über $\mathbb{K}$.

\subsection*{Definition}
Der $\operatorname{Raum}(E,\|\cdot\|)$ heißt reflexiv, wenn die lineare Isometrie aus 9.6 ,

$$
i_{E}: E \rightarrow E^{\prime \prime}, \quad \text { definiert durch } \quad i_{E}(x)(\varphi)=\varphi(x) \quad \text { für } x \in E, \varphi \in E^{\prime},
$$

surjektiv ist, also ein isometrischer Isomorphismus.

\subsection*{Beispiele und Bemerkung}
i). Ist $H$ ein Hilbertraum über $\mathbb{K}$, so ist $H$ reflexiv. Genauer gilt

$$
i_{H}=J_{H^{\prime}} \circ J_{H}: H \rightarrow H^{\prime} \rightarrow H^{\prime \prime}
$$

wobei $J_{H}: H \rightarrow H^{\prime}$ und $J_{H^{\prime}}: H^{\prime} \rightarrow H^{\prime \prime}$ die (antilinearen) isometrischen Isomorphismen aus dem Rieszschen Darstellungssatz 7.10 bzgl. der Hilberträume $H$ und $H^{\prime}$ (mit dem Skalarprodukt aus 7.10i) sind.\\
ii). Seien $p, q$ konjugierte Exponenten in $[1, \infty]$.

Die Abbildungen $S \in \mathcal{L}\left(l^{p},\left(l^{q}\right)^{\prime}\right)$ und $T \in \mathcal{L}\left(l^{q},\left(l^{p}\right)^{\prime}\right)$ definiert durch

$$
[S x](y)=\sum_{k \in \mathbb{N}} x_{k} y_{k}=[T y](x), \quad x=\left(x_{k}\right)_{k} \in l^{p}, y=\left(y_{k}\right)_{k} \in l^{q}
$$

sind Isometrien (Übungsaufgabe Blatt 10).\\
Mit $E:=l^{p}$ gilt dabei $T^{\prime} \circ i_{E}=S: l^{p} \rightarrow\left(l^{q}\right)^{\prime}$, denn für $x \in l^{p}$ und $y \in l^{q}$ ist

$$
\left[T^{\prime}\left(i_{E}(x)\right)\right](y)=i_{E}(x)(T y)=[T y](x)=[S x](y) .
$$

Sind num speziell $p, q \in(1, \infty)$, so sind $S$ und $T$ Isomorphismen (Übungsaufgabe Blatt 10), also ist auch $T^{\prime}$ eine Isomorphismus nach 9.11.\\
In diesem Fall ist also $i_{E}=\left(T^{\prime}\right)^{-1} S$ ein Isomorphismus, und somit ist der Raum $l^{p}$ reflexiv.\\
iii). Jeder reflexive normierte Raum ist isometrisch isomorph zu $E^{\prime \prime}$ und somit ein Banachraum, da $E^{\prime \prime}$ stets ein Banachraum ist. Allerdings ist nicht jeder

Banachraum $E$, welcher isometrisch isomorph zu $E^{\prime \prime}$ ist, auch reflexiv. Der sogenannte James-Raum $E$ ist isometrisch isomorph zu $E^{\prime \prime}$, aber nicht vermöge der Abbildung $i_{E}$. Genauer hat Bild $i_{E}$ Kodimension 1 in $E^{\prime \prime}$. Zur Definition von $E$ und dem Nachweis der genannten Eigenschaften siehe Robert James, A non-reflexive Banach space isometric with its second conjugate space, Proceedings National Academy of Sciences, Bd.37, 1951, S.174-177.

\subsection*{Hilfssatz}
i). Ist $E^{\prime}$ separabel, so auch $E$.\\
ii). Ist $E$ reflexiv und separabel, so ist auch $E^{\prime}$ separabel.

Beweis. (i) Sei $S:=\left\{\varphi \in E^{\prime}:\|\varphi\|=1\right\}$. Dann ist $S$ ebenfalls separabel, d.h. es existiert eine dichte Folge $\left(\varphi_{k}\right)_{k} \subset S$. Ferner existiert für jedes $k \in \mathbb{N}$ ein $x_{k} \in E$ mit $\left\|x_{k}\right\| \leq 1$ und $\varphi_{k}\left(x_{k}\right) \geq \frac{1}{2}$ (insbesondere reellwertig). Wir zeigen nun, dass die Menge $M:=\left\{x_{k}: k \in \mathbb{N}\right\}$ in $E$ total ist. Wäre $M^{\perp} \neq\{0\}$ und $\varphi \in S \cap M^{\perp}$, so würde $k \in \mathbb{N}$ existieren mit $\left\|\varphi-\varphi_{k}\right\|<\frac{1}{2}$ und somit

$$
\frac{1}{2} \leq \varphi_{k}\left(x_{k}\right)=\left|\varphi_{k}\left(x_{k}\right)-\varphi\left(x_{k}\right)\right| \leq\left\|x_{k}\right\|\left\|\varphi_{k}-\varphi\right\|<\frac{1}{2}
$$

Widerspruch. Es folgt also $M^{\perp}=\{0\}$ und somit $\overline{\operatorname{span} M}=M^{\perp \perp}=E$ nach 9.14 .\\
(ii) Da $E$ reflexiv ist, ist mit $E$ auch $E^{\prime \prime}$ separabel, also auch $E^{\prime}$ gemäß (i).

\subsection*{Bemerkung und Beispiele}
Aus der Separabilität von $E$ folgt im Allgemeinen nicht die Separabilität von $E^{\prime}$. Zum Beispiel ist $l^{1}$ separabel, aber $\left(l^{1}\right)^{\prime} \cong l^{\infty}$ ist nicht separabel (vgl. Übungsaufg. 2, Blatt 2). Insbesondere folgt mit 12.3 (ii), dass $l^{1}$ nicht reflexiv ist. Ist $K$ ein überabzählbarer kompakter metrischer Raum, so ist $C(K)$ separabel nach 4.28. Allerdings ist der Dualraum $C(K)^{\prime}$ nicht separabel, da die Diracmaße $\delta_{x} \in C(K)^{\prime}, x \in K$, definiert durch $\delta_{x}(f)=f(x)$ für $f \in C(K)$, eine überabzählbare diskrete Teilmenge von $C(K)^{\prime}$ bilden (vgl. Übungsblatt 4, Aufg. 4). Es folgt wiederum mit 12.3 (ii), dass $C(K)$ nicht reflexiv ist.

\subsection*{Satz}
Sei $E$ reflexiv und $F \subset E$ ein abgeschlossener Teilraum. Dann ist $F$ auch reflexiv.

Beweis. Sei $j: F \rightarrow E$ die Inklusion. Nach 9.10 ist das Diagramm\\
\includegraphics[max width=\textwidth, center]{2025_10_06_e8fc9b2ca9d05628f654g-116}\\
kommutativ. Wir bemerken zunächst: (*) Kern $j^{\prime \prime}=\{0\}$.\\
Ist nämlich $\Phi \in \operatorname{Kern} j^{\prime \prime} \subset F^{\prime \prime}$, so gilt

$$
0=\left[j^{\prime \prime} \Phi\right](\tau)=\Phi\left(j^{\prime} \tau\right)=\Phi\left(\left.\tau\right|_{F}\right) \quad \text { für alle } \tau \in E^{\prime}
$$

und somit $\Phi(\tau)=0$ für alle $\tau \in F^{\prime}$ gemäß dem Fortsetzungssatz 9.4, also $\Phi=0$.\\
Zum Beweis der Surjektivität von $i_{F}$ sei nun $\Psi \in F^{\prime \prime}$ beliebig. Nach Voraussetzung existiert ein $x \in E$ mit $i_{E}(x)=j^{\prime \prime} \Psi$. Wir zeigen zunächst: $x \in F^{\perp \perp}=\bar{F}=F$. Für $\varphi \in F^{\perp} \subset E^{\prime}$ ist nämlich $j^{\prime} \varphi=\left.\varphi\right|_{F} \in F^{\prime}$ die Nullabbildung und somit

$$
0=\Psi\left(j^{\prime} \varphi\right)=\left[j^{\prime \prime} \Psi\right](\varphi)=i_{E}(x)(\varphi)=\varphi(x) .
$$

Es folgt $x \in F^{\perp \perp}=\bar{F}=F$, also $x=j(x)$ und somit

$$
j^{\prime \prime}(\Psi)=i_{E}(j(x))=j^{\prime \prime}\left(i_{F}(x)\right), \quad \text { also } \quad j^{\prime \prime}\left(\Psi-i_{F}(x)\right)=0 .
$$

Dies liefert $\Psi=i_{F}(x)$ aufgrund von ( $*$ ). \(\square\)

\subsection*{Hauptsatz}
Ist $E$ reflexiv, so hat jede beschränkte Folge in $E$ eine schwach konvergente Teilfolge.

Beweis. Sei $\left(x_{n}\right)_{n} \subset E$ eine beschränkte Folge und $F:=\overline{\operatorname{span}\left\{x_{n}: n \in \mathbb{N}\right\}} \subset E$. Dann ist $F$ separabel, und $F$ ist auch reflexiv nach 12.5. Also ist auch $F^{\prime}$ separabel nach 12.3(ii). Sei nun $\Phi_{n}:=i_{F}\left(x_{n}\right) \in F^{\prime \prime}$ für $n \in \mathbb{N}$. Gemäß 11.11 können wir nach Übergang zu einer Teilfolge annehmen:

$$
\Phi_{n} \xrightarrow{w *} \Phi \quad \text { in } F^{\prime \prime}=\left(F^{\prime}\right)^{\prime}
$$

Da $F$ reflexiv ist, existiert $x \in F$ mit $i_{F}(x)=\Phi$, und für alle $\varphi \in E^{\prime}$ gilt dann mit $\psi=\left.\varphi\right|_{F} \in F^{\prime}$ :

$$
\varphi\left(x_{n}-x\right)=\psi\left(x_{n}-x\right)=\left[\Phi_{n}-\Phi\right](\psi) \rightarrow 0 \quad \text { für } n \rightarrow \infty .
$$ \(\square\)

\subsection*{Bemerkung}
i). Satz 12.6 liefert zusammen mit 11.5, dass in jedem reflexiven Raum die Einheitskugel "schwach folgenkompakt" ist.\\
ii). Ist $E$ separabel, so gilt sogar (ohne Beweis): $E$ ist genau dann reflexiv, wenn jede beschränkte Folge in $E$ eine schwach konvergente Teilfolge besitzt.\\
iii). Aus 12.6 folgt insbesondere, dass in jedem Hilbertraum jede beschränkte Folge eine schwach konvergente Teilfolge besitzt.

\subsection*{Bemerkung}
Ist $E$ reflexiv und $F$ topologisch isomorph zu $E$, so ist auch $F$ reflexiv. Dies folgt direkt aus 9.10(ii) und 9.11(ii). Insbesondere ist jeder endlich dimensionale normierte Raum reflexiv, da jeder solche Raum topologisch isomorph zu einem Hilbertraum ist.

\subsection*{Satz}
i). Ist $E$ reflexiv, so auch $E^{\prime}$.\\
ii). Ist $E$ ein Banachraum und $E^{\prime}$ reflexiv, so ist auch $E$ reflexiv.

Beweis. (i) Zu zeigen ist die Surjektivität von $i_{E^{\prime}}: E^{\prime} \rightarrow E^{\prime \prime \prime}$. Sei dazu $\Phi \in E^{\prime \prime \prime}$, und sei $\varphi:=\Phi \circ i_{E} \in E^{\prime}$. Dann gilt

$$
\Phi\left(i_{E}(x)\right)=\varphi(x)=\left[i_{E}(x)\right](\varphi) \quad \text { für alle } x \in E .
$$

Da $i_{E}$ nach Voraussetzung surjektiv ist, erhält man hieraus

$$
\Phi(\Psi)=\Psi(\varphi)=i_{E^{\prime}}(\varphi)(\Psi) \quad \text { für alle } \Psi \in E^{\prime \prime} .
$$

Es folgt $\Phi=i_{E^{\prime}}(\varphi)$, wie behauptet.\\
Randbemerkung: Dieser Beweis zeigt tatsächlich folgende Identität: $i_{E^{\prime}}=\left[i_{E}{ }^{\prime}\right]^{-1}: E^{\prime \prime \prime} \rightarrow E^{\prime}$.\\
(ii) Ist $E^{\prime}$ reflexiv, so ist nach (i) auch $E^{\prime \prime}$ reflexiv. Da $E$ vermöge $i_{E}$ zu einem (nach Voraussetzung) vollständigen und somit abgeschlossenen Teilraum von $E^{\prime \prime}$ topologisch isomorph ist, folgt mit 12.5 und 12.8, dass auch $E$ reflexiv ist.

\subsection*{Bemerkung}
Aus 12.4 und 12.9 (ii) folgt insbesondere, dass $l^{\infty}$ nicht reflexiv ist.

Im Folgenden wollen wir ein praktisches hinreichendes Kriterium für Reflexivität herleiten, die sogenannte gleichmäßige Konvexität eines Banachraums. Zuvor benötigen wir aber zwei Hilfssätze, welche auch anderweitig nützlich sind.

\subsection*{Hilfssatz ("komplexer" Trennungssatz in Ergänzung zu Kapitel 10)}
Sei $C \subset E$ absolut konvex und abgeschlossen, und sei $\eta \in E \backslash C$. Dann existiert $\psi \in E^{\prime}$ und $\alpha>0$ mit

$$
|\psi(y)|<\alpha \quad \text { für } y \in C \quad \text { und } \quad|\psi(\eta)|>\alpha .
$$

Beweis. Gemäß 10.3 (ii) existiert $\varphi \in \mathcal{L}(E, \mathbb{R})$ und $\alpha \in \mathbb{R}$ mit $\left.\varphi\right|_{C}<\alpha<\varphi(\eta)$. Wegen $0 \in C$ (aufgrund der absoluten Konvexität) erzwingt dies $\alpha>0$.\\
$\operatorname{Im}$ Fall $\mathbb{K}=\mathbb{R}$ erhält man dann auch $-\varphi(y)=\varphi(-y)<\alpha$ für $y \in C$, und die Behauptung folgt mit $\psi=\varphi$.\\
$\operatorname{Im}$ Fall $\mathbb{K}=\mathbb{C}$ betrachten wir nun die ( $\mathbb{C}$-lineare) Abbildung $\psi \in E^{\prime}$ definiert durch $\psi(y)=\varphi(y)-i \varphi(i y)$ (vgl. den Beweis von 9.3!). Für jedes $y \in C$ existiert dann ein $a \in \mathbb{C},|a|=1$ mit $|\psi(y)|=a \psi(y)=\psi(a y) \in \mathbb{R}$ und wegen $a y \in C$ somit $|\psi(y)|=\varphi(a y)<\alpha$. Ferner ist $|\psi(\eta)| \geq \operatorname{Re} \psi(\eta)=\varphi(\eta)>\alpha$.

Im Rest des Kapitels bezeichne stets $B:=B_{1}(0)$ die abgeschlossene Einheitskugel des normierten Raums $(E,\|\cdot\|)$.

\subsection*{Hilfssatz}
Seien $\Phi \in E^{\prime \prime}$ mit $\|\Phi\|_{E^{\prime \prime}} \leq 1$ und $\varphi_{1}, \ldots, \varphi_{n} \in E^{\prime}$ gegeben. Sei ferner $\varepsilon>0$. Dann existiert ein $x \in B$ mit

$$
\left|\Phi\left(\varphi_{i}\right)-\varphi_{i}(x)\right|<\varepsilon \quad \text { für } i=1, \ldots, n .
$$

Beweis. Es reicht, zu zeigen, dass $\eta:=\left(\Phi\left(\varphi_{1}\right), \ldots, \Phi\left(\varphi_{n}\right)\right) \in \mathbb{K}^{n}$ in der Menge

$$
C:=\overline{\left\{\left(\varphi_{1}(x), \ldots, \varphi_{n}(x)\right): x \in B\right\}} \subset \mathbb{K}^{n}
$$

liegt. Angenommen, es wäre $\eta \notin C$. Da $C$ absolut konvex und abgeschlossen in $\mathbb{K}^{n}$ ist, existiert gemäß 12.11 ein $\psi \in\left(\mathbb{K}^{n}\right)^{\prime}$ und $\alpha>0$ mit

$$
|\psi(y)|<\alpha \quad \text { für } y \in C \quad \text { und } \quad|\psi(\eta)|>\alpha .
$$

Dabei lässt sich $\psi$ schreiben als $\psi(y)=\sum_{i=1}^{n} \lambda_{i} y_{i}$ für $y \in \mathbb{K}^{n}$ mit geeigneten $\lambda_{1}, \ldots, \lambda_{n} \in \mathbb{K}$. Sei nun

$$
\widetilde{\psi}:=\psi \circ\left(\varphi_{1}, \ldots, \varphi_{n}\right)=\sum_{i=1}^{n} \lambda_{i} \varphi_{i} \in E^{\prime}
$$

Dann gilt $|\widetilde{\psi}(x)|=\left|\psi\left(\varphi_{1}(x), \ldots, \varphi_{n}(x)\right)\right|<\alpha$ für alle $x \in B$, also $\|\widetilde{\psi}\| \leq \alpha$. Andererseits ist

$$
|\Phi(\widetilde{\psi})|=\left|\sum_{i=1}^{n} \lambda_{i} \Phi\left(\varphi_{i}\right)\right|=|\psi(\eta)|>\alpha
$$

Dies widerspricht der Voraussetzung $\|\Phi\|_{E^{\prime \prime}} \leq 1$. Die Behauptung folgt.\\
Wir ergänzen einen alternativen (nicht ganz so anschaulichen) Beweis von 12.12, welcher ohne 12.11 auskommt und stattdessen 7.2 verwendet (kein Prüfungsstoff!).\\
2. Beweis. Sei $T \in \mathcal{L}\left(E, \mathbb{K}^{n}\right)$ definiert durch $T x=\left(\varphi_{1}(x), \ldots, \varphi_{n}(x)\right)$ und $C:= \overline{T(B)} \subset \mathbb{K}^{n}$. Dann ist $C$ konvex und abgeschlossen in $\mathbb{K}^{n}$. Sei ferner $\eta= \left(\Phi\left(\varphi_{1}\right), \ldots, \Phi\left(\varphi_{n}\right)\right)$. Es reicht wiederum, zu zeigen, dass $\eta \in C$ ist.\\
Sei dazu $P: \mathbb{K}^{n} \rightarrow C$ die abstandsminimierende Projektion aus 7.1 bezüglich des Standard-Skalarprodukts $\langle\cdot, \cdot\rangle$ auf $\mathbb{K}^{n}$ und $\delta:=\eta-P(\eta) \in \mathbb{K}^{n}$. Sei ferner $\varphi \in E^{\prime}$ definiert durch

$$
\sum_{i=1}^{n} \overline{\delta_{i}} \varphi_{i}, \quad \text { also } \quad \varphi(x)=\sum_{i=1}^{n} \varphi_{i}(x) \overline{\delta_{i}}=\langle T x, \delta\rangle \quad \text { für } x \in E .
$$

Dann gilt für alle $x \in B$ gemäß 7.2.\\
$0 \geq \operatorname{Re}\langle T x-P(\eta), \eta-P(\eta)\rangle \geq \operatorname{Re} \varphi(x)-\operatorname{Re}\langle P(\eta), \delta\rangle, \quad$ also $\operatorname{Re} \varphi(x) \leq \operatorname{Re}\langle P(\eta), \delta\rangle$.\\
Dies liefert (nach Betrachtung von $\alpha x$ anstelle von $x$ mit $\alpha \in \mathbb{K},|\alpha|=1$ ) auch $|\varphi(x)| \leq \operatorname{Re}\langle P(\eta), \delta\rangle$ für alle $x \in B$, also\\
$\|\varphi\| \leq \operatorname{Re}\langle P(\eta), \delta\rangle=\operatorname{Re}\langle\eta, \delta\rangle-|\delta|_{2}^{2} \leq \operatorname{Re}\langle\eta, \delta\rangle \leq|\langle\eta, \delta\rangle|=\left|\sum_{i=1}^{n} \overline{\delta_{i}} \Phi\left(\varphi_{i}\right)\right|=|\Phi(\varphi)| \leq\|\varphi\|$\\
Es folgt $|\delta|_{2}^{2}=0$, also $\delta=0$, wie gewünscht.

\subsection*{Definition}
Der Raum $E$ heißt gleichmäßig konvex, wenn für alle $\varepsilon>0$ ein $\delta>0$ existiert derart, dass für alle $x, y \in B$ die Implikation gilt:

$$
\left\|\frac{x+y}{2}\right\|>1-\delta \quad \Longrightarrow \quad\|x-y\|<\varepsilon
$$

\subsection*{Beispiele}
(i) Ist ( $H,\langle\cdot, \cdot\rangle$ ) ein Prähilbertraum, so ist $H$ gleichmäßig konvex, denn es gilt $\frac{1}{4}\|x-y\|^{2}=\frac{1}{2}\left(\|x\|^{2}+\|y\|^{2}\right)-\left\|\frac{x+y}{2}\right\|^{2} \leq 1-\left\|\frac{x+y}{2}\right\|^{2} \quad$ für alle $x, y \in B$.\\
(ii) Der Raum ( $\mathbb{R}^{N},|\cdot|_{2}$ ) ist gemäß (i) gleichmäßig konvex; allgemeiner ist ( $\mathbb{R}^{N},|\cdot|_{p}$ ) für $1<p<\infty$ gleichmäßig konvex. Die Räume ( $\mathbb{R}^{N},|\cdot|_{1}$ ) und ( $\mathbb{R}^{N},|\cdot|_{\infty}$ ) sind aber nicht gleichmäßig konvex.

\subsection*{Satz}
Sei $E$ gleichmäßig konvex und $\left(x_{n}\right)_{n}$ eine Folge in $E$.\\
i). Gilt $x_{n} \rightharpoonup x \in E$ und $\left\|x_{n}\right\| \rightarrow\|x\|$, so folgt $x_{n} \rightarrow x$ für $n \rightarrow \infty$.\\
ii). Gilt $\lim _{n \rightarrow \infty} \sup _{m \geq n}\left|\left\|x_{m}+x_{n}\right\|-2\right|=0$, so ist $\left(x_{n}\right)_{n}$ eine Cauchyfolge.

Beweis. (i) Ohne Einschränkung sei $x \neq 0$ (sonst folgt offensichtlich die starke Konvergenz) und $x_{n} \neq 0$ für alle $n \in \mathbb{N}$. Setze $z_{n}:=\frac{x_{n}}{\left\|x_{n}\right\|}$ und $z=\frac{x}{\|x\|}$. Dann gilt

$$
\varphi\left(z_{n}\right)=\frac{1}{\left\|x_{n}\right\|} \varphi\left(x_{n}\right) \rightarrow \frac{1}{\|x\|} \varphi(x)=\varphi(z) \quad \text { für } n \rightarrow \infty \text { und alle } \varphi \in E^{\prime},
$$

also $z_{n} \rightharpoonup z$. Nach 9.5 existiert $\varphi \in E^{\prime}$ mit $\|\varphi\|=1$ und $\varphi(z)=1$. Es folgt

$$
\left\|z_{n}+z\right\| \geq\left|\varphi\left(z_{n}+z\right)\right| \xrightarrow{n \rightarrow \infty} 2 \varphi(z)=2 .
$$

Da $E$ gleichmäßig konvex ist, folgt $\left\|z-z_{n}\right\| \rightarrow 0$ für $n \rightarrow \infty$, also

$$
\lim _{n \rightarrow \infty} x_{n}=\lim _{n \rightarrow \infty}\left\|x_{n}\right\| z_{n}=\|x\| z=x
$$

(ii) Nach Voraussetzung gilt $\sup _{m \geq n}\left\|\frac{x_{n}+x_{m}}{2}\right\| \rightarrow 1$ für $n \rightarrow \infty$.

\begin{enumerate}
  \item Spezialfall: $x_{n} \in B$ für alle $n \in \mathbb{N}$. Sei $\varepsilon>0$, und sei $\delta>0$ so gewählt, dass 12.1 gilt. Nach Voraussetzung existiert ein $n_{0} \in \mathbb{N}$ mit $\left\|\frac{x_{m}+x_{n}}{2}\right\|>1-\delta$ für $m \geq n \geq n_{0}$; also auch $\left\|x_{m}-x_{n}\right\|<\varepsilon$ für $m \geq n \geq n_{0}$. Dies zeigt, dass $\left(x_{n}\right)_{n}$ eine Cauchyfolge ist.
  \item Allgemeiner Fall: Die Voraussetzung liefert (mit $m=n$ ) insbesondere $\lim _{n \rightarrow \infty}\left\|x_{n}\right\|=1$, also ist ohne Einschränkung $\alpha_{n}:=\left\|x_{n}\right\| \neq 0$ für alle $n \in \mathbb{N}$. Setze $z_{n}:=\frac{x_{n}}{\alpha_{n}}$. Für $m, n \in \mathbb{N}$ ist dann
\end{enumerate}

$$
\left\|z_{n}+z_{m}\right\| \leq\left\|z_{n}\right\|+\left\|z_{m}\right\|=2
$$

und\\
$\left\|z_{n}+z_{m}\right\|=\left\|x_{n}-\left(\alpha_{n}-1\right) z_{n}+x_{m}-\left(\alpha_{m}-1\right) z_{m}\right\| \geq\left\|x_{n}+x_{m}\right\|-\left|\alpha_{n}-1\right|-\left|\alpha_{m}-1\right|$.\\
Es folgt also $\lim _{n \rightarrow \infty} \sup _{m \geq n}\left\|z_{n}+z_{m}\right\| \leq 2$ sowie

$$
\lim _{n \rightarrow \infty} \inf _{m \geq n}\left\|z_{n}+z_{m}\right\| \geq \lim _{n \rightarrow \infty} \inf _{m \geq n}\left\|x_{n}+x_{m}\right\| \geq 2
$$

nach Voraussetzung. Der Spezialfall liefert daher, dass $\left(z_{n}\right)_{n}$ eine Cauchyfolge ist. Da ferner\\
$\left\|x_{n}-x_{m}\right\|=\left\|z_{n}-\left(1-\alpha_{n}\right) z_{n}-z_{m}+\left(1-\alpha_{m}\right) z_{m}\right\| \leq\left\|z_{n}-z_{m}\right\|+\left|1-\alpha_{n}\right|+\left|1-\alpha_{m}\right|$\\
für $n, m \in \mathbb{N}$ gilt, ist auch $\left(x_{n}\right)_{n}$ eine Cauchyfolge.

\subsection*{Satz von Milman}
Ist $E$ ein gleichmäßig konvexer Banachraum, so ist $E$ reflexiv.

Beweis. Sei $\Phi \in E^{\prime \prime}$ mit $\|\Phi\|=1$. Es reicht, zu zeigen, dass $\Phi \in$ Bild $i_{E}$ ist. Sei dazu $\left(\varepsilon_{n}\right)_{n}$ eine Nullfolge positiver Zahlen und $\varphi \in E^{\prime}$ beliebig mit $\|\varphi\|=1$. Nach Definition der Norm in $E^{\prime \prime}$ existieren $\alpha_{n} \in E^{\prime}$ mit $\left\|\alpha_{n}\right\| \leq 1$ und

$$
\Phi\left(\alpha_{n}\right)=\left|\Phi\left(\alpha_{n}\right)\right|>1-\varepsilon_{n}
$$

Nach 12.12 existiert eine Folge $\left(x_{n}\right)_{n} \subset B$ derart, dass für alle $n \in \mathbb{N}$ gilt:

$$
\begin{aligned}
\left|\Phi(\varphi)-\varphi\left(x_{n}\right)\right| & <\varepsilon_{n} & & \text { und } \\
\left|\Phi\left(\alpha_{i}\right)-\alpha_{i}\left(x_{n}\right)\right| & <\varepsilon_{n} & & \text { für } i=1, \ldots, n .
\end{aligned}
$$

Wir zeigen: Als Konsequenz von $(*)$ ist $\left(x_{n}\right)_{n}$ eine Cauchyfolge.\\
Für $1 \leq i \leq n$ ist nämlich\\
$\operatorname{Re} \alpha_{i}\left(x_{n}\right)=\operatorname{Re} \Phi\left(\alpha_{i}\right)-\operatorname{Re}\left[\Phi\left(\alpha_{i}\right)-\alpha_{i}\left(x_{n}\right)\right]>1-\varepsilon_{i}-\left|\Phi\left(\alpha_{i}\right)-\alpha_{i}\left(x_{n}\right)\right|>1-\varepsilon_{i}-\varepsilon_{n}$.\\
Für $m, n \in \mathbb{N}$ mit $m \geq n$ folgt also

$$
2 \geq\left\|x_{m}\right\|+\left\|x_{n}\right\| \geq\left\|x_{m}+x_{n}\right\| \geq\left|\alpha_{n}\left(x_{m}+x_{n}\right)\right| \geq \operatorname{Re} \alpha_{n}\left(x_{m}+x_{n}\right)>2-\varepsilon_{m}-3 \varepsilon_{n}
$$

und somit

$$
\lim _{n \rightarrow \infty} \sup _{m \geq n}\left|\left\|x_{m}+x_{n}\right\|-2\right|=0
$$

Mit 12.15 (ii) folgt, dass $\left(x_{n}\right)_{n}$ eine Cauchyfolge in $E$ ist. Da $E$ vollständig ist, existiert $x:=\lim _{n \rightarrow \infty} x_{n}$, wobei gilt:

$$
\Phi(\varphi)=\varphi(x) \quad \text { und } \quad \Phi\left(\alpha_{i}\right)=\alpha_{i}(x) \quad \text { für } i \in \mathbb{N} .
$$

Jede andere Wahl von $\varphi \in E^{\prime}$ führt aber zu demselben $x \in E$, denn erfüllen $x, y \in E$ die Bedingungen

$$
\Phi\left(\alpha_{i}\right)=\alpha_{i}(x) \quad \text { und } \quad \Phi\left(\alpha_{i}\right)=\alpha_{i}(y) \quad \text { für } i \in \mathbb{N} \text {, }
$$

so erfüllt die Folge ( $x, y, x, y, \ldots$ ) Bedingung ( $*$ ) anstelle von ( $\left.x_{n}\right)_{n}$ und ist somit nach obigem Argument eine Cauchyfolge, d.h. es gilt $x=y$.\\
Fazit: Für das oben gefundene $x \in E$ gilt

$$
\Phi(\varphi)=\varphi(x) \quad \text { für alle } \varphi \in E^{\prime},
$$

also $\Phi=i_{E}(x) \in$ Bild $i_{E}$. \(\square\)

\section*{Anwendung auf die $L^{p}$-Räume}
Sei im Folgenden ( $\Omega, \mathcal{M}, \mu$ ) ein Maßraum.

\subsection*{Hilfssatz}
Für $p \in[2, \infty)$ und $u, v \in L^{p}(\Omega)$ gilt

$$
\|u+v\|_{p}^{p}+\|u-v\|_{p}^{p} \leq 2^{p-1}\left(\|u\|_{p}^{p}+\|v\|_{p}^{p}\right)
$$

Beweis. Es reicht offensichtlich, die Ungleichung

$$
|z+w|^{p}+|z-w|^{p} \leq 2^{p-1}\left(|z|^{p}+|w|^{p}\right) \quad \text { für } z, w \in \mathbb{C}
$$

zu zeigen. Wir verwenden dazu, dass für $a, b \geq 0$ gilt:

$$
\left(a^{p}+b^{p}\right)^{\frac{1}{p}} \leq\left(a^{2}+b^{2}\right)^{\frac{1}{2}} \quad \text { und } \quad a^{2}+b^{2} \leq 2^{\frac{p-2}{p}}\left(a^{p}+b^{p}\right)^{\frac{2}{p}}
$$

Damit erhält man nämlich

$$
\begin{aligned}
|z+w|^{p}+|z-w|^{p} & \leq\left(|z+w|^{2}+|z-w|^{2}\right)^{\frac{p}{2}}=\left[2\left(|z|^{2}+|w|^{2}\right)\right]^{\frac{p}{2}} \\
& =2^{\frac{p}{2}}\left(|z|^{2}+|w|^{2}\right)^{\frac{p}{2}} \leq 2^{\frac{p}{2}} 2^{\frac{p-2}{2}}\left(|z|^{p}+|w|^{p}\right)=2^{p-1}\left(|z|^{p}+|w|^{p}\right)
\end{aligned}
$$

für $z, w \in \mathbb{C}$, also (12.2). Nun zu (12.3): Wegen $p \geq 2$ ist die Funktion $[0, \infty) \rightarrow \mathbb{R}$, $t \mapsto t^{p / 2}$ konvex; somit gilt $\left(\frac{a^{2}+b^{2}}{2}\right)^{\frac{p}{2}} \leq \frac{a^{p}+b^{p}}{2}$ für $a, b \geq 0$. Dies zeigt die zweite Ungleichung in (12.3). In der ersten Ungleichung reicht es aufgrund der Homogenität, den Fall $a^{2}+b^{2}=1$ zu betrachten. Dann sind $a, b \in[0,1]$, und somit ist $a^{p}+b^{p} \leq a^{2}+b^{2}=1$, da $p \geq 2$ ist. Die Ungleichung folgt. \(\square\)

\subsection*{Satz}
Für $p \in[2, \infty)$ ist $L^{p}(\Omega)$ gleichmäßig konvex.\\
Beweis. Für $u, v \in L^{p}(\Omega)$ mit $\|u\|_{p},\|v\|_{p} \leq 1$ ist

$$
\|u-v\|_{p}^{p} \stackrel{\overline{12.17}}{\leq} 2^{p-1}\left(\|u\|_{p}^{p}+\|v\|_{p}^{p}\right)-\|u+v\|_{p}^{p} \leq 2^{p}\left(1-\left\|\frac{u+v}{2}\right\|_{p}^{p}\right)
$$

und dies liefert die gleichmäßige Konvexität.

\subsection*{Korollar}
Für $p \in(1, \infty)$ ist $L^{p}(\Omega)$ reflexiv, und mit $q=\frac{p}{p-1}$ hat man einen isometrischen Isomorphismus

$$
S: L^{q}(\Omega) \rightarrow\left(L^{p}(\Omega)\right)^{\prime}, \quad[S f](g)=\int_{\Omega} f g d \mu \quad \text { für } f \in L^{q}(\Omega), g \in L^{p}(\Omega)
$$

Beweis. Die Höldersche Ungleichung liefert

$$
|[S f](g)|=\left|\int_{\Omega} f g d \mu\right| \leq\|f\|_{q}\|g\|_{p} \quad \text { für } f \in L^{q}(\Omega), g \in L^{p}(\Omega)
$$

also ist $S f \in\left(L^{p}(\Omega)\right)^{\prime}$ mit $\|S f\| \leq\|f\|_{q}$. Ist speziell $g \in L^{p}(\Omega)$ definiert durch

$$
g(x)= \begin{cases}|f(x)|^{\frac{q}{p}-1} \overline{f(x)}=|f(x)|^{q-2} \overline{f(x)}, & \text { falls } f(x) \neq 0 \\ 0, & \text { falls } f(x)=0\end{cases}
$$

so erhält man zudem

$$
\|g\|_{p}=\|f\|_{q}^{q / p} \quad \text { und } \quad|[S f](g)|=\|f\|_{q}^{q}
$$

also $\|S f\| \geq \frac{|[S f](g)|}{\|g\|_{p}}=\|f\|_{q}^{q-q / p}=\|f\|_{q}$. Insgesamt folgt also $\|S f\|=\|f\|_{q}$, und somit ist $S$ eine Isometrie.\\
Wir betrachten num zunächst den Fall $p \in[2, \infty)$. Dann ist $E:=L^{p}(\Omega)$ reflexiv nach 12.18 und 12.16. Wir zeigen, dass $S$ in diesem Fall surjektiv ist. Sei dazu $F:=$ Bild $S \subset E^{\prime}$. Dann ist $F$ vollständig (als Bild eines Banachraums unter einer Isometrie) und somit abgeschlossen in $E^{\prime}$. Angenommen, es gäbe $\varphi \in E^{\prime} \backslash F$. Dann existiert auch $\varepsilon>0$ mit $U_{\varepsilon}(\varphi) \cap F=\varnothing$, und nach dem Trennungssatz von Mazur existiert $\Psi \in E^{\prime \prime}$ mit $\left.\Psi\right|_{F} \equiv 0$ und $\Psi(\varphi) \neq 0$ (dies kann man alternativ auch direkt aus 9.4 folgern, s.a. den Beweis von 9.14 (i)).

Aufgrund der Reflexivität von $E$ ist $\Psi=i_{E}(u)$ für ein $u \in E$. Für alle $v \in L^{q}(\Omega)$ ist dann

$$
0=\Psi(S v)=[S v](u)=\int_{\Omega} u v d \mu
$$

Durch Wahl von $v=|u|^{\frac{p}{q}-1} \bar{u}$ folgt leicht $u=0$, also $\Psi=0$. Widerspruch. Also ist $S$ surjektiv und damit ein isometrischer Isomorphismus.\\
Mit 12.9 (i) folgt nun auch, dass $L^{q}(\Omega) \cong E^{\prime}$ reflexiv ist, und mit dem gleichen Argument wie oben ist dann auch die Abbildung

$$
S: L^{p}(\Omega) \rightarrow\left(L^{q}(\Omega)\right)^{\prime}, \quad[S f](g)=\int_{\Omega} f g d \mu
$$

ein isometrischer Isomorphismus.\\
Man durchläuft also alle Zahlen $p \in[2, \infty)$ und $q \in(1,2]$ und erhält somit die Behauptung. \(\square\)

\subsection*{Bemerkung}
Man kann zeigen, dass $L^{p}(\Omega)$ auch für $1<p<2$ gleichmäßig konvex ist. Dazu verwendet man die elementare, aber nicht offensichtliche Ungleichung

$$
\|u+v\|_{p}^{q}+\|u-v\|_{p}^{q} \leq 2\left(\|u\|_{p}^{p}+\|v\|_{p}^{p}\right)^{q-1}, \quad u, v \in L^{p}(\Omega)
$$

für $1<p \leq 2$ und $q:=\frac{p}{p-1}$.\\
Beweis. Siehe z.B. Hirzebruch, Scharlau, Einführung in die Funktionalanalysis, Lemma 17.3. \(\square\)

\subsection*{Satz}
Das Maß $\mu$ sei $\sigma$-endlich, d.h. es gebe $\Omega_{k} \subset \Omega, k \in \mathbb{N}$ mit $\mu\left(\Omega_{k}\right)<\infty$ und $\Omega=\bigcup_{k=1}^{\infty} \Omega_{k}$. Dann ist ( $\left.L^{1}(\Omega)\right)^{\prime}$ isometrisch isomorph zu $L^{\infty}(\Omega)$; zu jedem $\varphi \in\left(L^{1}(\Omega)\right)^{\prime}$ existiert genau ein $f \in L^{\infty}(\Omega)$ mit $\|f\|_{\infty}=\|\varphi\|$ und

$$
\varphi(g)=\int_{\Omega} f g d \mu \quad \text { für alle } g \in L^{1}(\Omega) \text {. }
$$

Beweis. Ohne Einschränkung können wir annehmen, dass $\Omega_{k} \subset \Omega_{k+1}$ für alle $k \in \mathbb{N}$ gilt. Sei $\varphi \in\left(L^{1}(\Omega)\right)^{\prime}$, und sei $\varphi_{k} \in\left(L^{1}\left(\Omega_{k}\right)\right)^{\prime}$ definiert durch $\varphi_{k}(g)=\varphi(\tilde{g})$, wobei $\tilde{g} \in L^{1}(\Omega)$ die triviale Fortsetzung von $g \in L^{1}\left(\Omega_{k}\right)$ auf $\Omega$ bezeichne. Da $L^{2}\left(\Omega_{k}\right)$ stetig in $L^{1}\left(\Omega_{l}\right)$ eingebettet ist (gemäß 6.5(ii)), ist $\left.\varphi_{k}\right|_{L^{2}\left(\Omega_{k}\right)} \in\left(L^{2}\left(\Omega_{k}\right)\right)^{\prime}$. Somit existiert\\
nach 12.19 (oder dem Rieszschen Darstellungssatz, vgl. 7.10) genau ein $f_{k} \in L^{2}\left(\Omega_{k}\right)$ mit

$$
\varphi_{k}(g)=\int_{\Omega_{k}} f_{k} g d \mu \quad \text { für alle } g \in L^{2}\left(\Omega_{k}\right)
$$

Die Eindeutigkeit zeigt dabei $\left.f_{k}\right|_{\Omega_{j}}=f_{j}$ für $j \leq k$, und somit existiert eine $\mu$ messbare Funktion $f: \Omega \rightarrow \mathbb{K}$ mit $\left.f\right|_{\Omega_{k}}=f_{k}$ für $k \in \mathbb{N}$. Wir zeigen

$$
f \in L^{\infty}(\Omega) \quad \text { mit } \quad\|f\|_{\infty} \leq\|\varphi\| .
$$

Sei dazu

$$
\alpha: \Omega \rightarrow \mathbb{K}, \quad \alpha(x):= \begin{cases}\frac{\overline{f(x)}}{|f(x)|}, & f(x) \neq 0 \\ 0, & f(x)=0\end{cases}
$$

d.h. $\alpha \in L^{\infty}(\Omega)$ erfüllt $|\alpha| \leq 1$ und $\alpha f=|f|$ in $\Omega$. Sei ferner

$$
A_{n}:=\left\{x \in \Omega:|f(x)| \geq\|\varphi\|+\frac{1}{n}\right\} \quad \text { für } n \in \mathbb{N} .
$$

Wäre $\mu\left(A_{n}\right)>0$ für ein $n \in \mathbb{N}$, so würde $k \in \mathbb{N}$ existieren mit $\mu(A)>0$ für $A:=\Omega_{k} \cap A_{n}$, und mit $g:=\frac{\alpha 1_{A}}{\mu(A)} \in L^{2}\left(\Omega_{k}\right)$ wäre dann

$$
\|\varphi\|+\frac{1}{n} \leq \frac{1}{\mu(A)} \int_{A}|f| d \mu=\int_{\Omega_{k}} f g d \mu=\varphi_{k}(g)=\varphi(\tilde{g}) \leq\|\varphi\|\|\tilde{g}\|_{L^{1}(\Omega)}=\|\varphi\|
$$

Widerspruch. Es folgt $\mu\left(A_{n}\right)=0$ für alle $n \in \mathbb{N}$, und somit ist

$$
0=\mu\left(\bigcup_{n \in \mathbb{N}} A_{n}\right)=\mu(\{x \in \Omega:|f(x)|>\|\varphi\|\})
$$

Es folgt (12.6). Somit ist eine Linearform $\eta \in\left(L^{1}(\Omega)\right)^{\prime}$ wohldefiniert durch $\eta(g)= \int_{\Omega} f g d \mu$, und diese Linearform erfüllt

$$
\eta(\tilde{g})=\varphi_{k}(g)=\varphi(\tilde{g}) \quad \text { für alle } k \in \mathbb{N} \text { und } g \in L^{2}\left(\Omega_{k}\right)
$$

Da $L^{2}\left(\Omega_{k}\right) \subset L^{1}\left(\Omega_{k}\right)$ für $k \in \mathbb{N}$ dicht liegt (leichte Übung; man verwende $\mu\left(\Omega_{k}\right)< \infty)$, gilt (12.7) auch für $k \in \mathbb{N}$ und $g \in L^{1}\left(\Omega_{k}\right)$. Ist schließlich $g \in L^{1}(\Omega)$ und $g_{k}:=\left.g\right|_{\Omega_{k}} \in L^{1}\left(\Omega_{k}\right)$, so folgt $\tilde{g}_{k} \rightarrow g$ in $L^{1}(\Omega)$ für $k \rightarrow \infty$ und somit

$$
\varphi(g)=\lim _{k \rightarrow \infty} \varphi\left(\tilde{g}_{k}\right)=\lim _{k \rightarrow \infty} \eta\left(\tilde{g}_{k}\right)=\eta(g)
$$

Somit gilt $\varphi=\eta$, und man hat $\|\varphi\|=\|\eta\| \leq\|f\|_{\infty} \leq\|\varphi\|$, also $\|f\|_{\infty}=\|\varphi\|$.

\section*{§ 13 Fredholmoperatoren}
Seien stets $E, F$ normierte Räume über $\mathbb{K}=\mathbb{R}$ bzw. $\mathbb{K}=\mathbb{C}$. In diesem Kapitel wollen wir einige Aspekte der Klasse der Fredholmoperatoren studieren. Diese Operatorklasse verhält sich in Bezug auf das Lösen von Operatorgleichungen ähnlich gut wie Operatoren zwischen endlich-dimensionalen Unterräumen. Ferner hat die Klasse gute Stabilitätseigenschaften, und sie steht in einem engen Zusammenhang mit kompakten Operatoren. Wir benötigen zunächst ein paar vorbereitende Sätze. Wir erinnern daran, dass $\mathcal{K}(E)$ den Vektorraum der kompakten linearen Operatoren $E \rightarrow E$ bezeichne.

\subsection*{Hilfssatz}
Sei $W \subset E$ ein endlich dimensionaler Teilraum. Dann besitzt $W$ ein topologisches Komplement in $E$.

Beweis. Sei $n:=\operatorname{dim} W$ und $e_{1}, \ldots, e_{n}$ eine Basis von $W$. Sei ferner $\varphi_{1}, \ldots, \varphi_{n}$ die duale Basis von $W^{\prime}$ definiert durch

$$
\varphi_{i}\left(e_{j}\right)=\delta_{i j} \quad \text { für } i, j=1, \ldots, n .
$$

Nach 9.4 existieren Fortsetzungen $\psi_{1}, \ldots, \psi_{n} \in E^{\prime}$ von $\varphi_{1}, \ldots, \varphi_{n}$. Sei nun

$$
V:=\left\{x \in E: \psi_{i}(x)=0 \text { für } i=1, \ldots, n\right\} \text {. }
$$

Dann ist $V$ ein abgeschlossenes algebraisches Komplement zu $W$ in $E$, und mit 5.14(i) folgt $E=W \oplus_{\text {top }} V$.

\subsection*{Satz}
Sei $K \in \mathcal{K}(E)$ und $T=\operatorname{id}-K \in \mathcal{L}(E)$. Dann gilt:\\
i) $\operatorname{Kern} T$ ist endlichdimensional\\
ii) Bild $T$ ist abgeschlossen und hat endliche Kodimension.\\
iii) Ist $\operatorname{dim} E=\infty$, so ist $K \notin \operatorname{Iso}(E)$.

Beweis. i) Da $\left.K\right|_{\operatorname{Kern} T}=\left.\operatorname{id}\right|_{\operatorname{Kern} T} \in \mathcal{K}(\operatorname{Kern} T)$ ist, folgt $\operatorname{dimKern} T<\infty$ nach 4.7.\\
ii) Sei $W:=$ Kern $T$. Nach i) und 13.1 existiert ein topologisches Komplement $V$ in $E$ zu $W$, insbesondere ist $V \subset E$ abgeschlossen.\\
Sei nun $y \in \overline{\text { Bild } T}$; dann existiert eine Folge $\left(x_{n}\right)_{n}$ in $V$ mit $\lim _{n \rightarrow \infty} T x_{n}=y$.

\begin{enumerate}
  \item Fall: Die Folge $\left(x_{n}\right)_{n}$ ist beschränkt. Da $K$ ein kompakter Operator ist, können wir nach Übergang zu einer Teilfolge, wiederum $\left(x_{n}\right)_{n}$ genannt, annehmen, dass $z:=\lim _{n \rightarrow \infty} K x_{n} \in E$ existiert. Dann gilt aber
\end{enumerate}

$$
x_{n}=(T+K) x_{n} \rightarrow y+z \quad \text { für } n \rightarrow \infty,
$$

und aufgrund der Stetigkeit von $T$ folgt $T(y+z)=\lim _{n \rightarrow \infty} T x_{n}=y$.\\
2. Fall: Nach Übergang zu einer Teilfolge gilt $\left\|x_{n}\right\| \rightarrow \infty$ für $n \rightarrow \infty$. Setze dann $z_{n}:=\frac{x_{n}}{\left\|x_{n}\right\|} \in V$. Dann gilt $\left\|z_{n}\right\|=1$ für alle $n$ und $\lim _{n \rightarrow \infty} T z_{n}=0$. Nach erneutem Übergang zu einer Teilfolge existiert wiederum $z:=\lim _{n \rightarrow \infty} K z_{n}$, und somit folgt

$$
z_{n}=(T+K) z_{n} \rightarrow 0+z=z \quad \text { für } n \rightarrow \infty, \quad \text { also } T z=\lim _{n \rightarrow \infty} T z_{n}=0 .
$$

Aufgrund der Abgeschlossenheit von $V$ ist ferner $z \in V$ und somit $z=0$ aufgrund der Injektivität von $\left.T\right|_{V}$; dies widerspricht aber der Normierung $\left\|z_{n}\right\|=1$ für alle $n$. Also tritt dieser Fall nicht auf, und es folgt die Abgeschlossenheit von Bild $T$ in $E$. Da schließlich wegen 11.8 auch $K^{\prime} \in \mathcal{K}\left(E^{\prime}\right)$ gilt, ist wegen i) auch $\operatorname{Kern} T^{\prime}$ endlichdimensional. Ferner folgt mit 9.14.ii)

Bild $T=\overline{\operatorname{Bild} T}=(\operatorname{Bild} T)^{\perp \perp}=\left(\operatorname{Kern} T^{\prime}\right)^{\perp}=\left\{x \in E: \varphi(x)=0\right.$ für alle $\left.\varphi \in \operatorname{Kern} T^{\prime}\right\}$.\\
Dies zeigt codim Bild $T \leq m:=\operatorname{dimKern} T^{\prime}<\infty$, denn:\\
Ist $\varphi_{1}, \ldots, \varphi_{m}$ eine Basis von Kern $T^{\prime} \subset E^{\prime}$ und $X \subset F$ ein Unterraum mit $\operatorname{dim} X> m$, so existiert $x \in X \backslash\{0\}$ mit $\varphi_{j}(x)=0$ für $j=1, \ldots, m$, also $\varphi(x)=0$ für alle $\varphi \in \operatorname{Kern} T^{\prime}$ und somit $x \in \operatorname{Bild} T$. Die Behauptung folgt.\\
(iii) Wäre $K \in \operatorname{Iso}(E)$, und ist auch id $=K K^{-1} \in \mathcal{K}(E)$ gemäß 4.19. Mit 4.7 folgt $\operatorname{dim} E<\infty$, im Widerspruch zur Voraussetzung.

\subsection*{Definition und Satz}
Sei $V \subset E$ ein abgeschlossener Unterraum. Dann gilt:\\
(i) Auf dem Faktorraum $E / V:=\{x+V: x \in E\}$ wird eine Norm definiert durch

$$
\|x+V\|=\inf _{v \in V}\|x+v\|=\operatorname{dist}(x, V)
$$

(ii) Die Projektion $P: E \rightarrow E / V, x \mapsto x+V$ ist stetig mit $\|P\| \leq 1$.\\
(iii) Ist $T \in \mathcal{L}(E, F)$ mit $V \subset \operatorname{Kern} T$, so existiert $\tilde{T} \in \mathcal{L}(E / V, F)$ mit $T=\tilde{T} P$ und $\|\tilde{T}\|=\|T\|$.\\
(iv) Ist $E$ ein Banachraum, so auch $E / V$.

Beweis. Sei im Folgenden $\bar{x}=x+V \in E / V$ für $x \in E$.\\
(i) Die Seminormeigenschaften sind einfach zu sehen. Zur Definitheit: Ist $x \in E$ mit $\|\bar{x}\|=\operatorname{dist}(x, V)=0$, so folgt $x \in \bar{V}=V$ und somit $\bar{x}=0_{E / V}$.\\
(ii) Für $x \in E$ ist $\|P x\|=\|\bar{x}\| \leq\|x+0\|=\|x\|$ nach Definition der Norm. Dies zeigt die Stetigkeit von $P$.\\
(iii) Der Homomorphiesatz der linearen Algebra liefert eine lineare Abbildung mit $T=\tilde{T} P$. Sei $\varepsilon>0$. Ist $x \in E$ mit $\|\bar{x}\| \leq 1$, so existiert $y=x+v \in E$ mit $v \in V$ und $\|y\| \leq\|\bar{x}\|+\varepsilon \leq 1+\varepsilon$ und somit

$$
\|\tilde{T} \bar{x}\|=\|\tilde{T} \bar{y}\|=\|\tilde{T} P y\|=\|T y\| \leq\|T\|\|y\| \leq\|T\|(1+\varepsilon)
$$

Es folgt $\|\tilde{T}\| \leq\|T\|(1+\varepsilon)$. Da $\varepsilon>0$ beliebig gewählt war, folgt $\|\tilde{T}\| \leq\|T\|$. Ferner folgt auch

$$
\|T\|=\|\tilde{T} P\| \leq\|\tilde{T}\|\|P\| \leq\|\tilde{T}\|
$$

Insgesamt folgt $\|T\|=\|\tilde{T}\|$.\\
(iv) Wir verwenden 3.16. Seien dazu $x_{k} \in E, k \in \mathbb{N}$ mit $\sum_{k=1}^{\infty}\left\|\overline{x_{k}}\right\|<\infty$, und seien $y_{k} \in E, k \in \mathbb{N}$ mit $\overline{y_{k}}=\overline{x_{k}}$ und $\left\|y_{k}\right\| \leq\left\|\overline{x_{k}}\right\|+2^{-k}$ gewählt. Dann folgt auch $\sum_{k=1}^{\infty}\left\|y_{k}\right\|<\infty$. Da $E$ ein Banachraum ist, existiert somit $y:=\lim _{n \rightarrow \infty} \sum_{k=1}^{n} y_{k}$ in $E$ gemäß 3.16. Also existiert auch

$$
\lim _{n \rightarrow \infty} \sum_{k=1}^{n} \overline{x_{k}}=\lim _{n \rightarrow \infty} \sum_{k=1}^{n} \overline{y_{k}}=P\left(\lim _{n \rightarrow \infty} \sum_{k=1}^{n} y_{k}\right)=P y \quad \text { in } E / V
$$

\subsection*{Korollar}
Ist $T \in \mathcal{L}(E, F)$, so existiert ein weiterer normierter Raum $\tilde{E}$ und eine injektive Abbildung $\tilde{T} \in \mathcal{L}(\tilde{E}, F)$ mit Bild $\tilde{T}=$ Bild $T$.\\
Dabei kann $\tilde{E}$ als Banachraum gewählt werden, wenn $E$ ein Banachraum ist.\\
Beweis. Wähle $V=\operatorname{Kern} T, \tilde{E}=E / V$ und $\tilde{T}$ wie in 13.3 .

Im Folgenden seien nun $E, F$ Banachräume über $\mathbb{R}$ oder $\mathbb{C}$.

\subsection*{Satz}
Ist $T \in \mathcal{L}(E, F)$ mit codim Bild $T<\infty$, so ist Bild $T$ abgeschlossen.\\
Beweis. Gemäß 13.4 können wir ohne Einschränkung annehmen, dass $T$ injektiv ist. Sei $n:=\operatorname{codim}$ Bild $T$. Dann existieren $e_{1}, \ldots, e_{n} \in F$ derart, dass $F=\operatorname{Bild} T \oplus \operatorname{span}\left\{e_{1}, \ldots, e_{n}\right\}$ ist. Betrachte den Banachraum $E \times \mathbb{K}^{n}$, versehen mit der Norm $\left(x, \xi_{1}, \ldots, \xi_{n}\right) \mapsto\left\|\left(x, \xi_{1}, \ldots, \xi_{n}\right)\right\|=\|x\|_{E}+\sum_{i=1}^{n}\left|\xi_{i}\right|$. Sei $S: E \times \mathbb{K}^{n} \rightarrow F$ definiert durch $S\left(x, \xi_{1}, \ldots, \xi_{n}\right)=T x+\sum_{i=1}^{n} \xi_{i} e_{i}$. Dann ist $S \in \mathcal{L}\left(E \times \mathbb{K}^{n}, F\right)$ bijektiv und stetig, also $S \in \operatorname{Iso}\left(E \times \mathbb{K}^{n}, F\right)$ gemäß 5.7. Insbesondere ist Bild $T$ als Bild der abgeschlossenen Teilmenge $E \times\{0\} \subset E \times \mathbb{K}^{n}$ unter $S$ abgeschlossen in $F$.

\subsection*{Definition}
Eine Abbildung $T \in \mathcal{L}(E, F)$ heißt Fredholmoperator, falls $\operatorname{dimKern} T<\infty$ und codim Bild $T<\infty$ ist. In diesem Fall definieren wir den Fredholmindex durch $\operatorname{ind} T=\operatorname{dimKern} T-\operatorname{codim}$ Bild $T$. Wir schreiben:

$$
\begin{aligned}
\Phi(E, F) & =\{T \in \mathcal{L}(E, F): T \text { Fredholmoperator }\}, \\
\Phi_{k}(E, F) & =\{T \in \Phi(E, F): \text { ind } T=k\}, \\
\Phi(E) & :=\Phi(E, E), \Phi_{k}(E):=\Phi_{k}(E, E) .
\end{aligned}
$$

\subsection*{Beispiel}
i). Sind $E, F$ endlichdimensionale Banachräume mit $\operatorname{dim} E=m$ und $\operatorname{dim} F= n$, so ist jeder Operator $T \in \operatorname{Hom}(E, F)$ ein Fredholmoperator, und es gilt $\operatorname{ind} T=m-n$.\\
ii). Sei $K \in \mathcal{L}(E)$ ein kompakter Operator. Dann ist id $-K \in \Phi(E)$ gemäß 13.2 .\\
iii). Sei $p \in[1, \infty], E:=\ell^{p}(\mathbb{N}, \mathbb{K})$ und $T_{n} \in \mathcal{L}(E)$ für $n \in \mathbb{Z}$ wie folgt definiert: Im Fall $n \geq 0$ sei

$$
T_{n}(x)=\left(x_{n+1}, x_{n+2}, \ldots\right) \quad \text { für } x=\left(x_{j}\right)_{j} \in E, \quad \text { (Linksshift), }
$$

und $\operatorname{im}$ Fall $n<0$ sei

$$
T_{n}(x)=(\underbrace{0, \ldots, 0}_{-\mathrm{n}-\text { mal }}, x_{1}, x_{2}, \ldots) \quad \text { für } x=\left(x_{j}\right)_{j} \in E \quad \text { (Rechtsshift). }
$$

Dann ist $T_{n} \in \Phi_{n}(E)$ für alle $n \in \mathbb{Z}$.\\
iv). Sei $I \subset \mathbb{R}$ ein kompaktes (nichtdegeneriertes) Intervall und $T: C^{k}(I) \rightarrow C(I)$ gegeben durch $T f=f^{(k)}$. Dann ist $T$ ein Fredholmoperator vom Index $k$ (Übungsaufgabe, Blatt 13).

\subsection*{Bemerkung (Fredholm-Alternative)}
Ein Fredholmoperator $T \in \Phi(E, F)$ vom Index 0 erfüllt die sogenannte FredholmAlternative:\\
Entweder die Gleichung $T u=f$ hat für jedes $f \in F$ eine eindeutige Lösung, oder die Gleichung $T u=0$ hat eine nichttriviale Lösung $u \in E \backslash\{0\}$.

\subsection*{Lemma}
Sei $T \in \Phi(E, F)$. Dann existieren topologische Zerlegungen $E=\operatorname{Kern} T \oplus_{\text {top }} W$, $F=$ Bild $T \oplus_{\text {top }} Z$, und $\left.T\right|_{W}: W \rightarrow$ Bild $T$ ist ein topologischer Isomorphismus.

Beweis. Da $\operatorname{dimKern} T<\infty$, hat $\operatorname{Kern} T$ gemäß 13.1 ein topologisches Komplement $W$ in $E$, d.h. es gilt $E=\operatorname{Kern} T \oplus_{\text {top }} W$. Sei nun $Z$ ein algebraisches Komplement von Bild $T$ in $F$. Da nach 13.5 Bild $T \subset F$ abgeschlossen ist, gilt sogar $F=$ Bild $T \oplus_{\text {top }} Z$ gemäss 5.14. Nun ist $\left.T\right|_{W}: W \rightarrow$ Bild $T$ eine stetige bijektive Abbildung zwischen Banachräumen, da $W \subset E$ und Bild $T \subset F$ abgeschlossen sind. Also ist $\left.T\right|_{W}: W \rightarrow$ Bild $T$ ein topologischer Isomorphismus gemäß 5.7.

\subsection*{Satz (Indexformel)}
Seien $E, F, Z$ Banachräume, und sei $T \in \Phi(E, F), M \in \Phi(F, Z)$. Dann ist $M T \in \Phi(E, Z)$, und es gilt

$$
\operatorname{ind} M T=\operatorname{ind} M+\operatorname{ind} T
$$

Beweis. Der folgende Beweis basiert ausschließlich auf linearer Algebra. Wir benötigen zunächst eine spezielle Zerlegung von $F$. Dazu wählen wir zunächst einen Teilraum $\tilde{F} \subset F$ mit $F=($ Bild $T+\operatorname{Kern} M) \oplus \tilde{F}$. Ferner betrachten wir den nach Voraussetzung endlich dimensionalen Teilraum

$$
V:=\text { Bild } T \cap \text { Kern } M \subset F
$$

und wählen Teilräume $U \subset$ Bild $T, W \subset \operatorname{Kern} M$ mit

$$
\text { Bild } T=U \oplus V \quad \text { und } \quad \text { Kern } M=V \oplus W .
$$

Nach Konstruktion ist dann

$$
\text { Bild } T+\operatorname{Kern} M=U \oplus V \oplus W
$$

und

$$
F=U \oplus V \oplus W \oplus \tilde{F}
$$

Es folgt

$$
\text { Kern } M T=\{x \in E: T x \in \operatorname{Kern} M\}=T^{-1}(V) .
$$

Zerlegt man diesen Unterraum von $E$ nun in der Form

$$
\operatorname{Kern} M T=\operatorname{Kern} T \oplus \tilde{V},
$$

so ist $\left.T\right|_{\tilde{V}}: \tilde{V} \rightarrow V$ ein Isomorphismus und somit

$$
\text { (*) } \quad \operatorname{dimKern} M T=\operatorname{dimKern} T+\operatorname{dim} V=\operatorname{dimKern} T+\operatorname{dimKern} M-\operatorname{dim} W .
$$

Da $\left.M\right|_{V \oplus W}=0$ und $\left.M\right|_{U \oplus \tilde{F}}$ injektiv ist, ist Bild $M=M(U) \oplus M(\tilde{F})$. Schreibe nun $Z=M(U) \oplus M(\tilde{F}) \oplus \tilde{Z}$. Dann ist codim Bild $M=\operatorname{dim} \tilde{Z}$ und

$$
\operatorname{dim} M(\tilde{F})=\operatorname{dim} \tilde{F}=\operatorname{codim} \text { Bild } T-\operatorname{dim} W .
$$

Ferner ist Bild $M T=M($ Bild $T)=M(U)$, und somit gilt

$$
\begin{aligned}
\operatorname{codim} \operatorname{Bild} M T=\operatorname{codim} M(U) & =\operatorname{dim} M(\tilde{F})+\operatorname{dim} \tilde{Z} \\
(* *) & =\operatorname{codim} \operatorname{Bild} T-\operatorname{dim} W+\operatorname{codim} \operatorname{Bild} M .
\end{aligned}
$$

Subtraktion $(*)-(* *)$ liefert nun die Indexformel. \(\square\)

\subsection*{Satz}
Für $T \in \mathcal{L}(E, F)$ sind äquivalent:\\
(i) $T \in \Phi(E, F)$,\\
(ii) Es existieren $S_{1}, S_{2} \in \mathcal{L}(F, E)$ so, dass $\operatorname{id}_{E}-S_{1} T \in \mathcal{K}(E)$ und $\operatorname{id}_{F}-T S_{2} \in \mathcal{K}(F)$ ist.

Beweis. $(i) \Longrightarrow(i i)$ : Nach 13.9 gibt es topologische Zerlegungen $E=\operatorname{Kern} T \oplus_{\text {top }} W$, $F=$ Bild $T \oplus_{\text {top }} Z$ derart, dass $\left.T\right|_{W}: W \rightarrow$ Bild $T$ ein topologischer Isomorphismus ist. Sei $S_{0} \in \mathcal{L}$ (Bild $T, W$ ) die Inverse, und sei

$$
P \in \mathcal{L}(E) \text { die Projektion auf } W \text { längs } \operatorname{Kern} T
$$

sowie

$$
Q \in \mathcal{L}(F) \text { die Projektion auf Bild } T \text { längs } Z \text {. }
$$

Setze nun $S:=S_{0} Q \in \mathcal{L}(F, E)$. Dann gilt für $x \in E$

$$
S T x=S_{0} Q T x=S_{0} T x=P x
$$

und für $y \in F$

$$
T S y=T S_{0} Q y=Q y
$$

Also ist $\operatorname{id}_{E}-S T=\operatorname{id}_{E}-P \in \mathcal{L}(E)$ eine Projektion mit

$$
\operatorname{dim} \operatorname{Bild}\left(\operatorname{id}_{E}-S T\right)=\operatorname{dimKern} P=\operatorname{dimKern} T<\infty,
$$

und $\operatorname{id}_{F}-T S=\operatorname{id}_{F}-Q$ ist eine Projektion mit

$$
\operatorname{dim} \operatorname{Bild}\left(\operatorname{id}_{F}-T S\right)=\operatorname{dimKern} Q=\operatorname{dim} Z=\operatorname{codim} \text { Bild } T<\infty .
$$

Insbesondere ist $\operatorname{id}_{E}-S T \in \mathcal{K}(E)$ und $\operatorname{id}_{F}-T S \in \mathcal{K}(F)$.\\
$(i i) \Longrightarrow(i)$ : Nach 13.7 (ii) ist $T S_{2}=\operatorname{id}_{F}-\left(\operatorname{id}_{F}-T S_{2}\right) \in \Phi(F)$, also ist codim Bild $T \leq$ codim Bild $T S_{2}<\infty$ (da Bild $T \supset$ Bild $T S_{2}$ ). Zudem ist $S_{1} T= \operatorname{id}_{E}-\left(\operatorname{id}_{E}-S_{1} T\right) \in \Phi(E)$, also ist $\operatorname{dimKern} T \leq \operatorname{dimKern} S_{1} T<\infty$. Somit ist $T \in \Phi(E, F)$.

\subsection*{Hilfssatz}
i). Ist $T \in \mathcal{L}(E)$ und konvergiert die sogenannte Neumannsche Reihe $\sum_{k=0}^{\infty} T^{k}$ gegen ein $R \in \mathcal{L}(E)$, so ist $\mathrm{id}-T \in \operatorname{Iso}(E)$ mit $(\mathrm{id}-T)^{-1}=R$. Dies gilt insbesondere, wenn $\|T\|<1$ ist.\\
ii). Ist $T \in \operatorname{Iso}(E)$ und $S \in \mathcal{L}(E)$ mit $\|S\|<\frac{1}{\left\|T^{-1}\right\|}$, so ist $T-S \in \operatorname{Iso}(E)$ und $(T-S)^{-1}=\sum_{n=0}^{\infty}\left(T^{-1} S\right)^{n} T^{-1}$.\\
iii). Die Teilmenge $\operatorname{Iso}(E)$ ist offen in $\mathcal{L}(E)$.

Beweis. Sei $R_{n}:=\sum_{k=0}^{n} T^{k}$. Nach Voraussetzung ist $R:=\lim _{n \rightarrow \infty} R_{n}$, und dies erzwingt $\lim _{n \rightarrow \infty} T^{n}=\lim _{n \rightarrow \infty}\left(R_{n}-R_{n-1}\right)=0$. Es folgt

$$
(\mathrm{id}-T) R=\lim _{n \rightarrow \infty}(\mathrm{id}-T) R_{n}=\lim _{n \rightarrow \infty}\left(\mathrm{id}-T^{n+1}\right)=\mathrm{id}
$$

Genauso sieht man $T(\mathrm{id}-T)=R$.\\
(ii) Es ist $T-S=T$ (id $-T^{-1} S$ ), wobei nach Voraussetzung $\left\|T^{-1} S\right\| \leq\left\|T^{-1}\right\|\|S\|<$ 1 gilt. Mit i) folgt id $-T^{-1} S \in \operatorname{Iso}(E)$ und $\left(\mathrm{id}-T^{-1} S\right)^{-1}=\sum_{n=0}^{\infty}\left(T^{-1} S\right)^{n}$. Somit ist $T-S \in \operatorname{Iso}(E)$ und $(T-S)^{-1}=\left(\operatorname{id}-T^{-1} S\right)^{-1} T^{-1}=\sum_{n=0}^{\infty}\left(T^{-1} S\right)^{n} T^{-1}$.\\
iii) Da $\operatorname{dist}(T, \mathcal{L}(E) \backslash \operatorname{Iso}(E))>0$ für alle $T \in \operatorname{Iso}(E)$ gemäß i) gilt, ist $\operatorname{Iso}(E)$ offen in $\mathcal{L}(E)$.

\subsection*{Satz}
Für alle $n \in \mathbb{Z}$ ist $\Phi_{n}(E, F)$ eine offene Teilmenge von $\mathcal{L}(E, F)$. Mit anderen Worten: $\Phi(E, F)$ ist offen in $\mathcal{L}(E, F)$, und ind : $\Phi(E, F) \rightarrow \mathbb{Z}$ ist stetig, also konstant auf den Zusammenhangskomponenten von $\Phi(E, F)$.

\subsection*{Korollar(Stabilität des Index unter kompakten Störungen)}
(i) Sei $T \in \Phi(E, F), K \in \mathcal{K}(E, F)$. Dann ist $T+K \in \Phi(E, F)$, und es gilt $\operatorname{ind}(T+ K)=\operatorname{ind} T$.\\
(ii) Sei $K \in \mathcal{K}(E)$. Dann ist $\operatorname{id}_{E}-K \in \Phi_{0}(E)$.

Beweis. (i) Wegen $T \in \Phi(E, F)$ existieren nach 13.11 Operatoren $S_{1}, S_{2} \in \mathcal{L}(F, E)$ so, dass $\operatorname{id}_{E}-S_{1} T \in \mathcal{K}(E), \operatorname{id}_{F}-T S_{2} \in \mathcal{K}(F)$ gilt. Da ferner $S_{1} K \in \mathcal{K}(E)$ und $K S_{2} \in \mathcal{K}(F)$ ist, folgt $\operatorname{id}_{E}-S_{1}(T+K)=\operatorname{id}_{E}-S_{1} T-S_{1} K \in \mathcal{K}(E)$ und $\operatorname{id}_{F}-(T+ K) S_{2}=\operatorname{id}_{F}-T S_{2}-K S_{2} \in \mathcal{K}(F)$. Aus 13.11 folgt wiederum $T+K \in \Phi(E, F)$. Für $0 \leq t \leq 1$ ist auch $t K \in \mathcal{K}(E, F)$, also folgt $T+t K \in \Phi(E, F)$. Die Abbildung $[0,1] \rightarrow \mathbb{Z}, t \mapsto \operatorname{ind}(T+t K)$ ist stetig nach 13.13, also konstant. Daher folgt $\operatorname{ind}(T)= \operatorname{ind}(T+K)$.\\
(ii) Folgt direkt aus (i).

Beweis von 13.13.\\
Vormerkung: Wir bemerken zunächst, dass $\operatorname{Iso}(E, F) \subset \mathcal{L}(E, F)$ offen ist. Dies ist trivial im Fall $\operatorname{Iso}(E, F)=\varnothing$. Existiert andernfalls $L \in \operatorname{Iso}(E, F)$, so ist $\operatorname{Iso}(E, F)$ das Urbild von $\operatorname{Iso}(E)$ unter der stetigen Abbildung

$$
\mathcal{L}(E, F) \rightarrow \mathcal{L}(E), \quad T \mapsto L^{-1} T
$$

Dabei ist $\operatorname{Iso}(E) \subset \mathcal{L}(E)$ offen nach 13.12, und somit ist $\operatorname{Iso}(E, F)$ in $\mathcal{L}(E, F)$ offen. Nun zur Behauptung: Sei $T_{0} \in \Phi_{n}(E, F)$, und seien $W \subset E, Z \subset F$ gemäß 13.9 so gewählt, dass

$$
E=\operatorname{Kern} T_{0} \oplus_{\mathrm{top}} W, \quad F=\text { Bild } T_{0} \oplus_{\mathrm{top}} Z
$$

und $\left.T_{0}\right|_{W}: W \rightarrow$ Bild $T_{0}$ ein topologischer Isomorphismus ist. Sei $I \in \mathcal{L}(W, E)$ die Inklusion und $Q \in \mathcal{L}\left(F\right.$, Bild $\left.T_{0}\right)$ die 'Projektion' auf Bild $T_{0}$ mit Kern $Q=Z$; dann ist $Q T_{0} I=\left.T_{0}\right|_{W} \in \operatorname{Iso}\left(W\right.$, Bild $\left.T_{0}\right)$.\\
Wir betrachten nun die Abbildung

$$
\mu: \mathcal{L}(E, F) \rightarrow \mathcal{L}\left(W, \text { Bild } T_{0}\right), \quad \mu(T)=Q T I=\left.Q T\right|_{W}
$$

Man sieht leicht, dass $\mu$ stetig ist. Da Iso $\left(W\right.$, Bild $\left.T_{0}\right) \subset \mathcal{L}\left(W\right.$, Bild $\left.T_{0}\right)$ gemäß der Vorbemerkung offen ist, ist auch $\mathcal{M}:=\mu^{-1}\left(\operatorname{Iso}\left(W, \operatorname{Bild} T_{0}\right)\right)$ offen in $\mathcal{L}(E, F)$, wobei $T_{0} \in \mathcal{M}$ gilt. Es reicht nun, zu zeigen:

$$
\mathcal{M} \subset \Phi_{n}(E, F) .
$$

Für $T \in \mathcal{M}$ ist aber $W \cap \operatorname{Kern} T=\{0\}$, also $\operatorname{dimKern} T<\infty$. Zudem ist

$$
\text { (*) } \quad \text { Bild } T+Z=F .
$$

In der Tat: Ist $v=b+z \in F$ mit $b \in$ Bild $T_{0}, z \in Z$, so existiert wegen $Q T I \in$ Iso( $W$, Bild $T_{0}$ ) ein $w \in W$ mit $Q T w=b$, also $v=Q T w+z=T w+z-\left(\operatorname{id}_{F}-Q\right)(T w)$ mit $T w \in$ Bild $T$ und $z-\left(\operatorname{id}_{F}-Q\right)(T w) \in Z=\operatorname{Bild}\left(\operatorname{id}_{F}-Q\right)$. Es folgt $(*)$. Aus (*) folgt nun codim Bild $T \leq \operatorname{dim} Z=\operatorname{codim}$ Bild $T_{0}<\infty$, und insgesamt folgt $T \in \Phi(E, F)$. Ist ferner $\ell=\operatorname{dimKern} T_{0}$ und $k=\operatorname{codimBild} T_{0}=\operatorname{dim} Z$, so folgt $n=\operatorname{ind} T_{0}=\ell-k$ sowie

$$
I \in \Phi_{-\ell}(W, E) \quad \text { und } \quad Q \in \Phi_{k}\left(F, \text { Bild } T_{0}\right) .
$$

Ferner liefert die Indexformel 13.10

$$
0=\operatorname{ind}(Q T I)=\operatorname{ind} Q+\operatorname{ind} T+\operatorname{ind} I=-\ell+\operatorname{ind} T+k,
$$

d.h. ind $T=\ell-k=n$. Es folgt $M \subset \Phi_{n}(E, F)$, und damit folgt die Behauptung. \(\square\)

\subsection*{Beispiel}
Sei $\mathcal{C}_{2 \pi}$ bzw. $\mathcal{C}_{2 \pi}^{k}$ der Banachraum der stetigen bzw. $k$-mal stetig differenzierbaren $2 \pi$-periodischen Funktionen $\mathbb{R} \rightarrow \mathbb{R}$. Sei ferner $a \in \mathcal{C}_{2 \pi}$. Dann gilt:\\
Genau dann hat die Differentialgleichung

$$
u^{(k)}+a(x) u=f
$$

für jedes $f \in \mathcal{C}_{2 \pi}$ eine eindeutige Lösung $u \in \mathcal{C}_{2 \pi}^{k}$, wenn die homogene Gleichung

$$
u^{(k)}+a(x) u=0
$$

nur die triviale Lösung $u=0$ in $\mathcal{C}_{2 \pi}^{k}$ hat.\\
Dies sieht man so: Betrachte die Operatoren

$$
T_{0}, K \in \mathcal{L}\left(\mathcal{C}_{2 \pi}^{k}, \mathcal{C}_{2 \pi}\right), \quad\left[T_{0} u\right](x)=u^{(k)}(x), \quad[K u](x)=a(x) u(x) .
$$

Gemäß Übungsblatt 13 ist $T_{0}$ ein Fredholmoperator vom Index 0 ; ferner ist $K$ ein kompakter Operator, da die Einbettung $\mathcal{C}_{2 \pi}^{k} \hookrightarrow \mathcal{C}_{2 \pi}$ als Folge des Satzes von ArzelaAscoli kompakt und die lineare Abbildung

$$
\mathcal{C}_{2 \pi} \rightarrow \mathcal{C}_{2 \pi}, \quad u \mapsto a u
$$

stetig ist. Gemäß 13.14(i) ist also auch $T=T_{0}+K$ ein Fredholmoperator vom Index 0, und die obige Behauptung folgt.\\
Sei nun speziell $k=2 m$ gerade und $(-1)^{m} a$ eine positive Funktion; dann hat ( $* *$ ) nur die triviale Lösung in $\mathcal{C}_{2 \pi}^{k}$, denn für jede Lösung $u \in \mathcal{C}_{2 \pi}^{k}$ gilt mit partieller Integration

$$
(-1)^{m+1} \int_{0}^{2 \pi} a(x) u^{2}(x) d x=(-1)^{m} \int_{0}^{2 \pi} u(x) u^{(k)}(x) d x=\int_{0}^{2 \pi}\left[u^{(m)}(x)\right]^{2} d x \geq 0,
$$

und somit $u^{2} \equiv 0$ in $[0,2 \pi]$.\\
Folglich hat in diesem Fall (*) für jedes $f \in \mathcal{C}_{2 \pi}$ genau eine Lösung $u \in \mathcal{C}_{2 \pi}^{k}$.

\section*{§ 14 Die Fouriertransformation}
Im Folgenden sei stets $\mathbb{K}=\mathbb{C}$, und für $x, \xi \in \mathbb{R}^{N}$ sei $x \xi=x \cdot \xi=\sum_{k=1}^{N} x_{k} \xi_{k}$.

\subsection*{Definition}
$\mathcal{C}_{0}\left(\mathbb{R}^{N}\right)$ sei der Vektorraum der Funktionen $u \in C\left(\mathbb{R}^{N}\right)$ mit $\lim _{|x| \rightarrow \infty} u(x)=0$. Dieser Raum ist ein Banachraum bzgl. $\|\cdot\|_{\infty}$, da es ein abgeschlossener Unterraum des Banachraums $\mathcal{C}_{b}\left(\mathbb{R}^{N}\right)$ der beschränkten, stetigen Funktionen auf $\mathbb{R}^{N}$ ist.

\subsection*{Definition und Satz}
Für Funktionen $u, v \in L^{1}\left(\mathbb{R}^{N}\right)$ ist die Faltung $u * v \in L^{1}\left(\mathbb{R}^{N}\right)$ f.ü. definiert durch

$$
[u * v](x)=\int_{\mathbb{R}^{N}} u(x-y) v(y) d y=\int_{\mathbb{R}^{N}} u(y) v(x-y) d y
$$

Dabei gilt gemäß dem Satz von Fubini:

$$
\|[u * v]\|_{1} \leq \int_{\mathbb{R}^{N}} \int_{\mathbb{R}^{N}}|u(x-y)|\left\|v(y)\left|d y d x=\int_{\mathbb{R}^{N}}\right| v(y)\left|\int_{\mathbb{R}^{N}}\right| u(x-y) \mid d x d y=\right\| v\left\|_{1}\right\| u \|_{1}
$$

\subsection*{Beispiel}
Sei $u:=1_{[-1,1]}$ die Indikatorfunktion des Intervalls $[-1,1] \subset \mathbb{R}$. Dann ist

$$
\begin{aligned}
u * u(x) & =\int_{\mathbb{R}} u(x-y) u(y) d y=\int_{-1}^{1} 1_{[-1,1]}(x-y) d y=\int_{-1}^{1} 1_{[x-1, x+1]} d y \\
& =|[\max \{-1, x-1\}, \min \{1, x+1\}]|=\max \{0,2-|x|\}
\end{aligned}
$$

\subsection*{Definition}
Für $u \in L^{1}\left(\mathbb{R}^{N}\right)$ ist die Fouriertransformierte $\hat{u}: \mathbb{R}^{N} \rightarrow \mathbb{C}$ von $u$ definiert durch

$$
\hat{u}(\xi)=(2 \pi)^{-N / 2} \int_{\mathbb{R}^{N}} u(x) e^{-i x \xi} d x, \quad \xi \in \mathbb{R}^{N}
$$

Die Abbildung $u \mapsto \hat{u}$ heißt Fouriertransformation. Andere Bezeichnung: $\mathcal{F}(u)$ anstelle von $\hat{u}$.

\subsection*{Beispiel}
(i) Sei wieder $u:=1_{[-1,1]} \in L^{1}(\mathbb{R})$ die Indikatorfunktion des Intervalls $[-1,1] \subset \mathbb{R}$. Dann ist

$$
\sqrt{2 \pi} \hat{u}(\xi)=\int_{-1}^{1} e^{-i x \xi} d x=\int_{0}^{1}\left(e^{-i x \xi}+e^{i x \xi}\right) d x=2 \int_{0}^{1} \cos (x \xi) d x=2 \frac{\sin \xi}{\xi}
$$

und somit $\hat{u}(\xi)=\sqrt{\frac{2}{\pi}} \frac{\sin \xi}{\xi}$ (für $\xi=0$ erhält man $\hat{u}(0)=\sqrt{\frac{2}{\pi}}$ ).\\
(ii) Sei $u \in L^{1}(\mathbb{R})$ gegeben durch $u(x)=e^{-|x|}$. Dann ist

$$
\begin{aligned}
\sqrt{2 \pi} \hat{u}(\xi) & =\int_{\mathbb{R}} e^{-|x|} e^{-i x \xi} d x=\int_{0}^{\infty} e^{-x}\left(e^{-i x \xi}+e^{i x \xi}\right) d x \\
& =\lim _{R \rightarrow \infty}\left[\frac{e^{-x(1+i \xi)}}{-(1+i \xi)}+\frac{e^{-x(1-i \xi)}}{-(1-i \xi)}\right]_{0}^{R}=\frac{1}{1+i \xi}+\frac{1}{1-i \xi}=\frac{2}{1+\xi^{2}}
\end{aligned}
$$

und somit $\hat{u}(\xi)=\sqrt{\frac{2}{\pi}} \frac{1}{1+\xi^{2}}$.

\subsection*{Satz}
Seien $u, v \in L^{1}\left(\mathbb{R}^{N}\right)$. Dann gilt:\\
i). $\hat{u} \in C_{0}\left(\mathbb{R}^{N}\right)$ und $\|\hat{u}\|_{\infty} \leq(2 \pi)^{-N / 2}\|u\|_{1}$. Somit ist $\mathcal{F} \in \mathcal{L}\left(L^{1}\left(\mathbb{R}^{N}\right), C_{0}\left(\mathbb{R}^{N}\right)\right)$.\\
ii). Ist $a \in \mathbb{R}^{N}$ und $u_{a} \in L^{1}\left(\mathbb{R}^{N}\right)$ definiert durch $u_{a}(y)=u(y-a)$, so gilt

$$
\widehat{u_{a}}(\xi)=e^{-i a \xi} \hat{u}(\xi) \quad \text { für } \xi \in \mathbb{R}^{N} .
$$

iii). Ist $u$ stetig differenzierbar und auch $\partial_{k} u:=\frac{\partial u}{\partial x_{k}} \in L^{1}\left(\mathbb{R}^{N}\right)$ für ein $k$, so gilt

$$
\widehat{\partial_{k} u}(\xi)=i \xi_{k} \hat{u}(\xi) \quad \text { für } \xi \in \mathbb{R}^{N} .
$$

Die Fouriertransformation führt also Differenziationen in Multiplikationen über.\\
iv). Ist auch die Funktion $x \mapsto h(x):=x_{k} u(x)$ integrierbar für ein $k$, so ist $\hat{u}$ nach $\xi_{k}$ stetig partiell differenzierbar mit

$$
\hat{h}(\xi)=i \partial_{k} \hat{u}(\xi)
$$

Die Fouriertransformation führt also Multiplikationen in Differenziationen über.\\
v). $\widehat{u * v}(\xi)=(2 \pi)^{N / 2} \hat{u}(\xi) \hat{v}(\xi)$ für $\xi \in \mathbb{R}^{N}$.\\
vi). Für $\lambda>0$ gilt

$$
\int_{\mathbb{R}^{N}} \hat{u}(\xi) v(\lambda \xi) d \xi=\int_{\mathbb{R}^{N}} u(\lambda \xi) \hat{v}(\xi) d \xi
$$

Beweis. ii) Es ist

$$
\widehat{u_{a}}(\xi)=(2 \pi)^{-N / 2} \int_{\mathbb{R}^{N}} u(y-a) e^{-i y \xi} d y=(2 \pi)^{-N / 2} \int_{\mathbb{R}^{N}} u(z) e^{-i(z+a) \xi} d z=e^{-i a \xi} \hat{u}(\xi)
$$

für $\xi \in \mathbb{R}^{N}$.\\
iii) Wir betrachten zunächst den Fall $N=1$. Da nach Voraussetzung $u$ und $u^{\prime}$ in $L^{1}(\mathbb{R})$ liegen, gilt $\lim _{|x| \rightarrow \infty} u(x)=0$ (leichte Übung). Es folgt dann

$$
\begin{aligned}
& \sqrt{2 \pi} \xi \hat{u}(\xi)=\xi \int_{\mathbb{R}} u(x) e^{-i x \xi} d x=-\frac{1}{i} \int_{\mathbb{R}} u(x) \partial_{x}\left(e^{-i x \xi}\right) d x=-\frac{1}{i} \lim _{s \rightarrow \infty} \int_{-s}^{s} u(x) \partial_{x}\left(e^{-i x \xi}\right) d x \\
& =\frac{1}{i} \lim _{s \rightarrow \infty}\left(-\left[u(x) e^{-i x \xi}\right]_{x=-s}^{x=s}+\int_{-s}^{s} u^{\prime}(x) e^{-i x \xi} d x\right)=\frac{1}{i} \int_{\mathbb{R}} u^{\prime}(x) e^{-i x \xi} d \xi=\frac{\sqrt{2 \pi}}{i} \widehat{u^{\prime}}(\xi)
\end{aligned}
$$

Sei nun $N>1$, und ohne Einschränkung sei $k=N$. Für $x^{\prime}=\left(x_{1}, \ldots, x_{N-1}\right) \in \mathbb{R}^{N-1}$ sei ferner $U_{x^{\prime}} \in C^{1}(\mathbb{R})$ definiert durch $U_{x^{\prime}}(t)=u\left(x^{\prime}, t\right)$. Da $u, \partial_{N} u \in L^{1}\left(\mathbb{R}^{N}\right)$, folgt aus dem Satz von Fubini, dass für fast alle $x^{\prime} \in \mathbb{R}^{N-1}$ sowohl $U_{x^{\prime}}$ als auch $U_{x^{\prime}}^{\prime}$ in $L^{1}(\mathbb{R})$ liegen. Sei nun $\xi \in \mathbb{R}^{N}$ und $\xi^{\prime}=\left(\xi_{1}, \ldots, \xi_{n-1}\right)$; dann folgt mit dem Satz von Fubini:

$$
\begin{aligned}
& (2 \pi)^{N / 2} \xi_{N} \hat{u}(\xi)=\int_{\mathbb{R}} e^{-i x^{\prime} \xi^{\prime}} \xi_{N} \int_{\mathbb{R}} U_{x^{\prime}}\left(x_{N}\right) e^{-i x_{N} \xi_{N}} d x_{N} d x^{\prime}=\sqrt{2 \pi} \int_{\mathbb{R}^{N-1}} e^{-i x^{\prime} \xi^{\prime}} \xi_{N} \widehat{U_{x^{\prime}}}\left(\xi_{N}\right) d x^{\prime} \\
& \stackrel{\text { Fall } \mathrm{N}=1}{=} \sqrt{2 \pi} \frac{1}{i} \int_{\mathbb{R}^{N-1}} e^{-i x^{\prime} \xi^{\prime}} \widehat{U_{x^{\prime}}^{\prime}}\left(\xi_{N}\right) d x^{\prime}=\frac{1}{i} \int_{\mathbb{R}^{N-1}} e^{-i x^{\prime} \xi^{\prime}} \int_{\mathbb{R}} U_{x^{\prime}}^{\prime}\left(x_{N}\right) e^{-i x_{N} \xi_{N}} d x_{N} d x^{\prime} \\
& =\frac{1}{i} \int_{\mathbb{R}^{N}} \partial_{N} u(x) e^{-i x \xi} d x=(2 \pi)^{N / 2} \frac{1}{i} \widehat{\partial_{N} u}(\xi)
\end{aligned}
$$

(iv) Nach Voraussetzung werden die Funktionen

$$
x \mapsto u(x) \frac{\partial}{\partial \xi_{k}} e^{-i x \xi}=-i x_{k} u(x) e^{-i x \xi}, \quad \xi \in \mathbb{R}^{N}
$$

betragsmäßig durch die $L^{1}$-Funktion $x \mapsto\left|x_{k}\right||u(x)|$ majorisiert. Daher darf man im Folgenden unter dem Integral differenzieren und erhält

$$
\partial_{k} \hat{u}(\xi)=(2 \pi)^{-N / 2} \int_{\mathbb{R}^{N}} u(x) \frac{\partial}{\partial \xi_{k}} e^{-i x \xi} d x=-i(2 \pi)^{-N / 2} \int_{\mathbb{R}^{N}} x_{k} u(x) e^{-i x \xi} d x=-i \hat{h}(\xi)
$$

v) Mit dem Satz von Fubini ist

$$
\begin{aligned}
\widehat{u * v}(\xi) & =(2 \pi)^{-N / 2} \int e^{-i x \xi} \int u(y) v(x-y) d y d x=(2 \pi)^{-N / 2} \int u(y) \int e^{-i(x+y) \xi} v(x) d x d y \\
& =(2 \pi)^{-N / 2} \int u(y) e^{-i y \xi} d y \int e^{-i x \xi} v(x) d x=(2 \pi)^{N / 2} \hat{u}(\xi) \hat{v}(\xi)
\end{aligned}
$$

für $\xi \in \mathbb{R}^{N}$.\\
vi) Es ist

$$
\hat{u}(\xi)=(2 \pi)^{-N / 2} \int_{\mathbb{R}^{N}} u(x) e^{-i x \xi} d x=(2 \pi)^{-N / 2} \lambda^{N} \int_{\mathbb{R}^{N}} u(\lambda x) e^{-i \lambda x \xi} d x \quad \text { für } \xi \in \mathbb{R}
$$

mit dem Satz von Fubini also

$$
\begin{aligned}
\int_{\mathbb{R}^{N}} \hat{u}(\xi) v(\lambda \xi) d \xi & =(2 \pi)^{-N / 2} \lambda^{N} \int_{\mathbb{R}^{N}} u(\lambda x) \int_{\mathbb{R}^{N}} e^{-i \lambda x \xi} v(\lambda \xi) d \xi d x \\
& =(2 \pi)^{-N / 2} \int_{\mathbb{R}^{N}} u(\lambda x) \int_{\mathbb{R}^{N}} e^{-i x \xi} v(\xi) d \xi d x=\int_{\mathbb{R}^{N}} u(\lambda x) \hat{v}(x) d x
\end{aligned}
$$

Dies liefert (vi).\\
(i) Aus der Definition und dem Konvergenzsatz von Lebesgue folgt direkt, dass $\hat{u} \in \mathcal{C}_{b}\left(\mathbb{R}^{N}\right)$ ist und $\|\hat{u}\|_{\infty} \leq(2 \pi)^{-N / 2}\|u\|_{1}$ gilt. Somit ist $\mathcal{F} \in \mathcal{L}\left(L^{1}\left(\mathbb{R}^{N}\right), \mathcal{C}_{b}\left(\mathbb{R}^{N}\right)\right)$. Ist $u \in C_{c}^{\infty}\left(\mathbb{R}^{N}\right)$, so ist für $k=1, \ldots, N$ auch $\partial_{k} u \in L^{1}\left(\mathbb{R}^{N}\right)$. Gemäß iii) liegt also auch die Funktion $\xi \mapsto \xi_{k} \hat{u}(\xi)$ in $\mathcal{C}_{b}\left(\mathbb{R}^{N}\right)$ für $k=1, \ldots, N$. Dies erzwingt $\lim _{|\xi| \rightarrow \infty} \hat{u}(\xi)=0$, und somit ist $\hat{u} \in \mathcal{C}_{0}\left(\mathbb{R}^{N}\right)$. Da num $C_{c}^{\infty}\left(\mathbb{R}^{N}\right)$ gemäß 6.28 in $L^{1}\left(\mathbb{R}^{N}\right)$ dicht liegt und $\mathcal{F}$ als Abbildung von $L^{1}\left(\mathbb{R}^{N}\right)$ nach $\mathcal{C}_{b}\left(\mathbb{R}^{N}\right)$ stetig ist, folgt num

$$
\mathcal{F}\left(L^{1}\left(\mathbb{R}^{N}\right)\right)=\mathcal{F}\left(\overline{C_{c}^{\infty}\left(\mathbb{R}^{N}\right)}\right) \subset \overline{\mathcal{F}\left(C_{c}^{\infty}\left(\mathbb{R}^{N}\right)\right)}=\mathcal{C}_{0}\left(\mathbb{R}^{N}\right)
$$

\subsection*{Beispiel}
Sei $u:=1_{[-1,1]}$ die Indikatorfunktion des Intervalls $[-1,1] \subset \mathbb{R}$ und $v=u * u$, d.h. $v(x)=\max \{0,2-|x|\}$. Gemäß 14.5 und 14.6 (vi) ist dann $\hat{v}(\xi)=2 \sqrt{\frac{2}{\pi}}\left(\frac{\sin \xi}{\xi}\right)^{2}$.

\subsection*{Satz}
Sei $\varphi \in L^{1}\left(\mathbb{R}^{N}\right)$ definiert durch $\varphi(x)=e^{-\frac{|x|^{2}}{2}}$, vgl. 6.5. Dann gilt $\hat{\varphi}=\varphi$.\\
Beweis. Wir bemerken zunächst, dass $\int_{\mathbb{R}} e^{-\frac{x^{2}}{2}} d x=\sqrt{2 \pi}$ gilt. Tatsächlich ist

$$
\left(\int_{\mathbb{R}} e^{-\frac{x^{2}}{2}} d x\right)^{2}=\int_{\mathbb{R}} \int_{\mathbb{R}} e^{-\frac{x^{2}}{2}} e^{-\frac{y^{2}}{2}} d x d y=2 \pi \int_{0}^{\infty} r e^{-r^{2} / 2} d r=2 \pi \int_{0}^{\infty} e^{-\tau} d \tau=2 \pi
$$

Nun zum Beweis des Satzes; sei dazu zunächst $N=1$. Dann ist $\varphi^{\prime}(x)=-x e^{-\frac{x^{2}}{2}}$ für $x \in \mathbb{R}$, also offensichtlich $\varphi^{\prime} \in L^{1}(\mathbb{R})$. Somit liefert 14.6 iv ):

$$
\begin{aligned}
\xi \hat{\varphi}(\xi) & =-i \widehat{\varphi}^{\prime}(\xi)=\frac{i}{\sqrt{2 \pi}} \int_{\mathbb{R}} x e^{-\frac{x^{2}}{2}-i x \xi} d x=-\frac{1}{\sqrt{2 \pi}} \int_{\mathbb{R}} \partial_{\xi}\left(e^{-\frac{x^{2}}{2}-i x \xi}\right) d x \\
& =-\frac{1}{\sqrt{2 \pi}} \partial_{\xi} \int_{\mathbb{R}} e^{-\frac{x^{2}}{2}-i x \xi} d x=-\frac{d}{d \xi} \hat{\varphi}(\xi)
\end{aligned}
$$

für alle $\xi \in \mathbb{R}$. Integration dieser Differentialgleichung für $\hat{\varphi}$ liefert $\hat{\varphi}(\xi)=c e^{-\frac{\xi^{2}}{2}}= c \varphi(\xi)$ mit einem $c \in \mathbb{R}$. Da $\varphi(0)=1$ und gemäß der Vorbemerkung

$$
\hat{\varphi}(0)=(2 \pi)^{-1 / 2} \int_{\mathbb{R}} \varphi(x) d x \stackrel{[.5[i]}{=} 1
$$

ist, folgt $c=1$ und somit die Behauptung.\\
Im Fall $N>1$ verwendet man wieder den Satz von Fubini:\\
$\hat{\varphi}(\xi)=(2 \pi)^{-N / 2} \int_{\mathbb{R}^{N}} e^{-\frac{|x|^{2}}{2}-i x \xi} d x=\prod_{k=1}^{N}(2 \pi)^{-1 / 2} \int_{\mathbb{R}} e^{-\frac{x_{k}^{2}}{2}-i x_{k} \xi_{k}} d x_{k} \stackrel{\text { Fall } \stackrel{\mathrm{N}=1}{=}}{\prod_{k=1}^{N}} e^{-\frac{\xi_{k}^{2}}{2}}=\varphi(\xi)$.

\subsection*{Satz (Fouriersche Umkehrformel, 1. Version)}
Sei $u \in L^{1}\left(\mathbb{R}^{N}\right)$.\\
(i) Ist $u \in C_{b}\left(\mathbb{R}^{N}\right)$ und $\hat{u} \in L^{1}\left(\mathbb{R}^{N}\right)$, so ist\\
(14.2) $u(x)=(2 \pi)^{-N / 2} \int_{\mathbb{R}^{N}} \hat{u}(\xi) e^{i \xi x} d \xi=\hat{\hat{u}}(-x) \quad$ für alle $x \in \mathbb{R}^{N}$.\\
(ii) Ist $u \in C_{c}^{\infty}\left(\mathbb{R}^{N}\right)$, so ist $\hat{u} \in L^{1}\left(\mathbb{R}^{N}\right) \cap C_{0}\left(\mathbb{R}^{N}\right)$, und (14.2) gilt.\\
(iii) Für jedes $v \in C_{c}^{\infty}\left(\mathbb{R}^{N}\right)$ existiert ein $u \in L^{1}\left(\mathbb{R}^{N}\right) \cap C_{0}\left(\mathbb{R}^{N}\right)$ mit $\hat{u}=v$.

Beweis. (i) Sei $x \in \mathbb{R}^{N}$ und $u_{x} \in L^{1}\left(\mathbb{R}^{N}\right)$ definiert durch $u_{x}(y)=u(y-x)$, und sei $\varphi$ wie in 14.8 gegeben. Dann liefern 14.6 und 14.8 ;

$$
\int_{\mathbb{R}^{N}} \hat{u}(\xi) e^{-i \xi x} \varphi(\lambda \xi) d \xi=\int_{\mathbb{R}^{N}} \widehat{u_{x}}(\xi) \varphi(\lambda \xi) d \xi=\int_{\mathbb{R}^{N}} u_{x}(\lambda \xi) \varphi(\xi) d \xi
$$

Für $\lambda \rightarrow 0$ erhält man mit dem Konvergenzsatz von Lebesgue :

$$
\int_{\mathbb{R}^{N}} \hat{u}(\xi) e^{-i \xi x} d \xi=\int_{\mathbb{R}^{N}} u_{x}(0) \varphi(\xi) d \xi=u_{x}(0)(2 \pi)^{N / 2}=(2 \pi)^{N / 2} u(-x)
$$

Die Behauptung folgt durch Betrachtung von $-x$ anstelle von $x$.\\
(ii) Bekanntermaßen ist $\hat{u} \in C_{0}\left(\mathbb{R}^{N}\right)$. Mit $u$ ist ferner auch $v:=(-\Delta+1)^{N} u \in \mathcal{C}_{c}^{\infty}\left(\mathbb{R}^{N}\right)$. Also liefert 14.6 (iii), dass die Funktion

$$
\mathbb{R}^{N} \rightarrow \mathbb{R}, \quad \xi \mapsto \hat{v}(\xi)=\left(1+|\xi|^{2}\right)^{N} \hat{u}(\xi)
$$

stetig und beschränkt ist. Mit $C:=\|\hat{v}\|_{\infty}$ folgt also $|\hat{u}(\xi)| \leq \frac{C}{\left(1+|\xi|^{2}\right)^{N}}$ für alle $\xi \in \mathbb{R}^{N}$, und somit ist $\hat{u} \in L^{1}\left(\mathbb{R}^{N}\right)$. Gemäß (i) gilt also (14.2).\\
(iii) Sei $w \in C_{c}^{\infty}\left(\mathbb{R}^{N}\right)$ definiert durch $w(x)=v(-x)$. Gemäß (ii) ist dann $u:= \hat{w} \in L^{1}\left(\mathbb{R}^{N}\right) \cap C_{0}\left(\mathbb{R}^{N}\right)$, und 14.2 liefert $v(x)=w(-x)=\hat{\hat{w}}(x)=\hat{u}(x)$ für alle $x \in \mathbb{R}^{N}$.

\subsection*{Beispiel}
Sei $u:=1_{[-1,1]}$ die Indikatorfunktion des Intervalls $[-1,1] \subset \mathbb{R}$ und $v=u * u$, d.h. $v(x)=\max \{0,2-|x|\}$. Gemäß 14.7 ist dann $\hat{v}(\xi)=2 \sqrt{\frac{2}{\pi}}\left(\frac{\sin \xi}{\xi}\right)^{2}$, d.h. $\hat{v} \in L^{1}(\mathbb{R})$. Mit 14.9 folgt also\\
$\max \{0,2-|x|\}=\frac{1}{\sqrt{2 \pi}} \int_{\mathbb{R}} \hat{v}(\xi) e^{i \xi x} d \xi=\frac{1}{\sqrt{2 \pi}} \int_{\mathbb{R}} \hat{v}(\xi) \cos (\xi x) d \xi=\frac{2}{\pi} \int_{\mathbb{R}}\left(\frac{\sin \xi}{\xi}\right)^{2} \cos (\xi x) d \xi$.\\
Speziell mit $x=0$ folgt

$$
\int_{0}^{\infty}\left(\frac{\sin \xi}{\xi}\right)^{2} d \xi=\frac{\pi}{2}
$$

\subsection*{Hauptsatz}
Für alle $u \in L^{1}\left(\mathbb{R}^{N}\right) \cap L^{2}\left(\mathbb{R}^{N}\right)$ ist $\hat{u} \in L^{2}\left(\mathbb{R}^{N}\right)$ und $\|\hat{u}\|_{2}=\|u\|_{2}$. Somit lässt sich die Fouriertransformation in eindeutiger Weise zu einem isometrischen Isomorphismus (bzw. einem unitären Operator, siehe 8.12)

$$
\mathcal{F}: L^{2}\left(\mathbb{R}^{N}\right) \rightarrow L^{2}\left(\mathbb{R}^{N}\right)
$$

fortsetzen.\\
Um dies zu beweisen, benötigen wir noch folgenden

\subsection*{Hilfsatz}
Sei $\Omega \subset \mathbb{R}^{N}$ offen, $1 \leq p<q<\infty$, und sei $u \in L^{p}(\Omega) \cap L^{q}(\Omega)$. Dann existiert eine Folge $\left(u_{n}\right)_{n}$ in $\mathcal{C}_{c}^{\infty}(\Omega)$ mit $\left\|u_{n}-u\right\|_{s} \rightarrow 0$ für $n \rightarrow \infty$ und $s \in[p, q]$.

Beweis. Sei $\Omega_{n}:=\Omega \cap U_{n}(0)$ und $v_{n}:=u \cdot 1_{\Omega_{n}}$ für $n \in \mathbb{N}$. Mit der Interpolationsungleichung 6.5(i) und dem Satz von Lebesgue folgt dann

$$
(*) \quad\left\|v_{n}-u\right\|_{s} \rightarrow 0 \quad \text { für } n \rightarrow \infty \text { und } s \in[p, q]
$$

Sei $C_{n}:=\min \left\{1,\left|\Omega_{n}\right|^{\frac{1}{q}-\frac{1}{p}}\right\}$ für $n \in \mathbb{N}$. Gemäß 6.28 existiert zu jedem $n \in \mathbb{N}$ ein $u_{n} \in \mathcal{C}_{c}^{\infty}(\Omega)$ mit $\left\|u_{n}-v_{n}\right\|_{q} \leq \frac{C_{n}}{n}$, gemäß der Interpolationsungleichung 6.5(i) also

$$
\left\|u_{n}-v_{n}\right\|_{s} \leq \frac{1}{n} \quad \text { für } n \in \mathbb{N} \text { und } s \in[p, q] \text {. }
$$

Die Kombination von $(*)$ und $(* *)$ ergibt, dass die Folge $\left(u_{n}\right)_{n}$ die gewünschte Eigenschaft hat. \(\square\)

Beweis von 14.11. Ist $u \in \mathcal{C}_{c}^{\infty}\left(\mathbb{R}^{N}\right)$, so liefert 14.9 (ii), dass $\hat{u} \in L^{1}\left(\mathbb{R}^{N}\right) \cap L^{2}\left(\mathbb{R}^{N}\right)$ und $\hat{\hat{u}}(x)=u(-x)$ für alle $x \in \mathbb{R}^{N}$ gilt. Mit 14.9 folgt ferner

$$
\hat{\hat{\bar{u}}}(x)=(2 \pi)^{-\frac{N}{2}} \int_{\mathbb{R}^{N}} \overline{\hat{u}}(\xi) e^{-i \xi x} d \xi=\overline{(2 \pi)^{-\frac{N}{2}}} \int_{\mathbb{R}^{N}} \hat{u}(\xi) e^{i \xi x} d \xi=\bar{u}(x) .
$$

Anwendung von 14.6 (ii) mit $\lambda=1$ liefert nun

$$
\|\hat{u}\|_{2}^{2}=\int_{\mathbb{R}^{N}} \hat{u}(\xi) \overline{\hat{u}(\xi)} d \xi=\int_{\mathbb{R}^{N}} u(\xi) \hat{\overline{\hat{u}}}(\xi) d \xi=\int_{\mathbb{R}^{N}} u(\xi) \bar{u}(\xi) d \xi=\|u\|_{2}^{2}
$$

Es folgt, dass die Abbildung $\mathcal{C}_{c}^{\infty}\left(\mathbb{R}^{N}\right) \rightarrow L^{2}\left(\mathbb{R}^{N}\right), u \mapsto \hat{u}$ eine Isometrie bzgl. $\|\cdot\|_{2}$ ist und sich somit nach 6.28 und 5.1 in eindeutiger Weise zu einer linearen Isometrie $\mathcal{F}: L^{2}\left(\mathbb{R}^{N}\right) \rightarrow L^{2}\left(\mathbb{R}^{N}\right)$ fortsetzen lässt. Für $u \in L^{1}\left(\mathbb{R}^{N}\right) \cap L^{2}\left(\mathbb{R}^{N}\right)$ gilt dabei $\mathcal{F}(u)= \hat{u}$, denn: Gemäß 14.12 existiert eine Folge $\left(u_{n}\right)_{n} \subset \mathcal{C}_{c}^{\infty}\left(\mathbb{R}^{N}\right)$ mit $u_{n} \rightarrow u$ in $L^{1}\left(\mathbb{R}^{N}\right)$ und $L^{2}\left(\mathbb{R}^{N}\right)$, insbesondere also

$$
\hat{u}_{n}=\mathcal{F}\left(u_{n}\right) \rightarrow \mathcal{F}(u) \quad \text { bzgl. }\|\cdot\|_{2} .
$$

Wegen 6.8 dürfen wir (nach Übergang zu einer Teilfolge) auch annehmen, dass $\hat{u}_{n}$ punktweise f.ü. gegen $\mathcal{F}(u)$ konvergiert. Allerdings gilt für alle $\xi \in \mathbb{R}^{N}$ auch

$$
\left|\hat{u}(\xi)-\hat{u}_{n}(\xi)\right| \leq(2 \pi)^{-\frac{N}{2}} \int_{\mathbb{R}^{N}}\left|\left(u(x)-u_{n}(x)\right) e^{-i x \xi}\right| d x=(2 \pi)^{-\frac{N}{2}}\left\|u-u_{n}\right\|_{1} \rightarrow 0
$$

für $n \rightarrow \infty$, also folgt $\mathcal{F}(u)=\hat{u}$.\\
Zu zeigen bleibt noch die Surjektivität von $\mathcal{F}$. Aus 14.9(iii) folgt zunächst die Inklusion $\mathcal{C}_{c}^{\infty}\left(\mathbb{R}^{N}\right) \subset$ Bild $\mathcal{F}$. Da Bild $\mathcal{F}$ als Bild von $L^{2}\left(\mathbb{R}^{N}\right)$ unter einer Isometrie vollständig und damit abgeschlossen in $\mathcal{L}^{2}\left(\mathbb{R}^{N}\right)$ ist, folgt dann $L^{2}\left(\mathbb{R}^{N}\right)=\overline{\mathcal{C}_{c}^{\infty}\left(\mathbb{R}^{N}\right)} \subset$ Bild $F$ und somit die Surjektivität von $\mathcal{F}$. \(\square\)

\subsection*{Bemerkung}
i). Die Fouriertransformierte einer Funktion $u \in L^{2}\left(\mathbb{R}^{N}\right)$ lässt sich als Limes

$$
\mathcal{F}(u)=\lim _{R \rightarrow \infty} \mathcal{F}_{R}(u) \quad \text { in } L^{2}\left(\mathbb{R}^{N}\right)
$$

darstellen mit

$$
\mathcal{F}_{R}(u)(\xi):=(2 \pi)^{-\frac{N}{2}} \int_{B_{R}(0)} e^{-i x \xi} u(x) d x, \quad \xi \in \mathbb{R}^{N}
$$

Dies folgt mit dem Hauptsatz 14.11 aus der Tatsache, dass $\mathcal{F}_{R}(u)=\mathcal{F}\left(u_{R}\right)$ mit $u_{R}:=u \cdot 1_{B_{R}(0)} \in L^{2}\left(\mathbb{R}^{N}\right) \cap L^{1}\left(\mathbb{R}^{N}\right)$ für $R>0$ gilt und die Funktionen $u_{R}$ für $R \rightarrow \infty$ in $L^{2}\left(\mathbb{R}^{N}\right)$ gegen $u$ konvergieren.\\
ii). Als Abbildung $\mathcal{F}: L^{1}\left(\mathbb{R}^{N}\right) \rightarrow C_{0}\left(\mathbb{R}^{N}\right)$ ist die Fouriertransformation injektiv, aber nicht surjektiv (ohne Beweis).

\subsection*{Definition und Satz}
Für $n \in \mathbb{N}_{0}$ seien die Funktionen $h_{n}, H_{n}: \mathbb{R} \rightarrow \mathbb{C}$ definiert durch

$$
h_{n}(x)=(-1)^{n} e^{x^{2}}\left(\frac{d}{d x}\right)^{n} e^{-x^{2}}, \quad H_{n}(x)=\left(\sqrt{\pi} 2^{n} n!\right)^{-1 / 2} h(x) e^{-\frac{x^{2}}{2}} .
$$

Die Funktionen $h_{n}$ heißen Hermite-Polynome, und die Funktionen $H_{n}$ heißen Hermite-Funktionen. Man kann zeigen (Übungsaufgabe, Blatt 13):\\
(i) Die Funktionen $H_{n} n \in \mathbb{N}_{0}$ bilden ein Orthonormalsystem von $L^{2}(\mathbb{R}, \mathbb{C})$.\\
(ii) Für alle $n \in \mathbb{N}_{0}$ gilt $\widehat{H_{n}}=(-i)^{n} H_{n}$.

Wir zeigen nun, dass die Menge der Funktionen $H_{n}$ auch total in $L^{2}(\mathbb{R})$ und somit eine Orthonormalbasis ist. Sei dazu $f \in \overline{\operatorname{span}\left\{H_{n}: n \in \mathbb{N}_{0}\right\}} \subset L^{2}(\mathbb{R})$; zu zeigen ist $f \equiv 0$. Sei dazu $\xi \in \mathbb{R}$ beliebig. Da für $n \in \mathbb{N}_{0}$ die Funktion

$$
g_{n} \in L^{1}(\mathbb{R}) \cap L^{2}(\mathbb{R}), \quad g_{n}(x)=\frac{(i \xi x)^{n}}{n!} e^{-\frac{x^{2}}{2}}
$$

im Erzeugnis der Funktionen $H_{k}, k \leq n$ liegt, erhält man

$$
\int_{\mathbb{R}} f(x) \bar{g}_{n}(x) d x=0 \quad \text { für alle } n \in \mathbb{N} \cup\{0\} .
$$

Da die Reihe $\sum_{n=0}^{\infty} g_{n}$ in $L^{2}(\mathbb{R})$ gegen die Funktion

$$
g \in L^{1}(\mathbb{R}) \cap L^{2}(\mathbb{R}), \quad g(x)=e^{i x \xi-\frac{x^{2}}{2}}
$$

konvergiert, folgt mit (14.3) also

$$
0=\int_{\mathbb{R}} f(x) \overline{g(x)} d x=\int_{\mathbb{R}} e^{-i x \xi} f(x) e^{-\frac{x^{2}}{2}} d x=\sqrt{2 \pi} \hat{h}(\xi)
$$

für $h \in L^{1}(\mathbb{R}) \cap L^{2}(\mathbb{R})$ definiert durch $h(x)=f(x) e^{-\frac{x^{2}}{2}}$. Da $\xi \in \mathbb{R}$ beliebig gewählt war, folgt $\hat{h} \equiv 0$, und 14.11 liefert $h \equiv 0$. Es folgt $f \equiv 0$, was zu zeigen war.

\section*{§ 15 Der Rieszsche Darstellungssatz für Linearformen auf $\mathcal{C}_{c}(X)$}
Sei stets $(X, d)$ ein lokal kompakter und $\sigma$-kompakter metrischer Raum. Ziel dieses Kapitels ist eine Charakterisierung des Dualraums von ( $\mathcal{C}_{c}(X),\|\cdot\|_{\infty}$ ).

\subsection*{Definition}
Eine $\mathbb{R}$-lineare Abbildung $\Phi: \mathcal{C}_{c}(X, \mathbb{R}) \rightarrow \mathbb{R}$ heißt positiv, wenn gilt:

$$
\Phi(f) \geq 0 \text { für alle nichtnegativen Funktionen } f \in \mathcal{C}_{c}(X) \text {. }
$$

Mit $\Lambda_{+}(X)$ bezeichnen wir im Folgenden die Menge der positiven Linearformen auf dem Raum $\mathcal{C}_{c}(X, \mathbb{R})$.\\
Bemerkung: (15.1) ist offensichtlich äquivalent ist zur Monotonieeigenschaft

$$
\Phi(f) \leq \Phi(g) \quad \text { für } f, g \in \mathcal{C}_{c}(X) \text { mit } f \leq g .
$$

\subsection*{Satz und Definition}
Sei $\Phi \in \Lambda_{+}(X)$, und sei

$$
\mu(U)=\sup \left\{\Phi(f): f \in \mathcal{C}_{c}(X,[0,1]), \operatorname{supp} f \subset U\right\}
$$

für $U \subset X$ offen und

$$
\mu(A)=\inf \{\mu(U): U \subset X \text { offen mit } A \subset U\}
$$

für beliebiges $A \in \mathcal{P}(X)$. Dann ist $\mu$ ein äußeres Maß auf $X$, d.h. es gilt:\\
(i) $\mu(\varnothing)=0$.\\
(ii) Sind $A, A_{k} \subset X, k \in \mathbb{N}$ mit $A \subset \bigcup_{k=1}^{\infty} A_{k}$, so ist $\mu(A) \leq \sum_{k=1}^{\infty} \mu\left(A_{k}\right)$.

Ferner ist die Einschränkung von $\mu$ auf die Borelalgebra $\mathcal{B}(X)$ ein Radon-Maß. Wir nennen diese Einschränkung das von $\Phi$ induzierte Radon-Maß $\mu$.

Bemerkung: Für offene Teilmengen $A \subset B \subset X$ folgt aus (15.3) die Monotonieeigenschaft $\mu(A) \subset \mu(B)$. Daher stimmt für offene Teilmengen $A \subset X$ die Definition (15.4) mit (15.3) überein, d.h. $\mu$ ist wohldefiniert.

Wir verschieben den etwas längeren maßtheoretischen Beweis ans Ende des Kapitels.

\subsection*{Bemerkung}
Ist $\Phi \in \Lambda_{+}(X)$, so gilt für das von $\Phi$ induzierte Borelmaß $\mu$ auch

$$
\mu(U)=\sup \left\{\Phi(h): h \in \mathcal{C}_{c}(X, \mathbb{R}), 0 \leq h \leq 1_{U}\right\} \quad \text { für } U \subset X \text { offen. }
$$

Dabei ist ' $\leq$ ' offensichtlich. Zum Beweis von ' $\geq$ ' sei $h \in \mathcal{C}_{c}(X,[0,1])$ beliebig mit $h \leq 1_{U}$. Sei ferner $h_{\varepsilon}:=(h-\varepsilon)^{+} \in \mathcal{C}_{c}(X,[0,1])$ für $\varepsilon>0$. Dann ist $\operatorname{supp} h_{\varepsilon} \subset U$. Wählt man zudem $u \in \mathcal{C}_{c}(X,[0,1])$ mit $u \equiv 1$ auf $\operatorname{supp} h$, so gilt $h-h_{\varepsilon} \leq \varepsilon u$ in $X$ und somit $\Phi\left(h-h_{\varepsilon}\right) \leq \Phi(\varepsilon u)=\varepsilon \Phi(u)$. Es folgt

$$
\mu(U) \geq \Phi\left(h_{\varepsilon}\right)=\Phi(h)-\Phi\left(h-h_{\varepsilon}\right) \geq \Phi(h)-\varepsilon \Phi(u) \quad \text { für alle } \varepsilon>0
$$

und somit $\mu(U) \geq \Phi(h)$. Die Behauptung folgt.

\subsection*{Satz}
Sei $\Phi \in \Lambda_{+}(X)$ und $\mu$ das induzierte Radon-Maß auf $X$. Dann gilt

$$
\Phi(f)=\int_{X} f d \mu \quad \text { für alle } f \in \mathcal{C}_{c}(X, \mathbb{R})
$$

Beweis. Es reicht aufgrund der Linearität beider Seiten, (15.5) für eine gegebene nichtnegative Funktion $f \in \mathcal{C}_{c}(X, \mathbb{R})$ zu zeigen. Sei $\varepsilon>0$. Wir wählen $0=t_{0}< t_{1}<\cdots<t_{n}=2\|f\|_{\infty}$ mit $t_{i}-t_{i-1}<\varepsilon$ und $\mu\left(f^{-1}\left(t_{i}\right)\right)=0$ für $i=1, \ldots, n$. Sei $U_{i}:=f^{-1}\left(\left(t_{i-1}, t_{i}\right)\right)$ für $i=1, \ldots, n$. Dann ist $U_{i}$ offen und $\mu\left(U_{i}\right)<\infty$. Da $\mu$ ein Radon-Maß ist, existieren kompakte Mengen $K_{i} \subset U_{i}$ mit

$$
\mu\left(U_{i} \backslash K_{i}\right)<\frac{\varepsilon}{n} \quad \text { für } i=1, \ldots, n .
$$

Nach Definition von $\mu$ existieren ferner $g_{i} \in \mathcal{C}_{c}(X,[0,1]), i=1, \ldots, n$ mit $\operatorname{supp} g_{i} \subset U_{i}$ und $\mu\left(U_{i}\right)-\frac{\varepsilon}{n} \leq \Phi\left(g_{i}\right)$ für $i=1, \ldots, n$. Wähle nun wie im Beweis von 6.23(ii) Funktionen $h_{i} \in \mathcal{C}_{c}(X,[0,1])$ mit $\operatorname{supp} h_{i} \subset U_{i}$ und $h_{i} \equiv 1$ auf $\operatorname{supp} g_{i} \cup K_{i}$ für $i=1, \ldots n$. Dann ist $g_{i} \leq h_{i}$ für $i=1, \ldots, n$, und aufgrund der Monotonie von $\Phi$ und der Definition von $\mu$ folgt

$$
\mu\left(U_{i}\right)-\frac{\varepsilon}{n} \leq \Phi\left(h_{i}\right) \leq \mu\left(U_{i}\right) \quad \text { für } i=1, \ldots, n .
$$

Betrachte num $\tilde{f}:=f\left(1-\sum_{i=1}^{n} h_{i}\right) \in \mathcal{C}_{c}(X,[0, \infty))$ und die offene Menge $A:=\{x \in X: \tilde{f}(x)>0\}$ Es gilt $A \subset \bigcup_{i=1}^{n}\left[\left(U_{i} \backslash K_{i}\right) \cup f^{-1}\left(t_{i}\right)\right]$ und damit

$$
\mu(A) \leq \sum_{i=1}^{n} \mu\left(U_{i} \backslash K_{i}\right)<\varepsilon
$$

gemäß 15.6. Ferner folgt mit 15.3 .

$$
\begin{aligned}
& \Phi(\tilde{f}) \leq \sup \left\{\Phi(h): h \in \mathcal{C}_{c}(X, \mathbb{R}), 0 \leq h \leq\|f\|_{\infty} 1_{A}\right\} \\
& =\|f\|_{\infty} \sup \left\{\Phi(h): h \in \mathcal{C}_{c}(X, \mathbb{R}), 0 \leq h \leq 1_{A}\right\}=\|f\|_{\infty} \mu(A) \leq \varepsilon\|f\|_{\infty}
\end{aligned}
$$

Mit (15.7) folgt nun

$$
\begin{aligned}
\Phi(f)=\Phi\left(f \sum_{i=1}^{n} h_{i}\right)+\Phi(\tilde{f}) \leq \Phi\left(\sum_{i=1}^{n} t_{i} h_{i}\right)+\varepsilon\|f\|_{\infty} & =\sum_{i=1}^{n} t_{i} \Phi\left(h_{i}\right)+\varepsilon\|f\|_{\infty} \\
& \leq \sum_{i=1}^{n} t_{i} \mu\left(U_{i}\right)+\varepsilon\|f\|_{\infty}
\end{aligned}
$$

und

$$
\begin{aligned}
\Phi(f) \geq \Phi\left(f \sum_{i=1}^{n} h_{i}\right) \geq \Phi\left(\sum_{i=1}^{n} t_{i-1} h_{i}\right)=\sum_{i=1}^{n} t_{i-1} \Phi\left(h_{i}\right) & \geq \sum_{i=1}^{n} t_{i-1}\left(\mu\left(U_{i}\right)-\frac{\varepsilon}{n}\right) \\
& \geq \sum_{i=1}^{n} t_{i-1} \mu\left(U_{i}\right)-t_{n} \varepsilon
\end{aligned}
$$

Da weiterhin aufgrund der Monotonie des Integrals

$$
\sum_{i=1}^{n} t_{i-1} \mu\left(U_{i}\right) \leq \int_{X} f d \mu \leq \sum_{i=1}^{n} t_{i} \mu\left(U_{i}\right)
$$

gilt, erhält man

$$
\left|\Phi(f)-\int_{X} f d \mu\right| \leq \sum_{i=1}^{n}\left(t_{i}-t_{i-1}\right) \mu\left(U_{i}\right)+\varepsilon\left(\|f\|_{\infty}+t_{n}\right) \leq \varepsilon \mu(\operatorname{supp} f)+3 \varepsilon\|f\|_{\infty}
$$

Durch Grenzübergang $\varepsilon \rightarrow 0$ folgt nun (15.5), wie behauptet.

Im Folgenden sei $\mathcal{C}_{c}(X)=\mathcal{C}_{c}(X, \mathbb{K})$ mit $\mathbb{K}=\mathbb{R}$ oder $\mathbb{K}=\mathbb{C}$. Wir bemerken, dass $\mathcal{C}_{c}(X)$, versehen mit $\|\cdot\|_{\infty}$, im Allgemeinen nicht vollständig ist; die Vervollständigung ist gegeben durch den Raum $\mathcal{C}_{0}(X)$ der stetigen Funktionen $f: X \rightarrow \mathbb{K}$ derart, dass für jedes $\varepsilon>0$ eine kompakte Menge $K \subset X$ existiert mit $|f| \leq \varepsilon$ auf $X \backslash K$. Ist $X$ selbst kompakt, so gilt offensichtlich $\mathcal{C}_{c}(X)=\mathcal{C}(X)=\mathcal{C}_{0}(X)$, und man hat somit bereits einen Banachraum.

\subsection*{Satz}
Sei $\Psi \in \mathcal{C}_{c}(X)^{\prime}$. Dann existiert ein $\Phi \in \Lambda_{+}(X)$ mit

$$
\Phi(f):=\sup \left\{|\Psi(g)|: g \in \mathcal{C}_{c}(X),|g| \leq f\right\} \quad \text { für alle } f \in \mathcal{C}_{c}(X,[0, \infty))
$$

Ist ferner $\mu$ das von $\Phi$ induzierte Radon-Maß, so gilt für offene Teilmengen $U \subset X$ :

$$
\mu(U)=\sup \left\{|\Psi(g)|: g \in \mathcal{C}_{c}(X), \operatorname{supp} g \subset U,\|g\|_{\infty} \leq 1\right\}
$$

Beweis. Für $f \in \mathcal{C}_{c}(X,[0, \infty))$ sei $\Phi(f)$ wie in (15.8) definiert; dann gilt $\Phi(f) \geq 0$ und $\Phi(c f)=c \Phi(f)$ für $c \geq 0$. Ferner gilt $\Phi\left(f_{1}\right) \leq \Phi\left(f_{2}\right)$ für $f_{1}, f_{2} \in \mathcal{C}(X,[0, \infty))$, $f_{1} \leq f_{2}$. Wir zeigen nun:

$$
\Phi\left(f_{1}+f_{2}\right)=\Phi\left(f_{1}\right)+\Phi\left(f_{2}\right) \quad \text { für } f_{1}, f_{2} \in \mathcal{C}(X,[0, \infty)) .
$$

Seien dazu $g_{1}, g_{2} \in \mathcal{C}_{c}(X)$ beliebig mit $\left|g_{i}\right| \leq f_{i}$ für $i=1,2$, und seien $\alpha_{i} \in \mathbb{K}$ mit $\left|\alpha_{i}\right|=1$ und $\alpha_{i} \Psi\left(g_{i}\right)=\left|\Psi\left(g_{i}\right)\right|$ für $i=1,2$ gewählt. Dann ist $\left|\alpha_{1} g_{1}+\alpha_{2} g_{2}\right| \leq f_{1}+f_{2}$ und

$$
\left|\Psi\left(g_{1}\right)\right|+\left|\Psi\left(g_{2}\right)\right|=\Psi\left(\alpha_{1} g_{1}+\alpha_{2} g_{2}\right) \leq\left|\Psi\left(\alpha_{1} g_{1}+\alpha_{2} g_{2}\right)\right| \leq \Phi\left(f_{1}+f_{2}\right)
$$

Durch Supremumsbildung folgt also $\Phi\left(f_{1}\right)+\Phi\left(f_{2}\right) \leq \Phi\left(f_{1}+f_{2}\right)$. Zum Beweis von ' $\geq$ ' sei $g \in \mathcal{C}_{c}(X)$ beliebig mit $|g| \leq f_{1}+f_{2}$. Für $i=1,2$ sei ferner

$$
g_{i}: X \rightarrow \mathbb{K}, \quad g_{i}(x)= \begin{cases}\frac{g(x) f_{i}(x)}{f_{1}(x)+f_{2}(x)}, & f_{1}(x)+f_{2}(x)>0 \\ 0, & f_{1}(x)+f_{2}(x)=0\end{cases}
$$

Mittels der Ungleichung $|g| \leq f_{1}+f_{2}$ sieht man dann leicht, dass $g_{i} \in \mathcal{C}_{c}(X)$ für $i=1,2$ gilt. Ferner ist $\left|g_{i}\right| \leq f_{i}$ in $X$ und somit

$$
|\Psi(g)|=\left|\Psi\left(g_{1}+g_{2}\right)\right| \leq\left|\Psi\left(g_{1}\right)\right|+\left|\Psi\left(g_{2}\right)\right| \leq \Phi\left(f_{1}\right)+\Phi\left(f_{2}\right) .
$$

Durch Supremumsbildung folgt $\Phi\left(f_{1}+f_{2}\right) \leq \Phi\left(f_{1}\right)+\Phi\left(f_{2}\right)$, und insgesamt folgt (15.10). Wir setzen nun

$$
\Phi(f):=\Phi\left(f^{+}\right)-\Phi\left(f^{-}\right) \quad \text { für } f \in \mathcal{C}_{c}(X, \mathbb{R}) .
$$

Unter Verwendung von (15.10) sieht man dann leicht, dass $\Phi$ linear ist.\\
Die Charakterisierung (15.9) folgt direkt aus (15.3) und (15.8).

\subsection*{Hauptsatz (Rieszscher Darstellungssatz)}
Zu jedem $\Psi \in \mathcal{C}_{c}(X)^{\prime}$ existiert ein Radon-Maß $\mu$ und eine $\mu$-messbare Funktion $\alpha: X \rightarrow \mathbb{K}$ mit $|\alpha|=1 \mu$-f.ü. auf $X$ und derart, dass

$$
\Psi(f)=\int_{X} f \alpha d \mu \quad \text { für alle } f \in \mathcal{C}_{c}(X)
$$

gilt. Ist ferner $X$ kompakt, so gilt

$$
\|\Psi\|=\mu(X)
$$

Ferner ist $\mu$ durch (15.11) auf Borelmengen eindeutig bestimmt, und $\alpha$ ist $\mu$-f.ü. eindeutig bestimmt.

Beweis. Zur Existenz: Sei $\Phi \in \Lambda_{+}(X)$ wie in 15.5 gewählt, und sei $\mu$ das von $\Phi$ induzierte Radon-Maß. Dann gilt gemäß (15.8) und (15.5)

$$
|\Psi(f)| \leq \Phi(|f|)=\int_{X}|f| d \mu=\|f\|_{L^{1}(X, d \mu)} \quad \text { für } f \in \mathcal{C}_{c}(X) \text {. }
$$

Sei im Folgenden $L^{p}(X)=L^{p}(X, d \mu)$ für $1 \leq p \leq \infty$. Da $\mathcal{C}_{c}(X)$ gemäß 6.25 in $L^{1}(X, d \mu)$ dicht liegt, lässt sich $\Psi$ zu einer stetigen Linearform auf $L^{1}(X, d \mu)$ mit $\|\Psi\|_{\left(L^{1}(X, d \mu)\right)^{\prime}} \leq 1$ fortsetzen. Nach 12.21 existiert also ein $\alpha \in L^{\infty}(X, d \mu)$ mit $\|\alpha\|_{L^{\infty}(X, d \mu)}=\|\Psi\|_{\left(L^{1}(X)\right)^{\prime}} \leq 1$ und derart, dass (15.11) gilt. Es bleibt $|\alpha|=1 \mu$ f.ü. in $X$ zu zeigen. Sei dazu $U \subset X$ eine beliebige offene Menge mit $\mu(U)<\infty$. Nach (15.9) existiert dann eine Folge von Funktionen $g_{n} \in \mathcal{C}_{c}(X)$ mit $\operatorname{supp} g_{n} \subset U$, $\left\|g_{n}\right\|_{\infty} \leq 1$ für $n \in \mathbb{N}$ und

$$
\mu(U)=\lim _{n \rightarrow \infty}\left|\Psi\left(g_{n}\right)\right| \stackrel{15.11}{=} \lim _{n \rightarrow \infty}\left|\int_{U} \alpha g_{n} d \mu\right| \leq \int_{U}|\alpha| d \mu
$$

und somit

$$
\int_{U}(1-|\alpha|) d \mu=\mu(U)-\int_{U}|\alpha| d \mu \leq 0 .
$$

Da $1-|\alpha| \geq 0 \mu$-f.ü. in $U$ gilt, folgt $|\alpha|=1 \mu$-f.ü. in $U$. Aufgrund der beliebigen Wahl von $U$ folgt $|\alpha|=1 \mu$-f.ü. in $X$.\\
Ist schließlich $X$ kompakt, so folgt mit (15.13)

$$
|\Psi(f)| \leq\|f\|_{L^{1}(X, d \mu)} \leq \mu(X)\|f\|_{\infty} \quad \text { für alle } f \in \mathcal{C}_{c}(X)
$$

und somit $\|\Psi\| \leq \mu(X)$. Ferner gilt (15.14) in diesem Fall auch für $U=X$, und somit folgt

$$
\mu(X) \leq \limsup _{n \rightarrow \infty}\|\Psi\|\left\|g_{n}\right\|_{\infty} \leq\|\Psi\| .
$$

Insgesamt folgt also (15.12).\\
Zur Eindeutigkeit: Sei $\mu$ ein Radon-Maß und $\alpha: X \rightarrow \mathbb{K}$ eine $\mu$-messbare Funktion mit $|\alpha|=1$ f.ü. auf $X$ und derart, dass (15.11) gilt. Wir zeigen, dass dann

$$
\mu(U)=\sup \left\{|\Psi(g)|: g \in \mathcal{C}_{c}(X), \operatorname{supp} g \subset U,\|g\|_{\infty} \leq 1\right\}
$$

für $U \subset X$ offen mit $\mu(U)<\infty$ gilt; damit ist $\mu$ auf Borelmengen von $X$ eindeutig bestimmt. Wegen (15.11) gilt

$$
|\Psi(g)| \leq \int_{X}|g| d \mu \leq \mu(U) \quad \text { für } g \in \mathcal{C}_{c}(X) \text { mit } \operatorname{supp} g \subset U \text { und }\|g\|_{\infty} \leq 1 \text {, }
$$

also folgt ' $\geq$ ' in (15.15). Zum Beweis von ' $\leq$ ' bemerken wir zunächst, dass $U$ als metrischer Raum ebenfalls lokal kompakt und die Einschränkung von $\mu$ auf die Borel-Teilmengen von $U$ wieder ein Radon-Maß ist. Anwendung des Satzes von Lusin 6.23(ii) mit $U$ anstelle von $X$ auf die Funktion $\bar{\alpha}$ liefert für jedes $n \in \mathbb{N}$ somit eine Funktion $g_{n} \in \mathcal{C}_{c}(X)$ mit $\operatorname{supp} g_{n} \subset U,\left\|g_{n}\right\|_{\infty} \leq\|\alpha\|_{\infty} \leq 1$ und

$$
\mu\left(\left\{x \in U: g_{n}(x) \neq \bar{\alpha}\right\}\right) \leq \frac{1}{n} .
$$

Insbesondere folgt $g_{n} \rightarrow \bar{\alpha}$ in $L^{1}(U)$ und nach Übergang zu einer Teilfolge auch punktweise in $U$. Der Konvergenzsatz von Lebesgue liefert dann

$$
\mu(U)=\int_{U}|\alpha|^{2} d \mu=\lim _{n \rightarrow \infty} \int_{X} \alpha g_{n} d \mu \stackrel{\text { 15.11 }}{=} \lim _{n \rightarrow \infty} \Psi\left(g_{n}\right)
$$

Somit folgt ' $\leq$ ' in (15.15), und insgesamt folgt Gleichheit.\\
Sei schließlich $\beta: X \rightarrow \mathbb{K}$ eine weitere $\mu$-messbare Funktion mit $|\beta|=1$ f.ü. auf $X$ und

$$
\Psi(f)=\int_{X} f \alpha d \mu=\int_{X} f \beta d \mu \quad \text { für alle } f \in \mathcal{C}_{c}(X)
$$

Durch Betrachtung von $U \subset X$ offen mit $\mu(U)<\infty$ und $f=g_{n}, n \in \mathbb{N}$ wie oben folgt dann wiederum

$$
\mu(U)=\lim _{n \rightarrow \infty} \Psi\left(g_{n}\right)=\lim _{n \rightarrow \infty} \int_{X} g_{n} \beta d \mu=\int_{U} \bar{\alpha} \beta d \mu
$$

Da $|\bar{\alpha} \beta|=1$ f.ü. in $X$ gilt, folgt $\bar{\alpha} \beta=1 \mu$-f.ü. in $U$ und somit auch in $X$ aufgrund der beliebigen Wahl von $U$. Dies zeigt $\alpha=\beta \mu$-f.ü. in $X$. \(\square\)

Zum Abschluss des Kapitels liefern wir nun den

\section*{Beweis von 15.2 ,}
\begin{enumerate}
  \item Schritt: $\mu$ ist ein äußeres Maß auf $X$.\\
$\mathrm{Zu}(\mathrm{i}):$ Aus (15.3) folgt offensichtlich die Eigenschaft $\mu(\varnothing)=0$.\\
Zu (ii): Seien $A_{k} \subset X, k \in \mathbb{N}$ und $A \subset \bigcup_{k \in \mathbb{N}} A_{k}$ gegeben. Sei zunächst angenommen, dass $A_{k}$ und $A$ offene Mengen sind, und sei $f \in \mathcal{C}_{c}(X,[0,1])$ eine Funktion mit $K:=\operatorname{supp} f \subset A$. Da $K$ kompakt ist, existiert $l \in \mathbb{N}$ mit $K \subset \bigcup_{k=1}^{l} A_{k}$. Wähle nun
\end{enumerate}

Funktionen $\xi_{k} \in \mathcal{C}_{c}(X,[0,1]), k=1, \ldots, l$ mit $\operatorname{supp} \xi_{k} \subset A_{k}$ und $\sum_{k=1}^{l} \xi_{k} \equiv 1$ auf $K$ (endliche Zerlegung der Eins). Dann ist $f \leq \sum_{k=1}^{l} f \xi_{k}$ und somit wegen 15.2

$$
\Phi(f) \leq \sum_{k=1}^{l} \Phi\left(f \xi_{k}\right) \leq \sum_{k=1}^{l} \mu\left(A_{k}\right) \leq \sum_{k \in \mathbb{N}} \mu\left(A_{k}\right) .
$$

Supremumsbildung über $f$ liefert also

$$
\mu(A) \leq \sum_{k \in \mathbb{N}} \mu\left(A_{k}\right) .
$$

Seien nun $A$ und $A_{k}, k \in \mathbb{N}$ beliebig (d.h. nicht notwendig offen), sei $\varepsilon>0$, und seien $V_{k}, k \in \mathbb{N}$ offene Mengen mit $A_{k} \subset V_{k}$ und $\mu\left(V_{k}\right) \leq \mu\left(A_{k}\right)+\varepsilon 2^{-k}$. Dann ist nach Definition und dem bisher Gezeigten

$$
\mu(A) \leq \mu\left(\bigcup_{k \in \mathbb{N}} V_{k}\right) \leq \sum_{k \in \mathbb{N}} \mu\left(V_{k}\right) \leq \varepsilon+\sum_{k \in \mathbb{N}} \mu\left(A_{k}\right) .
$$

Da $\varepsilon>0$ beliebig gewählt war, folgt $\mu(A) \leq \sum_{k \in \mathbb{N}} \mu\left(A_{k}\right)$.\\
Wir haben somit gezeigt, dass $\mu$ ein äußeres Maß ist.\\
2. Schritt: Die Einschränkung von $\mu$ auf die Borelalgebra $\mathcal{B}(X)$ ist ein Borelmaß. Dies folgt mit dem sogenannten\\
Kriterium von Carathéodory: Sei ( $X, d$ ) ein metrischer Raum und $\mu$ ein äußeres Maß auf X. Gilt

$$
\mu(A \cup B)=\mu(A)+\mu(B)\left\{\begin{array}{l}
\text { für } A, B \subset X \text { mit } \\
\operatorname{dist}(A, B):=\inf \{d(x, y): x \in A, y \in B\}>0
\end{array}\right.
$$

so ist jede Borelmenge $A \subset X \mu$-messbar im Sinne von Carathéodory, d.h. es gilt

$$
\mu(Z)=\mu(Z \cap A)+\mu(Z \backslash A) \quad \text { für } Z \subset X \text { beliebig. }
$$

Wir prüfen zunächst die Voraussetzung 15.17):\\
Die Ungleichung " $\leq$ " ist aufgrund der Subadditivität klar, zu zeigen bleibt also " $\geq$ ". Sei dazu zunächst angenommen, dass $A, B \subset X$ offen sind, und sei $\delta>0$. Dann existieren Funktionen $f, g \in \mathcal{C}_{c}(X,[0,1])$ mit $\operatorname{supp} f \subset A, \operatorname{supp} g \subset B$ und

$$
\mu(A) \leq \Phi(f)+\frac{\delta}{2} \quad \text { sowie } \quad \mu(B) \leq \Phi(g)+\frac{\delta}{2} .
$$

Da $A \cap B=\varnothing$ gilt, ist $h:=f+g \in \mathcal{C}_{c}(X,[0,1])$ mit $\operatorname{supp} h \subset A \cup B$, und somit gilt

$$
\mu(A \cup B) \geq \Phi(h)=\Phi(f)+\Phi(g)=\mu(A)+\mu(B)+\delta .
$$

Da $\delta>0$ beliebig war, folgt die Behauptung in diesem Spezialfall.\\
Im Allgemeinen Fall sei $\varepsilon:=\operatorname{dist}(A, B)>0, U \subset X$ eine beliebige offene Menge mit $A \cup B \subset U$, und sei

$$
\tilde{A}:=U_{\varepsilon / 4}(A) \cap U, \quad \tilde{B}:=U_{\varepsilon / 4}(B) \cap U .
$$

Dann $\operatorname{sind} \tilde{A}, \tilde{B} \subset X$ offen mit $\operatorname{dist}(\tilde{A}, \tilde{B}) \geq \frac{\varepsilon}{2}$, und somit folgt mit dem Spezialfall und der Monotonie

$$
\mu(A)+\mu(B) \leq \mu(\tilde{A})+\mu(\tilde{B})=\mu(\tilde{A} \cup \tilde{B}) \leq \mu(U)
$$

Aufgrund der beliebigen Wahl von $U$ folgt also $\mu(A)+\mu(B) \leq \mu(A \cup B)$. Die zeigt Voraussetzung (15.17) im vorliegenden Fall.\\
3. Schritt: Die Einschränkung von $\mu$ auf $\mathcal{B}(X)$ ist ein Radon-Maß.

Nach Definition ist $\mu$ von außen regulär. Ferner gilt

$$
\mu(K)<\infty \quad \text { für jede kompakte Menge } K \subset X,
$$

denn: Gemäß 6.15(i) existiert $u \in \mathcal{C}_{c}(X,[0,1])$ mit $u \equiv 1$ in einer offenen Umgebung $U$ von $K$. Für alle $f \in \mathcal{C}_{c}(X,[0,1])$ mit $\operatorname{supp} f \subset U$ gilt dann $\Phi(f) \leq \Phi(u)$, und somit folgt $\mu(U) \leq \Phi(u)$. Somit ist auch $\mu(K) \leq \mu(U) \leq \Phi(u)<\infty$.\\
Gemäß der Bemerkung in 6.19 ist $\mu$ somit lokal endlich, und wegen 6.20(i) ist $\mu$ somit ein Radon-Maß.\\
4. Schritt: Beweis des obigen Kriteriums von Carathéodory.

Dazu reicht es, die Eigenschaft (15.18) für abgeschlossene Teilmengen $A \subset X$ zu beweisen (vgl. das Lemma von Carathéodory aus der Analysis II). Sei dazu $Z \subset X$ beliebig. Es reicht dann, die Ungleichung " $\geq$ " in (15.20) zu zeigen (" $\leq$ " folgt direkt aus den Eigenschaften eines äußeren Maßes):

$$
\mu(Z) \geq \mu(Z \cap A)+\mu(Z \backslash A) .
$$

Dies ist trivial, falls $\mu(Z)=\infty$ ist. Sei nun $\mu(Z)<\infty$ angenommen, und sei $A_{n}:= \left\{x \in X: \operatorname{dist}(x, A) \leq \frac{1}{n}\right\}$ für $n \in \mathbb{N}$. Dann ist $\operatorname{dist}\left(Z \backslash A_{n}, Z \cap A\right)>0$ und somit

$$
\mu\left(Z \cap A_{n}\right)+\mu(Z \backslash A)=\mu\left(\left[Z \cap A_{n}\right] \cup[Z \backslash A]\right) \leq \mu(Z),
$$

nach Voraussetzung und aufgrund des Monotonie von $\mu$. Es reicht also, zu zeigen:

$$
\lim _{n \rightarrow \infty} \mu\left(Z \cap A_{n}\right)=\mu(Z \cap A) .
$$

Sei dazu

$$
R_{k}:=\left\{x \in Z: \frac{1}{k+1}<\operatorname{dist}(x, A) \leq \frac{1}{k}\right\} \subset Z \quad \text { für } k \in \mathbb{N} \text {. }
$$

Aufgrund der Abgeschlossenheit von $A$ ist dann $Z \backslash A=\left[Z \backslash A_{n}\right] \cup \bigcup_{k=n}^{\infty} R_{k}$ und somit

$$
\mu\left(Z \backslash A_{n}\right) \leq \mu(Z \backslash A) \leq \mu\left(Z \backslash A_{n}\right)+\sum_{k=n}^{\infty} \mu\left(R_{k}\right)
$$

Es reicht also, zu zeigen:

$$
\lim _{n \rightarrow \infty} \sum_{k=n}^{\infty} \mu\left(R_{k}\right)=0, \quad \text { bzw. äquivalent: } \quad \sum_{k=1}^{\infty} \mu\left(R_{k}\right)<\infty .
$$

Nach Voraussetzung (und per Induktion) gelten dabei für $m \in \mathbb{N}$ die Identitäten

$$
\sum_{k=1}^{m} \mu\left(R_{2 k}\right)=\mu\left(\bigcup_{k=1}^{m} R_{2 k}\right) \quad \text { und } \quad \sum_{k=0}^{m} \mu\left(R_{2 k+1}\right)=\mu\left(\bigcup_{k=1}^{m} R_{2 k+1}\right),
$$

da die hier jeweils betrachten Mengenfamilien paarweise positiven Abstand haben. Aufgrund der Monotonie folgt somit

$$
\sum_{k=1}^{2 m+1} \mu\left(R_{k}\right) \leq 2 \mu(Z)<\infty \quad \text { für alle } m \in \mathbb{N} \text {. }
$$

Dies zeigt (15.22), wie benötigt.



\pagebreak
\printbibliography
\end{document}